% !TeX program = xelatex
% XeLaTeX 编译两遍以更新目录。全书在单个源文件中。
\documentclass[UTF8,a4paper,zihao=-4,fontset=fandol]{ctexrep}
\usepackage[left=23mm,right=23mm,top=23mm,bottom=24mm,headheight=15pt,headsep=7mm,footskip=11mm]{geometry}
\usepackage{amsmath,amssymb,bm}
\usepackage{booktabs,tabularx,longtable,array,multirow}
\usepackage{enumitem,xcolor,needspace,fancyhdr,newunicodechar,tocloft}
\usepackage[unicode,colorlinks,linkcolor=navy,urlcolor=navy]{hyperref}
\usepackage{bookmark}
\definecolor{navy}{HTML}{183A53}
\definecolor{muted}{HTML}{51616C}
\newunicodechar{★}{\texorpdfstring{\ensuremath{\bigstar}}{★}}
\hypersetup{pdftitle={哈尔滨工业大学模式识别与机器学习A复习考点手册},bookmarksdepth=1,pdfdisplaydoctitle=true}
\setcounter{tocdepth}{0}
\setlength{\cftbeforechapskip}{.4em}
\setcounter{secnumdepth}{0}
\ctexset{
  chapter={format=\Huge\bfseries\color{navy},name={第,章},number=\arabic{chapter},beforeskip=0pt,afterskip=1.1em},
  section={format=\Large\bfseries\color{navy},beforeskip=1.25em plus .2em,afterskip=.55em},
  subsection={format=\large\bfseries\color{navy},beforeskip=.9em plus .15em,afterskip=.28em},
  paragraph={format=\normalsize\bfseries\color{navy},beforeskip=.65em,afterskip=.7em,runin=true}
}
\setlength{\parindent}{0pt}
\setlength{\parskip}{.32em plus .07em minus .05em}
\setlength{\emergencystretch}{2em}
\setlength{\LTleft}{0pt}
\setlength{\LTright}{0pt plus 1fill}
\setlength{\LTpre}{.35em}
\setlength{\LTpost}{.65em}
\renewcommand{\arraystretch}{1.16}
\allowdisplaybreaks[2]
\setlist[itemize]{leftmargin=1.65em,topsep=.35em,itemsep=.43em,parsep=.11em}
\setlist[enumerate]{leftmargin=2em,topsep=.35em,itemsep=.48em,parsep=.11em}
\setlist[description]{leftmargin=0pt,labelindent=0pt,font=\bfseries,itemsep=.48em,style=nextline}
\clubpenalty=10000
\widowpenalty=10000
\raggedbottom
\pagestyle{fancy}
\fancyhf{}
\fancyhead[L]{\small\textcolor{muted}{模式识别与机器学习 A}}
\fancyhead[R]{\small\textcolor{muted}{期末复习手册}}
\fancyfoot[C]{\small\textcolor{muted}{\thepage}}
\renewcommand{\headrulewidth}{.36pt}
\fancypagestyle{plain}{\fancyhf{}\fancyfoot[C]{\small\textcolor{muted}{\thepage}}\renewcommand{\headrulewidth}{0pt}}

\begin{document}
\begin{titlepage}
\thispagestyle{empty}
\vspace*{.15\textheight}
{\Large\color{muted}哈尔滨工业大学\par}
\vspace{1.2em}
{\Huge\bfseries\color{navy}模式识别与机器学习 A\par}
\vspace{.45em}
{\Huge\bfseries\color{navy}期末复习考点手册\par}
\vspace{1.5em}
{\color{navy}\rule{.77\linewidth}{.9pt}\par}
\vspace{1em}
{\large 刘远超、刘扬课件与本地往年题\par}
{\large 公式推导、手算例题、真题得分点与章末自测\par}
\vfill
{\small 课件出处中的 p 采用 PDF 物理页。往年题与模拟题分别编号。\par}
\end{titlepage}
\renewcommand{\contentsname}{目录}
\tableofcontents
\clearpage
\setcounter{chapter}{-1}

\chapter{总览与复习路线}

\section{本章考点清单}
\begin{itemize}
\item ★★★ 读出设定、目标与所求，区分概率、误差、风险、距离和损失。常见题会把「定义」「手算」「解释原因」放在同一道题里。 【2024秋真题 大题2、4–6】【2025秋补缓考 简答1、4】【刘远超·模拟试题范本 p2】
\item ★★★ 先通读最常重现的计算链：后验与风险、回归残差与梯度、最大间隔、聚类与 EM、深层网络的前向和反传。 【2019秋真题 第4、7题】【2022秋真题 第3、5题】【2024秋真题 大题2–5、7、9】【2025秋补缓考 简答1–4、6–8】
\item ★★ 用两位老师的课件互补。先按刘远超的应用例子代入数字，再看刘扬的符号推导。遇到不同约定写出含义和维度。 【刘远超·课程简介 p6–10】【刘扬·课程与实验大纲】【模拟卷三 大题1–8】【模拟卷四 大题1–8】
\item ★ 看时间分配、可见分值与题号标注。回忆中的缺项保留缺项。试题范本和模拟卷分别用其自己的题号。 【2020春真题 p1–4】【2024秋真题 回忆版 DOCX】【2025秋补缓考 回忆 PDF p1–3】【刘远超·模拟试题范本 p1–2】
\end{itemize}

\section{核心概念与定义}
\subsection{课程知识地图}
机器学习（machine learning）围绕任务、经验与可度量表现安排训练。模式识别（pattern recognition）从所测对象获取模式表示，再进行分类、检测、分组或决策。深度学习（deep learning）使用多层可学习的表示与任务目标共同训练。刘远超引论以学习、模式识别和人工智能（Artificial Intelligence, AI）的关系开篇，课程概述转向数据、评价设置和半监督方法。刘扬的课件从决策树与概率模型进入统计学习和深度模型。【刘远超·1 引论 p5–10、p17–21】【刘远超·2 机器学习概述 p17–47】【刘扬·1 决策树】【刘扬·2 概率论速通】

地图中的缩写包括支持向量机（Support Vector Machine, SVM）、期望最大化（Expectation-Maximization, EM）、隐马尔可夫模型（Hidden Markov Model, HMM）、反向传播（Backpropagation, BP）、卷积神经网络（Convolutional Neural Network, CNN）、循环神经网络（Recurrent Neural Network, RNN）、自编码器（Autoencoder, AE）与变分自编码器（Variational Autoencoder, VAE）。MSE 是均方误差（Mean Squared Error），KL 是库尔贝克–莱布勒散度（Kullback–Leibler Divergence）。

刘远超从已有数据、模型假设和参数求解来解释机器学习，又用多层特征组合解释深度学习。刘扬的决策树导论引用 Mitchell 的任务 $T$、表现 $P$、经验 $E$。以垃圾邮件分类为例，$T$ 是给新邮件预测标签，$E$ 是已有标注邮件，$P$ 可以是验证集 F1。刘远超还对照形式规则的符号主义、连接网络的连接主义、通过感知和动作联系解释智能的行为主义。计算平台侧，CPU 擅长通用控制，GPU 适合批量并行算子。CUDA 提供 GPU 计算平台，cuDNN 提供部分深度学习算子，PyTorch 和 TensorFlow 面向模型与训练实验，虚拟环境隔离依赖。这里掌握分工和用途即可。【刘远超·1 引论 p5–10、p20】【刘扬·1 决策树 p7】【刘远超·2 机器学习概述 p4–14】

\begin{tabularx}{\linewidth}{@{}p{.20\linewidth}XX@{}}
\toprule
入口 & 向后连接 & 考场上要会做的事\\
\midrule
数据与评价 & 第 1、2 章：拆分、概率、信息 & 写明独立性和对数底数，找到正确分母，读指标而不混用\\
监督决策 & 第 3–9 章：树、风险、近邻、回归、感知机、SVM & 从目标写决策规则，按题设估参，说明边界与误判代价\\
无监督表示 & 第 10、11 章：聚类、EM、降维 & 按总误差或似然优化，写清硬／软分配、投影、重构\\
概率序列 & 第 12 章：HMM & 用状态和观测连乘，决定求和、取最大、后验或更新\\
深层网络 & 第 13–17 章：BP、卷积、序列网络、AE/VAE & 画信息流和梯度流，以维度、参数共享和目标分解作检查\\
综合应用 & 第 18、19 章：实验、速查 & 把模型、目标、实现、验证和结论串起来\\
\bottomrule
\end{tabularx}

\subsection{读题时区分的四种产物}
\begin{enumerate}
\item \textbf{分数}可以是超平面带符号值、神经网络得分（logit）或未归一化的后验权重。分数的尺度由模型定义。
\item \textbf{概率与密度}有不同归一方式。连续点密度允许超过 1。条件后验是给定样本后对类别归一，符号 $p(x\mid y)$ 的含义要按离散或连续特征决定。
\item \textbf{预测与决策}除模型输出外还依赖阈值、损失表和任务目标。最大后验对应零一损失，不同误判代价下比较条件风险。
\item \textbf{训练与评价}的目标可能不同。训练用于更新权重，验证用于选择，测试用于评估固定后的流程。例如网络用交叉熵训练，取 F1 或准确率评价预测分类。
\end{enumerate}
【刘扬·3 贝叶斯判别 p1–20】【刘扬·4 朴素贝叶斯 p25–28】【刘扬·5 逻辑回归与SVM】【刘远超·2 机器学习概述 p28–34】

\section{公式与完整推导}
\subsection{共通的三条推理链}
\textbf{从联合到决策。} 条件定义给 $P(x,y)=P(x\mid y)P(y)$，对标签求和得到 $P(x)=\sum_jP(x\mid y_j)P(y_j)$，故 $P(y_i\mid x)=P(x\mid y_i)P(y_i)/P(x)$。选动作 $a$ 时对真实标签加权求损失 $\sum_i L(a,y_i)P(y_i\mid x)$，并选风险最小动作。样本密度相同时也需乘类别先验。对应第 2、4、6、12 章。

\textbf{从目标到更新。} 给模型参数 $\theta$，先写预测或似然，再写 $J(\theta)=\frac1n\sum_iL_i(\theta)+\lambda\Omega(\theta)$。若各部分可微，和式求导得 $\nabla J=\frac1n\sum_i\nabla L_i+\lambda\nabla\Omega$。梯度下降 $\theta'=\theta-\eta\nabla J$。闭式法常令梯度为零求解，但是否为最优还要看凸性或二阶条件。对应第 7、9、13、14 章。

\textbf{从不可见量到可计算目标。} 观察概率 $p(x)=\sum_zp(x,z)$。引入归一化的分布 $q(z)$，把 $\ln p(x)$ 分解为 $\mathbb E_q[\ln p(x,z)-\ln q(z)]+D_{\mathrm{KL}}(q\Vert p(z\mid x))$。KL 非负，得到下界。EM 的 E 步让后验贴合旧模型，VAE 则用可训练的编码器近似后验。推导、参数顺序和实现见第 10、12、17 章。 【刘扬·8GMM-EM p19–33】【刘扬·14VAE】【刘远超·10 生成式神经网络】

\section{计算题型与解题模板}
\paragraph{通用模板。}
把条件记成五栏：已知数据、指定符号、假设、要求的量、结果单位。先把通式写在等号左侧，下一行才代入。留能核验的中间值，例如叶子类别计数、$t=2$ 的前向向量、每个 CNN 的有效核与通道数。最后检查概率和为 1、总分母是否含样本数、参数维度是否和题面一致。

\paragraph{原创例：从密度到最小风险。}
两类先验为 $P(A)=1/4$、$P(B)=3/4$，某一点的类条件密度分别为 $0.6,0.2$。\(P(A)p(x\mid A)=0.15\)，\(P(B)p(x\mid B)=0.15\)，归一后后验都为 $1/2$。若错把 $A$ 判成 $B$ 的损失是 4，错把 $B$ 判成 $A$ 的损失是 1，正确为 0，则预测 $A$ 的风险 $1/2$，预测 $B$ 的风险 $2$，选择 $A$。题目若只问最小错误率，则两个动作最优错误率相同。密度与损失方向分别检查。【补充：原创复习例】【见第4章】

\paragraph{如何按题面分值安排速度。}
同一张卷上先统计题型权重，再给主观问答保留列条件与验算的时间。例如独立练习一套 100 分卷，如果用 100 分钟做题，预留 10 分钟回看图与边界条件，剩余 90 分钟按 20/10/70 分比例分为 18、9、63 分钟。遇到有公共输入的大题，先把中间向量或标签计数写完整，后续小问可以继续沿用自己的清晰步骤。这个时间例只演示比例分配，各次实际用时依本次安排。

\section{易混辨析}
\begin{tabularx}{\linewidth}{@{}p{.24\linewidth}XX@{}}
\toprule
容易混的概念 & 第一项 & 第二项\\
\midrule
历年回忆／模拟卷 & 对应年份和留存题号，缺失条件须明示 & 有完整题面、用来练特定技能，单独编号\\
刘远超／刘扬课件 & 中文应用例、流程和代表模型 & 英文推导较多、概率与数学结构较细\\
题面给定值／假设补充 & 可直接代入，优先遵循题面约定 & 为演示新造数据，不能当成原卷信息\\
定义／解释 & 写条件和数学对象 & 写哪一步起作用、何时成立、会产生什么结果\\
\bottomrule
\end{tabularx}

\section{常见错误与陷阱}
\begin{itemize}
\item ★★★ 表示明确出现或多次复现、应能手算或写推导。★★ 表示常见配套知识，至少能说明并做基本题。★ 表示先掌握术语和适用场景。在没有正式考试细则时，这些星级表示复习顺序。
\item 同一年有不同课程卷或回忆版。正文的「2021 秋」与「2021 深度」分别引用当地两份不同记录。「2024 大题」按回忆 DOCX 编号，PDF 片段顺序不同。页码用 PDF 从第 1 页开始的物理页。
\item 单选题先检查问的是正确还是错误。判断题先检查「必然、只、所有」的适用条件。主观题遇到方差定义、度量、对数底数、是否带偏置时直接写定采用的约定。
\item 若真题题干缺数据，不借参考答案的演示数当原题数据。原样保留可确认的短问，按知识点组织答法。
\end{itemize}

\section{往年真题与参考答案}
\subsection{题型、分值与题风}
\begin{tabularx}{\linewidth}{@{}p{.23\linewidth}p{.24\linewidth}X@{}}
\toprule
材料 & 已核实结构 & 训练重点\\
\midrule
2019 秋、2021 两个版本、2022 春秋、2023 & 多为回忆或题目清单。有的记录小问分值、缺全卷 & 定义、原理、推导、与实验联系。2022 春涉及模型结构与工程判断\\
2020 春试题 A & 完整 4 页、100 分。含规定选答 & 地形识别、样本与算力限制、影像去噪与生成超分辨率方案。复习本书范围内的分类、验证、CNN 部分\\
2024 秋回忆版 & 60 分。单选 5、判断 10、计算分析 9 & MSE、SVM 软硬间隔、KL、从右向左的 RNN 前向、AE、GMM。常把结构定义和算式混在同题\\
2025 秋补缓考 & 100 分。5 道选择 10 分、10 道判断 10 分、8 道简答 80 分 & KNN、聚类、逻辑回归、LSTM、RNN、Transformer。简答第 5 题题干缺失\\
2026 课件试题范本 & 可见选择和判断样例。题类说明给判断 10 分、计算问答 40 分 & 可用计算器一栏，典型问 KL（底 10）、RNN 分类、Iris 高斯 NB。按样例标注，不从样例题数反推正式总卷\\
\bottomrule
\end{tabularx}

【2019秋至2025秋本地往年题及回忆原件】【刘远超·模拟试题范本 p1–2】统计以本地九套材料为基础，同一套的多个版本只记一套。复现最明显的一组是回归/MSE、SVM、RNN、聚类/GMM。随后把 CNN/BP、HMM、VAE、KNN 和贝叶斯基础补齐。

\subsection{各章复习优先级}
\begin{tabularx}{\linewidth}{@{}p{.13\linewidth}p{.30\linewidth}X@{}}
\toprule
顺序 & 章节 & 第一遍结束能做到\\
\midrule
优先一 & 2、4、6、7、9、10、12、13、15、16、17 & 分别用正确底数、损失表、对数项、回归分母、间隔约束、责任度、状态路径、链式梯度、尺寸和潜变量串起一题\\
优先二 & 1、3、5、8、11、14、18 & 用混淆矩阵和验证划分解释指标，算树切分、密度与投影，说明感知机更新及优化器差别，再组织实验答案\\
最后检查 & 0、19 与各章标 ★ 的延伸 & 自己闭卷写公式，按易混表纠错，再复核符号方向和实验细节\\
\bottomrule
\end{tabularx}

每个实质章节的真题节给出按得分点组织的参考回答。比如 2024 大题 4–5 应先写分离超平面、距离、松弛约束，再解释核如何替换内积。2025 简答 6 应先从输入、状态、方向把递推画出来，最后说明输出。2022 秋第 4 题的实验答法要写数据、现象、原因、改法和对照验证。以上给的是核心得分点，留存材料没有逐点官方评分时不加造分值。

\section{本章自测}
\begin{enumerate}
\item 题目给了三个 $p(x\mid y_i)$ 与先验，且只求类别，归一化一定要计算吗？\\
\textbf{答} 零一损失时比较 $P(y_i)p(x\mid y_i)$ 即可，公共归一项不改比较。求数值后验则必须归一化。
\item 计算题指定了底数为 10，但公式表常用 $\ln$。哪一个写在数字代入式？\\
\textbf{答} 题目指定的 $\log_{10}$，相关熵项也用同一个底数。
\item 一段回忆只留到「第三步换成向右读」和一个矩阵。怎样写复盘？\\
\textbf{答} 标清缺失的初始状态、输入与所求，只给适用递推和明确缺口，完整数值另设例题。
\item 一道题给了真实类别并要求对未标特征聚类，这些类别最稳妥用在何处？\\
\textbf{答} 在普通无监督设置中用于事后比较或外部评价，训练使用不含标签的特征。
\item 9 道分析题每题 5 分，另有 5 道选择和 10 道判断各 1 分，总分多少？\\
\textbf{答} $9\cdot5+5+10=60$。这是本地 2024 秋回忆卷的可知分值结构。
\item 写公式时看到 $W$ 有两个可能的维度，最简单的核对办法是什么？\\
\textbf{答} 明确样本用行还是用列，在矩阵乘法里对照输入维与输出维。批量反传再检查梯度和对应参数形状一致。
\end{enumerate}

\chapter{机器学习概述与性能评价}
\section{本章考点清单}
\begin{itemize}[leftmargin=2em]
\item \(\bigstar \bigstar \bigstar\) 监督学习、无监督学习、半监督学习，训练集、验证集、测试集，留出法与分层 \(K\) 折，调参与最终评估的职责。【刘远超·2 机器学习概述 p17--20、p28--41】【刘扬·1机器学习（1）决策树 p16】【2022春真题 三(1)、八(1)】【2022秋真题 第1题】【2023真题 第10题】【2024秋真题 判断1、单选5】【模拟卷一 大题27(3)】【模拟卷三 单选1】
\item \(\bigstar \bigstar \bigstar\) 欠拟合、过拟合和学习曲线，数据泄露及分组切分。【刘远超·7 网络优化 p62--65】【刘扬·1机器学习（1）决策树 p17、p54、p70--80】【2022春真题 三(1)--(2)】【2022秋真题 第4题】【2024秋真题 单选1】【2025秋补缓考真题 简答8(3)】【模拟卷二 单选1、判断13、大题7(4)(10)】【模拟卷四 单选8、判断12】组间或时间重叠造成的泄露规则【补充】。
\item \(\bigstar \bigstar \bigstar\) 混淆矩阵、准确率、精确率、召回率、\(F_\beta\)、宏微平均、ROC 与 AUC。【刘远超·3 分类及回归问题 p9--18】【模拟卷一 填空4】【模拟卷三 单选2--3、判断19】【模拟卷四 单选9】同分 AUC 的半胜计数及下文的多类手算【补充】。
\item \(\bigstar \bigstar \bigstar\) 均方误差、平均绝对误差，特征变换与统计量估计的先后。【刘远超·3 分类及回归问题 p53--58】【2024秋真题 大题2】【2025秋补缓考真题 单选3、简答4(1)、8(1)】【模拟卷一 判断11、大题27(3)】【模拟卷三 大题7(5)--(6)】【模拟卷四 单选1--2】标准化、最小最大缩放的分折应用【补充】。
\item \(\bigstar \bigstar\) 平方误差的偏差--方差--噪声分解。【刘扬·1机器学习（1）决策树 p17】【模拟卷三 判断11】完整分解及独立性条件【补充】。
\item \(\bigstar\) 主动学习、伪标签、标注一致性和混淆矩阵的行列约定。【刘远超·2 机器学习概述 p20--27、p40--46】【模拟卷一 填空10】
\end{itemize}

\section{核心概念与定义}
\begin{description}[leftmargin=2em,style=nextline]
\item[学习范式] 监督学习（supervised learning）从 \((x_i,y_i)\) 拟合映射 \(x\mapsto y\)，分类的输出是类别，回归的输出是数值。无监督学习（unsupervised learning）从无外部类别标签的 \(x_i\) 寻找簇、分布或低维结构。半监督学习（semi-supervised learning）让有标签集和无标签集共同影响训练。自训练（self-training）先用有标签集拟合，再选高置信无标签样本作伪标签，须防止错误伪标签被迭代强化。主动学习（active learning）从无标签池选最值得请人标注的样本，如最大熵或投票分歧大的样本。自监督学习（self-supervised learning）从样本自身构造预测目标，弱监督学习（weak supervision）使用不精确、不完整或不可靠标注，小样本学习（few-shot learning）强调目标任务的可用标注少。这几个条件可以同时成立。
\item[样本与标注] 样本 \(x_i\) 是特征向量，\(y_i\) 为监督信号。真实状态（ground truth）是想预测的对象，金标准（gold standard）是实践中可信的参考标注。两位标注者的 Cohen \(\kappa\) 及多位标注者的 Fleiss \(\kappa\) 通过减去机会一致率评价一致性。\(\kappa=(p_o-p_e)/(1-p_e)\) 适用 \(p_e<1\) 的情形，\(p_o,p_e\) 分别为观测一致率与由类别边际估算的机会一致率。【刘远超·2 机器学习概述 p20--27】
\item[评估数据] 训练集拟合参数和一切依赖数据的变换。验证集比较超参数、选阈值、选模型和停止轮次。完成选择后，在独立测试集报告泛化性能（generalization）。复用测试分数做选择相当于泄露（data leakage）。同患者、同原图的变体宜按组划分。预测未来时宜按时间切分。各组类不均衡且样本可随机分时可用分层留出或分层交叉验证。
\item[泛化与诊断] 训练误差和验证误差都高，提示欠拟合（underfitting）、优化不足或较大数据噪声，需参考可达基线。训练误差明显较低而验证误差较高，提示过拟合（overfitting），也要检查分布差异和评估频数。学习曲线（learning curve）通常把训练样本量置于横轴，画同一损失口径下的训练和验证表现。迭代轮数横轴的图用于看优化与早停，诊断时先检查横轴。
\item[分类指标] 先指定正类。以真实类别为行、预测类别为列时，正类行中的 \((\mathrm{TP},\mathrm{FN})\) 给出检出与漏检，负类行中的 \((\mathrm{FP},\mathrm{TN})\) 给出误报与正确排除。精确率（precision）问预测为正的人中谁真为正，召回率（recall）问真正为正的人中多少被检出。ROC（receiver operating characteristic）在不同阈值下画假阳性率与真阳性率，AUC（area under the curve）是面积。宏平均（macro-average）先按类计算再等权平均，微平均（micro-average）先加总各类分子、分母再计算。
\item[回归指标] 均方误差（mean squared error, MSE）对误差平方取均值，单位是目标量单位的平方。均方根误差（RMSE）开平方后与目标同单位。平均绝对误差（mean absolute error, MAE）对误差绝对值取均值，同样与目标同单位。平方损失更重罚大误差。拟合优度 \(R^2=1-\mathrm{SSE}/\sum_i(y_i-\bar y)^2\)，分母为正时定义，独立数据上的 \(R^2\) 可以小于零。【刘远超·3 分类及回归问题 p53--58】
\end{description}

\section{公式与完整推导}
\subsection*{划分、交叉验证与缩放}
先预留测试集。其余样本分为 \(K\) 个互斥折 \(D_k\)。第 \(k\) 次只在 \(\bigcup_{j\ne k}D_j\) 内拟合预处理及模型，在 \(D_k\) 计算损失 \(L_k\)。对每组超参数重复，等大小折用 \(K^{-1}\sum_kL_k\) 选最小者。若折大小不同且 \(L_k\) 为样本均值，按样本汇总的总体误差是 \(\sum_k|D_k|L_k/\sum_k|D_k|\)。若评价 \(F_1\) 等非线性指标，应先确认题设要求折指标均值还是合并预测后计算。选参后用非测试数据重拟合整条流水线。
\[
 \mu=\frac1n\sum_i x_i,\quad s^2=\frac1n\sum_i(x_i-\mu)^2,\quad
 z(x)=\frac{x-\mu}{s},\qquad m(x)=\frac{x-x_{\min}}{x_{\max}-x_{\min}}.
\]
上述均值、标准差（约定 \(1/n\)）、极值均从当前训练部分估计。同一变换作用到验证和测试点，故新点可得 \(m(x)>1\) 或 \(m(x)<0\)。零方差或极差为零的特征应说明删除、置零或使用约定的稳定项。每列分别处理。

\subsection*{平方损失的偏差--方差分解}
固定 \(x\)，令 \(f(x)=\mathbb{E}[Y\mid x]\)、\(Y=f(x)+\epsilon\)、\(\mathbb{E}[\epsilon\mid x]=0\)，固定 \(x\) 时测试噪声与训练集 \(D\) 独立。记 \(m(x)=\mathbb{E}_D\hat f_D(x)\)。将误差展开为
\[
 Y-\hat f_D(x)=\epsilon+(f(x)-m(x))+(m(x)-\hat f_D(x)).
\]
平方展开后，有 \(\mathbb{E}\epsilon=0\)、\(\mathbb{E}_D[m-\hat f_D]=0\) 以及 \(\mathbb{E}[\epsilon(m-\hat f_D)]=\mathbb{E}\epsilon\cdot\mathbb{E}_D(m-\hat f_D)=0\)。因此
\[
 \mathbb{E}[(Y-\hat f_D)^2\mid x]
 = (m-f)^2+\mathbb{E}_D[(\hat f_D-m)^2]+\operatorname{Var}(Y\mid x).
\]
提高模型灵活度常可降低偏差并增加估计方差，但实际取舍要看数据、正则与拟合过程。

\subsection*{从混淆矩阵到曲线}
\[
 \begin{aligned}
 \mathrm{Acc}&=(\mathrm{TP}+\mathrm{TN})/N,&
 P&=\frac{\mathrm{TP}}{\mathrm{TP}+\mathrm{FP}},&
 R&=\frac{\mathrm{TP}}{\mathrm{TP}+\mathrm{FN}},\\
 \mathrm{TPR}&=R,&
 \mathrm{FPR}&=\frac{\mathrm{FP}}{\mathrm{FP}+\mathrm{TN}},&
 F_\beta&=\frac{(1+\beta^2)PR}{\beta^2P+R}.
 \end{aligned}
\]
把 \(P,R\) 代入 \(F_1=2PR/(P+R)\)，同乘 \((\mathrm{TP}+\mathrm{FP})(\mathrm{TP}+\mathrm{FN})\) 得 \(F_1=2\mathrm{TP}/(2\mathrm{TP}+\mathrm{FP}+\mathrm{FN})\)。分母为零要按题设约定，不能自行把无定义都记为 \(1\)。多类每次把一类当正，其余当负。每个样本恰有一个预测类时，各类合计的 \(\mathrm{TP}\) 等于正确样本数，总 \(\mathrm{FN}\) 和总 \(\mathrm{FP}\) 均等于错分数，所以微平均 \(P,R,F_1\) 等于准确率。
 
按分数从大到小遍历阈值，同分样本一起判正。起止点为 \((0,0)\)、\((1,1)\)，每遇正样本纵增 \(1/n_+\)，每遇负样本横增 \(1/n_-\)。梯形积分为
\[
 \mathrm{AUC}=\sum_k (x_{k+1}-x_k)\frac{y_k+y_{k+1}}2
 =\frac{\#(s_+>s_-)+\tfrac12\#(s_+=s_-)}{n_+n_-}.
\]
第二个等号来自每次横向跨过一个负样本时，当前高度等于排在其前面的正样本占比。同分组用起止高度平均，故同分配对计半次。两类都应有非零样本。

\subsection*{误差聚合}
令 \(e_i=\hat y_i-y_i\)，则 \(\mathrm{SSE}=\sum e_i^2\)、\(\mathrm{MSE}=\mathrm{SSE}/n\)、\(\mathrm{MAE}=\sum|e_i|/n\)。若变量只允许常数 \(a\)，平方误差和的导数为 \(2(na-\sum y_i)\)，最优为均值。绝对误差和在非样本点的斜率为左侧样本数减右侧样本数，最优在中位数处，偶数样本时两个中间次序统计量之间都取到最小。

\subsection*{课件中的 PR 曲线、标注一致性与搜索}
精确率—召回率曲线（Precision–Recall curve, PR curve）以召回率为横轴、精确率为纵轴。按预测分数降序扫描阈值，同分同组加入，更新 $TP,FP,FN$ 后依次得到点。随机排序的基线精确率与正例比例相应。PR 曲线面积依插值约定，平均精确率（Average Precision, AP）常用各组召回增量乘该组精确率求和。写 AUC 时说明是 ROC 还是 PR。多类中一对其余（One-vs-Rest, OvR）也可用于训练，即每次把一类作为正类训练一个二分类器，推理比较各类分数。分类报告的 support 是该真实类的样本数。【刘远超·3 分类及回归问题 p4、p11–12、p18】

设两位标注者的交叉表为 $n_{ij}$，Cohen 的 $\kappa$ 中 $p_o=\sum_i n_{ii}/N$，$p_e=\sum_i n_{i\cdot}n_{\cdot i}/N^2$，$\kappa=(p_o-p_e)/(1-p_e)$。Pearson 相关只反映线性关系，两位评分者相差固定常数时可仍高度相关，绝对评分却可能不一致。多人每项各给 $m$ 次标注时，Fleiss 的同项一致率为 $P_i=\sum_j n_{ij}(n_{ij}-1)/[m(m-1)]$，对各项平均得 $p_o$。所有标注中的类别比例为 $p_j$，故 $p_e=\sum_jp_j^2$，再用同样的 $\kappa$ 公式。它考虑偶然一致，$p_e=1$ 时该式分母为零。【刘远超·2 机器学习概述 p21–27】

主动学习（active learning）的查询策略是在标注预算内选最有价值的未标样本。最低置信度用 $1-\max_jp_j$，边距法挑前两项概率差最小者，熵采样挑 $-\sum_jp_j\ln p_j$ 最大者。委员会查询（Query by Committee）比较多个模型的票数熵或比较成员分布到平均预测的平均 KL。意见分歧大的点优先。训练中的超参数通过网格列出候选组合，在相同的训练折划分上比较验证平均得分，选定后重训，独立测试集最后使用。直推式（transductive）设置可利用给定目标池的无标签特征，评估泛化到新样本时须明确该权限，使用归纳式（inductive）流程。【刘远超·2 机器学习概述 p33–37、p44–46】

\section{计算题型与解题模板}

\paragraph{补充题型：PR 曲线和 Fleiss 一致性。}
原创例：正例得分 $.9,.55$，负例得分 $.8,.58,.2$。从高到低扫阈值，召回上升的两步分别为正例 $.9$ 和 $.55$，对应精确率 $1$、$2/4=1/2$，两次召回增量都为 $1/2$，所以 AP 为 $3/4$。随机基线精确率为 $2/5$。
另有三个项目，每项恰有三人给二分类标签，类别计数是 $(3,0),(1,2),(2,1)$。各项目同项一致率为 $1,1/3,1/3$，故 $p_o=5/9$。两类总频率为 $6/9,3/9$，$p_e=5/9$，所以 $\kappa=0$，观察到的集体一致与按边际频率预计的一致相同。
\paragraph{例 1 选参和缩放。}
\textbf{步骤} 先固定测试集，逐折拟合变换，记录同一指标，再全量重拟合。四个等大小折的准确率，方案 A 为 \(0.82,0.76,0.80,0.74\)，方案 B 为 \(0.78,0.83,0.79,0.80\)。某折训练特征四值为 \(1,5,9,13\)，求验证点 \(15\) 的两种缩放。
\textbf{解} A、B 四折均值分别为 \(0.78,0.80\)，选 B。训练均值 \(7\)，总体方差 \((36+4+4+36)/4=20\)，故 \(z=(15-7)/\sqrt{20}=4/\sqrt5\)，最小最大变换为 \((15-1)/(13-1)=7/6\)。不截断且不把 \(15\) 放进统计量。

\paragraph{例 2 偏差方差与曲线。}
\textbf{步骤} 定位测试噪声、估计均值与估计方差，再比较不同规模训练曲线。固定 \(x\) 时 \(f(x)=3,\ \operatorname{Var}(Y\mid x)=1/4\)。训练出的预测为 \(2,4,5\) 的概率依次为 \(1/4,1/2,1/4\)。
\textbf{解} 预测均值 \(15/4\)，平方期望 \(61/4\)，故偏差平方 \((3/4)^2=9/16\)，估计方差 \(61/4-(15/4)^2=19/16\)，测试误差为 \(9/16+19/16+1/4=2\)。若训练误差很小且验证误差高，可试增加代表性样本或正则。若两者高且接近，先排除训练未收敛，再试增加表达能力或减弱过强的正则。

\paragraph{例 3 二类与多类混淆矩阵。}
\textbf{步骤} 标明行列，二类先填四格，多类逐类使用一对其余，再明确平均方式。二类四格 \((\mathrm{TP},\mathrm{FP},\mathrm{FN},\mathrm{TN})=(18,6,2,14)\)。
\textbf{解} \(\mathrm{Acc}=32/40=0.8,\ P=18/24=0.75,\ R=18/20=0.9,\ F_1=36/44=9/11,\ \mathrm{FPR}=6/20=0.3\)。另有真实类别作行的三类矩阵
\[
 \begin{pmatrix}9&1&0\\2&3&0\\1&1&3\end{pmatrix}.
\]
三行合计 \(10,5,5\)，列合计 \(12,5,3\)。各类精确率 \(3/4,3/5,1\)，各类召回率 \(9/10,3/5,3/5\)。宏 \(P=47/60\)、宏 \(R=7/10\)。微平均三个指标都为 \(15/20=3/4\)，宏 \(F_1=(9/11+3/5+3/4)/3=159/220\)。

\paragraph{例 4 ROC 扫阈值与同分。}
\textbf{步骤} 合并排序，记录同分组，补端点后作梯形面积。也可逐对计胜。三个正样本分数依次为 \(0.87\)、\(0.64\)、\(0.41\)，三个负样本分数依次为 \(0.72\)、\(0.64\)、\(0.23\)。
\textbf{解} 点依次为 \((0,0)\)、\((0,1/3)\)、\((1/3,1/3)\)、\((2/3,2/3)\)、\((2/3,1)\)、\((1,1)\)。只在有横向跨度的三段积累面积 \(1/9+1/6+1/3=11/18\)。三个正例胜数依次为 \(3\)、\(1.5\)、\(1\)，合计 \(5.5/9\)，与面积一致。

\paragraph{例 5 误差聚合与异常值。}
\textbf{步骤} 先列有符号残差，要求 MSE 时每项平方，MAE 时逐项取绝对值，都除当前样本数。残差为 \((-1,2,-4)\)。
\textbf{解} \(\mathrm{SSE}=1+4+16=21\)，\(\mathrm{MSE}=7\)，\(\mathrm{MAE}=(1+2+4)/3=7/3\)。若最后一个残差变为 \(-8\)，MSE 变为 \(23\)，MAE 变为 \(11/3\)。不要把误差之和先平方。

\section{易混辨析}
\noindent\begin{tabularx}{\linewidth}{@{}p{.22\linewidth}XX@{}}
\toprule
概念 & 判别依据 & 考试中怎么用\\
\midrule
训练／验证／测试 & 拟合参数／选择方案／最终评估 & 预处理拟合、模型选择、报告三步分别看各自使用了哪些数据\\
欠拟合／过拟合 & 两端误差偏高／训练好而验证明显差 & 先核对数据分布、收敛性及基线，再建议调整容量、样本或正则\\
精确率／召回率 & 预测正作分母／真实正作分母 & 有漏诊代价的题强调召回，还要评估误报\\
ROC／P--R & 横轴 FPR／横轴召回率 & 类别稀少时，P--R 显示精确率随召回的变化，ROC 不直接表达预测正类中的真例比例\\
宏／微平均 & 每类等权／全样本计数汇总 & 单标签任务的微 \(F_1\) 等于准确率，宏 \(F_1\) 先逐类再平均\\
MSE／MAE & 平方后平均／取绝对值后平均 & 前者权重大误差，后者随残差绝对值线性增长\\
\bottomrule
\end{tabularx}

\section{常见错误与陷阱}
\begin{itemize}[leftmargin=2em]
\item 特征标准化、缺失值填补、特征筛选即使不看标签，也要在训练段拟合。验证或测试数据不得提前进入这些统计量。
\item 增广前按原图或患者编组选集合。同一观察对象的重复数据到不同集合会把识别个体特征误当成泛化。
\item 看到低准确率不能立刻诊断方差大，看到两条误差曲线相近也不能证明泛化已达到最佳。
\item 对调混淆矩阵的行列会对调 \(\mathrm{FP},\mathrm{FN}\)。空类分母为零须按题设处理。用总样本数作分类 \(F_1\) 的分母、用算术平均冒充调和平均，都会改变结果。
\item ROC 阈值与 AUC 是两层问题。阈值相同时，单调变换可能改变具体分类，遍历全阈值时排序与整体曲线保持。ROC 积分用横坐标变化量，遇同分时两类一起入列。
\item MSE 既不除自由度也不带默认 \(1/2\)。训练目标写成 \(\mathrm{SSE}/2\) 时，梯度少一个因子 \(2\)，报告 MSE 仍按定义计算。
\end{itemize}

\section{往年真题与参考答案}
\begin{enumerate}[leftmargin=2em]
\item \textbf{2019 秋第 3 题，过拟合。} 先给训练损失分别与独立验证表现的定义，再写树可用停止与后剪枝限制划分，SVM 可选合适 \(C\) 与核、控制复杂度。用验证数据选择超参数，不把所有模型一律用同一招。
\item \textbf{2022 春三(1)，人脸数据划分。} 按人群完全分割的方案把分布变化与泛化分数混在一起。按年龄、性别直接分组也要交代交叉群体及同一个体。建议先按身份归组，再按目标部署人群和交叉属性分层，使训练、验证、测试代表目标分布。若专门考跨群体迁移，另留可解释的域外测试并分别汇报。
\item \textbf{2022 春三(2)，学习曲线。} 回忆记录为高方差情形。得分点是读出训练好／验证差，说明落差，提出扩充代表性数据、减小容量或增强正则，说明哪项影响估计方差。原图未保存，不能报坐标值。
\item \textbf{2022 春八(1) 与 2023 第 10 题，监督范式。} 按标注成本和数量展开。语义分割的全监督通常要逐像素标签。弱监督可使用图像级标签或有限区域提示。半监督结合部分像素级标注及无标签图。自监督目标从图像本身产生。小样本学习强调目标类别每类只有少量标注。各方法的标注粒度、目标任务和数据量应分别说明。
\item \textbf{2022 秋第 1、4 题。} 第 1 题定义两类学习，监督例举逻辑回归、SVM、带标签分类树，无监督例举 K-means、PCA、GMM。第 4 题可用 lab1 拟合正弦说明，按阶数记录训练和验证曲线，解释高阶在采样点间随噪声摆动。可减阶或加岭惩罚，验证后定方案。
\item \textbf{2024 秋单选 1、5、判断 1。} 减少训练样本一般不能作为缓解高方差的常规办法。监督与否取决于训练是否用到外部标签，具体待分类模型名单未恢复。验证集可报告模型选择时的指标、早停与调参，梯度更新使用训练集，最后独立报告用测试集。
\item \textbf{2024 秋大题 2，MSE。} 原模型 \(\hat y=0.8x^2\)，两点输入 \(4,5\) 得预测 \(12.8,20.0\)，真值 \(14.8,22.0\)，故损失 \((4+4)/2=4\)。写报值本身、误差、分母三个得分点。
\item \textbf{2025 秋补缓考单选 1、3 与判断 3。} \(k\) 小的 KNN 更依赖局部样本，通常方差较大。\(k\) 增大一般更平滑，偏差可能增大。特征量纲差距大时缩放可减轻大尺度坐标支配距离或优化梯度的影响，标准化不替代验证。
\item \textbf{2025 秋补缓考简答 4(1)、8(1)(3)。} 先写 MSE，再按样本预测列残差。第 4(1) 的原数据缺失，参考答案给的 \((1,2,3)\)、\((1,2.5,2)\) 仅为示例。高阶多项式可把训练噪声一起拟合，使训练好而测试差。
\end{enumerate}

\section{本章自测}
\begin{enumerate}[leftmargin=2em]
\item 模型参数按训练集更新，超参数按同一独立测试集挑过三次，最后报告最高测试分数。流程如何修正？\textbf{答} 把选择移到验证数据，完成后另用真正未参与选择的测试数据估计泛化。
\item 训练特征为 \(3,7,11\)。新点 \(1\) 的 min-max 为多少？\textbf{答} \((1-3)/(11-3)=-1/4\)。
\item 混淆计数为 \(8,5,2,15\)，按 \((\mathrm{TP},\mathrm{FP},\mathrm{FN},\mathrm{TN})\) 排列。求召回率、FPR、\(F_1\)。\textbf{答} \(4/5,1/4,16/23\)。
\item 正例分数 \(0.3,0.9\)，负例 \(0.4,0.6\)。AUC 为多少？\textbf{答} \(0.9\) 赢两次，\(0.3\) 赢零次，故 \(1/2\)。
\item 验证曲线和训练曲线均已稳定在较大误差，差别很小。写两个后续动作。\textbf{答} 先确认基线、标注和收敛情况。验证更强的模型或较弱的正则能否同时降低两种误差。
\item 误差绝对值为 \(2,2,5,7\)。MSE 与 MAE 各为多少？\textbf{答} \(82/4=20.5\) 与 \(16/4=4\)。
\item 校验集合中的实际一致率 \(0.82\)，机会一致率 \(0.55\)，求 \(\kappa\)。\textbf{答} \((0.82-0.55)/(1-0.55)=0.60\)。
\end{enumerate}

\chapter{概率论与信息论}
\section{本章考点清单}
\begin{itemize}[leftmargin=2em]
\item \(\bigstar \bigstar \bigstar\) 条件概率、全概率、贝叶斯公式，先验、似然、证据、后验，条件独立与边缘化。【刘扬·2机器学习（2）概率论速通 p24--42】【刘远超·3 分类及回归问题 p39、p49】【2019秋真题 第2、4题】【2021深度真题 第1题】【2022春真题 一(1)】【2023真题 第1题】【模拟卷一 填空6、大题26(3)】【模拟卷二 单选3】【模拟卷四 单选4】
\item \(\bigstar \bigstar \bigstar\) 极大似然 MLE、最大后验 MAP、伯努利与一元正态的估计、Beta 与 Dirichlet 共轭，固定高斯噪声与最小二乘。【刘扬·2机器学习（2）概率论速通 p45--58】【刘远超·3 分类及回归问题 p48--50、p61--63】【2021深度真题 第2题】【2021秋真题 第1题】【模拟卷一 填空1、大题26(3)】【模拟卷三 大题7(7)】【模拟卷四 大题7(2)--(3)】
\item \(\bigstar \bigstar \bigstar\) 信息熵、交叉熵和 KL 的方向、底数与支持集，数值保留两位小数。【刘远超·7 网络优化 p4--10、p15--17、p23--25】【2024秋真题 大题6】【模拟卷一 填空2】【模拟卷三 单选4】【刘远超·模拟试题范本-2026-11-3 p2 第1计算问】KL 的非负性证明与相等条件【补充】。
\item \(\bigstar \bigstar \bigstar\) 联合熵、条件熵、互信息、决策树的信息增益及对称性。【刘远超·7 网络优化 p11--14】【刘扬·1机器学习（1）决策树 p45--49】【2019秋真题 第1题】【2021秋真题 第2题】【2022秋真题 第2题】【模拟卷一 填空3】【模拟卷二 单选6】【模拟卷四 单选10】
\item \(\bigstar \bigstar\) 随机变量、离散概率质量与连续密度，伯努利、二项、分类、多项、均匀、正态、Beta、Dirichlet 和常用矩。【刘扬·2机器学习（2）概率论速通 p13--24、p47--54、p57】【刘远超·3 分类及回归问题 p40--50】【模拟卷一 填空1】【模拟卷三 判断20】拉普拉斯密度及其尺度计算【补充】。
\end{itemize}

\section{核心概念与定义}
\begin{description}[leftmargin=2em,style=nextline]
\item[概率表示] 概率质量函数（probability mass function, PMF）对离散点给概率。概率密度函数（probability density function, PDF）需积分后给区间概率。若 \(X\) 有密度 \(p\)，则 \(P(l<X<u)=\int_l^u p(x)\,\mathrm dx\)，单点概率为零，密度值可以超过 \(1\)。随机变量的期望（expectation）是对取值乘概率求和或对 \(xp(x)\) 积分。方差是 \(\mathbb E[X^2]-(\mathbb E X)^2\)，协方差为 \(\mathbb E[(X-\mathbb EX)(Y-\mathbb EY)]\)。
\item[条件与边缘] 联合分布（joint distribution）描述共现，边缘分布（marginal distribution）对其他变量求和或积分。独立对应联合是边缘的积。条件独立（conditional independence）对应在给定第三变量后各条件分布相乘，不能由两者边缘独立直接推出。零概率条件事件不能用离散比值定义。
\item[先验和参数估计] 看到数据前对参数或类归属的概率为先验（prior）。似然（likelihood）把已观测数据固定、把参数当自变量，比较数据对各参数的支持。后验（posterior）把似然乘先验再归一。极大似然估计（maximum likelihood estimation, MLE）仅最大化似然，最大后验估计（maximum a posteriori, MAP）最大化后验密度或质量。后验均值、后验众数和后验预测有不同用途。
\item[两种老师的记法] 类别可记 \(y\) 或 \(\omega_i\)。\(\Pr(Y=y\mid X=x)\) 对应 \(P(\omega_i\mid x)\)。若 \(x\) 连续，写 \(p(x\mid y)\) 表示类条件密度。以后验比较作分类时 \(p(x)\) 是全类共享的归一因子。参数估计时通常写 \(p(D\mid\theta)\)、\(p(\theta\mid D)\)，角色与类别贝叶斯公式相同。
\item[信息量] 单次自信息（self-information）为 \(-\log_b p\)，稀有事件信息大。信息熵（entropy）为自信息的期望。交叉熵（cross-entropy）在真实分布 \(P\) 下平均按 \(Q\) 编码的成本。KL 散度（Kullback--Leibler divergence）扣除真实熵后的代价差，非对称，不能当作距离使用。互信息（mutual information）是已知另一个变量平均减少的不确定性。
\end{description}

\noindent\begin{tabularx}{\linewidth}{@{}p{.23\linewidth}X X@{}}
\toprule
分布 & 形式或产生方式 & 常用矩或参数\\
\midrule
伯努利（Bernoulli） & \(P(X=x)=p^x(1-p)^{1-x},\ x=0,1\) & 均值 \(p\)，方差 \(p(1-p)\)\\
二项（binomial） & \(n\) 个独立同参数伯努利之和，概率 \(\binom nkp^k(1-p)^{n-k}\) & 均值 \(np\)，方差 \(np(1-p)\)\\
分类／多项（categorical／multinomial） & 单次按 \(K\) 类概率 \(\theta_j\) 抽样，重复 \(n\) 次的计数质量为 \(n!\prod_j(\theta_j^{n_j}/n_j!)\) & \(n_j\) 的均值 \(n\theta_j\)，单次给出一个类别结果\\
均匀（uniform） & \(a<b\)，区间内密度 \(1/(b-a)\)，区间外为零 & 均值 \((a+b)/2\)，方差 \((b-a)^2/12\)\\
正态（normal/Gaussian） & \(p(x)=\exp(-(x-\mu)^2/(2\sigma^2))/(\sqrt{2\pi}\sigma)\) & \(\mu\) 是均值，\(\sigma^2>0\) 是方差\\
Beta／Dirichlet & Beta 在 \(0<\theta<1\) 上正比于 \(\theta^{a-1}(1-\theta)^{b-1}\)。Dirichlet 在概率单纯形上正比于 \(\prod_j\theta_j^{a_j-1}\) & 所有形状参数为正，均值分别 \(a/(a+b)\)、\(a_j/\sum_\ell a_\ell\)\\
拉普拉斯（Laplace） & \(p(x)=e^{-|x-\mu|/b}/(2b)\)，尺度 \(b>0\)【补充】 & 均值 \(\mu\)，方差 \(2b^2\)\\
\bottomrule
\end{tabularx}

\section{公式与完整推导}
\subsection*{全概率与贝叶斯公式}
设 \(A_1,\ldots,A_m\) 互斥穷尽。因 \(B\) 由互斥的 \(B\cap A_j\) 组成，
\[
 P(B)=\sum_jP(B\cap A_j)=\sum_jP(B\mid A_j)P(A_j).
\]
又 \(P(A_i\cap B)=P(B\mid A_i)P(A_i)=P(A_i\mid B)P(B)\)。对 \(P(B)>0\) 相除得
\[
 P(A_i\mid B)=\frac{P(B\mid A_i)P(A_i)}{\sum_jP(B\mid A_j)P(A_j)}.
\]
先算各类的未归一得分，然后除分数之和。连续 \(X=x\) 用相应密度作分数并对类求和，不能把密度直接当作单点概率。

\subsection*{MLE、MAP 与三个常见推导}
独立同分布样本 \(D=(x_1,\ldots,x_n)\) 下，\(L(\theta)=\prod_i p(x_i\mid\theta)\)，故 \(\ell(\theta)=\sum_i\ln p(x_i\mid\theta)\)。对正似然，取对数保留最优点。由于证据 \(p(D)\) 不依赖 \(\theta\)，
\[
 \theta_{\mathrm{MLE}}\in\operatorname*{arg\,max}_\theta\ell(\theta),
 \qquad
 \theta_{\mathrm{MAP}}\in\operatorname*{arg\,max}_\theta\{\ell(\theta)+\ln p(\theta)\}.
\]
似然乘了参数范围内恒定的先验时，两种最优点一致。非均匀先验的 \(-\ln p(\theta)\) 是最小化负对数后验时新增的惩罚，平方范数来自独立零均值正态先验，绝对值和来自独立零均值拉普拉斯先验。正则系数需随损失整体缩放。

设 \(n\) 次伯努利中有 \(s\) 次为 \(1\)。写成
\[
 \ell(p)=s\ln p+(n-s)\ln(1-p),\quad
 \ell'(p)=\frac{s-np}{p(1-p)}.
\]
若 \(0<s<n\)，导数在 \(s/n\) 前正后负，最大为 \(s/n\)。若 \(s=0\) 或 \(s=n\)，在允许的端点 \(0\) 或 \(1\) 达最大。Beta\((a,b)\) 先验乘上似然后，\(p\) 的两端指数为 \(s+a-1,n-s+b-1\)，故得到 Beta\((a+s,b+n-s)\)。当两个后验形状参数均大于 \(1\)，按同一求导法得 MAP \((a+s-1)/(a+b+n-2)\)，后验预测成功率 \((a+s)/(a+b+n)\)。边界形状需另查众数，多项情形指数按计数逐类相加。

一元高斯的对数似然，记 \(v=\sigma^2>0\)，为
\[
 \ell(\mu,v)=-\frac n2\ln(2\pi v)-\frac1{2v}\sum_i(x_i-\mu)^2 .
\]
对 \(\mu\) 求导得 \(\sum_i(x_i-\mu)/v\)，二阶导为 \(-n/v<0\)，故每个固定 \(v\) 下均取 \(\bar x\)。以 \(\bar x\) 代入后对 \(v\) 求导，
\[
 \ell_v=-\frac n{2v}+\frac{\sum_i(x_i-\bar x)^2}{2v^2}.
\]
当平方和非零，此导数先正后负，取到 \(v_{\mathrm{MLE}}=\sum_i(x_i-\bar x)^2/n\)。若观测完全相同且要求 \(v>0\)，似然在 \(v\downarrow0\) 时可无界。回归噪声固定 \(v\) 时，把 \(\mu\) 改为 \(f_\theta(x_i)\)，优化前面的正系数平方误差和等价于最小二乘。

\subsection*{熵、KL 与互信息}
\[
 \begin{aligned}
 H_b(P)&=-\sum_i p_i\log_b p_i,&
 H_b(P,Q)&=-\sum_i p_i\log_b q_i,\\
 D_{\mathrm{KL},b}(P\|Q)
 &=\sum_i p_i\log_b(p_i/q_i)
 =H_b(P,Q)-H_b(P).
 \end{aligned}
\]
所有式子使用同一底数 \(b>1\)。以 2 为底用 bit，以自然对数用 nat，以 10 为底可直接报数值和对数底。约定 \(0\log 0=0\)，\(p_i>0,q_i=0\) 时 KL 和交叉熵为 \(+\infty\)。当该情形未出现，对每个 \(p_i>0\)，\(-\ln(q_i/p_i)\ge 1-q_i/p_i\)，故
\[
 D_{\mathrm{KL},e}(P\|Q)\ge \sum_{p_i>0}(p_i-q_i)=1-\sum_{p_i>0}q_i\ge 0.
\]
两次等号要求在 \(P\) 的支持上 \(p_i=q_i\) 且剩余坐标 \(q_i=0\)，等价于 \(P=Q\)。以其他底计算只需除正数 \(\ln b\)。JS 散度取 \(M=(P+Q)/2\) 并平均两项 \(D(P\|M),D(Q\|M)\)。两项交换位置对值无影响，因 \(M_i\ge P_i/2,Q_i/2\)，以 2 为底时范围为 \([0,1]\)。【刘远超·7 网络优化 p7--10】

将 \(p(x,y)=p(y)p(x\mid y)\) 代入联合熵，拆对数后对 \(x\) 求和，分别得 \(H(Y)\) 和 \(\sum_y p(y)H(X\mid Y=y)\)。因此
\[
 \begin{aligned}
 H(X,Y)&=H(Y)+H(X\mid Y)=H(X)+H(Y\mid X),\\
 I(X;Y)&=H(X)-H(X\mid Y)=H(Y)-H(Y\mid X)\\
 &=H(X)+H(Y)-H(X,Y)\\
 &=\sum_{x,y}p(x,y)\log_b\frac{p(x,y)}{p(x)p(y)}.
 \end{aligned}
\]
最后的式子是联合与边缘积之间的 KL，所以互信息非负，等于零当且仅当 \(X,Y\) 独立。离散情形，下式亦成立 \(0\le I(X;Y)\le\min\{H(X),H(Y)\}\)。决策树的信息增益具体是属性与标签的互信息，不能把任意两个分布的 KL 直接当作属性信息增益。

\subsection*{概率公理、分布函数与正态共轭}
概率测度非负，全样本空间的概率为 1，对两两不交事件可加。因 $A,A^c$ 不相交且并为样本空间，$P(A^c)=1-P(A)$。将并集拆成 $A$ 与 $B\setminus A$，得 $P(A\cup B)=P(A)+P(B)-P(A\cap B)$。分布函数（Cumulative Distribution Function, CDF）$F(x)=P(X\le x)$ 总存在，单调、右连续。连续且有密度时，$F(x)=\int_{-\infty}^x f(t)dt$，可微处 $f=F'$。离散 CDF 的跃点高度是该点的概率质量。独立同分布、方差有限且非零时，中心极限定理（Central Limit Theorem, CLT）给 $\sqrt n(\bar X-\mu)/\sigma\Rightarrow\mathcal N(0,1)$，说的是标准化样本均值的极限分布。【刘扬·2 概率论速通 p3–11、p16–23】

若 $x_i$ 独立于 $\mathcal N(\mu,\sigma^2)$，$\sigma^2$ 已知，先验为 $\mathcal N(\mu_0,\tau^2)$，后验仍是正态。取对数后与 $\mu$ 有关的项为
\[
-\frac{1}{2\sigma^2}\sum_i(x_i-\mu)^2-\frac{(\mu-\mu_0)^2}{2\tau^2}
=-\frac12(a\mu^2-2b\mu)+C .
\]
其中 $a=n/\sigma^2+1/\tau^2,\ b=n\bar x/\sigma^2+\mu_0/\tau^2$。完全平方为 $-\tfrac a2(\mu-b/a)^2+C'$，故后验方差 $1/a$、均值和 MAP 均为 $b/a$。样本量增大时似然部分的精度增大，这解释了后验均值为何更靠近样本均值。【刘扬·2 概率论速通 p57】

\section{计算题型与解题模板}
\paragraph{例 1 先验到后验。}
\textbf{步骤} 将事件名和列概率写清，按类乘先验后做全概率归一。工件来自设备 A 的概率 \(3/20\)。A 的报警率 \(4/5\)，其余设备统一记 B，报警率 \(1/10\)。
\textbf{解} 两类未归一分数分别为 \(3/25\) 与 \(17/200\)，报警总概率 \(41/200\)，故 \(P(A\mid\text{报警})=24/41\)。若只比较报警率 \(4/5\) 与 \(1/10\)，无法给出这个后验数值。

\paragraph{例 2 伯努利、Beta 与多项更新。}
\textbf{步骤} 数成功失败，写似然，更新先验的两个形状参数，区分后验 mode 和预测均值。9 次有 7 次成功，先验 Beta\((2,4)\)。
\textbf{解} MLE 为 \(7/9\)，后验 Beta\((9,6)\)，MAP 为 \(8/13\)，下一次成功的概率为 \(9/15=3/5\)。另一三类样本计数 \((2,0,5)\)，Dirichlet 先验参数 \((1,3,2)\)，则后验参数 \((3,3,7)\)，下一次为第 3 类的概率 \(7/13\)。零计数的第二类没有消失，因为先验参数仍为正。

\paragraph{例 3 高斯估计与 MAP 平衡。}
\textbf{步骤} 先求均值、中心化平方和。同估均值方差时用 MLE 分母 \(n\)，已知方差和正态先验时对负对数后验求导。样本为 \((-2,0,4,6)\)。
\textbf{解} \(\bar x=2\)，平方和 \(16+4+4+16=40\)，方差 MLE 为 \(10\)。若改设观测方差固定为 \(4\)，并且 \(\mu\sim\mathcal N(1,1)\)，最小化 \(\sum_i(x_i-\mu)^2/8+(\mu-1)^2/2\)。导数 \((4\mu-8)/4+(\mu-1)=2\mu-3\)，二阶导为正，所以 \(\mu_{\mathrm{MAP}}=3/2\)。

\paragraph{例 4 以 2 为底的四数联算。}
\textbf{步骤} 先算两边熵，再由 \(P\) 给 \(-\log_2 Q\) 加权，最后减 \(H(P)\)。取 \(P=(1/2,3/8,1/8)\)，\(Q=(1/4,1/2,1/4)\)。
\textbf{解} \(H(P)=\tfrac12-\tfrac38\log_2(3/8)+\tfrac38\approx1.4056\)，\(H(Q)=1.5\)。交叉熵为 \(1+3/8+1/4=1.625\)，KL 为 \(1.625-1.4056\approx0.2194\)。保留两位分别为 \(1.41,1.50,1.63,0.22\)（\(1.625\) 按通常四舍五入）。

\paragraph{例 5 以 10 为底与对数换底。}
\textbf{步骤} 逐项保留比值，检查底和方向。取 \(P=(0.6,0.3,0.1)\)，\(Q=(0.3,0.3,0.4)\)。
\textbf{解} \(D_{\mathrm{KL},10}(P\|Q)=0.6\log_{10}2+0+0.1\log_{10}(1/4)=0.4\log_{10}2\approx0.1204\)。以 2 为底时结果是 \(0.4\)，用 \(\ln\) 时结果是 \(0.4\ln 2\)。直接在最后一行才舍入。

\paragraph{例 6 列联表到条件熵与互信息。}
\textbf{步骤} 求边缘与行内条件分布，按行权重平均行熵，再求差。联合计数按行 \(X=0,1\)、列 \(Y=0,1\) 排为 \(\begin{pmatrix}3&1\\2&6\end{pmatrix}\)。
\textbf{解} 两行权重 \(1/3,2/3\)，两行条件 \(P(Y=1\mid X)\) 分别为 \(1/4,3/4\)，二元熵同为 \(h_2(1/4)\approx0.8113\)。\(H(Y)=h_2(5/12)\approx0.9799\)，故 \(I(X;Y)\approx0.1686\)。按另一方向，\(H(X\mid Y)=h_2(1/3)-I(X;Y)\approx0.7497\)。两个条件熵一般不同。

\section{易混辨析}
\noindent\begin{tabularx}{\linewidth}{@{}p{.23\linewidth}XX@{}}
\toprule
概念 & 区分要点 & 容易失分的例子\\
\midrule
先验／似然／后验 & 未见数据／按参数比较数据适配性／见数据后更新 & \(P(D\mid\theta)\) 与 \(P(\theta\mid D)\) 的条件反过来了\\
MLE／MAP／预测 & 最大化似然／最大化似然乘先验／对参数后验求平均 & Beta 的预测均值不必等于 Beta 的 mode\\
质量／密度 & 单点概率／积分给概率 & 密度 \(1.5\) 不表示事件概率 \(1.5\)\\
交叉熵／KL & 编码代价／多出的代价 & \(P=Q\) 时 KL 为零，交叉熵一般仍大于零\\
条件熵／联合熵 & 知道一个后还剩多少／共同不确定性 & 链式法则用条件项相加，不能假设两个变量独立\\
KL／互信息 & 用来比较指定的两个分布／联合与边缘积之间的 KL & KL 一般非对称，互信息关于两个变量对称\\
\bottomrule
\end{tabularx}

\section{常见错误与陷阱}
\begin{itemize}[leftmargin=2em]
\item 看见条件概率先确认条件发生的对象，反向概率需要先验与证据。条件独立仅能在其指定的条件下拆乘积。
\item 观察数据固定时似然随参数变化，不能按参数当成已归一概率。先验为零的区域在 MAP 中不可取。
\item 极值候选须查边界或二阶符号。对高斯方差不能从样本方差直接抄 \(n-1\)。后验形状参数不大于 \(1\) 时不可照抄内部众数公式。
\item 用 \(P\) 为权重时减 \(H(P)\)。写清 \(P\|Q\)，底数和小数位数，不在中间结果逐项舍入。
\item \(H(X\mid Y)\) 先按 \(Y=y\) 分组，每个条件分布内概率要归一，再按 \(p(y)\) 加权。联合熵与交叉熵虽然都常写两个参数，前者输入的是一个联合分布，后者比较同一字母表上的两个分布。
\end{itemize}

\section{往年真题与参考答案}
\begin{enumerate}[leftmargin=2em]
\item \textbf{2019 秋第 1、2、4 题。} 第 1 题写 \(I(Y;A)=H(Y)-H(Y\mid A)\)，将属性子集频率作条件熵权重并说明切分后不确定性减少。第 2 题先写 \(P(Y=i\mid x)\propto p(x\mid i)P(Y=i)\)，再以条件错误率 \(1-P(Y=i\mid x)\) 得最大后验类别。第 4 题的两类高斯共协方差，在对数后验比中比较两类均值项与先验项，\(x^\top\Sigma^{-1}x\) 相减为零，边界为仿射方程。
\item \textbf{2021 秋第 1、2 题。} 正则化对应负对数先验，与数据的负对数似然相加得到负对数后验。条件熵回答“已知 \(A\) 后标签还有多不确定”，写分组熵及加权形式。决策树中用与父节点熵的差衡量切分信息量。
\item \textbf{2021 深度第 1、2 题。} 癌症判别例的可核得分点是先验、条件概率的乘积及两个后验比较，资料缺参数。MLE 例仅回忆为课件同类题，通用步骤为写联合似然、取对数、求驻点、检查边界与最大性，不能补造分布或样本。
\item \textbf{2022 春一(1)、2023 第 1 题。} 在 0--1 损失下比较类别后验。能从类别与观测的联合概率恢复后验时，应先统一分母。2022 春原条件数值和 2023 具体所求均缺失。
\item \textbf{2022 秋第 2 题。} 给出熵、条件熵和互信息定义，然后针对 \(X_1,X_2\) 分别计算按子节点比例加权的条件熵并比较信息增益。分类结果图或计数未出现在现存 DOCX，不能宣布某属性胜出。
\item \textbf{2024 秋大题 6。} 原 \(P=(0.1,0.3,0.6),Q=(0.1,0.2,0.7)\)。以 2 为底，可得 \(H(P)\approx1.30,H(Q)\approx1.16,D_{\mathrm{KL}}(P\|Q)\approx0.04\)。得分点依次为两组各自加权，对应坐标写比值，用 \(P\) 加权，并按要求保留两位。相对熵的方向由同套回忆 PDF 补明。
\item \textbf{本次课程范本计算第 1 题。} 题设非本地往年回忆卷。原 \(P=(0.8,0.1,0.1),Q=(0.1,0.1,0.8)\)，要求以 10 为底。\(D(P\|Q)=0.8\log_{10}8+0.1\log_{10}(1/8)=0.7\log_{10}8\approx0.63\)。中间同概率坐标贡献为零。
\end{enumerate}

\section{本章自测}
\begin{enumerate}[leftmargin=2em]
\item \(P(A)=1/5,\ P(E\mid A)=3/4,\ P(E\mid A^\mathrm c)=1/8\)，求 \(P(A\mid E)\)。\textbf{答} 分子 \(3/20\)，分母 \(3/20+1/10=1/4\)，后验 \(3/5\)。
\item 12 次二值试验得到 3 次成功，先验 Beta\((3,2)\)。\textbf{答} MLE \(1/4\)，后验 Beta\((6,11)\)，MAP \(1/3\)，后验预测 \(6/17\)。
\item \(X\) 的取值为 \(1,2,6\)，在一元高斯模型中求均值和方差 MLE。 \textbf{答} \(3,14/3\)。
\item \(P=(3/4,1/4),\,Q=(1/2,1/2)\)。以 2 为底求交叉熵和 KL，写出精确表达式。\textbf{答} \(H(P,Q)=1\)，\(D(P\|Q)=1-h_2(1/4)\)。
\item \(X\) 等概率取 \(0,1\)，且 \(Y=1-X\)。以 2 为底求两个条件熵和互信息。 \textbf{答} 条件熵都为 \(0\)，\(I=1\)。
\item 若密度 \(p(x)=2\) 对 \(x\in[0,1/2]\) 成立，区间外为零，求 \(P(X<1/8)\)。 \textbf{答} \(\int_0^{1/8}2\,\mathrm dx=1/4\)，密度虽大于 1，分布有效。
\item 若 \(P=(0.4,0.6)\) 且 \(Q=(1,0)\)，两个方向的 KL 哪个为无穷？ \textbf{答} \(D(P\|Q)=+\infty\)，反方向为 \(\log_b(1/0.4)\)。
\end{enumerate}

\chapter{决策树}
\label{ch:prml-tree}

\section{本章考点清单}
\begin{itemize}
\item \ensuremath{\bigstar\bigstar\bigstar} 信息熵、条件熵、信息增益（information gain）及节点切分。要按子节点样本占比加权，并能解释为什么最大化增益有意义。【刘扬·1机器学习（1）决策树 p44–50】【2019秋真题 第1题】【2021秋机器学习真题 第2题（条件熵应用）】【2022秋真题 第2题】【模拟卷一 填空3】【模拟卷二 选择6】
\item \ensuremath{\bigstar\bigstar\bigstar} 单变量切分、连续属性阈值、决策边界。会按路径写出叶节点区域，并区分单步边界与整棵树的边界。【刘扬·1机器学习（1）决策树 p23–35、38–42、57】【2019秋真题 第1题】【模拟卷四 判断11】
\item \ensuremath{\bigstar\bigstar\bigstar} 过拟合（overfitting）、停止规则、预剪枝与后剪枝。【刘扬·1机器学习（1）决策树 p51、54、60–64】【2019秋真题 第3题】【模拟卷四 选择3】
\item \ensuremath{\bigstar\bigstar} ID3、C4.5、分类与回归树（classification and regression tree, CART）的切分准则。会算信息增益率（gain ratio）和加权基尼指数（Gini index）。课件第 56 页提及 ID3、C4.5。增益率、CART 剪枝的公式详见教材笔记。【刘扬·1机器学习（1）决策树 p44–56】【补充·李航教材笔记“五、决策树”】【模拟卷三 选择5】【模拟卷四 选择3】
\item \ensuremath{\bigstar} 缺失值的加权传播、最小描述长度（minimum description length, MDL）。【刘扬·1机器学习（1）决策树 p62、65–68】
\end{itemize}

\section{核心概念与定义}
决策树（decision tree）是一组递归的特征测试规则。一个内部节点实施一个测试，每条分支对应测试结果，叶节点输出预测。分类树可以用叶内类频率估计后验，再按多数类预测。回归树在平方损失下通常用叶内目标均值预测。预测时从根沿满足条件的分支下降到叶。【刘扬·1机器学习（1）决策树 p23–35】【补充·李航教材笔记“五、决策树”】

训练从当前节点的样本集出发，按准则选切分，把子集交给各子节点递归处理。节点已纯、无有效切分或达到了停止条件时成为叶节点。该过程每轮只优化眼前的切分，属于贪心（greedy）搜索。多个结构可能同样拟合训练集。课件中的 $P,S$ 指父样本集或当前子集，下面统一用 $D$ 表示当前节点数据。【刘扬·1机器学习（1）决策树 p32–43、51–56、59】

名义属性可以一值一分支，或将若干取值组成两个集合。顺序属性可按顺序切分。对连续特征 $x_j$，常用 $x_j\le t$ 与 $x_j>t$ 两路切分。单变量的切分超平面 $x_j=t$ 与其余坐标轴平行。一个叶子的区域是沿路径各条件的交集。多个叶子的预测区域可构成阶梯状或分离的区域。【刘扬·1机器学习（1）决策树 p38–42、57】

\section{公式与完整推导}
\subsection*{从条件熵推到信息增益}
设当前节点有 $n$ 个样本，类别变量为 $Y$，第 $c$ 类样本数为 $n_c$。令 $q_c=n_c/n$。默认以 2 为底，约定 $0\log_2 0=0$。一次切分的分支变量为 $A$。第 $v$ 个分支样本数为 $n_v$，该分支内第 $c$ 类样本数为 $n_{vc}$。于是
\[
 H(Y)=-\sum_c q_c\log_2q_c,\qquad
 H(Y\mid A=v)=-\sum_c\frac{n_{vc}}{n_v}\log_2\frac{n_{vc}}{n_v}.
\]
按分支概率取均值，得到
\[
 H(Y\mid A)
   =\sum_v\frac{n_v}{n}H(Y\mid A=v).
\]
利用 $p(c,v)=p(v)p(c\mid v)$，把 $\log_2 p(c,v)$ 拆成两项并对 $c$ 求和，得到 $H(Y,A)=H(A)+H(Y\mid A)$。故最大化分裂前后不确定性的降幅，等价于最大化类别与分支的经验互信息：
\begin{align*}
 \operatorname{Gain}(D,A)
 &=H(Y)-H(Y\mid A)\\
 &=H(Y)+H(A)-H(Y,A)=I(Y;A).
\end{align*}
若切分后每个叶都纯，则条件熵为 0。若每个子节点的类比例与父节点相同，则条件熵等于父节点熵，切分未提供类别信息。【刘扬·1机器学习（1）决策树 p44–50】

\subsection*{增益率与基尼指数}
原始信息增益可能偏好多值属性。例如样本编号可把每条训练样本单独分出去，训练子节点都纯。C4.5 用分支熵
\[
 H(A)=-\sum_v\frac{n_v}{n}\log_2\frac{n_v}{n},
 \qquad
 \operatorname{GainRatio}(D,A)=\frac{\operatorname{Gain}(D,A)}{H(A)}
\]
衡量每单位分支信息对应的类别信息。$H(A)=0$ 表示没有实际切分，跳过该候选。若题目明确采用含平均信息增益筛选的 C4.5 流程，先过滤信息增益过低的候选，再在剩余项中比较增益率。只问增益率的计算时，按给定候选逐项计算即可。【补充·李航教材笔记“五、决策树”。C4.5 候选筛选按指定口径】

节点中按照类频率随机预测时，真实类与随机预测独立。两者都等于 $c$ 的概率是 $q_c^2$，故基尼不纯度
\[
 G(D)=1-\sum_c q_c^2,\qquad
 G_{\mathrm{split}}=\sum_v\frac{n_v}{n}G(D_v).
\]
分类 CART 按二叉切分比较 $G_{\mathrm{split}}$，取最小者。父节点固定，等价于让 $G(D)-G_{\mathrm{split}}$ 最大。二分类下 $G(D)=2q(1-q)$。回归 CART 可比较两子节点的总平方残差。对一个叶节点，$\sum_i(y_i-u)^2$ 的导数是 $2\sum_i(u-y_i)$，为零时 $u=\bar y$。二阶导数 $2n>0$，因此叶均值确实最小化平方残差。【补充·李航教材笔记“五、决策树”】

\subsection*{预剪枝、后剪枝与代价复杂度}
预剪枝（pre-pruning）在生长时限制最小叶样本数、最大深度或最小增益。后剪枝（post-pruning）先生成较大树，再比较保留子树与将其替换为叶子的表现。验证集后剪枝可固定叶预测并比较独立验证集的误分类数。【刘扬·1机器学习（1）决策树 p60–64】

【补充·李航教材笔记“五、决策树”】在同一个节点 $t$，令 $R(T_t)$ 为子树的训练风险，$R(t)$ 为将其替成单叶的风险，叶数为 $L_t>1$。定义 $R_\alpha(T)=R(T)+\alpha |L(T)|$。比较两种局部代价，
\[
 R(t)+\alpha\le R(T_t)+\alpha L_t
 \quad\Longleftrightarrow\quad
 \alpha\ge\frac{R(t)-R(T_t)}{L_t-1}.
\]
分子是剪枝增加的拟合误差，分母是减少的叶数。CART 可以按最小的临界值逐步得到一列嵌套子树，再用独立验证评估选树。误分类数与误分类率两种风险定义会使 $\alpha$ 的数值单位不同。

\subsection*{树的搜索空间、MDL 与缺失值}
有 $n$ 个布尔属性时，$2^n$ 种输入组合各可独立指定两个标签之一，因而可表示 $2^{2^n}$ 个不同的布尔分类函数。等价函数还可由不同树表达，穷举树的代价通常很高。决策树采用局部最优分裂，配合剪枝控制增长。最小描述长度（Minimum Description Length, MDL）把方案看成 $L(T)+L(D\mid T)$ 两段编码，前者度量模型复杂度，后者度量在树已知时还需说明的数据和错误信息，选择和较短的树。对缺失属性的一个课程处理是估计切分收益时用已知样本及其比例，训练时按已知分支占比将未知值样本的权重分到多个分支，预测时再按分支权重聚合类别后验。【刘扬·1 决策树 p59、p61–62、p65–68】

\section{计算题型与解题模板}
\subsection*{题型一：用列联计数选择切分}
\textbf{模板} 写出父节点类比例。对每个候选逐分支计算 $n_v/n$ 和局部纯度。信息增益用父熵减加权熵，增益率再除分支熵，CART 按合法的二分比较加权基尼。

\textbf{原创例题} 父节点有 6 正、6 负。属性 $A$ 的两个分支为 $(5,1),(1,5)$。属性 $B$ 的三个取值对应 $(3,0),(0,2),(3,4)$。比较信息增益、增益率，并按二叉规则比较最佳基尼切分。这里计算增益率时先不筛选候选。

\textbf{解答} 父熵为 1，令 $h(p)=-p\log_2p-(1-p)\log_2(1-p)$。属性 $A$ 两边等大，$H(Y\mid A)=h(5/6)\approx0.6500$，增益约为 $0.3500$。分支熵为 1，增益率也约为 $0.3500$。属性 $B$ 的前两支纯，$H(Y\mid B)=\frac7{12}h(3/7)\approx0.5747$，增益约为 $0.4253$。分支熵由 $(3/12,2/12,7/12)$ 给出，约为 $1.3844$，增益率约为 $0.3072$。所以增益选 $B$，直接比较增益率选 $A$。若按“增益至少达到候选平均值”筛选，只有 $B$ 保留，再选 $B$。

对于 CART，$A$ 的加权基尼为 $2(5/6)(1/6)=5/18$。将 $B$ 的第一个、第二个、第三个取值分别单独分出，三个合法二分的加权基尼依次为 $1/3,\,2/5,\,17/35$，故 CART 选择 $A$。不能把三个分支同时当成 CART 的一次二叉切分。

\subsection*{题型二：枚举连续属性的阈值}
\textbf{模板} 排序不同的取值。在相邻不同值之间取代表阈值，并算每种阈值产生的两路类计数。同一取值的样本必须留在同一侧。按题目指定准则和并列规则选择。

\textbf{原创例题} 五条记录为 $(1,-),(2,-),(4,+),(6,+),(8,-)$，采用 $x\le t$ 分支，按信息增益选阈值。父计数 $(2,3)$。
\[
\begin{array}{c|c|c|c}
t&\text{左侧（正，负）}&\text{右侧（正，负）}&\text{增益（bit）}\\\hline
1.5&(0,1)&(2,2)&0.1710\\
3&(0,2)&(2,1)&0.4200\\
5&(1,2)&(1,1)&0.01997\\
7&(2,2)&(0,1)&0.1710
\end{array}
\]
例如 $t=3$ 时，左侧纯，右侧权重 $3/5$，增益是 $h(2/5)-\frac35 h(2/3)$。四者中应取 $t=3$。任意严格位于 2 和 4 之间的阈值都会给这些训练样本相同分配。

\subsection*{题型三：验证剪枝和临界代价}
\textbf{模板} 固定节点被压成叶时的预测。验证剪枝先比较两种树对同一验证数据的误判数。代价复杂度剪枝则列出两种训练风险加叶数惩罚并解不等式。

\textbf{原创例题} 一棵三叶子树在训练集上错 2 个，合成单叶会错 6 个。独立验证集中，保留子树错 3 个，合并后错 2 个。按验证错误数应剪枝。若风险取训练错误个数，则保留代价 $2+3\alpha$，剪去代价 $6+\alpha$。所以 $\alpha>2$ 时剪枝，$\alpha<2$ 时保留，$\alpha=2$ 时同值并按给定并列规则处理。

\section{易混辨析}
\noindent\begin{tabularx}{\linewidth}{@{}l>{\raggedright\arraybackslash}X>{\raggedright\arraybackslash}X@{}}
\toprule
概念 & 准则 & 作答时写出的区别\\
\midrule
ID3 & 最大信息增益 & 原始做法常对离散属性多路分裂，注意高基数偏好。\\
C4.5 & 用增益率选分裂 & 引入分支熵，能处理连续属性，分母为 0 的划分不可比较。\\
分类 CART & 最小加权基尼 & 二叉树。基尼不需要对数，不能照抄多路候选的结果。\\
回归 CART & 最小加权平方残差总和 & 输出数值，平方损失下的叶预测是均值。\\
预剪枝与后剪枝 & 生长时停止与生长后删枝 & 过早停止可能错过有用的后续交互。后剪枝可评估完整子树。\\
\bottomrule
\end{tabularx}

\section{常见错误与陷阱}
\begin{itemize}
\item 计算子节点熵时用本分支样本数作分母，求总条件熵时才用父节点样本数确定权重。两个叶的大小不同，不能简单平均。
\item 分支熵用各分支的大小，节点熵用各类别的大小。一个纯节点还可沿与标签无关的特征分支，但增益为零。
\item 相邻不同真实值之间才存在有效阈值。把重复值一条切到左侧、一条切到右侧违反测试规则。
\item 节点切分的边界轴对齐，整树可以组合多条这样的边界，覆盖 XOR 等单一线性阈值无法表示的类别区域。
\item 用验证数据选深度或剪枝后，最终泛化评估仍需独立测试数据。课件还讨论了缺值时按可见样本的分支比例分配权重，避免把未知直接当成随意选中的某分支。【刘扬·1机器学习（1）决策树 p65–68】
\end{itemize}

\section{往年真题与参考答案}
\paragraph{2019 秋第 1 题。} 问信息增益的定义、它为什么适合选分裂、整树边界是否线性。
\begin{enumerate}
\item 写 $I(Y;A)=H(Y)-\sum_vP(A=v)H(Y\mid A=v)$。
\item 举平衡根节点例：一条属性若把正负样本完全分开，分裂后条件熵为 0，增益 1 bit。若两支正负比例仍各半，增益是 0。最大增益挑选最能降低标签不确定性的候选。
\item 可用二维 XOR 举例：先按 $x_1$ 的阈值分，再在两支内分别按 $x_2$ 的阈值分，给相反的象限配同一标签。每次切分是轴对齐直线的一部分，合成的类别区域没有单一线性超平面。
\end{enumerate}
\paragraph{2019 秋第 3 题。} 分别说明树与 SVM 怎样缓解过拟合。树的给分点是区分训练拟合好和未见数据误差高，指出过深的树可把少量异常训练点划为纯叶，再举限深或设最小叶样本数的预剪枝、以及用验证集选择剪后的子树。SVM 部分见第 \ref{ch:prml-svm} 章。
\paragraph{2021 秋机器学习第 2 题。} 写条件熵并举课程中的用途。可取决策树为例：先写 $H(Y\mid X)=\sum_xp(x)H(Y\mid X=x)$，说明知道属性取值后计算剩余的不确定性。以取值分支占比加权，避免把只含少量样本的纯分支夸大为整个划分的高质量。固定同一父节点时，最小条件熵等价于最大信息增益。
\paragraph{2022 秋第 2 题。} 题目要求解释熵、条件熵、互信息并比较 $X_1,X_2$ 两个候选，但现存回忆未附“如下”所指的分类计数。可评分的步骤是写三项定义、按各分支人数加权、在相同父节点上选条件熵更小或增益更大的属性。若两者相等，应按另有给定的准则处理，存稿无法判定哪一个属性胜出。

\Needspace{24\baselineskip}
\section{本章自测}
\begin{enumerate}
\item 当前节点有 5 正、5 负，两个子节点为 $(4,0),(1,5)$。求信息增益。
\item 节点三类频率为 $(1/2,1/4,1/4)$，求基尼不纯度。
\item 一个候选属性在当前节点只出现同一个取值，能否计算有效的增益率？说明原因。
\item 给定二元 XOR。写出一棵深度为 2 的树的分裂和叶预测，并说明为什么它能表达交互。
\item 当前子树有 4 个叶、训练错误数 4。替换成单叶后错误数 6。使用训练错误数定义代价复杂度，求剪枝的临界 $\alpha$。
\item 依次观察到特征值 $1,1,3,5$，第一和第二个样本标签不同。列出所有按相邻不同值中点得到的候选阈值。能否单靠该特征在一次分裂中将前两个样本分离？
\end{enumerate}
\textbf{参考答案}
\begin{enumerate}
\item 条件熵为 $\tfrac6{10}h(1/6)$，增益 $1-\tfrac6{10}h(1/6)\approx0.6100$ bit。
\item $1-(1/4+1/16+1/16)=5/8$。
\item 分支熵为零，没有发生有效划分，跳过这个候选。
\item 根按 $x_1=0$ 或 $1$ 分。左支输出 $x_2$，右支输出 $1-x_2$，二者均由再按 $x_2$ 切成的两个叶实现。对 $x_2$ 的决策取决于 $x_1$ 的分支。
\item 解 $6+\alpha=4+4\alpha$ 得 $\alpha=2/3$。大于临界值时单叶代价更低。
\item 中点为 2、4。相同特征值总进入同一分支，不能用这项特征分离前两条。
\end{enumerate}

\chapter{贝叶斯决策}
\label{ch:prml-bayes}

\section{本章考点清单}
\begin{itemize}
\item \ensuremath{\bigstar\bigstar\bigstar} 先验、类条件密度、后验和最小错误率判别。能从点态风险证明全局最优、写似然比和判别函数。【刘扬·2机器学习（2）概率论速通 p26–42】【刘扬·3机器学习（3）贝叶斯判别+KNN分类器（非参数估计） p4–12】【2019秋真题 第2题】【2021深度学习与模式识别真题 第1题】【2022春真题 一(1)】【2023真题 第1题】【模拟卷四 选择4】
\item \ensuremath{\bigstar\bigstar\bigstar} 高斯类条件密度（Gaussian class-conditional density）下的等后验边界与后验。相同协方差时二次项抵消，不同时一般保留二次项。【刘扬·3机器学习（3）贝叶斯判别+KNN分类器（非参数估计） p5–9】【2019秋真题 第4题】【2022秋真题 第5题】【模拟卷三 选择6】【模拟卷四 大题6(6、8)】
\item \ensuremath{\bigstar\bigstar} 不对称误判代价和最小风险（minimum risk）。课件先给最小错误率。一般损失矩阵和阈值是其延伸。【刘扬·3机器学习（3）贝叶斯判别+KNN分类器（非参数估计） p8–11】【补充·一般代价矩阵】【模拟卷一 填空6】【模拟卷二 选择3】
\item \ensuremath{\bigstar\bigstar} 贝叶斯错误率（Bayes error rate）是已知真实分布和给定观测特征时的最低可能期望错误率，能把条件错误与总体错误区分开。【刘扬·3机器学习（3）贝叶斯判别+KNN分类器（非参数估计） p8–12、18】【2019秋真题 第2题】【模拟卷四 选择4、5】
\end{itemize}

\section{核心概念与定义}
设类别为 $\omega_1,\ldots,\omega_K$，其先验概率（prior probability）为 $\pi_k=P(Y=\omega_k)$，类条件密度为 $p(x\mid Y=\omega_k)$。观测到 $x$ 后，后验概率（posterior probability）是 $\eta_k(x)=P(Y=\omega_k\mid X=x)$。贝叶斯公式给出
\[
 \eta_k(x)=\frac{\pi_k p(x\mid\omega_k)}
                  {\sum_j\pi_j p(x\mid\omega_j)}.
\]
对离散特征使用概率质量，对连续特征使用密度。只比较类别时，分母与 $k$ 无关。题目要求后验数值时需要归一化。刘扬课件写 $P(Y\mid X)$、$P(X\mid Y)$。刘远超朴素贝叶斯一节写 $P(y)\prod_iP(x_i\mid y)$ 的最大项，这对应在条件独立假设下把类条件项分解后的形式。【刘扬·3机器学习（3）贝叶斯判别+KNN分类器（非参数估计） p6–11】【刘远超·3 分类及回归问题 p39】

条件风险（conditional risk）是给定当前观测时选择某个动作的期望损失。损失矩阵以动作作行、真实类别作列，$\lambda_{ak}$ 表示动作 $a$、真实为 $k$ 时的代价：
\[
 R(a\mid x)=\sum_k\lambda_{ak}\eta_k(x).
\]
选择条件风险最小的动作称为贝叶斯决策。在 $0$--$1$ 损失（zero–one loss）下，相当于选后验概率最大的类别。若题目给出拒绝等动作，应将其单独作为一行比较。【刘扬·3机器学习（3）贝叶斯判别+KNN分类器（非参数估计） p8–11】【补充·一般损失矩阵】

\section{公式与完整推导}
\subsection*{最小错误率的全局最优证明}
分类规则 $\delta(x)$ 给出预测类别。对给定总体，设其真实后验已知，则
\begin{align*}
 R(\delta)
 &=P(\delta(X)\ne Y)\\
 &=\mathbb E_X\!\left[
   \sum_{k=1}^K\mathbf 1\{\delta(X)\ne \omega_k\}\eta_k(X)
   \right]\\
 &=\mathbb E_X[1-\eta_{\delta(X)}(X)].
\end{align*}
对每个 $x$，$\eta_{\delta(x)}(x)\le\max_k\eta_k(x)$。右侧逐点给出最低条件错误，所以取最大后验的规则达到全局最优：
\[
 R^*=\mathbb E_X[1-\max_k\eta_k(X)].
\]
任何平局的最大后验类别都最优。有限样本学出的模型应当根据数据估计后验，实际测试错误率不保证达到 $R^*$。【刘扬·3机器学习（3）贝叶斯判别+KNN分类器（非参数估计） p8–12】

\subsection*{不对称代价、阈值和似然比}
【补充】由全期望公式，任意决策规则的总体风险等于 $\mathbb E_X[R(\delta(X)\mid X)]$。在每个 $x$ 选择最小条件风险便得到总体最优。二分类正确代价为零，令 $c_{10}$ 表示把真类 0 判为 1 的代价，$c_{01}$ 表示把真类 1 判为 0 的代价，两者均为正。记 $\eta=P(Y=1\mid x)$，则
\[
 R(1\mid x)=c_{10}(1-\eta),\qquad
 R(0\mid x)=c_{01}\eta.
\]
判 1 的严格条件为 $c_{10}(1-\eta)<c_{01}\eta$，即
\[
 \eta>\frac{c_{10}}{c_{10}+c_{01}}.
\]
在两类似然为正的区域，代入后验赔率，有
\[
 \frac{p(x\mid 1)}{p(x\mid0)}>
 \frac{c_{10}\pi_0}{c_{01}\pi_1}.
\]
增加假阳性的代价 $c_{10}$ 会提高判正阈值。代价相等时可令阈值为 $1/2$。如果正确预测也有代价，直接比较两行期望，不应套用零对角的阈值公式。【模拟卷一 填空6】【模拟卷二 选择3】

\subsection*{高斯判别函数的展开}
设 $x\in\mathbb R^d$，$\mu_k$、$\Sigma_k$ 分别是第 $k$ 类均值与正定协方差。省去所有类都相同的 $-d\ln(2\pi)/2$，取对数判别函数（discriminant function）
\[
 g_k(x)=\ln\pi_k-\tfrac12\ln|\Sigma_k|
        -\tfrac12(x-\mu_k)^\top\Sigma_k^{-1}(x-\mu_k).
\]
因对数严格递增，最大 $g_k$ 等价于最大后验。把平方型逐项展开：
\begin{align*}
 g_k(x)
 &=-\tfrac12 x^\top\Sigma_k^{-1}x
   +\mu_k^\top\Sigma_k^{-1}x
   -\tfrac12\mu_k^\top\Sigma_k^{-1}\mu_k\\
 &\quad-\tfrac12\ln|\Sigma_k|+\ln\pi_k .
\end{align*}
若 $\Sigma_k=\Sigma$ 对所有类相同，则二次项和行列式项在两类评分相减时消去。可用 $w_k^\top x+b_k$ 代替，其中
\[
 w_k=\Sigma^{-1}\mu_k,\qquad
 b_k=-\tfrac12\mu_k^\top\Sigma^{-1}\mu_k+\ln\pi_k.
\]
两类均值不同且先验为正时，边界为 $(w_1-w_0)^\top x+b_1-b_0=0$。共同协方差不必为对角或单位矩阵。若两类协方差不同，则差式的二次系数矩阵为 $-\tfrac12(\Sigma_1^{-1}-\Sigma_0^{-1})$，一般得到二次判别面（quadratic discriminant boundary）。若 $\Delta=g_1-g_0$，则
\[
 P(1\mid x)=\frac{\exp(g_1)}{\exp(g_0)+\exp(g_1)}
           =\frac{1}{1+\exp(-\Delta)}.
\]
该式也说明等协方差的二类高斯生成模型会诱导线性对数赔率，但与直接拟合逻辑回归的训练过程不同。【刘扬·3机器学习（3）贝叶斯判别+KNN分类器（非参数估计） p5–9】【2019秋真题 第4题】【2022秋真题 第5题】

\section{计算题型与解题模板}
\subsection*{题型一：离散候选动作和代价风险}
\textbf{模板} 计算每类 $\pi_kp(x\mid k)$。若问概率，将各项除以总和。若问最小风险，对损失矩阵逐行动作求加权和。

\textbf{原创例题} 三类的先验为 $(1/2,3/10,1/5)$，在某观测处的类条件值为 $(1/4,1/2,1/4)$。损失以动作作行为
\[
 \Lambda=\begin{pmatrix}0&1&5\\2&0&5\\1&1&0\end{pmatrix}.
\]
求最大后验类和最小条件风险动作。

\textbf{解答} 乘以先验的三项为 $(5,6,2)/40$，总和为 $13/40$，所以后验为 $(5,6,2)/13$。按最小错误率判第 2 类。三行分别点乘后验，得 $(16,20,11)/13$，故按最小风险选动作 3。

\subsection*{题型二：共协方差高斯的边界与后验}
\textbf{模板} 明确正类记号，求 $\Sigma^{-1}(\mu_1-\mu_0)$，单独算由均值和先验得到的截距。令差等于零得边界，实际预测看差的符号，计算后验再用 Sigmoid。

\textbf{原创例题} 设 $\mu_0=(0,0)^\top,\mu_1=(2,1)^\top$，两类共享
$\Sigma=\left(\begin{smallmatrix}2&1\\1&2\end{smallmatrix}\right)$，先验为 $\pi_0=3/4,\pi_1=1/4$。求边界和 $x=(2,-1)^\top$ 的类别 1 后验。

\textbf{解答} $\Sigma^{-1}=\frac13\left(\begin{smallmatrix}2&-1\\-1&2\end{smallmatrix}\right)$，$\Sigma^{-1}\mu_1=(1,0)^\top$，$\mu_1^\top\Sigma^{-1}\mu_1=2$。所以 $\Delta=x_1-1-\ln 3$。边界为 $x_1=1+\ln 3$，给定点的差为 $1-\ln 3<0$，判 0。类别 1 后验为 $e/(3+e)\approx0.4754$。

\subsection*{题型三：异协方差高斯和两个阈值}
\textbf{模板} 保留正态密度归一化常数中的行列式。写 $\Delta(x)$，解其零点，检查各区间内的符号。

\textbf{原创例题} 一维两类均值均为 0，$\sigma_0^2=1,\sigma_1^2=4$，两类先验相等。求判成类 1 的区域和 $x=2$ 的后验。

\textbf{解答} 方差为 4 的密度归一化因子是方差为 1 的一半。因此
\[
 \Delta(x)=-\ln2-x^2/8+x^2/2=-\ln2+3x^2/8.
\]
当 $|x|>\sqrt{(8/3)\ln2}$ 判 1，中间判 0，边界可任择其中一类。$x=2$ 时，后验为 $e^{3/2}/(e^{3/2}+2)$，约为 $0.6914$。

\section{易混辨析}
\noindent\begin{tabularx}{\linewidth}{@{}l>{\raggedright\arraybackslash}X>{\raggedright\arraybackslash}X@{}}
\toprule
对象 & 衡量什么 & 决策和推导常用形式\\
\midrule
先验与后验 & 看见 $x$ 前和看见 $x$ 后的类概率 & 将先验乘以类条件密度才得到可归一化的类得分。\\
最小错误率 & 相同的误判代价 & 取最大后验。最低条件错误是 $1-\max_k\eta_k(x)$。\\
最小风险 & 可为不等的各种代价 & 逐动作算条件期望，取风险最小的一行。\\
线性高斯判别 & 类间共用协方差 & 比较评分时共享二次项相消，先验改变截距。\\
二次高斯判别 & 类别特有协方差 & 类间二次项和行列式项一般都要保留。\\
\bottomrule
\end{tabularx}

\section{常见错误与陷阱}
\begin{itemize}
\item $\eta_k(x)$ 的分母要对所有候选类求和。只写似然或只归一化似然而漏乘先验，均可能改变类别。
\item 风险的代价索引容易颠倒。先用句子写清“真实 0 却判 1”，再填矩阵，最后与对角为零的约定核对。
\item 共同协方差导致的是成对判别式为仿射函数。多类每对边界可以是不同的超平面。若两类均值相同且先验不等，则会退化为恒定判别。
\item 对异方差题只比较马氏距离（Mahalanobis distance）会漏掉行列式和先验。密度数值可大于 1，最后归一化的后验才是类别概率。
\item 贝叶斯最优依赖真实分布、给定特征和指定损失。不能据此保证某个有限训练集学出的模型总体最优。
\end{itemize}

\section{往年真题与参考答案}
\paragraph{2019 秋第 2 题。} 求最小错误率的最优分类方法并证明。给分点是：写条件错误 $1-P(Y=\delta(x)\mid x)$，逐点选择最大后验，最后对 $X$ 取期望。应明确这是 $0$--$1$ 损失。
\paragraph{2019 秋第 4 题。} 二元类别服从伯努利分布，两类为等协方差高斯。记 $\pi_1=\rho,\pi_0=1-\rho$。写 $\Delta=(\mu_1-\mu_0)^\top\Sigma^{-1}x-\frac12(\mu_1^\top\Sigma^{-1}\mu_1-\mu_0^\top\Sigma^{-1}\mu_0)+\ln[\rho/(1-\rho)]$，令 $\Delta=0$ 给出边界。后验为 $1/(1+e^{-\Delta})$。给分点包括保留先验、说明二次项为什么抵消和写出归一化后验。
\paragraph{2021 深度学习与模式识别第 1 题。} 回忆指向患癌或未患癌的贝叶斯判别例题，未保留数据。所能确认的解答框架是：将癌症先验乘以对应观测在患癌时的概率，另算未患癌的联合项，二者相加作分母。若题目规定按最小错误率，则比较后验。若另有不同误判代价，再比较条件风险。没有原始数值，无法给出确定类别。
\paragraph{2022 春第一部分第 (1) 问。} 给定类条件和先验求最小错误率分类器。按 $g_k(x)=\pi_kp(x\mid\omega_k)$ 比大小，若求概率则除以所有 $g_k$ 的和。存稿未保留概率表，故只列通式。
\paragraph{2022 秋第 5 题。} 两类高斯何时产生线性边界。标准条件是共享同一正定协方差。写出二次项系数差为零，剩下均值产生的一次项和先验、均值产生的常数项。两类均值相异时边界是超平面。
\paragraph{2023 第 1 题。} 回忆仅记最小错误率贝叶斯决策。可确认的答题步骤与 2019 秋第 2 题相同。题目数据和更细的问法未保存。

\section{本章自测}
\begin{enumerate}
\item $P(1)=1/5$，$p(x\mid1)=4/5$，$p(x\mid0)=1/10$。求 $P(1\mid x)$。当误报代价为 3、漏报代价为 1 时判哪一类？
\item 某点三个后验为 $(0.2,0.5,0.3)$。写出 $0$--$1$ 损失下最优预测和条件错误率。
\item 一维高斯同方差为 2、等先验，均值分别为 $-1,3$。求边界。
\item 一维两类方差均为 1，$\mu_0=0,\mu_1=2$，$\pi_0=3/4,\pi_1=1/4$。求边界并指出先验造成的偏移方向。
\item 两类具有完全相同的高斯密度，先验为 $0.3,0.7$。是否出现取决于 $x$ 的等后验边界？
\item 若给损失矩阵的第 $k$ 列所有元素同时加上同一个常数 $b_k$，贝叶斯最优动作是否改变？证明。
\end{enumerate}
\textbf{参考答案}
\begin{enumerate}
\item 未归一化项为 $4/25$ 与 $2/25$，得类别 1 后验 $2/3$。阈值是 $3/4$，判 0。
\item 判第 2 类，条件错误率 $1-0.5=0.5$。
\item 等先验同方差令两边距离相等，边界为中点 $x=1$。
\item $\Delta=2x-2-\ln3$，边界 $x=1+\frac12\ln3$。先验偏向类 0，类 1 从更大的 $x$ 开始预测。
\item 不会。每处后验都等于先验，恒判先验较大的第 2 类。
\item 不变。每个动作的风险都增加 $\sum_k b_k\eta_k(x)$，此增量与动作无关，两两排序不变。
\end{enumerate}

\chapter{非参数估计与近邻分类}
\label{ch:prml-knn}

\section{本章考点清单}
\begin{itemize}
\item \ensuremath{\bigstar\bigstar\bigstar} $k$ 近邻（$k$-nearest neighbors, KNN）分类：排序、投票、距离和量纲、$k$ 的选择。【刘扬·3机器学习（3）贝叶斯判别+KNN分类器（非参数估计） p13–35】【2025秋补缓考 真题单选1（仅见一项）及判断3】【模拟卷二 选择1、大题3(3)】【模拟卷三 选择7、大题3(2)】【模拟卷四 选择8】
\item \ensuremath{\bigstar\bigstar\bigstar} Parzen 窗（Parzen window）和 KNN 密度估计、窗宽与体积。能区分“固定局部窗体积”和“固定样本数”这两种估计。【刘扬·3机器学习（3）贝叶斯判别+KNN分类器（非参数估计） p19–22】【模拟卷一 大题23(1)】【模拟卷二 大题3(1–2)】【模拟卷三 大题3(1、3)】【模拟卷四 大题3(1)】
\item \ensuremath{\bigstar\bigstar} 欧氏、曼哈顿、Minkowski、最大坐标差和汉明距离。会解释特征缩放对近邻的影响。【刘扬·3机器学习（3）贝叶斯判别+KNN分类器（非参数估计） p22–23】【刘扬·6机器学习（6）层次聚类+K-means初步 p11–19】【模拟卷三 选择7、大题3(2)】
\item \ensuremath{\bigstar\bigstar} 维数灾难（curse of dimensionality）及 1-NN 的渐近错误率上界。掌握结论的适用条件。【刘扬·3机器学习（3）贝叶斯判别+KNN分类器（非参数估计） p18、34】【模拟卷四 选择5】
\item \ensuremath{\bigstar} 窗口估计的一致性直觉、加权投票和零距离约定。【刘扬·3机器学习（3）贝叶斯判别+KNN分类器（非参数估计） p19–22】【补充·李航教材笔记“三、k 近邻法”】
\end{itemize}

\section{核心概念与定义}
参数估计（parametric estimation）先选一个有限维分布族，再估计参数，例如由样本估计高斯均值和协方差。非参数估计（nonparametric estimation）不强设这种有限维形式，直接利用附近样本估计分布或预测。窗宽 $h$、近邻数 $k$ 仍是需要选择的超参数。

\textbf{KNN 分类。} 保存带标签的训练样本。预测时按已选距离找出最近的 $k$ 条，用等权多数票或预先约定的权重投票。$k=1$ 时每个训练点形成一个最近点区域（Voronoi region）。权重若取 $1/d$ 或 $1/d^2$，有零距离时要事先定义规则，例如只在所有零距离训练样本间投票。【刘扬·3机器学习（3）贝叶斯判别+KNN分类器（非参数估计） p13–17、21–28】【补充·李航教材笔记“三、k 近邻法”】

\textbf{密度估计。} Parzen 窗先定形状和大小，再数样本贡献。KNN 密度估计先定要包含的样本数，再根据第 $k$ 个近邻决定窗口大小。课件中的 $V$ 表示体积。若一维窗半宽为 $r$，则 $V=2r$。若为 $d$ 维欧氏球，$V=v_dr^d$。边长为 $h$ 的 $d$ 维立方窗体积为 $h^d$。【刘扬·3机器学习（3）贝叶斯判别+KNN分类器（非参数估计） p19–20】

\section{公式与完整推导}
\subsection*{从局部频率得到窗口估计}
给定查询点 $x$、窗口 $R$，其中包含 $k_R$ 个训练样本，总样本数为 $n$。若训练样本独立同分布，窗口命中率的经验估计为
\[
 \widehat P(X\in R)=\frac{k_R}{n}.
\]
另一方面，$P(X\in R)=\int_Rp(u)\,du$。若窗足够小、$p$ 在窗内变化不大，则积分约为 $p(x)V$。结合两式得 $\widehat p(x)=k_R/(nV)$。固定 $V$ 估计 $k_R$ 是矩形 Parzen，固定 $k_R=k$ 估计 $V$ 是 KNN 密度。如果数据连续且距离没有平局，最小闭窗口恰含 $k$ 条样本。若边界有并列，按题目指定的计数或打破并列方式处理，不应默认球内恰有 $k$ 条。【刘扬·3机器学习（3）贝叶斯判别+KNN分类器（非参数估计） p19–20】

\subsection*{从矩形核到高斯核}
令基核 $K$ 非负且 $\int_{\mathbb R^d}K(z)\,dz=1$。带宽为 $h>0$ 时，对 $\int h^{-d}K((x-x_i)/h)\,dx$ 做变换 $z=(x-x_i)/h$，有 $dx=h^d\,dz$，故每个平移缩放的核积分仍为 1。取它们的均值，
\[
 \widehat p_h(x)=\frac{1}{nh^d}\sum_{i=1}^nK((x-x_i)/h).
\]
若 $K(z)=\mathbf 1\{\|z\|_\infty\le1/2\}$，非零项恰是边长 $h$ 的立方窗中的样本，退化为 $k_R/(nh^d)$。若 $K(z)=(2\pi)^{-d/2}\exp(-\|z\|_2^2/2)$，则
\[
 \widehat p_h(x)=\frac1{n(2\pi)^{d/2}h^d}
 \sum_i\exp\!\left[-\frac{\|x-x_i\|_2^2}{2h^2}\right].
\]
高斯核为近点赋更大权重，远点仍有很小的正贡献。矩形窗若给的是半宽 $a$，应以 $2a$ 代入边长 $h$。【刘扬·3机器学习（3）贝叶斯判别+KNN分类器（非参数估计） p19–20】

\subsection*{投票与局部后验}
令整体第 $k$ 近邻确定共用窗体积 $V$。第 $c$ 类共有 $n_c$ 条，其中 $k_c$ 条在窗内。
\[
 \widehat p(x\mid c)=\frac{k_c}{n_c V},\qquad
 \widehat P(c)=\frac{n_c}{n}.
\]
两者相乘为 $k_c/(nV)$，再按类别归一化，得到 $\widehat P(c\mid x)=k_c/k$。因此取局部后验最大等价于等权投票。这里所有类别共用同一个窗口。如果每类各自放大或缩小窗口，就不能将计算简化为一个共同邻域中的票数。【刘扬·3机器学习（3）贝叶斯判别+KNN分类器（非参数估计） p20–22】

\subsection*{距离、平滑与维数灾难}
对向量 $u,v$，Minkowski 距离为 $d_p(u,v)=(\sum_j|u_j-v_j|^p)^{1/p}$，$p\ge1$。$p=1$ 是曼哈顿距离（Manhattan distance），$p=2$ 是欧氏距离（Euclidean distance），$p\to\infty$ 时为最大坐标绝对差。相同长度的离散序列可按不同位置的个数定义汉明距离（Hamming distance）。若特征分别除以标准差 $s_j>0$，变换后的欧氏距离平方就是 $\sum_j(u_j-v_j)^2/s_j^2$。均值平移在相减时相消。用训练数据估计标准差时，两边要用同一组统计量。【刘扬·3机器学习（3）贝叶斯判别+KNN分类器（非参数估计） p22–23】【刘扬·6机器学习（6）层次聚类+K-means初步 p11–19】【补充·标准化距离计算】

【补充】设均匀分布在单位立方体。一个完整落在其中、边长 $h$ 的立方窗的预期样本比例为 $h^d$。若固定要覆盖 $q$ 的质量，则 $h=q^{1/d}$。以 $q\in(0,1)$ 固定时，此量随 $d$ 增大趋于 1，邻域已经失去局部性。这解释了高维稀疏和样本需求上升的关系。对 KNN 来说，小 $k$ 通常偏差较小、对抽样和噪声更敏感，方差较大。大 $k$ 通常更平滑、方差较小、可能跨类别边界而增大偏差。选 $k$ 用验证集或训练集内部交叉验证。对 Parzen，相似的平滑调节由 $h$ 控制。【刘扬·3机器学习（3）贝叶斯判别+KNN分类器（非参数估计） p34】【2025秋补缓考 真题判断3】

\subsection*{1-NN 渐近错误率上界}
课件给出 Cover–Hart 结论。在独立同分布训练和测试、样本量趋于无穷、最近邻位置趋向查询点，且此处类后验连续等通常条件下，记 $\eta_c=P(Y=c\mid x)$。最近邻标签与查询标签条件上趋于对同一后验分布独立抽取。两次同类的概率为 $\sum_c\eta_c^2$，故渐近条件错误率为 $1-\sum_c\eta_c^2$。设最大后验为 $m$，则
\[
 1-\sum_c\eta_c^2\le 1-m^2=(1-m)(1+m)\le 2(1-m).
\]
对 $X$ 取期望得到 $R_{1\mathrm{NN}}\le2R^*$。二类时，设较小后验为 $r$，条件错误为 $2r(1-r)$。利用 $\mathbb E[r^2]\ge(\mathbb E[r])^2$，还有 $R_{1\mathrm{NN}}\le2R^*(1-R^*)$。【刘扬·3机器学习（3）贝叶斯判别+KNN分类器（非参数估计） p18】【补充·条件风险推导】

\subsection*{相关特征与方向相似度}
马氏距离（Mahalanobis distance）$d_M(x,y)=\sqrt{(x-y)^\top\Sigma^{-1}(x-y)}$，要求 $\Sigma$ 正定。若 $\Sigma=U\Lambda U^\top$，写白化（whitening）矩阵 $\Lambda^{-1/2}U^\top$，其作用后的欧氏平方正好等于上述二次型。大方差方向被弱化。文档非零向量先归一化为 $\bm u,\bm v$ 后，$\|\bm u-\bm v\|^2=2-2\bm u^\top\bm v=2(1-\cos\theta)$，因此最小单位化欧氏距离与最大余弦相似度给相同的邻居排序。零向量的余弦需另设处理。【刘扬·3 贝叶斯判别与KNN p14、p23、p34】

\section{计算题型与解题模板}
\subsection*{题型一：同一点的矩形窗与近邻密度}
\textbf{模板} 先算每点到查询点的距离。矩形窗按给定边长或半宽计数，近邻密度按第 $k$ 小距离找半径。最后分别除以总样本数与各自的窗体积。

\textbf{原创例题} 样本为 $\{-3,-1,0,2,5\}$，查询 $x=1/4$。求半宽为 2 的一维闭矩形窗和 $k=2$ 的近邻密度。五个距离依次是 $13/4,5/4,1/4,7/4,19/4$。矩形窗计入 $-1,0,2$，体积 4，所得为 $3/(5\cdot4)=3/20$。第二近邻距离 $5/4$，其闭窗体积 $5/2$，估计为 $2/[5(5/2)]=4/25$。

\subsection*{题型二：高斯核的带宽与归一化}
\textbf{原创例题} 一维样本 $-2,0,3$，求 $x=1$ 处的高斯窗估计，取 $h=1$。先算平方距离 $9,1,4$，核贡献分别为 $e^{-9/2}/\sqrt{2\pi}$、$e^{-1/2}/\sqrt{2\pi}$ 和 $e^{-2}/\sqrt{2\pi}$。再平均，得
\[
 \widehat p(1)=\frac{e^{-9/2}+e^{-1/2}+e^{-2}}{3\sqrt{2\pi}}
 \approx0.1001.
\]
若新加一条远处同分布样本，计算要在分子加入其核贡献，并把分母的样本数由 3 改为 4。

\subsection*{题型三：多数投票和距离加权票}
\textbf{原创例题} 查询 $q=(1,1)$，训练样本 $(1,2),(3,1),(1,4),(5,2)$ 的标签依次为 $+,-,-,+$。采用欧氏距离，$k=3$。求多数票与以 $1/d^2$ 加权时的结果。

\textbf{解答} 距离平方依次为 $1,4,9,17$，所以近邻是前三点。等权时负票 2 张，判负。加权后正类票重为 $1$，负类票重为 $1/4+1/9=13/36$，判正。先确定邻域，再累加同类权重，不能只比较单个样本权重。

\section{易混辨析}
\noindent\begin{tabularx}{\linewidth}{@{}l>{\raggedright\arraybackslash}X>{\raggedright\arraybackslash}X@{}}
\toprule
方法 & 固定什么 & 使用什么作分母或预测\\
\midrule
矩形 Parzen 窗 & 带宽与窗形 & 窗内个数除以 $nV$，不同位置的计数可变。\\
高斯 Parzen 窗 & 核的带宽 & 对归一化后的连续距离权重取平均。\\
KNN 密度估计 & 近邻数 $k$ & $k/(nV(x))$，局部窗口体积变动。\\
KNN 分类 & 近邻数和距离规则 & 对近邻标签投票，局部后验可取 $k_c/k$。\\
参数模型 & 分布族或函数形式 & 存和估参数。KNN 常保存训练样本以便查询。\\
\bottomrule
\end{tabularx}

\section{常见错误与陷阱}
\begin{itemize}
\item 一维窗长是两倍半径，二维圆窗面积为 $\pi r^2$。给边长时不要再乘两倍。
\item KNN 密度的分子是近邻计数，KNN 分类分母若用于票数归一化，应是邻域大小而非全训练集大小。
\item 距离平方对排序与距离相同。但权重 $1/d$ 与 $1/d^2$ 是不同规则。
\item $k$ 奇数只减少二分类等权计票平局，多分类或距离边界仍可能并列。留一验证在检索时需排除被留出的样本本身。
\item 不把渐近风险上界理解为任何小样本集的保证。不要在用来报告性能的测试集上回头选 $k$ 或带宽。
\end{itemize}

\section{往年真题与参考答案}
\paragraph{2025 秋补缓考单选第 1 题。} 回忆仅保留一条选项，声称“$k$ 越小，模型越复杂、越容易过拟合”。在通常的局部平滑比较中该项成立。给分依据是少数近邻让单个异常训练样本更容易改变预测，模型对训练数据的波动更敏感。整题及其他选项缺失，不能恢复该选项的字母。
\paragraph{2025 秋补缓考判断第 3 题。} 回忆题认为 $k$ 增大时偏差减小、方差增大。按课程中的一般趋势，此说法错误。增大 $k$ 导致更大邻域内的信息被平均，方差趋小。附近的分布变化更容易被平滑掉，偏差趋大。
\paragraph{2020 春 A 卷第 1.1(2) 问（可选解法）。} 题目要求按小天体窗口特征做快速地形分类，允许自行选算法。若选 KNN，答题可写清按训练集统计量标准化、距离和投票、验证选 $k$、留出未见区域验证。还应评估飞行器查询阶段的时间与存储，在训练样本较少时再论证其适用。原题没有指定用 KNN。

\section{本章自测}
\begin{enumerate}
\item 二维查询点的第 3 个欧氏近邻距离为 2。训练集有 8 条，无平局。求近邻密度。
\item 一维固定窗长为 3，含 2 条样本，训练总数 9。求矩形窗密度。
\item 假设满足课件的渐近条件，贝叶斯错误率为 $0.14$。给出 1-NN 的二倍上界。
\item 两坐标的差分别为 2、4，训练标准差为 $1/2,4$。求标准化后欧氏距离的平方。
\item 对某查询点，两个训练样本与它完全重合，标签相同，其余点均有正距离。若希望按距离加权，怎样避免分母为零？
\item 等权 KNN 用 $k=n$，训练集中 B 类占全部样本的 $60\%$。为什么不论查询点移到哪里都预测同一类？
\end{enumerate}
\textbf{参考答案}
\begin{enumerate}
\item 圆窗面积为 $4\pi$，估计为 $3/(32\pi)$。
\item $2/(9\cdot3)=2/27$。
\item 上界为 $0.28$。若已知为二分类，进一步上界 $2\cdot0.14\cdot0.86=0.2408$。
\item $(2/(1/2))^2+(4/4)^2=17$。
\item 可规定有零距离时只在所有零距离点间投票，于是预测这两个样本的标签。
\item 此时邻域包含全部样本，票数分布总与全体训练样本相同，故恒预测 B。
\end{enumerate}

\chapter{朴素贝叶斯}
\label{ch:prml-nb}

\section{本章考点清单}
\begin{itemize}
\item \ensuremath{\bigstar\bigstar\bigstar} 朴素贝叶斯（naive Bayes）的建模假设、分类式、生成模型与判别模型的区别，以及何时等价于真实的贝叶斯最优规则。【刘扬·4机器学习（4）朴素贝叶斯 p2–14、18、33–34】【刘远超·3 分类及回归问题 p39–40】【2021秋机器学习真题 第3题】【模拟卷二 大题6(1)】【模拟卷三 大题6(8)】
\item \ensuremath{\bigstar\bigstar\bigstar} 多项式朴素贝叶斯（multinomial naive Bayes）与伯努利朴素贝叶斯（Bernoulli naive Bayes）。区分词元总数与文档个数，记住重复词和缺席词各自的处理。【刘远超·3 分类及回归问题 p40–47】【刘扬·4机器学习（4）朴素贝叶斯 p19–23】【模拟卷一 大题26(1–2)】【模拟卷二 大题6(2–7)】【模拟卷三 大题6(3–5)】【模拟卷四 大题6(1–3)】
\item \ensuremath{\bigstar\bigstar\bigstar} 高斯朴素贝叶斯（Gaussian naive Bayes, GNB）的均值、极大似然方差、联合评分和归一化后验。【刘远超·3 分类及回归问题 p48–50】【刘扬·4机器学习（4）朴素贝叶斯 p24–32】【模拟试题范本 2026-11-3 p2 第三部分第3题】【模拟卷一 大题26(3)】【模拟卷二 大题6(8)】【模拟卷三 大题6(6–7)】
\item \ensuremath{\bigstar\bigstar} 名义类别列计数、拉普拉斯（Laplace）与 Lidstone 平滑、对数域计算。明确每个分母对应的全部可能事件。【刘扬·4机器学习（4）朴素贝叶斯 p6–18】【刘远超·3 分类及回归问题 p39–47】【模拟卷二 大题6(3、5、7)】【模拟卷三 大题6(1–5、8)】【模拟卷四 大题6(4)】
\item \ensuremath{\bigstar\bigstar} 类间共享各维方差时为线性对数赔率，异方差一般保留平方项。【刘扬·4机器学习（4）朴素贝叶斯 p26–28】【刘远超·3 分类及回归问题 p48–50】【模拟卷四 大题6(5–8)】
\end{itemize}

\section{核心概念与定义}
给定类别 $Y=c$ 后，各特征\textbf{联合条件独立}。亦即对所有允许的特征取值
\[
 p(x_1,\ldots,x_d\mid c)=\prod_{j=1}^{d}p(x_j\mid c).
\]
因此只需估计类先验和每类下的一维条件分布。对离散分类，若每列有 $m_j$ 种取值，直接拟合类条件联合表需约 $\prod_jm_j-1$ 个自由参数。条件独立假设下需 $\sum_j(m_j-1)$ 个。这可以节省训练样本与估计成本。对于三个及以上特征，两两条件独立不足以推出这个联合分解。即便条件独立，边缘化掉类别后，各特征也可能相关。【刘扬·4机器学习（4）朴素贝叶斯 p2–10、18】

多项式朴素贝叶斯的文本特征是每个词的\textbf{出现次数}，按类估计词元分布。伯努利模型的文本特征是每个词的\textbf{出现指示}，每个词在一篇文档中只取 $0,1$。对于非文本的结构化类别数据，每列用自己的类别分布建模。高斯朴素贝叶斯则假设每个类中的各连续维独立且服从正态分布，不要求不同类别使用相同方差。【刘远超·3 分类及回归问题 p39–50】【刘扬·4机器学习（4）朴素贝叶斯 p19–28】

\section{公式与完整推导}
\subsection*{分类式和参数计数}
贝叶斯公式与条件独立给出
\[
 P(c\mid x)=
 \frac{P(c)\prod_{j=1}^d p(x_j\mid c)}
      {\sum_{c'}P(c')\prod_{j=1}^d p(x_j\mid c')}.
\]
分类只需比较分子。真正的后验需要分母。含标签训练集共有 $N$ 条，类 $c$ 有 $N_c$ 条，则频数先验 $\widehat P(c)=N_c/N$。若条件独立成立且各概率恰为真实值，最大后验就是 $0$--$1$ 损失下的贝叶斯最优决策。有限样本的频率估计还存在统计误差。【刘扬·4机器学习（4）朴素贝叶斯 p8–14、33–34】【刘远超·3 分类及回归问题 p39】

\subsection*{离散似然、加性平滑和符号}
对某个多类别分布 $\theta_v$，各事件计数为 $n_v$、总计数 $n$。忽略与参数无关的常数，对数似然是 $\sum_v n_v\ln\theta_v$。用拉格朗日乘子 $\lambda$ 加入 $\sum_v\theta_v=1$，一阶条件为 $n_v/\theta_v=\lambda$。当所有计数为正时求和得到 $\lambda=n$，故 $\widehat\theta_v=n_v/n$。零频时其极大似然值为 0。若每个事件的参数有对称 Dirichlet 先验 $\alpha>0$，后验参数为 $n_v+\alpha$，下一次事件的预测概率是后验均值
\[
 \widehat\theta_v=\frac{n_v+\alpha}{n+\alpha m}\qquad(m\text{ 为可能事件总数}).
\]
$\alpha=1$ 是加一的拉普拉斯平滑，$0<\alpha<1$ 常称 Lidstone 平滑。上式是后验预测。在同一个 Dirichlet 先验参数下，MAP 的内点表达式会使用 $n_v+\alpha-1$，边界情形需另行考虑。【刘扬·4机器学习（4）朴素贝叶斯 p10–18】【刘远超·3 分类及回归问题 p41–47】

\begin{itemize}
\item \textbf{名义类别列。} 对第 $j$ 列的某个取值 $v$，分子为类 $c$ 中该列等于 $v$ 的记录数加 $\alpha$，分母为 $N_c+\alpha m_j$。$m_j$ 是该列的可能取值数，不能用词表或其他列的种类数。
\item \textbf{多项式文本。} 设 $n$ 为词表大小（刘远超记法），$N_{yi}$ 为类 $y$ 中词 $i$ 出现的总次数，$N_y=\sum_iN_{yi}$ 为该类总词元数。记测试文本中词 $i$ 的次数为 $t_i$，使用
\[
 \theta_{yi}=P(x_i\mid y)=\frac{N_{yi}+\alpha}{N_y+\alpha n},
 \qquad
 s_y=\ln P(y)+\sum_i t_i\ln\theta_{yi}.
\]
计数向量的多项式系数是 $(\sum_it_i)!/\prod_i t_i!$，只依赖测试文本，可从所有类别得分中约去。缺席词 $t_i=0$，其项为零。
\item \textbf{伯努利文本。} 用 $D_y$ 表示类 $y$ 的文档数，用 $D_{yi}$ 表示该类至少出现一次词 $i$ 的文档数。课件的加一式为 $q_{yi}=P(x_i=1\mid y)=(D_{yi}+1)/(D_y+2)$。令 $z_i=\mathbf1\{t_i>0\}$，则
\[
 s_y=\ln P(y)+\sum_i\bigl[z_i\ln q_{yi}+(1-z_i)\ln(1-q_{yi})\bigr].
\]
缺席词有独立的补概率因子。分母加 2 是因为每个词的状态只有出现与未出现两种。
\end{itemize}

\subsection*{高斯极大似然与共方差边界}
对一个类别的一个特征，样本 $u_1,\ldots,u_m$。令 $\tau=\sigma^2>0$，其对数似然为
\[
 \ell(\mu,\tau)=-\frac m2\ln(2\pi\tau)
                -\frac1{2\tau}\sum_i(u_i-\mu)^2.
\]
对 $\mu$ 求偏导，得到 $\partial\ell/\partial\mu=\tau^{-1}\sum_i(u_i-\mu)=0$，所以 $\widehat\mu=\bar u$。在固定 $\tau$ 下，最大化 $\ell$ 等价于最小化平方残差，$\bar u$ 是唯一最优。代入后对 $\tau$ 求导，
\[
 \frac{\partial\ell}{\partial\tau}=
 -\frac{m}{2\tau}+
 \frac{\sum_i(u_i-\bar u)^2}{2\tau^2}=0,
 \qquad
 \widehat\sigma^2=\frac1m\sum_i(u_i-\bar u)^2.
\]
若平方和为正，导数在此点前正后负。如果类内某维所有值相同，正方差模型没有内点 MLE。实现时可规定方差下限等平滑策略。【刘扬·2机器学习（2）概率论速通 p45–49】【刘扬·4机器学习（4）朴素贝叶斯 p24–28】【刘远超·3 分类及回归问题 p48–50】

逐维相加得到
\[
 s_c=\ln P(c)-\frac12\sum_j
 \left[\ln(2\pi\sigma_{cj}^2)+\frac{(x_j-\mu_{cj})^2}{\sigma_{cj}^2}\right].
\]
若 $\sigma_{1j}^2=\sigma_{0j}^2=v_j$ 对每一维成立，令差值 $\Delta=s_1-s_0$，展开成
\[
 \Delta=\ln\frac{P(1)}{P(0)}
 +\sum_j\frac{\mu_{1j}-\mu_{0j}}{v_j}x_j
 -\frac12\sum_j\frac{\mu_{1j}^2-\mu_{0j}^2}{v_j}.
\]
这就是贝叶斯线性判别的对角共同协方差情形。类间方差不同，$\Delta$ 一般留下 $x_j^2$ 项。不同特征可以有不同的共享方差 $v_j$，共享指的是同一维在不同类别之间相等。【刘扬·4机器学习（4）朴素贝叶斯 p26–28】【模拟卷四 大题6(5–8)】

\subsection*{归一化和防下溢}
许多小概率直接相乘可能在浮点数中变成 0。对数把连乘变成求和，排序仍相同。要从对数分数求后验，用 $M=\max_c s_c$ 得
\[
 P(c\mid x)
 =\frac{e^{s_c-M}}{\sum_{c'}e^{s_{c'}-M}}.
\]
该变换来自分子、分母同除以 $e^M$。分母至少有一项为 1，指数最大为 0。两类时也可直接用对数优势 $s_1-s_0$ 代入 Sigmoid。【补充·稳定的对数归一化】【模拟卷三 大题6(4、7–8)】

\subsection*{条件独立的假设强度与方差约束}
把同一个特征复制一份会在对数得分中出现同一项两次，即使类别仍预测正确，也可能让后验概率过于极端。可删掉重复特征、对少数依赖特征合并设联合条件项，或以不同概率模型处理各列。离散与连续列可同用，只需为各列定义对应的条件质量或密度。高斯版本可以共享同一特征的跨类方差、共享同一类的跨特征方差，或让全部类和特征使用同一标量方差。共享参数能减少估计量数。若每个坐标的方差跨类别相同，对数判别差的对应二次项才会消掉。【刘扬·4 朴素贝叶斯 p18、p26、p34】

\section{计算题型与解题模板}
\subsection*{题型一：表格里的类别特征和平滑}
\textbf{模板} 逐列确定取值空间，按类计数，列出先验与单属性似然，最后乘起来再按类归一化。

\textbf{原创例题} 两类各 4 条记录，两个属性分别各有 3 种可能取值。待测记录的第一属性取值在两类的计数为 $3,1$，第二属性取值在两类的计数为 $0,2$。先验相等，取 $\alpha=1/2$。第一属性在两类的平滑概率分别为 $7/11,3/11$，第二属性为 $1/11,5/11$。于是条件似然为 $7/121,15/121$，第二类后验为 $15/22$。若不平滑，第一类似然会因第二属性的零频而成为 0。

\subsection*{题型二：同一词表下并排算多项式与伯努利}
\textbf{原创例题} 词表依次为“水、温、阀、泵”。A 类有 3 篇文档，B 类有 2 篇。
\[
\begin{array}{c|cc}
 &\text{各词总次数}&\text{含各词文档数}\\\hline
 A&(6,2,0,2)&(3,2,0,1)\\
 B&(1,3,3,1)&(1,2,2,1)
\end{array}
\]
先验用频率，平滑取 $\alpha=1$。对“水、水、阀”分别用两模型分类。

\textbf{多项式解答} 两类先验为 $3/5,2/5$，词向量是 $(7,3,1,3)/14$、$(2,4,4,2)/12$。测试得分为
\[
 S_A=\tfrac35(7/14)^2(1/14)=3/280,\qquad
 S_B=\tfrac25(2/12)^2(4/12)=1/270.
\]
故判 A，归一化的 A 后验为 $81/109$。

\textbf{伯努利解答} 测试存在向量为 $(1,0,1,0)$。两类出现率是 $(4,3,1,2)/5$、$(2,3,3,2)/4$。A 得分为 $\frac35\frac45(1-\frac35)\frac15(1-\frac25)=72/3125$。B 得分为 $\frac25\frac24(1-\frac34)\frac34(1-\frac24)=3/160$。也判 A，后验为 $768/1393$。这里二者类别相同，但打分不同，来源是计数和缺席事件不同。

\subsection*{题型三：多特征高斯似然及后验}
\textbf{原创例题} A 类样本为 $(0,1),(2,1),(1,-1),(1,3)$，B 类为 $(2,2),(4,2),(2,4),(4,4)$。同先验，用高斯朴素贝叶斯判别 $(2,2)$。

\textbf{解答} A 均值为 $(1,1)$，各维平方残差和为 $(2,8)$，方差为 $(1/2,2)$。B 均值为 $(3,3)$，各维平方残差和为 $(4,4)$，方差为 $(1,1)$。两类方差乘积均为 1，归一化常数相同。A 的标准化平方残差和为 $1/(1/2)+1/2=5/2$，B 为 $1+1=2$。所以 $s_B-s_A=1/4$，判 B，B 后验为 $1/(1+e^{-1/4})\approx0.5622$。

\section{易混辨析}
\noindent\begin{tabularx}{\linewidth}{@{}l>{\raggedright\arraybackslash}X>{\raggedright\arraybackslash}X@{}}
\toprule
模型或量 & 用什么统计量 & 易混点\\
\midrule
多项式 NB & 类内每词出现次数与词元总数 & 重复词进入幂次，缺席词指数为 0。\\
伯努利 NB & 类内含词的文档数与文档总数 & 重复不增加存在指示，缺席词乘 $1-q$.\\
表格类别 NB & 每属性各水平的计数 & 平滑项按该列可能水平数确定。\\
高斯 NB & 每类每维均值和 MLE 方差 & 分母用 $m$，各类可异方差。\\
生成模型与判别模型 & 前者拟合 $P(Y)$ 与 $p(X\mid Y)$，后者直接拟合 $P(Y\mid X)$ 或分类分数 & 即使都给出线性边界，NB 与逻辑回归的建模过程不同。\\
\bottomrule
\end{tabularx}

\section{常见错误与陷阱}
\begin{itemize}
\item 多项式模型不能把“同一篇是否出现”当词元次数。伯努利模型不能把一个词重复出现的次数当作文档数。
\item 类先验用各类文档数，词条件用词元数或含词文档数。不要以类内总词元数替代类先验样本数。
\item 未归一化的类得分可用于比较，写“属于该类的概率”时得除以全体得分之和。
\item 对高斯模型，用类内样本估计参数，方差除以 $m$。零方差不可直接代进密度分母。
\item 伯努利遗漏不出现的词、对数打分遗漏 $\ln P(c)$、高斯异方差遗漏 $\ln\sigma_{cj}^2$，都会潜在改变预测。
\item “条件独立就一定最优”需补真分布或一致估计等条件，不能从假设本身推出任意训练集上的最优性。
\end{itemize}

\section{往年真题与参考答案}
\paragraph{2021 秋机器学习第 3 题。} 问基本假设、参数估计好处，以及实际满足假设时能否最优。
\begin{enumerate}
\item 写 $p(x\mid y)=\prod_jp(x_j\mid y)$，指出是在给定类别下的联合条件独立。
\item 于是只拟合一维类条件分布而不用估计整个高维联合表。离散参数量由取值数连乘级降为连加级。
\item 若假设为真且所用的类先验及条件概率都是真实值，用贝叶斯公式得到真实后验，最大后验给 $0$--$1$ 损失下的最优分类器。有限样本的估计不能作同样保证。
\end{enumerate}
\paragraph{课程模拟范本 2026-11-3 第三部分第 3 题。} 此题是官方课程范本，单列供复习。鸢尾花三类各 50 条，共 150 条，特征向量是 $(4.8,2.0,2.5,2.9)$。第一得分点是写 $P(Y=c)=50/150=1/3$，三类相同。第二得分点是按 $Y=2$ 的 50 条数据对第三个特征估计 $\mu_{23}$、$\sigma_{23}^2$，写密度
\[
 p(X_3=2.5\mid Y=2)
 =\frac{1}{\sqrt{2\pi\sigma_{23}^2}}\,
  \exp\!\left[-\frac{(2.5-\mu_{23})^2}{2\sigma_{23}^2}\right].
\]
这里 $p$ 指密度，题目仅要求写公式，无需猜测每类内部的均值或方差数值。【模拟试题范本 2026-11-3 p2】

\Needspace{32\baselineskip}
\section{本章自测}
\begin{enumerate}
\item 某类训练文本共有 20 个词元，词表含 5 词，词 $w$ 从未出现。加一平滑后该词概率多少？
\item 某类有 6 篇文档，其中 2 篇含某词。用加一伯努利模型求这个词在测试文档中缺席时应乘的因子。
\item 两类对数得分为 $-103,-105$。第二类的后验是多少？怎样避免求两个极小指数后的除法失准？
\item 某类内一属性的取值范围为三种，4 条记录中目标取值出现两次。$\alpha=1/2$ 时其条件概率是多少？
\item 某类一维数据是 $-2,0,2$。求高斯 MLE 均值及方差。
\item 设二类高斯 NB 在第一维共享方差 2，正负类这维均值分别为 3 和 1。求这维在“正类减负类”的对数赔率中的一次与平方项系数。
\item 若两个特征在总体中相关，能否直接否定朴素贝叶斯的条件独立假设？同一词在伯努利模型的一篇文档中重复五次，该特征记为什么？
\end{enumerate}
\textbf{参考答案}
\begin{enumerate}
\item $(0+1)/(20+5)=1/25$。
\item 出现率为 $3/8$，缺席因子 $5/8$。
\item 两分数都减去 $-103$，第二类后验为 $e^{-2}/(1+e^{-2})\approx0.1192$。
\item $(2+1/2)/(4+3/2)=5/11$。
\item 均值为 0，MLE 方差为 $(4+0+4)/3=8/3$。
\item 一次项系数 $(3-1)/2=1$，平方项的系数为 0。
\item 条件独立应准确表述为“给定类别后，各特征联合独立”。对整个总体逐维观察，特征仍可相关。伯努利特征把同一文档中多次出现的词记为 1。
\end{enumerate}

\chapter{线性回归、多项式回归与逻辑回归}
\section{本章考点清单}
\begin{itemize}[leftmargin=2em]
\item \(\bigstar \bigstar \bigstar\) 线性回归、平方误差、正规方程、高斯噪声的 MLE 与最小二乘。【刘远超·3 分类及回归问题 p59--66】【刘扬·1机器学习（1）决策树 p70--73】【2024秋真题 大题1】【2025秋补缓考真题 简答8(1)--(2)】【模拟卷一 大题27(1)】【模拟卷二 大题7(1)--(2)】【模拟卷三 大题7(1)--(3)(7)】
\item \(\bigstar \bigstar \bigstar\) 多项式特征、验证集选阶、误差指标与过拟合。【刘远超·3 分类及回归问题 p69--80】【刘扬·1机器学习（1）决策树 p70--80】【实验 lab1 多项式拟合正弦】【2024秋真题 大题2】【2025秋补缓考真题 简答4(1)、8(3)】【模拟卷一 大题27(2)--(3)】【模拟卷二 大题7(3)--(5)】【模拟卷三 大题7(4)--(6)】
\item \(\bigstar \bigstar \bigstar\) 岭回归、Lasso、MAP 先验、截距是否惩罚及损失缩放。【刘远超·3 分类及回归问题 p77--80】【刘远超·7 网络优化 p21--22、p64】【刘扬·5机器学习（5）逻辑回归+SVM p15--17】【2021秋真题 第1题】【模拟卷一 填空8、大题27(2)】【模拟卷二 判断14、大题7(2)(10)】【模拟卷三 大题7(7)(10)】【模拟卷四 判断18、大题7(2)--(5)】Lasso 零点次梯度和软阈值推导【补充】。
\item \(\bigstar \bigstar \bigstar\) 二分类逻辑回归、Sigmoid、条件似然与二元交叉熵梯度、线性边界。【刘远超·3 分类及回归问题 p83--85】【刘扬·5机器学习（5）逻辑回归+SVM p1--18、p31--33】【实验 lab2 逻辑回归】【2019秋真题 第5题】【2021秋真题 第7题】【2022秋真题 第6题】【2024秋真题 大题3】【2025秋补缓考真题 简答1、4(2)】【模拟卷一 大题27(4)】【模拟卷二 大题7(6)--(8)】【模拟卷三 大题7(8)--(10)】严格可分时有限极小点不存在的证明【补充】。
\item \(\bigstar \bigstar \bigstar\) Softmax 回归、类别对数比、梯度 \(\bm p-\bm y\)、共同平移不变性。【刘远超·7 网络优化 p24--25】【刘扬·10前馈神经网与反向传递算法 p10、p17--18、p26】【模拟卷三 判断18】【模拟卷四 大题7(6)--(10)】Softmax 参数的不可辨识性【补充】。
\end{itemize}

\section{核心概念与定义}
\begin{description}[leftmargin=2em,style=nextline]
\item[输出与参数] 线性回归（linear regression）给实数预测 \(\hat y=\bm w^\top\bm x+b\)。把特征加首列 \(1\) 后可将截距并入 \(\bm\theta=(b,\bm w^\top)^\top\)。多项式回归（polynomial regression）改用 \(1,x,\ldots,x^m\) 等已知特征，预测关于 \(x\) 可弯曲，关于系数仍线性。因此仍可用正规方程（normal equations）。Lab1 的实验对象为带噪正弦曲线，用多项式和误差曲线比较容量。
\item[损失与正则] 最小二乘（ordinary least squares, OLS）最小化残差平方和。岭回归（ridge regression）加系数平方和，Lasso（least absolute shrinkage and selection operator）加绝对值和。截距是否入罚项、数据是否中心化由题意决定。正则化（regularization）控制容量，依独立验证选择强度，不能只求训练误差最小。
\item[二类与多类分类] 逻辑回归（logistic regression）拟合条件概率而非实值响应，二类用 \(p=\sigma(\bm w^\top\bm x+b)\)，多类的 softmax 回归（softmax regression）以每类一个线性分数 \(a_k\) 得 \(p_k=e^{a_k}/\sum_\ell e^{a_\ell}\)。两者通常以交叉熵（cross-entropy）训练。实验 lab2 的重点是分出训练和测试、用特征得到概率或判别边界，并用分类指标评价。
\item[老师的符号对应] 刘远超课件记第 \(i\) 个输入 \(\bm x^{(i)}\) 和参数 \(\bm\theta\)，偏置常并进系数向量。刘扬逻辑回归课件用输出为 \(Y\in\{0,1\}\) 的判别模型与条件对数似然 MCLE，最大化条件对数似然等价于最小化本章的交叉熵和。系数整体除 \(n\) 或加入 \(\tfrac12\) 会改变梯度尺度。正则系数和学习率的口径要与目标相配套。
\end{description}

\section{公式与完整推导}
\subsection*{从平方损失推出正规方程}
设 \(X\in\mathbb R^{n\times d}\) 由样本行组成，若含常数列则 \(d\) 已计进截距。定义
\[
 J=\frac12\|X\bm\theta-\bm y\|^2
 =\frac12\bm\theta^\top X^\top X\bm\theta-\bm y^\top X\bm\theta
   +\frac12\bm y^\top\bm y .
\]
第一项系数矩阵对称，偏导为 \(X^\top X\bm\theta\)，第二项偏导为 \(-X^\top\bm y\)。于是 \(\nabla J=X^\top(X\bm\theta-\bm y)\)。设梯度为零得正规方程 \(X^\top X\hat{\bm\theta}=X^\top\bm y\)。因为任意 \(\bm h\) 有 \(\bm h^\top X^\top X\bm h=\|X\bm h\|^2\ge0\)，驻点即全局最小。满列秩时 \(X^\top X\) 正定且可逆，\(\hat{\bm\theta}=(X^\top X)^{-1}X^\top\bm y\) 唯一。秩亏时，最小二乘解可用广义逆（pseudoinverse）\(X^+\bm y\)，也可写全部满足正规方程的解。不同解可能得到相同训练预测。
 
对单独的截距 \(b\)，残差定义 \(e_i=\hat y_i-y_i\)，有 \(\partial J/\partial b=\sum_i e_i\)，\(\partial J/\partial\bm w=\sum_i e_i\bm x_i\)。全批量梯度下降同步更新 \(\bm\theta\gets\bm\theta-\eta\nabla J\)。若题设采用 MSE，梯度为 \(2X^\top(X\bm\theta-\bm y)/n\)。若题设采用 SSE，梯度多一倍于 \(J=\mathrm{SSE}/2\) 的梯度。

给定输入 \(x_i\) 时，若 \(y_i=\bm\theta^\top\bm x_i+\epsilon_i\)，噪声独立且 \(\epsilon_i\sim\mathcal N(0,\sigma^2)\)，则
\[
 -\ln p(\bm y\mid X,\bm\theta)=\frac n2\ln(2\pi\sigma^2)
                       +\frac1{2\sigma^2}\|X\bm\theta-\bm y\|^2 .
\]
固定正方差，忽略常数并乘 \(\sigma^2\)，即得 \(\mathrm{SSE}/2\)。此等价依赖噪声假设，异方差情况下要相应按方差倒数加权。

\subsection*{岭、Lasso 与先验}
取对角矩阵 \(R\)，一般系数为 \(1\)，不罚截距时第一个元素为 \(0\)。岭目标
\[
 J_2=\tfrac12\|X\bm\theta-\bm y\|^2+\tfrac\lambda2\bm\theta^\top R\bm\theta
\]
逐项求导得 \((X^\top X+\lambda R)\bm\theta-X^\top\bm y\)。若新矩阵正定，直接解线性方程。全罚且 \(\lambda>0\)，任一非零 \(\bm h\) 有 \(\|X\bm h\|^2+\lambda\|\bm h\|^2>0\)，因此可逆。截距不罚时要看两个半正定项是否存在共同零空间。有非空的常数列且其余系数均被罚，结论仍成立。
 
若斜率先验独立为 \(\mathcal N(0,\tau^2)\)，其负对数关于参数部分为 \(\sum_jw_j^2/(2\tau^2)\)。把似然和此项相加，再乘 \(\sigma^2\)，对应上式的 \(\lambda=\sigma^2/\tau^2\)。若先验改为尺度 \(s>0\) 的零均值 Laplace，各项的负对数为 \(|w_j|/s\)，乘同样的 \(\sigma^2\) 后，Lasso
\[
 J_1=\tfrac12\mathrm{SSE}+\lambda\|\bm w\|_1,\qquad \lambda=\sigma^2/s.
\]
如果将 \(\mathrm{SSE}/2\) 改为 \(\mathrm{SSE}/(2n)\)，保持同一极值时应将惩罚系数同步除以 \(n\)。

给定一维目标 \(g(w)=\tfrac{a}{2}(w-z)^2+\lambda|w|\)，\(a>0\)。在 \(w>0\) 令导数为零得 \(w=z-\lambda/a\)，需 \(z>\lambda/a\)。在 \(w<0\) 得 \(w=z+\lambda/a\)，需 \(z<-\lambda/a\)。在 \(w=0\)，次梯度（subgradient）为 \(-az+\lambda[-1,1]\)，包含 \(0\) 当且仅当 \(|z|\le\lambda/a\)。据凸性得到
\[
 w^\star=\operatorname{sign}(z)\max\{|z|-\lambda/a,\,0\}.
\]
这解释了 Lasso 为什么可能精确置零。等价约束的 \(L_1\) 区域在二维是菱形，轴上角点有较宽的支撑法向集合，损失等高线常先接触角点。一般相关特征下各坐标不能独立用同一个软阈值式，必须结合其他坐标或满足最优性条件。

\subsection*{二元交叉熵梯度与线性边界}
令 \(a_i=\bm w^\top\bm x_i+b\)，\(p_i=(1+e^{-a_i})^{-1}\)。对独立标签的条件似然，
\[
 -\ln\mathcal L = L
   =-\sum_i\{y_i\ln p_i+(1-y_i)\ln(1-p_i)\}.
\]
单样本 \(\partial L_i/\partial p_i=-y_i/p_i+(1-y_i)/(1-p_i)\)，而 \(\mathrm dp_i/\mathrm da_i=p_i(1-p_i)\)，相乘为 \(-y_i(1-p_i)+(1-y_i)p_i=p_i-y_i\)。再乘 \(\partial a_i/\partial \bm w=\bm x_i,\ \partial a_i/\partial b=1\)，得
\[
 \nabla_{\bm w}L=\sum_i(p_i-y_i)\bm x_i,\qquad
 \partial_bL=\sum_i(p_i-y_i).
\]
若用均值损失则两式再除样本数。用 \(p\ge t\) 判正，移项得 \(a\ge\ln(t/(1-t))\)，\(t=1/2\) 就是 \(a=0\) 的超平面。线性分数经过非线性 Sigmoid 后仍有线性阈值边界。若输入已经做非线性特征变换，原输入空间的边界可以是曲线。
 
若存在 \(\bm\theta_0\) 使有符号间隔 \(m_i=y_i'\bm\theta_0^\top\tilde{\bm x}_i>0\)，且 \(y_i'\in\{-1,+1\}\)，将参数放大到 \(t\bm\theta_0\) 时
\[
 L(t)=\sum_i\ln(1+e^{-tm_i}),\quad
 \frac{\mathrm dL}{\mathrm dt}=-\sum_i\frac{m_i}{1+e^{tm_i}}<0 .
\]
\(L(t)\downarrow0\)，所有有限分数对应的损失均为正，因此无正则问题无有限的全局最优参数。加入给所有方向正惩罚的二次项后，沿该射线的惩罚随 \(t^2\) 增长。这个证明描述目标，不足以单独推出任意优化器、任意步长的迭代轨迹。

\subsection*{Softmax 交叉熵}
令 \(a_k=\bm w_k^\top\bm x+b_k\)，标签 \(\bm y\) 为 one-hot，\(\sum_k y_k=1\)。把 \(-\sum_k y_k\ln p_k\) 展开成
\[
 L=\ln\sum_k e^{a_k}-\sum_k y_ka_k.
\]
对 \(a_j\) 求导，第一项给 \(e^{a_j}/\sum_k e^{a_k}=p_j\)，第二项为 \(-y_j\)。令 \(W\) 的第 \(j\) 行是 \(\bm w_j^\top\)，则
\[
 \nabla_{\bm a}L=\bm p-\bm y,\quad
 \nabla_WL=(\bm p-\bm y)\bm x^\top,\quad
 \nabla_{\bm b}L=\bm p-\bm y.
\]
同一样本所有 \(a_k\) 加同一个常数时，分子分母同乘同一因子，概率不变。实现时减 \(\max_k a_k\) 以防指数溢出。类别 \(j,k\) 的对数比为 \(a_j-a_k=(\bm w_j-\bm w_k)^\top\bm x+b_j-b_k\)。同时给所有行加同一个向量、所有偏置加同一常数不会改变条件概率。

\subsection*{向量范数、矩阵范数与惩罚对象}
向量 $\|w\|_1=\sum_j|w_j|$、$\|w\|_2=\sqrt{\sum_j w_j^2}$、$\|w\|_\infty=\max_j|w_j|$。常用 $L_2$ 目标罚的是平方范数。矩阵诱导 $1$ 范数是绝对列和的最大值，诱导 $\infty$ 范数是绝对行和的最大值，诱导 $2$ 范数为最大奇异值。Frobenius 范数 $\|W\|_F=\sqrt{\sum_{ij}W_{ij}^2}$ 即把权重展成向量取 $L_2$ 范数。课件介绍的 $\min_j|w_j|$ 在非零向量出现零坐标时会取 0，因此不满足范数要求。线性回归的截距是 $\bar y-\hat w_1\bar x$，斜率是 $\sum(x_i-\bar x)(y_i-\bar y)/\sum(x_i-\bar x)^2$。若课件公式的系数名字次序与代码不同，以 $w_0+w_1x$ 的明确约定为准。【刘远超·3 分类及回归问题 p64、p75–76】

\section{计算题型与解题模板}
\paragraph{例 1 手算线性正规方程。}
\textbf{步骤} 写增广设计矩阵，合计列乘积，解联立式，代回求误差。样本为 \((-2,0)\)、\((0,1)\)、\((1,5)\)。
\textbf{解} 左边和右边分别为
\[
 X^\top X=\begin{pmatrix}3&-1\\-1&5\end{pmatrix},\qquad
 X^\top\bm y=(6,5)^\top .
\]
联立 \(3b-w=6\) 与 \(-b+5w=5\)，得 \(b=5/2\)、\(w=3/2\)。三个残差为 \(-1/2\)、\(3/2\)、\(-1\)。故 \(\mathrm{SSE}=7/2\)，\(\mathrm{MSE}=7/6\)。

\paragraph{例 2 多项式选阶与留出误差。}
\textbf{步骤} 利用对称输入简化正规方程，算训练误差，再用验证集比较。样本 \((-2,4),(-1,2),(1,4),(2,8)\)，另给验证 \((-3,2),(3,8)\)。
\textbf{解} 训练集均值 \(\bar x=0,\bar y=9/2\)，\(\sum x_iy_i=10,\sum x_i^2=10\)，线性为 \(9/2+x\)。二次模型 \(2+x+x^2\) 逐点成立，因设计满秩，零残差解唯一。训练 SSE 为 \(9,0\)。验证处两模型分别预测 \((3/2,15/2)\) 与 \((8,14)\)，MSE 为 \(1/4,36\)，故只看给定验证集时选择线性。

\paragraph{例 3 岭与 Lasso 的不同收缩。}
\textbf{步骤} 写 \(X^\top X\)、\(X^\top y\)，岭求线性方程。各斜率方向独立时，改写二次项并用软阈值。固定截距为零，\(X=\operatorname{diag}(\sqrt2,2)\)，\(\bm y=(\sqrt2,6)^\top\)。
\textbf{解} \(X^\top X=\operatorname{diag}(2,4),X^\top\bm y=(2,12)^\top\)，OLS 点 \((1,3)\)。岭 \(\lambda=2\) 时得到 \((2/4,12/6)=(1/2,2)\)。Lasso \(\lambda=5/2\) 时，阈值是 \(5/4\) 和 \(5/8\)，故 \((0,19/8)\)。若观测方差 \(3\)、正态先验方差 \(12\)，岭的 MAP 对应 \(\lambda=1/4\)。同一观测方差配尺度 \(2\) 的 Laplace 先验时，Lasso 的 MAP 对应 \(\lambda=3/2\)。

\paragraph{例 4 逻辑回归的一次同步更新。}
\textbf{步骤} 算各点 logit 和概率，累积交叉熵，按同一旧参数求两组梯度，最后同更新。二类样本 \((-2,0),(1,1)\)，\(w=b=\ln 3\)，用交叉熵和及学习率 \(1/5\)。
\textbf{解} logit 为 \(-\ln3,2\ln3\)，故 \(p=(1/4,9/10)\)，\(L=-\ln(3/4)-\ln(9/10)=\ln(40/27)\)。\(p-y=(1/4,-1/10)\)，故 \(\partial_wL=-1/2-1/10=-3/5,\ \partial_bL=1/4-1/10=3/20\)。结果 \(w^+=\ln3+3/25,\ b^+=\ln3-3/100\)。

\paragraph{例 5 Softmax 概率与矩阵梯度。}
\textbf{步骤} 指数归一，计算真实类负对数，求 logit 梯度及外积。设 \(a=(\ln2,\ln5,\ln3)\)，真实类为 1，\(\bm x=(-1,4)^\top\)。
\textbf{解} \(\bm p=(1/5,1/2,3/10)\)，\(L=\ln 5\)，\(\nabla_{\bm a}L=(-4/5,1/2,3/10)^\top\)。乘 \(\bm x^\top\) 得
\[
 \nabla_WL=
 \begin{pmatrix}4/5&-16/5\\-1/2&2\\-3/10&6/5\end{pmatrix}.
\]
行数是类别数，列数是特征数。第一类与第三类的对数概率比为 \(\ln(2/3)\)。

\section{易混辨析}
\noindent\begin{tabularx}{\linewidth}{@{}p{.22\linewidth}XX@{}}
\toprule
方法 & 输出或约束 & 目标与边界\\
\midrule
线性／多项式回归 & 实数。输入分别用原变量或其幂次 & 可用同一正规方程，目标通常是平方误差\\
二元逻辑回归 & Sigmoid 概率 & 二元交叉熵。固定概率阈值对应线性特征空间的超平面\\
Softmax 回归 & \(K\) 个相加为 1 的互斥概率 & 多类交叉熵。任两类分数相等定义一条成对边界\\
L1／L2 & 绝对值和／平方和 & L1 可能精确置零。L2 通常连续缩小系数\\
感知机／逻辑回归／SVM & 符号输出／概率输出／间隔得分 & 误分驱动更新／交叉熵拟合／最大间隔或合页损失加正则\\
\bottomrule
\end{tabularx}

\section{常见错误与陷阱}
\begin{itemize}[leftmargin=2em]
\item \(X\) 的每行是训练样本，每列是参数对应特征。含截距时显式加常数列。独立写 \(b\) 时不要再重复加常数。
\item 写正规方程以后确认满列秩。实际实现通常解线性系统，不显式求逆。高阶 Vandermonde 可能病态，可缩放特征或用数值稳定的分解。
\item MSE、SSE、半平方和的分母不同。先验换算只在定义一致时有效，不可以在复制代码时同时换损失归一而留下原 \(\lambda\)。
\item 阶数高只能保证包含低阶模型的最优训练误差不增，验证或测试误差不保证下降。Lab1 不可拿同一路径上的训练拟合作为泛化证明。
\item 逻辑回归给条件概率，预测阈值不必等于 \(1/2\)。交叉熵无漏项，\(y=0\) 时仍有 \(-\ln(1-p)\)。计算 softmax 时先减最大 logit 以减少溢出。
\end{itemize}

\section{往年真题与参考答案}
\begin{enumerate}[leftmargin=2em]
\item \textbf{2019 秋第 5 题。} 写二分类逻辑回归 \(p=\sigma(w^\top x+b)\) 与条件似然。写交叉熵中的两个分布：单样本目标为真实标签的一点质量 \((1-y,y)\)，模型为 \((1-p,p)\)，按样本累加。
\item \textbf{2021 秋第 1、7 题。} 第 1 题从贝叶斯后验取负对数，先验变成惩罚，高斯配 L2、拉普拉斯配 L1，正则强度要对应数据损失口径。第 7 题写 \(p\ge1/2\Leftrightarrow w^\top x+b\ge0\)，说明概率映射非线性而判别面线性。
\item \textbf{2022 秋第 6 题。} 共同点是在不做非线性特征映射时可用线性判别函数。逻辑回归的 log 损失在正间隔一侧仍变化，给概率。SVM 以合页损失惩罚间隔过小的点并追求大间隔，得分本身不必为概率。
\item \textbf{2024 秋大题 1，线性回归。} 定义 \(e_i=\bm w^\top\bm x_i+b-y_i\)。若题目说平方和，则给 \(L=\sum_i e_i^2\)，写梯度 \(2\sum_i e_i\bm x_i\) 和 \(2\sum_i e_i\)。两参数同时做旧值减学习率乘梯度。说明使用半平方和时整组常数改动。
\item \textbf{2024 秋大题 2，多项式与 MSE。} 写高次特征后按系数线性拟合。已知 \(\hat y=0.8x^2\)。对测试样本 \((4,14.8)\) 和 \((5,22.0)\)，预测分别为 \(12.8\) 和 \(20.0\)，故 MSE 为 \(4\)。
\item \textbf{2024 秋大题 3，逻辑回归。} 写 \(1/(1+e^{-a})\)，说明线性回归的输出可以是任意实数。逻辑回归将线性分数转成正类概率，用交叉熵训练，可按阈值判别。
\item \textbf{2025 秋补缓考简答 1。} 感知机输出符号且更新误分点，逻辑回归学习概率用交叉熵。按课程算法，硬间隔 SVM 显式优化最大几何间隔，感知机与标准逻辑回归训练本身没有这个有限间隔目标。
\item \textbf{2025 秋补缓考简答 4、8。} 第 4 题写 MSE、给二元交叉熵。原预测和真值缺失，参考答案的数字是样例。第 8 题列 \(J=n^{-1}\sum_i(f_{\bm\omega}(x_i)-y_i)^2\)，链式求导 \(2n^{-1}\sum_i(f_{\bm\omega}(x_i)-y_i)\nabla_{\bm\omega}f_{\bm\omega}(x_i)\)。高阶在训练集低误差却在测试集高误差，结合噪声及独立验证解释。
\end{enumerate}

\section{本章自测}
\begin{enumerate}[leftmargin=2em]
\item 样本 \((1,2),(2,6)\) 的最小二乘直线是什么？\textbf{答} 两个不同横坐标可以插值，\(w=4,b=-2\)。
\item 哪种情况下 \(X^\top X\) 不可逆？给例。 \textbf{答} 列线性相关时，如并列同一个特征两次。可用伪逆求最小范数解。
\item 半平方误差加岭中 \(a=3,z=-2,\lambda=1\)，目标为 \(a(w-z)^2/2+\lambda w^2/2\)。\textbf{答} \(w=az/(a+\lambda)=-3/2\)。
\item \(g(w)=2(w-1/3)^2+2|w|\) 的最优 \(w\) 是什么？\textbf{答} \(a=4\) 的阈值为 \(1/2\)，故 \(w=0\)。
\item 二类某点 \(a=\ln4,y=0,x=2\)，求单样本关于 \(w,b\) 的梯度。 \textbf{答} \(p=4/5\)，梯度分别 \(8/5,4/5\)。
\item 三类未归一化分数的指数为 \(1,2,7\)，真实类第三，求其交叉熵及第二类 logit 梯度。 \textbf{答} \(\ln(10/7),1/5\)。
\item 线性可分的标签是否保证无正则逻辑回归有有限 MLE？\textbf{答} 不保证。全体间隔同时放大可使训练损失趋零，有限参数时损失仍大于零。
\end{enumerate}

\chapter{感知机}
\section{本章考点清单}
\begin{itemize}[leftmargin=2em]
\item \(\bigstar \bigstar \bigstar\) 感知机判别式、增广表示、误分驱动随机更新、手算迭代。【刘远超·6 感知机模型 p5、p16--19】【刘扬·10前馈神经网与反向传递算法 p4】【李航教材笔记·二 感知机】【2025秋补缓考真题 简答1(1)】【模拟卷二 大题7(9)】【模拟卷三 单选10】
\item \(\bigstar \bigstar \bigstar\) 对偶计数、Gram 内积表、Novikoff 收敛定理和条件。【李航教材笔记·二 感知机】【模拟卷三 单选10】以增广向量写出的错误次数上界完整证明【补充】。
\item \(\bigstar \bigstar \bigstar\) XOR 不能用单个线性超平面严格分类，阈值逻辑单元与多层构造。【刘远超·6 感知机模型 p8--13】【2025秋补缓考真题 判断1】【模拟卷一 大题22(4)】【模拟卷四 大题2(1)--(2)】
\item \(\bigstar \bigstar\) 和二类逻辑回归、硬间隔 SVM、全连接网络的区别。【刘远超·6 感知机模型 p16--27】【刘扬·10前馈神经网与反向传递算法 p4--14】【刘扬·5机器学习（5）逻辑回归+SVM p31--33】【2022春真题 四(1)】【2025秋补缓考真题 简答1(1)--(2)】【模拟卷二 单选10、大题7(9)】
\end{itemize}

\section{核心概念与定义}
感知机（perceptron）是二类线性判别模型。对 \(y\in\{-1,+1\}\)，得分 \(f(\bm x)=\bm w^\top\bm x+b\)，输出其符号。令 \(z_i=y_if(\bm x_i)\)，则 \(z_i>0\) 代表严格分对。刘远超课件判别阶跃在 \(f=0\) 处取 \(+1\)，在组成逻辑门时常把输出改为 \(0,1\)。本章训练统一把 \(z_i\le0\) 视作更新事件，保证在分离面的点仍可推动边界。这条训练规则与零分数时如何输出类别是两项独立的约定。

超平面（hyperplane）为 \(f(\bm x)=0\)。当 \(\bm w\ne0\)，样本的有符号几何间隔 \(z_i/\|\bm w\|\)，到平面的无符号距离 \(|f(\bm x_i)|/\|\bm w\|\)。增广形式取 \(\tilde{\bm x}_i=(\bm x_i^\top,1)^\top\)、\(\bm u=(\bm w^\top,b)^\top\)，则 \(f=\bm u^\top\tilde{\bm x}\)。Novikoff 证明里的半径、单位向量与间隔均可在这个增广空间定义，不能把两个空间的规范混算。

对偶形式（dual form）将权重写成已触发更新的带标签训练点的线性组合。李航笔记分开写 \(b\)，偏置可单独维护，或用增广向量一并表达。两种写法只是表示不同，无需像 SVM 一样分别优化一套有盒约束的拉格朗日乘子。

\section{公式与完整推导}
\subsection*{误分损失和原始更新}
令 \(L=\sum_i\max(0,-z_i)\)。对 \(z_i<0\) 的点，该损失为 \(-y_i(\bm w^\top\bm x_i+b)\)，与不含 \(\|\bm w\|\) 分母的函数间隔相关。求梯度得 \(-y_i\bm x_i\)、\(-y_i\)，故随机梯度更新
\[
 \bm w^+=\bm w+\eta y_i\bm x_i,\qquad b^+=b+\eta y_i.
\]
对 \(z_i>0\) 梯度为零，无需更新。对 \(z_i=0\)，可在次梯度集合中选择误分侧梯度并按题设更新。目标不等同于真几何距离之和，不能在任意参数点把它们互换或额外除 \(\|\bm w\|\)。对正在更新的点，之后其有符号函数间隔增加 \(\eta(\|\bm x_i\|^2+1)>0\)：
\[
 y_i((\bm w+\eta y_i\bm x_i)^\top\bm x_i+b+\eta y_i)
   =z_i+\eta(\|\bm x_i\|^2+1).
\]
该等式只保证该样本的间隔增加，之后重访别的点时仍可能改变此前判决。

\subsection*{从原始更新推出对偶形式}
从零参数开始，按固定步长 \(\eta\) 处理训练序列。第 \(j\) 点触发 \(m_j\) 次更新，总贡献是 \(m_j\eta y_j\bm x_j\) 和 \(m_j\eta y_j\)。令 \(\alpha_j=m_j\eta\ge0\)，则
\[
 \bm w=\sum_j\alpha_jy_j\bm x_j,\quad
 b=\sum_j\alpha_jy_j,\quad
 y_if(\bm x_i)=y_i\left(\sum_j\alpha_j y_j(\bm x_j^\top\bm x_i)+b\right).
\]
预计算 Gram 矩阵（Gram matrix）\(G_{ji}=\bm x_j^\top\bm x_i\)。若 \(y_if(\bm x_i)\le0\)，只改 \(\alpha_i\gets\alpha_i+\eta\) 与 \(b\gets b+\eta y_i\)。若采用增广内积 \(\tilde G_{ji}=G_{ji}+1\)，就只维护 \(\alpha\)，不能再另加同一个 \(b\)。

\subsection*{Novikoff 结论与证明}
假设有限训练集严格线性可分，存在增广单位向量 \(\bm u_\star\) 使 \(y_i\bm u_\star^\top\tilde{\bm x}_i\ge\gamma>0\)，并且 \(\|\tilde{\bm x}_i\|\le R\)。令 \(\bm u_0=0\)，步长固定 \(\eta>0\)。只对有符号间隔不大于零的样本更新。若第 \(t\) 次更新使用第 \(i\) 点，
\[
 \bm u_{t+1}^\top\bm u_\star
 =\bm u_t^\top\bm u_\star+\eta y_i\tilde{\bm x}_i^\top\bm u_\star
 \ge\bm u_t^\top\bm u_\star+\eta\gamma.
\]
更新 \(M\) 次后左侧至少是 \(M\eta\gamma\)。再看范数增长：
\[
 \|\bm u_{t+1}\|^2
 =\|\bm u_t\|^2+2\eta y_i\bm u_t^\top\tilde{\bm x}_i
                  +\eta^2\|\tilde{\bm x}_i\|^2
 \le\|\bm u_t\|^2+\eta^2 R^2,
\]
因为这是需要更新的点，中项非正。故 \(\|\bm u_M\|\le\eta R\sqrt M\)。柯西--施瓦茨给 \(M\eta\gamma\le\bm u_M^\top\bm u_\star\le\eta R\sqrt M\)，解得 \(M\le (R/\gamma)^2\)。这是总更新次数的上界。循环遍历训练集时最终有一整轮无需再更新。严格不可分时这个证明失去正间隔假设，不能据此保证终止或最大间隔。

\subsection*{XOR 线性不可分及阈值网络}
输入 \(r_1,r_2\in\{0,1\}\)，异或（exclusive OR, XOR）的正例为 \((0,1),(1,0)\)，负例是 \((0,0),(1,1)\)。若仿射得分把正例判为非负，负例判为严格负，则两正得分的和非负，两负得分的和严格负。但仿射得分在这两对相同中点上都等于各自均值，两个和必相等，矛盾。参数不可能存在。
 
用 \(H(s)=1\) 当 \(s\ge0\)，否则为 \(0\)。单阈值能实现与门或或门，例如 \(H(r_1+r_2-1.5)\) 与 \(H(r_1+r_2-0.5)\)。取隐藏输出 \(h_{\mathrm{OR}}\)、\(h_{\mathrm{AND}}\)，外层 \(H(h_{\mathrm{OR}}-h_{\mathrm{AND}}-0.5)\) 仅在和为 \(1\) 时输出 \(1\)。若所有隐藏层都只有线性操作，矩阵相乘以后仍是单个线性映射，不会获得 XOR 的非线性表达能力。

\section{计算题型与解题模板}
\paragraph{例 1 原始更新和对偶核对。}
\textbf{步骤} 先写点序与触发规则，依次算当前有符号间隔。更新后再算下一个点，最后验证全体间隔。对偶计算用 Gram 表逐项核对。点 \(A=(-1,0)\) 为负，\(B=(1,2),C=(2,-1)\) 为正，点序为 \(A,B,C\)，步长 \(1/2\)，初始全零。
\textbf{解} A 间隔 0，更新成 \((\bm w,b)=((1/2,0),-1/2)\)。B 的得分 \(1/2-1/2=0\)，更新成 \(((1,1),0)\)。C 得分 \(2-1=1\)，不更新。再查三点间隔 \(1,3,1\)，下一轮全跳过。Gram 表为
\[
 G=\begin{pmatrix}1&-1&-2\\-1&5&0\\-2&0&5\end{pmatrix}.
\]
\(\alpha=(1/2,1/2,0)\)，故 \(b=-1/2+1/2=0\)，并且 \(-\tfrac12A+\tfrac12B=(1,1)\)。算测试点新得分时只需与 A、B 做内积，C 的系数为零。

\paragraph{例 2 误分上界的口径。}
\textbf{步骤} 增广、验单位分离向量、取最小正间隔及最大范数平方。样本 \((-3,1)\) 标签 \(-1\)，\((2,0),(1,2)\) 标签 \(+1\)。
\textbf{解} 用 \(\bm u_\star=(1,0,0)^\top\) 在增广空间分离，三个间隔 \(3,2,1\)，所以 \(\gamma=1\)。三个增广范数平方为 \(11,5,6\)，故 \(R^2=11\)，零初值的常步长算法总更新次数不超过 \(11\)。这不是实际一定更新 \(11\) 次。

\paragraph{例 3 重设物理输入阈值。}
\textbf{步骤} 枚举四组输入，找仅覆盖“恰有一个高档”的带状区间。用两个阈值单元夹住区间，再检查真值表。传感器一低、高档分别读 \(1,4\)，传感器二低、高档分别读 \(4,8\)。恰有一个传感器在高档时报警。
\textbf{解} 四种两读数和为 \(5,8,9,12\)。定义 \(h_1=H(x_1+x_2-6.5)\)、\(h_2=H(x_1+x_2-10.5)\)，输出 \(H(h_1-h_2-0.5)\)。两端得到 \((0,0)\)、\((1,1)\) 的隐藏输出，均判 \(0\)。中间两种为 \((1,0)\)，均判 \(1\)。从两类读数矩形对角线中点重合也可证明无单个仿射阈值能实现。

\section{易混辨析}
\noindent\begin{tabularx}{\linewidth}{@{}p{.25\linewidth}XX@{}}
\toprule
方法 & 训练条件或目标 & 标准预测\\
\midrule
感知机原始／对偶 & 同一误分更新，分别存 \((\bm w,b)\) 或存各样本累积系数 & 对偶只需样本内积和偏置，判别与原始相同\\
感知机／逻辑回归 & 误分时更新／条件似然或交叉熵 & 符号类别／概率再作阈值判决\\
感知机／硬间隔 SVM & 找到任一分离平面／标准化间隔并显式最大化几何间隔 & 两者可线性分类。感知机通常随训练顺序变化\\
单位阶跃／Sigmoid & 离散阈值输出／光滑概率映射 & 阶跃常用于逻辑表示，Sigmoid 可端到端求导\\
\bottomrule
\end{tabularx}

\section{常见错误与陷阱}
\begin{itemize}[leftmargin=2em]
\item 更新时同时保留标签 \(y_i\) 和学习率。\(b\) 的更新项只含 \(\eta y_i\)，不再乘特征中的某一坐标。
\item \(0,1\) 标签直接代入 \(\eta y_i\bm x_i\) 会让负类永不更新。先转换成 \(-1,+1\)，或改用与 \(0,1\) 约定一致的误差更新。
\item 增广 Gram 已含常数列，不能两次计入偏置。对偶计数与 SVM 的 KKT 乘子不要混用。
\item 使用 Novikoff 时要写明严格可分、正间隔、有界输入、零初值与固定正步长。该结论不保证解唯一或最大间隔。一次无更新完整遍历才是实际停止判据。
\item 两层网络的表达需要非线性。解决 XOR 时还须按实际正负定义核对四个角，不能只给出隐藏两个感知机的个数。
\end{itemize}

\section{往年真题与参考答案}
\begin{enumerate}[leftmargin=2em]
\item \textbf{2022 春四(1)，从浅层到深层。} 多层结构扩展表达能力，训练仍可能受连乘梯度衰减、放大与优化困难限制。得分点宜写三种机制及其影响，例如 ReLU 减少饱和区，适当方差初始化控制层间尺度，残差路径保留短梯度通路。单个感知机本身只能给线性边界，叠层还须插入非线性。
\item \textbf{2025 秋补缓考判断 1。} 单个未做非线性映射的感知机只有一个线性判别面，不能拟合 XOR 这类线性不可分问题。可用两个对角线中点重合说明反例。
\item \textbf{2025 秋补缓考简答 1(1)。} 共同部分是仿射打分。感知机以阶跃或符号直接判类，对误分点按规则纠正。逻辑回归把分数送入 Sigmoid、估计条件概率，通常用交叉熵并按阈值输出类。
\item \textbf{2025 秋补缓考简答 1(2)。} 原题围绕最大几何间隔。标准硬间隔 SVM 约束 \(y_if(\bm x_i)\ge1\) 并最小化 \(\|\bm w\|^2/2\)，显式得到最大间隔。若只在经典感知机和逻辑回归之间选，两者的训练目标都不是该最大间隔规划。
\end{enumerate}

\section{本章自测}
\begin{enumerate}[leftmargin=2em]
\item 在 \(\bm w=(0,1),b=0\) 处，样本 \((2,3)\) 的标签为 \(-1\)，步长 \(1/3\)。\textbf{答} 当前有符号间隔 \(-3\) 应更新，\(\bm w^+=(-2/3,0),b^+=-1/3\)。
\item \(\bm x=(2,-1),y=1\) 在原始参数中触发 3 次更新，步长均为 \(1/4\)。只问这一点的贡献。\textbf{答} \(\alpha=3/4\)，\(\Delta\bm w=(3/2,-3/4)\)，\(\Delta b=3/4\)。
\item \(R=3,\gamma=3/4\)，从零开始对严格可分集按题给规则更新。上界是多少？\textbf{答} \((3/(3/4))^2=16\) 次。
\item \(H(-r_1-r_2+1.5)\) 对二值输入实现哪个逻辑门？\textbf{答} 与非（NAND），仅 \(r_1=r_2=1\) 时输出 \(0\)。
\item 若两个互异样本的增广内积为 \(7\)，原输入内积是多少？\textbf{答} \(7-1=6\)。
\item 感知机得到一个能严格分类的平面，就能确定它是硬间隔 SVM 平面吗？\textbf{答} 不能由可分性确定。感知机做到零训练错误即可，SVM 还需优化几何间隔。
\end{enumerate}

\chapter{支持向量机}
\label{ch:prml-svm}

\section{本章考点清单}
\begin{itemize}
\item \ensuremath{\bigstar\bigstar\bigstar} 支持向量机（support vector machine, SVM）、超平面、函数间隔、几何间隔和最大间隔。会从近邻两类约束求硬间隔解并标明支持向量。【刘远超·3 分类及回归问题 p22–30】【刘扬·5机器学习（5）逻辑回归+SVM p23–30、44–46】【2021秋机器学习真题 第5题】【2024秋真题 大题4】【2025秋补缓考 真题简答3(1)】【模拟卷一 大题22(1)】【模拟卷二 大题2(1–4)】【模拟卷三 大题2(4–7)】
\item \ensuremath{\bigstar\bigstar\bigstar} 拉格朗日对偶（Lagrange duality）、KKT 条件（Karush–Kuhn–Tucker conditions）和偏置求法。要写出非负或盒约束、等式约束、驻点和互补松弛。【刘远超·3 分类及回归问题 p26–32】【刘扬·5机器学习（5）逻辑回归+SVM p37–49】【2021秋机器学习真题 第5题】【模拟卷一 大题22(2–3)】【模拟卷二 大题2(4、7)】【模拟卷三 大题2(7–10)】【模拟卷四 大题2(4–9)】
\item \ensuremath{\bigstar\bigstar\bigstar} 软间隔（soft margin）、松弛变量和合页损失（hinge loss）。会从样本位置判断 $\xi_i$，会用 KKT 认证数值最优解。【刘远超·3 分类及回归问题 p31–32】【刘扬·5机器学习（5）逻辑回归+SVM p31–36、47–49】【2024秋真题 大题5】【2025秋补缓考 真题简答3(2)】【模拟卷一 大题22(2–3)】【模拟卷二 判断15、大题2(5–7)】
\item \ensuremath{\bigstar\bigstar\bigstar} 核技巧（kernel trick），线性、多项式和高斯核。会算核矩阵和测试点得分，能区分在特征空间线性与原空间非线性。【刘远超·3 分类及回归问题 p33–35】【刘扬·5机器学习（5）逻辑回归+SVM p50–57】【2019秋真题 第8、9题】【2024秋真题 单选2、大题5】【2025秋补缓考 真题判断6】【模拟卷一 大题22(4)】【模拟卷二 大题2(8)】【模拟卷四 大题2(3–9)】
\item \ensuremath{\bigstar\bigstar} SVM 与逻辑回归、感知机的目标与边界比较，判断何者显式优化最大间隔。【刘扬·5机器学习（5）逻辑回归+SVM p23–36】【刘远超·3 分类及回归问题 p22–30、83–85】【2019秋真题 第3题】【2022秋真题 第6题】【2025秋补缓考 真题简答1(2)】【模拟卷二 选择10】
\end{itemize}

\section{核心概念与定义}
二分类标签为 $y_i\in\{-1,+1\}$。设得分 $f(x)=w^\top x+b$。当 $w\ne0$，超平面（hyperplane）为 $f(x)=0$，法向量是 $w$。函数间隔为 $m_i=y_i f(x_i)$，带符号的几何间隔为 $m_i/\|w\|_2$，点到平面的无符号距离为 $|f(x_i)|/\|w\|_2$。函数间隔检验标签与得分符号是否一致。刘远超课件 p24 先从绝对距离引入几何量，训练约束用含标签的带符号量。刘扬课件的最小允许间隔可直接用 $y_i f(x_i)$ 表示。硬间隔正确分类时二者的距离数值一致，软间隔误分点上带符号量为负。【刘远超·3 分类及回归问题 p23–27】【刘扬·5机器学习（5）逻辑回归+SVM p24–31】

线性可分数据可以规范化到所有 $y_i f(x_i)\ge1$。两条间隔边界是 $f(x)=+1$ 与 $f(x)=-1$，宽度为 $2/\|w\|$。到中心平面的最小训练距离是 $1/\|w\|$。对偶中 $\alpha_i>0$ 的训练样本参与 $w$ 或核得分的展开，称为支持向量（support vector）。非退化硬间隔情形下，靠在间隔边界上的点是支持向量。一般 KKT 只能保证 $\alpha_i>0\Longrightarrow y_i f(x_i)=1$（硬间隔），贴边点可因线性冗余得到零乘子。【刘远超·3 分类及回归问题 p24–30】【刘扬·5机器学习（5）逻辑回归+SVM p41–46】

\section{公式与完整推导}
\subsection*{点到平面距离和硬间隔原问题}
沿法向找投影 $x'=x-\lambda w$，代入 $w^\top x'+b=0$ 得 $\lambda=(w^\top x+b)/\|w\|^2$。所以 $\|x-x'\|=|w^\top x+b|/\|w\|$。同理，将一个间隔边界上的点沿法向投影到另一边界时，分子为两条得分水平之差 2。

严格线性可分时，对任何将训练集分开的候选 $(w,b)$，先求 $m_{\min}=\min_i y_i f(x_i)>0$。用正数 $m_{\min}$ 同除 $w,b$ 不改变决策面，且使最小函数间隔变为 1。此时最大化几何间隔等价于最小化 $\|w\|$，通常使用可微的等价目标：
\[
 \min_{w,b}\frac12\|w\|^2,\qquad
 y_i(w^\top x_i+b)\ge1.
\]
【刘远超·3 分类及回归问题 p22–28】【刘扬·5机器学习（5）逻辑回归+SVM p23–30】

\subsection*{硬间隔对偶逐步消元}
把约束改写成 $1-y_i(w^\top x_i+b)\le0$，引入 $\alpha_i\ge0$。拉格朗日函数为
\[
 L=\tfrac12\|w\|^2+\sum_i\alpha_i
  -w^\top\sum_i\alpha_i y_i x_i-b\sum_i\alpha_i y_i.
\]
对 $w$ 求导令零得 $w=\sum_i\alpha_i y_ix_i$。由于 $b$ 无二次惩罚，若 $\sum_i\alpha_i y_i\ne0$，关于 $b$ 的下确界为负无穷。因此对偶需要 $\sum_i\alpha_i y_i=0$。记 $v=\sum_i\alpha_i y_ix_i$，代入后 $\|v\|^2/2-v^\top v=-\|v\|^2/2$。将平方内积展开，得到
\[
 \max_\alpha \sum_i\alpha_i
 -\frac12\sum_{i,j}\alpha_i\alpha_j y_i y_j x_i^\top x_j,
 \qquad
 \alpha_i\ge0,\quad\sum_i\alpha_i y_i=0.
\]
硬间隔严格可分且两类均出现时，可增大候选的尺度得到所有约束严格成立。凸目标与仿射约束满足 Slater 严格可行条件，原始与对偶最优值相等。KKT 互补松弛为 $\alpha_i[y_i f(x_i)-1]=0$。若已知某个 $\alpha_s>0$，用 $b=y_s-w^\top x_s$ 恢复偏置。【刘远超·3 分类及回归问题 p28–30】【刘扬·5机器学习（5）逻辑回归+SVM p37–46】

\subsection*{软间隔、合页损失、对偶与 KKT}
加入 $\xi_i\ge0$ 允许违反间隔约束：
\[
 \min_{w,b,\xi}\tfrac12\|w\|^2+C\sum_i\xi_i,\qquad
 y_i f(x_i)\ge1-\xi_i,\quad \xi_i\ge0,\quad C>0.
\]
当 $w,b$ 固定，目标随每个 $\xi_i$ 严格增加，故最优 $\xi_i=\max(0,1-y_i f(x_i))$。代回即得 L2 权重惩罚加合页损失：
\[
 \tfrac12\|w\|^2+C\sum_i\max(0,1-y_i f(x_i)).
\]
课件也会把这个目标按 $C$ 或训练集大小作整体缩放。只要正则项与经验损失的相对系数匹配，缩放不改变最优解。

分别对 $1-\xi_i-y_i f(x_i)\le0$、$-\xi_i\le0$ 引入 $\alpha_i,\beta_i\ge0$，
\[
 L=\tfrac12\|w\|^2+\sum_i\alpha_i
    -\sum_i\alpha_i y_i f(x_i)
    +\sum_i(C-\alpha_i-\beta_i)\xi_i.
\]
原变量 $\xi_i$ 现在的非负性已进入拉格朗日约束。对 $\xi_i$ 求下确界有限要求 $C-\alpha_i-\beta_i=0$，从而 $0\le\alpha_i\le C$。对 $w,b$ 的消元与硬间隔相同，所以对偶目标不变，只把 $\alpha_i$ 的非负约束替换为盒约束。取 $w=b=0,\ \xi_i>1$ 即可严格可行。KKT 还要求原始与对偶可行，驻点满足 $w=\sum_i\alpha_i y_i x_i$ 及 $\sum_i\alpha_i y_i=0$，两组互补松弛为
\[
 \alpha_i[y_i f(x_i)-1+\xi_i]=0,\qquad
 (C-\alpha_i)\xi_i=0.
\]
因此，当 $\alpha_i<C$ 时 $\xi_i=0$。若 $0<\alpha_i<C$，必有 $y_i f(x_i)=1$，可由该点恢复偏置。若 $\alpha_i=C$，则 $y_i f(x_i)=1-\xi_i\le1$。乘子刚到上界的点也可能在间隔边界上，不一定已经误分。若 $\alpha_i=0$，则 $y_i f(x_i)\ge1$，等号在退化情形允许。无自由支持向量时，由 KKT 不等式合并所有对 $b$ 的限制。【刘远超·3 分类及回归问题 p31–32】【刘扬·5机器学习（5）逻辑回归+SVM p40、47–49】

【补充】令 $A=\|w\|^2/2$，$L=\sum_i\xi_i$，不同 $C$ 的最优量记 $(A_1,L_1)$、$(A_2,L_2)$。若 $C_2>C_1>0$，分别用两个解代入对方目标得到
\[
 A_1+C_1L_1\le A_2+C_1L_2,\qquad
 A_2+C_2L_2\le A_1+C_2L_1.
\]
相加得 $(C_2-C_1)(L_2-L_1)\le0$，所以最优合页损失和不增。再由第一式得 $A_2\ge A_1$。这不等于误分类个数必然随 $C$ 单调变化。某些数据下 $w=0$，此时不存在用 $2/\|w\|$ 度量的几何带宽。

\subsection*{核技巧及常见核}
把输入 $\phi(x)$ 代入上述推导，训练数据仅以 $\phi(x_i)^\top\phi(x_j)$ 成对出现。令
\[
 K(x,z)=\langle\phi(x),\phi(z)\rangle,\qquad
 f(x)=\sum_i\alpha_i y_iK(x_i,x)+b.
\]
用 $K$ 直接计算成对内积，无须显式构造高维 $\phi$。核矩阵 $G_{ij}=K(x_i,x_j)$ 必须对称半正定。对任意向量 $a$，
\[
 a^\top G a=\left\|\sum_i a_i\phi(x_i)\right\|^2\ge0.
\]
常用线性核为 $x^\top z$。刘远超的多项式核写为 $(x^\top z+1)^p$，一般形式为 $(\gamma x^\top z+r)^p$，对 $\gamma,r\ge0$、正整数 $p$ 可用多项式展开表示为内积。高斯径向基核（radial basis function kernel, RBF）为
\[
 K(x,z)=\exp[-\|x-z\|^2/(2\sigma^2)],\qquad \sigma>0.
\]
也常写成 $\exp[-\gamma\|x-z\|^2]$，其中 $\gamma=1/(2\sigma^2)$。很小的 $\sigma$ 让不同训练点的核相似度迅速趋近零，模型可能只在很小的邻域内拟合，容易过拟合。核空间仍可使用软间隔，使用核本身不保证训练数据严格可分。【刘远超·3 分类及回归问题 p33–35】【刘扬·5机器学习（5）逻辑回归+SVM p50–57】

\subsection*{弱对偶、强对偶与显式核特征}
对于约束 $g_i(u)\le0$ 用 $\alpha_i\ge0$，等式约束的乘子可以取任意实数。对任一原可行点 $u$，$L(u,\alpha)\le f(u)$，所以 $\inf_u L(u,\alpha)\le \inf_{u:\,g(u)\le0}f(u)$。再对乘子取上确界得到弱对偶，即对偶值不会超过原问题值。写成极大极小顺序是 $\sup_\alpha\inf_u L\le\inf_u\sup_\alpha L$。凸问题满足严格可行条件（Slater condition）时取强对偶。标准软间隔 SVM 可将松弛量取足够大而严格可行，线性可分的硬间隔可放大可分参数得到严格间隔。因此解对偶并用 KKT 可得到原问题的最优解。二维二次非齐次核 $(1+x^\top z)^2$ 的显式映射可写为 $(1,\sqrt2x_1,\sqrt2x_2,x_1^2,x_2^2,\sqrt2x_1x_2)$。各分量点积后常数、一次与二次项正好是展开式。刘扬展示的高斯核中单位带宽版本是 $\sigma=1$ 特例，刘远超用一般 $\sigma$。【刘扬·5 逻辑回归与SVM p37–39、p55–56】【刘远超·3 分类及回归问题 p35】

\section{计算题型与解题模板}
\subsection*{题型一：两侧最近支持点、硬间隔和对偶核对}
\textbf{模板} 用规范化约束给出 $\|w\|$ 的下界。找能同时达到该下界并满足所有样本约束的 $w,b$。用 $w=\sum\alpha_i y_ix_i$ 和 $\sum\alpha_i y_i=0$ 恢复乘子并核对。

\textbf{原创例题} 负点 $n_1=(-1,0)$、$n_2=(-3,-2)$，正点 $p_1=(2,1)$、$p_2=(4,2)$。求最优超平面、间隔和对偶乘子。
\begin{enumerate}
\item 支持候选对 $p_1,n_1$ 的差向量为 $u=(3,1)$。任何原问题可行解都有 $w^\top(p_1-n_1)\ge2$。由 Cauchy–Schwarz，$\|w\|\ge2/\sqrt{10}$。
\item 令 $w=2u/\|u\|^2=(3/5,1/5)$。使 $f(n_1)=-1$，得 $b=-2/5$，此时 $f(p_1)=1$。$n_2,p_2$ 的函数间隔分别是 $13/5,12/5$，均大于 1，故此解可行且达到下界。
\item 间隔总宽为 $\sqrt{10}$。两支持向量为 $n_1,p_1$。将两个乘子取为 $1/5$，其余为 0，满足等式约束且恢复 $w$。原始目标 $1/5$，对偶目标 $2/5-1/5=1/5$。
\end{enumerate}

\subsection*{题型二：给定评分算松弛量和损失}
\textbf{原创例题} 取 $w=(1,-1)$，$b=1/2$，$C=1/5$。四条带标签记录是正例 $(2,0),(0,0)$、负例 $(1,0),(0,1/2)$，只评价这组已给定的参数。
\[
\begin{array}{c|rrrr}
 & (2,0)_+&(0,0)_+&(1,0)_-&(0,1/2)_-\\\hline
 f(x)&5/2&1/2&3/2&0\\
 m_i&5/2&1/2&-3/2&0\\
 \xi_i&0&1/2&5/2&1
\end{array}
\]
第二点被正确分类但进了间隔，第三点被分错，第四点位于决策面并按规定的平局处理。$\sum_i\xi_i=4$，目标为 $1+(1/5)4=9/5$。给定 $w,b$ 后的这一数值不是整个优化问题的最优值。

\subsection*{题型三：软间隔数值认证}
\textbf{原创例题} 一维样本依次为 $(-3,-1),(0,+1),(2,-1)$，取 $C=2/5$。给出一组原始和对偶最优解。

\textbf{解答} 中间正例无法被一条线性得分同时从两侧负例之间隔离。尝试两个负点恰在负间隔边界：解 $-3w+b=-1$、$2w+b=-1$ 得 $w=0,b=-1$，故 $\xi=(0,2,0)$。对偶取 $\alpha_2=C$、$\alpha_1+\alpha_3=C$。驻点 $w=0$ 还要求 $3\alpha_1-2\alpha_3=0$，得
\[
 \alpha=(4/25,\,2/5,\,6/25).
\]
前后两点乘子严格在盒子内部，与其零松弛、间隔等于 1 吻合。中间点 $\alpha_2=C$、函数间隔为 $-1=1-\xi_2$。可行、驻点及互补松弛都成立，确为最优。也可用原始、对偶目标值均为 $4/5$ 验证。此时 $w=0$，决策得分恒为负。

\subsection*{题型四：从核特征空间求分类得分}
\textbf{原创例题} 一维训练输入 $-2,0,2$，标签依次为 $+,-,+$。使用 $K(x,z)=(xz)^2$，按硬间隔训练。求测试点 $1,\sqrt 3$ 的预测。
\begin{enumerate}
\item 核对应 $\phi(x)=x^2$，因此三个输入映射成 $4,0,4$。硬间隔要求 $b\le-1$、$4w+b\ge1$，得 $w\ge1/2$，取 $w=1/2,b=-1$ 恰可行。
\item 以两个正点等权选取一组乘子，令 $\alpha_1=\alpha_3=1/16,\ \alpha_2=1/8$。等式约束满足，恢复 $w=4(\alpha_1+\alpha_3)=1/2$，三个点都贴边，KKT 成立。
\item 测试得分为 $f(x)=x^2/2-1$。两点得分分别为 $-1/2,1/2$，分别判负和正。原空间的边界是 $\pm\sqrt2$。
\end{enumerate}

\section{易混辨析}
\noindent\begin{tabularx}{\linewidth}{@{}l>{\raggedright\arraybackslash}X>{\raggedright\arraybackslash}X@{}}
\toprule
概念 & 对应条件 & 作答时的判别\\
\midrule
函数与几何间隔 & $y_i f(x_i)$ 与 $y_i f(x_i)/\|w\|$ & 同比例缩放 $w,b$ 只改变函数间隔。\\
硬间隔与软间隔 & 强制所有间隔至少为 1，或用 $\xi_i$ 放松 & 软间隔罚的是违反量之和，包含正确但靠得太近的点。\\
自由与盒上界支持向量 & $0<\alpha<C$ 与 $\alpha=C$ & 自由者必贴边。上界者可能贴边、进带或误分。\\
感知机、LR、SVM & 按误分更新、最大条件似然、正则与最大间隔 & 线性特征下都可用线性分数，标准硬间隔 SVM 显式最大化几何间隔。\\
核映射与软间隔 & 改变特征空间与放宽硬约束 & 二者可以组合，处理不同的问题。\\
\bottomrule
\end{tabularx}

\section{常见错误与陷阱}
\begin{itemize}
\item 用 $1/\|w\|$ 时说的是归一化后的单侧间隔，$2/\|w\|$ 才是两条间隔边界间的距离。$|f(x)|$ 是得分绝对值，求欧氏距离还需除以法向量的范数。
\item 对偶等式来自偏置的驻点。漏 $b$ 或擅自罚 $b^2$ 会改变约束。
\item 软间隔最优松弛量 $0<\xi<1$ 表示有正确符号但落在间隔内。$\xi=1$ 表示在决策面，$\xi>1$ 才是严格落在错误一侧。
\item 只有自由支持向量可无条件用 $m_i=1$ 求软间隔偏置。上界 $\alpha_i=C$ 不能推出误分，也不能无条件拿来代入 $b=y_s-w^\top x_s$。
\item 使用核后仍需要 $y_i$ 以及偏置。训练仍需计算样本间的核相似度。用相似度代作核之前，还需检查它是否为有效的半正定核。
\end{itemize}

\section{往年真题与参考答案}
\paragraph{2019 秋第 3 题。} 过拟合表现为训练误差低而未见数据误差较高。SVM 通过增大规范化间隔、限制权重范数控制模型复杂度。软间隔通过验证选 $C$，在拟合与平滑之间折中。核 SVM 还需选择核的参数，过窄的 RBF 可能导致只拟合训练点周围。决策树部分见第 \ref{ch:prml-tree} 章。
\paragraph{2019 秋第 8 题。} 解释升维后为什么有时可线性分开。给分点是指出线性空间的线性组合在原输入里可对应非线性函数，如 $\phi(x)=x^2$ 时 $w\phi(x)+b=0$ 在原坐标可有两个阈值。选一个自己确实完成过的实验说明特征变换前后输入和模型，若没有使用，直接说明并用该例解释原理。
\paragraph{2019 秋第 9 题及 2021 秋机器学习第 8 题。} 各要求三个内积的例子及其含义。可列：SVM 的 $w^\top x$ 将样本投影到平面法向。核 SVM 的 $\phi(x)^\top\phi(z)$ 衡量特征空间的内积。PCA 的 $u^\top(x-\bar x)$ 得主轴投影坐标。也可用 KNN 的 $\|x-z\|^2=x^\top x-2x^\top z+z^\top z$ 替换其中一例，说明内积用于算距离。
\paragraph{2021 秋机器学习第 5 题。} 问线性可分目标、约束、解的形式和对偶好处。列出 $\min\|w\|^2/2$、$y_i(w^\top x_i+b)\ge1$，再写 $w=\sum_i\alpha_i y_ix_i$ 和对偶目标。应说明对偶凸二次规划可用 KKT 判断最优、零乘子样本不出现在得分展开，且训练样本只以内积出现，便于核化。
\paragraph{2022 秋第 6 题。} 在线性特征下，SVM 和逻辑回归都可用 $w^\top x+b=0$ 作中心边界。标准 SVM 采用合页损失，达到间隔后对单样本损失为零。逻辑回归采用对数损失 $\ln(1+e^{-y f(x)})$，直接拟合条件概率，不显式最大化几何间隔。SVM 原始得分没有概率校准。
\paragraph{2024 秋单选第 2 题及大题第 4 题。} 单选描述为用特征映射在高维找到分离超平面，对应核 SVM。大题配图的中心线为 $H_1$，两侧虚线是间隔边界。得分点为标明 $w^\top x+b=0$、指出贴两侧间隔线的最近样本在非退化硬间隔下为支持向量，再写有符号几何间隔 $y_if(x_i)/\|w\|$ 和无符号距离 $|f(x_i)|/\|w\|$。
\paragraph{2024 秋大题第 5 题。} 软间隔给分点为写 $\xi_i\ge0$、$y_if(x_i)\ge1-\xi_i$ 和 $\|w\|^2/2+C\sum_i\xi_i$，说明允许少量点侵入间隔或分错。核技巧给分点为写 $\phi$、$K(x,z)=\phi(x)^\top\phi(z)$ 及“对偶和预测只需内积”，指出特征空间线性可在原空间形成非线性边界。PDF 简记第 8 项与 DOCX 回忆的大题第 5 题相对应。
\paragraph{2025 秋补缓考判断第 6 题。} 核 SVM 在核所定义的特征空间中寻找线性超平面，该决策在原空间可以非线性。按照这一含义，回忆中的判断正确，不能把它理解成总能完全分开任意训练集。
\paragraph{2025 秋补缓考简答第 1(2) 问和第 3 题。} 感知机和逻辑回归的标准训练目标都未显式最大化几何间隔。标准硬间隔 SVM 才直接优化该量。第 3 题的第一问以硬间隔图示解释最近点，并可补充 $\alpha_i>0$ 的一般定义。第二问对当前特征中无法严格线性分开的数据使用软间隔，写出引入 $\xi_i$ 的约束和目标。若希望改变分界形状，可再选核。

\section{本章自测}
\begin{enumerate}
\item $w=(2,-1)$，$b=-1$，样本 $x=(2,1)$ 的标签为 $-1$。求 $f(x)$、函数间隔、带符号的几何间隔及最小松弛量。
\item 某规范化硬间隔解的 $\|w\|=3$。求单侧间隔与总间隔。将 $w,b$ 同乘 3 后，原决策平面会移动吗？
\item 软间隔中 $\alpha_i=C$、函数间隔为 $0.7$。在最优解中该样本的 $\xi_i$ 是多少？是否误分？
\item 令 $K(x,z)=(x^\top z+1)^2$，$x=(1,1),z=(1,2)$。求核值。
\item 为什么对偶问题需要 $\sum_i\alpha_i y_i=0$？写出含 $b$ 的拉格朗日项。
\item 软间隔数据为 $(-1,-1),(1,+1)$。给定 $C=1/4,\alpha_1=\alpha_2=C$。求 $w$ 及全部最优 $b$。
\item 若训练样本互异，RBF 的 $\sigma\to0^+$ 时核矩阵对角和非对角元素各趋于什么？这对分类泛化有什么可能影响？
\end{enumerate}
\textbf{参考答案}
\begin{enumerate}
\item 得分 $2$，函数间隔 $-2$，几何间隔 $-2/\sqrt5$，$\xi_i=3$。
\item 单侧 $1/3$，总宽 $2/3$。同乘正数不移动分数为零的平面，只会改变函数间隔尺度。
\item $\xi_i=0.3$，得分与标签同符号，正确但落在间隔内。
\item 内积为 $3$，所以是 $16$。
\item 含 $b$ 的项是 $-b\sum_i\alpha_i y_i$。若系数非零，对无约束 $b$ 求下确界会得到负无穷。
\item $w=C+C=1/2$。因两乘子都为 $C>0$，KKT 给两点函数间隔均不超过 1。由 $1/2-b\le1$、$1/2+b\le1$，得 $b\in[-1/2,1/2]$。
\item 对角项恒为 1，非对角项趋于 0。模型将不同训练点视作几乎彼此无关，可能只记住孤立样本，需用验证数据选择尺度。
\end{enumerate}

\chapter{聚类}

\section{本章考点清单}

\begin{itemize}
\item ★★★ $K$ 均值的目标函数、分配与更新、初始化和局部解。能复算一轮中心、簇内平方误差，解释簇形状与尺度的限制。 【刘远超·4 聚类问题 p18–22】【刘扬·7 K-means 符号化表达 p2–7】【2022秋真题 第3题】【2025秋补缓考 简答2】【模拟卷一 大题23(2)】【模拟卷二 大题3(4–6)】【模拟卷三 大题3(4–6)】【模拟卷四 大题3(2)】
\item ★★★ 高斯混合模型、完全数据似然、$Q$ 函数、责任度和一轮 EM。会分别更新权重、均值、协方差，解释似然不减、协方差奇异及小方差极限。 【刘远超·4 聚类问题 p36–46】【刘扬·8GMM-EM p9–40】【2019秋真题 第7题】【2021秋真题 第4、6题】【2024秋真题 大题9】【模拟卷一 大题23(4)】【模拟卷二 大题3(9–10)】【模拟卷三 大题3(10)】【模拟卷四 大题3(3–5)】
\item ★★★ DBSCAN 的核心点、边界点、噪声与题境选型。 【2025秋补缓考 单选2、简答2】【模拟卷一 大题23(5)】【模拟卷二 大题3(11)】【模拟卷三 大题3(9)】【模拟卷四 判断17】【补充】
\item ★★ 单链、全链、平均链与树状图。能每轮更新距离，按合并高度读指定簇数。 【刘远超·4 聚类问题 p25】【刘扬·6 层次聚类 p21–29】【模拟卷一 大题23(3)】【模拟卷二 大题3(8)】【模拟卷三 大题3(7–8)】【模拟卷四 大题3(10)】
\item ★★ 聚类的纯度、兰德指数、配对精确率、召回率、F 值、互信息、同质性、完整性、V-measure 和轮廓系数。 【刘远超·4 聚类问题 p8–15】【模拟卷二 大题3(7)】【模拟卷四 大题3(6–9)】
\item ★★ 谱聚类为何在谱嵌入后使用 $K$ 均值。 【2021深度真题 第4题】【2022春真题 二(3)】【补充】
\item ★ 均值偏移、自组织映射、混合伯努利、AIC 与 BIC。 【刘远超·4 聚类问题 p28、p32–33、p47–50】【刘扬·8GMM-EM p41–47】【模拟卷四 大题3(11)】
\end{itemize}

\section{核心概念与定义}

\paragraph{目标与表示。}
聚类（clustering）从没有类别标签的输入中寻找组内相似、组间有区分的划分。划分式方法直接给簇，层次方法保留由小簇合并为大簇的树，密度方法利用稠密区域的连通性，概率混合模型给每点属于各成分的后验。真实标签可以在训练后专门用来评价。一个数据集可能在颜色、大小等不同属性上产生不同且合理的聚类。距离的尺度和任务的相似性定义应先说明。【刘远超·4 聚类问题 p4–5】【刘扬·6 层次聚类 p6–20】

距离满足非负性、同一性、对称性、三角不等式。闵可夫斯基距离（Minkowski distance）
$d_p(x,y)=(\sum_j|x_j-y_j|^p)^{1/p}$ 在 $p=1,2,\infty$ 时对应曼哈顿、欧氏、切比雪夫距离。等长离散串的汉明距离（Hamming distance）统计不同位，编辑距离（edit distance）统计最少编辑代价。相关系数（correlation coefficient）衡量形状随动，常数向量的相关系数需要另行约定，不能直接用于有零方差的变量。把输入按训练数据缩放等于改变不同坐标在距离中的权重。【刘扬·6 层次聚类 p11–19】

\paragraph{$K$ 均值及初始化。}
$K$ 均值（$K$-means）用 $K$ 个中心代表 $N$ 个向量，按到中心的平方欧氏距离作硬分配。Lloyd 迭代交替执行最近中心分配与簇内均值更新。$K$-means++ 是初始化算法。首中心从样本均匀抽取，已有中心集合 $M$ 后，以 $D(x)^2/\sum_i D(x_i)^2$ 抽取下一个中心，$D(x)=\min_{\mu\in M}\|x-\mu\|$。小批量 $K$ 均值（Mini-Batch $K$-means）只抽一批样本更新，常用累计计数均值。$K$ 可通过肘部法观察 SSE 随 $K$ 的边际下降，或结合轮廓系数选用。多个起点能降低选中较差局部解的概率。【刘远超·4 聚类问题 p18–22】【刘扬·6 层次聚类 p30–39】【刘扬·8GMM-EM p2–8】

\paragraph{凝聚层次与密度聚类。}
凝聚层次聚类（agglomerative hierarchical clustering）以单点簇为初值，每轮合并距离最小的簇对。树状图（dendrogram）的高度是发生合并时的簇间距离。单链看最近点对，全链看最远点对，平均链看全部跨簇点对。按题设的簇数或高度切树，再读叶节点集合。【刘远超·4 聚类问题 p25】【刘扬·6 层次聚类 p21–28】

DBSCAN（Density-Based Spatial Clustering of Applications with Noise）给定半径 $\varepsilon$ 与阈值 $\mathrm{MinPts}$。通常闭邻域包含中心点自身。邻居数足够的是核心点。邻居数不足但进入某个核心点邻域的是边界点。其余为噪声。核心点通过核心邻居扩展，边界点不会继续扩展。两个核心连通分量若接触同一边界点，该点的最终簇归属可随访问次序变化，题面可规定。$\varepsilon$ 缩小时核心判据更严，$\mathrm{MinPts}$ 增大时同样更严。密度相差很大的簇可能难以共用一套参数。【2025秋补缓考 简答2】【补充】

\paragraph{模态与拓扑表示。}
均值偏移（mean shift）反复把查询点移动到局部邻域加权中心，收敛到相同密度模态的点可归为一组。刘远超课件先用半径球内等权平均讲偏移向量，再用邻域反复移动示意。自组织映射（Self-Organizing Map, SOM）把原型排在预定网格，每次先选最匹配单元（Best Matching Unit, BMU），再更新该原型及网格邻居。相邻网格原型有机会保留输入空间近邻关系，训练后可用 BMU 表示样本或再聚类原型。【刘远超·4 聚类问题 p28、p32–33】

\paragraph{概率混合模型。}
设观测 $\bm x_i\in\mathbb R^d$，隐变量 $\bm z_i$ 用 one-hot 指示成分。高斯混合模型（Gaussian Mixture Model, GMM）
\[
p(\bm x_i\mid\theta)=\sum_{k=1}^K\pi_k\mathcal N(\bm x_i\mid\bm\mu_k,\Sigma_k),
\quad \pi_k\ge0,\quad\sum_k\pi_k=1,\quad \Sigma_k\succ0 .
\]
其中 $\pi_k$ 是混合权重，$\bm\mu_k$ 是均值，$\Sigma_k$ 是协方差。刘远超在一维中用观测 $y_j$、指示 $\gamma_{jk}$、权重 $\alpha_k$、响应度 $\widehat\gamma_{jk}$。刘扬用 $\bm x_n,z_{nk},\pi_k,\gamma(z_{nk})$ 和多维 $\Sigma_k$。记号不同，指示量的后验均值都是责任度（responsibility）。日常说的“软分配”是 $K$ 个非负数和为 1 的责任度向量。【刘远超·4 聚类问题 p36–45】【刘扬·8GMM-EM p9–17、p34–40】

\paragraph{评价与选择。}
外部指标有参考类别，内部指标只用样本几何和预测簇。设 $n_{ck}$ 是真实类别 $c$ 与预测簇 $k$ 的交叉计数，$n_{c\cdot}$、$n_{\cdot k}$ 为边际，$N=\sum_{c,k}n_{ck}$。课件的类别划分 $C$、簇划分 $K$，分别对应以下随机变量 $Y$ 与 $Z$。同质性问「给定簇，类别是否单一」，完整性问「给定类别，是否集中在同一簇」。V-measure 调和这两者。轮廓系数（silhouette coefficient）用簇内平均距离与最近其他簇的平均距离衡量边界位置。AIC 与 BIC 用拟合度和参数复杂度评价统计模型，参数和样本量的含义需一致。 【刘远超·4 聚类问题 p8–15、p47–50】

\section{公式与完整推导}

\subsection{$K$ 均值交替下降}
以 $r_{ik}\in\{0,1\}$ 表示是否分到第 $k$ 簇，且 $\sum_k r_{ik}=1$。
\[
J(R,M)=\sum_{i=1}^N\sum_{k=1}^K r_{ik}\|\bm x_i-\bm\mu_k\|_2^2.
\]
固定中心，每个样本只涉及自己的 $K$ 个候选代价。单独取最小项能使整项和最小。固定分配，令 $n_k=\sum_i r_{ik}$，第 $k$ 簇梯度为
\[
\begin{aligned}
\nabla_{\bm\mu_k}J
&=\sum_i r_{ik}\nabla_{\bm\mu_k}
   (\bm x_i^\top\bm x_i-2\bm x_i^\top\bm\mu_k+\bm\mu_k^\top\bm\mu_k)\\
&=-2\sum_i r_{ik}\bm x_i+2 n_k\bm\mu_k .
\end{aligned}
\]
当 $n_k>0$，零梯度解 $\bm\mu_k=\sum_i r_{ik}\bm x_i/n_k$ 是最小值，因为 Hessian 为 $2n_k I$。两次操作各自不增 $J$，目标下界为 0，可稳定在依赖初始中心的结果。空簇应按题设补种中心或保留旧中心，计算时不能除以零。一次 Mini-Batch 更新中，原中心的累计样本数为 $m$，新分进该中心的点为 $\bm u_1,\ldots,\bm u_b$，累计均值
$\bm\mu'=(m\bm\mu+\sum_j\bm u_j)/(m+b)$，它来自旧样本和加新样本和。【刘远超·4 聚类问题 p18–22】【刘扬·7 K-means 符号化表达 p2–7】

\subsection{链接准则与更新}
设两个互不相交的簇为 $A,B$，点间距离为 $d$。
\[
\begin{aligned}
d_{\mathrm{single}}(A,B)&=\min_{i\in A,\,j\in B}d(i,j),\\
d_{\mathrm{complete}}(A,B)&=\max_{i\in A,\,j\in B}d(i,j),\\
d_{\mathrm{average}}(A,B)&=\frac{1}{|A||B|}\sum_{i\in A,j\in B}d(i,j).
\end{aligned}
\]
若先合并 $A,B$ 再计算到 $C$，单链的候选点对来自 $A\times C$ 或 $B\times C$，故新距离为原来两个距离的最小值。全链取最大值。平均链把两张跨簇距离表相加，其分母是 $(|A|+|B|)|C|$，因此
\[
d_{\mathrm{average}}(A\cup B,C)
=\frac{|A|d_{\mathrm{average}}(A,C)+|B|d_{\mathrm{average}}(B,C)}{|A|+|B|}.
\]
每次合并先写出选中簇对、距离与更新表，再进行下一次比较。全链通常更偏向紧凑簇，单链更容易被一串近邻桥接。【刘远超·4 聚类问题 p25】【刘扬·6 层次聚类 p24–29】

\subsection{配对计数、信息与轮廓}
记 $B(m)=\binom m2$。同类同簇的无序对数可对每个单元格分别计数，因为一对样本若同类且同簇只能落在同一格。因此
\[
\begin{aligned}
TP&=\sum_{c,k}B(n_{ck}),&
FP&=\sum_k B(n_{\cdot k})-TP,\\
FN&=\sum_c B(n_{c\cdot})-TP,&
TN&=B(N)-TP-FP-FN .
\end{aligned}
\]
$\mathrm{RI}=(TP+TN)/B(N)$，配对精确率 $TP/(TP+FP)$，配对召回率 $TP/(TP+FN)$，$F_1=2TP/(2TP+FP+FN)$。纯度为 $\sum_k \max_c n_{ck}/N$。把每个样本单独成簇可拿到纯度 1，因此纯度常需和完整性、簇数一起阅读。分母为零时按题面或实现约定，不能直接除零。【刘远超·4 聚类问题 p8–11】

设 $p_{ck}=n_{ck}/N$，边际为 $p_c,p_k$。任一非零格的条件概率 $p(c\mid k)=p_{ck}/p_k$。将其代入
\[
\begin{aligned}
H(Y\mid Z)&=-\sum_{c,k}p_{ck}\log(p_{ck}/p_k),\\
H(Y)-H(Y\mid Z)
&=-\sum_c p_c\log p_c+\sum_{c,k}p_{ck}\log(p_{ck}/p_k)\\
&=\sum_{c,k}p_{ck}\log\frac{p_{ck}}{p_cp_k}=I(Y;Z).
\end{aligned}
\]
同质性 $h=1-H(Y\mid Z)/H(Y)$，完整性 $c=1-H(Z\mid Y)/H(Z)$。分母所对应的熵为 0 时，该项按课件取 1。V-measure 为 $V_\beta=(1+\beta)hc/(\beta h+c)$，$\beta=1$ 时为 $2hc/(h+c)$，两项同为 0 时定为 0。评价过程用统一的对数底数。【刘远超·4 聚类问题 p12、p14–15】

对于点 $i$，$a_i$ 是到同簇其余点的平均距离。对每个别的簇分别算均距并取最小为 $b_i$。$s_i=(b_i-a_i)/\max(a_i,b_i)$，整体为点的平均。分母为 0 或单点簇时常约定 $s_i=0$。$s_i<0$ 表明有别的簇平均更近，不等同于已知其真实标签。 【刘远超·4 聚类问题 p13】【补充：零分母约定】

\subsection{均值偏移与自组织映射}
若以 $x$ 为中心的半径球内有 $m$ 个样本，课件的位移
$m_h(x)=m^{-1}\sum_{x_i\in S_h(x)}(x_i-x)$。加回 $x$ 得 $x'=m^{-1}\sum_{x_i\in S_h(x)}x_i$。用高斯权重时，令 $w_i(x)=\exp(-\|x-x_i\|^2/(2h^2))$。高斯核密度 $\widehat p$ 的梯度与 $\sum_i w_i(\bm x_i-\bm x)/h^2$ 成正比。把它除以 $\sum_iw_i/h^2$ 得 $\sum_iw_i\bm x_i/\sum_iw_i-\bm x$，即局部加权中心减当前点。SOM 更新可写 $\bm w_j'=\bm w_j+\eta(t)H_{j,j_*}(t)(\bm x-\bm w_j)$。$H$ 用网格的拓扑距离，获胜节点的计算则比较输入与原型在特征空间的距离。【刘远超·4 聚类问题 p28、p32–33】【补充：高斯加权推导】

\subsection{EM 的完全数据分解与两步}
已观测 $\bm X$、未观测 $\bm Z$。观察数据的边际对数似然为 $\ell(\theta)=\ln\sum_Z p(X,Z\mid\theta)$。其对数里含对成分的和。若知道 one-hot 指示量，$p(X,Z\mid\theta)=\prod_{i,k}[\pi_k p(x_i\mid k;\theta_k)]^{z_{ik}}$，由乘积取对数得
\[
\ln p(X,Z\mid\theta)=\sum_{i,k} z_{ik}\bigl[\ln\pi_k+\ln p(x_i\mid k;\theta_k)\bigr].
\]
期望最大化（Expectation-Maximization, EM）第 $t$ 轮固定旧参数 $\theta^{(t)}$。E 步先由贝叶斯公式算
\[
\gamma_{ik}=\mathbb E[z_{ik}\mid x_i;\theta^{(t)}]
=\frac{\pi_k^{(t)}p(x_i\mid k;\theta_k^{(t)})}
{\sum_j\pi_j^{(t)}p(x_i\mid j;\theta_j^{(t)})}.
\]
期望对每个指标线性作用，得到
$Q(\theta,\theta^{(t)})=\sum_{i,k}\gamma_{ik}[\ln\pi_k+\ln p(x_i\mid k;\theta_k)]$。
M 步固定这些旧后验，最大化的变量是新的候选参数 $\theta$。为什么可以这样做？设 $q$ 是 $Z$ 上的分布，展开
\[
\begin{aligned}
\mathcal F(q,\theta)
&=\sum_Z q(Z)\ln\frac{p(X,Z\mid\theta)}{q(Z)}\\
&=\sum_Z q(Z)\ln p(X\mid\theta)
 +\sum_Z q(Z)\ln\frac{p(Z\mid X,\theta)}{q(Z)}\\
&=\ell(\theta)-D_{\mathrm{KL}}(q\Vert p(Z\mid X,\theta)).
\end{aligned}
\]
因此 $\mathcal F\le\ell$。E 步令 $q=q_t=p(Z\mid X,\theta^{(t)})$，下界在旧参数处取等。M 步增大 $\mathcal F$，因为固定 $q_t$ 时熵项不变，等价于增大 $Q$。于是
\[
\ell(\theta^{(t+1)})\ge\mathcal F(q_t,\theta^{(t+1)})
\ge\mathcal F(q_t,\theta^{(t)})=\ell(\theta^{(t)}).
\]
这条链证明似然不减，不能推断全局收敛。刘扬将 EM 讲成证据下界的坐标更新。刘远超的推导从完全数据似然进入 $Q$ 函数，二者可用上面的分解衔接。【刘远超·4 聚类问题 p39–45】【刘扬·8GMM-EM p19–33】

\subsection{GMM 的 M 步三类参数}
设 $N_k=\sum_i\gamma_{ik}$。权重子目标为 $\sum_k N_k\ln\pi_k$。拉格朗日式减去 $\lambda(\sum_k\pi_k-1)$，对正权重求导得 $N_k/\pi_k-\lambda=0$，求和得 $\lambda=\sum_kN_k=N$，故 $\pi_k'=N_k/N$。

对第 $k$ 个均值，只有平方距离项含 $\mu_k$。令 $R_k=\Sigma_k^{-1}$，
\[
\nabla_{\mu_k}Q=\sum_i\gamma_{ik}R_k(\bm x_i-\bm\mu_k)
=R_k\Big(\sum_i\gamma_{ik}\bm x_i-N_k\bm\mu_k\Big).
\]
因为 $R_k$ 可逆，令梯度为 0 得 $\bm\mu_k'=\sum_i\gamma_{ik}\bm x_i/N_k$。
把新均值代回，定义 $S_k=\sum_i\gamma_{ik}(\bm x_i-\bm\mu_k')(\bm x_i-\bm\mu_k')^\top$。利用 $\ln|\Sigma_k|=-\ln|R_k|$ 和二次型的迹形式，协方差相关的目标是
\[
Q_k(R_k)=\frac{N_k}{2}\ln|R_k|-\frac12\operatorname{tr}(R_k S_k)+\text{常数}.
\]
由 $d\ln|R_k|=\operatorname{tr}(R_k^{-1}dR_k)$，
$dQ_k=\frac12\operatorname{tr}[(N_kR_k^{-1}-S_k)dR_k]$。
令微分为零，正定且 $N_k>0$ 时得 $\Sigma_k'=S_k/N_k$。
一维记方差为 $v_k$，也可直接求导
$Q_k=-N_k\ln v_k/2-\sum_i\gamma_{ik}(x_i-\mu_k')^2/(2v_k)+C$，
令 $-N_k/(2v_k)+\sum_i\gamma_{ik}(x_i-\mu_k')^2/(2v_k^2)=0$ 得相同公式。
若所有成分共用球形方差 $\Sigma_k=vI_d$，对共用 $v$ 求导得到 $v'=\sum_{i,k}\gamma_{ik}\|x_i-\mu_k'\|^2/(Nd)$。【刘远超·4 聚类问题 p43–45】【刘扬·8GMM-EM p14–17、p34–38】

若固定所有混合权重为正并使用相同球形 $v$，比两个成分的责任度时公共因子抵消。设 $k_*$ 是唯一最近中心，则对于 $k\ne k_*$，
\[
\frac{\gamma_{ik}}{\gamma_{i k_*}}=\frac{\pi_k}{\pi_{k_*}}
\exp\!\left[-\frac{\|x_i-\mu_k\|^2-\|x_i-\mu_{k_*}\|^2}{2v}\right]\longrightarrow0
\quad (v\downarrow0).
\]
最近成分的责任度趋于 1。均值更新随之趋向簇内普通平均。有限方差下，方差、相关方向和权重仍会影响后验。协方差奇异、空成分和初始化的处理见本章陷阱及实验三。【刘扬·8GMM-EM p39–40】【2021秋真题 第4、6题】

\subsection{混合伯努利、模型选择与谱嵌入}
若 $x_{ij}\in\{0,1\}$，成分 $k$ 的条件独立坐标概率为 $p_{kj}$，有
$p(x_i\mid k)=\prod_j p_{kj}^{x_{ij}}(1-p_{kj})^{1-x_{ij}}$。沿用相同责任度归一化。固定责任度，令 $a=\sum_i\gamma_{ik}x_{ij}$，$b=N_k-a$，关于 $p_{kj}$ 的 M 步子目标为 $a\ln p_{kj}+b\ln(1-p_{kj})$。导数为 $a/p_{kj}-b/(1-p_{kj})$，解得 $p_{kj}'=a/N_k$。当 $a$ 或 $b$ 为零，未平滑解在边界。外层仍按混合权重加和，单个伯努利分量的特征独立并不蕴含边际仍独立。【刘扬·8GMM-EM p41–47】

同样本、同可比较似然刻度下，令 $d_f$ 为自由参数数，$\hat\ell$ 为极大对数似然。
$\mathrm{AIC}=2d_f-2\hat\ell$，
$\mathrm{BIC}=d_f\ln N-2\hat\ell$，越小越好。
全协方差 $d$ 维 $K$ 成分 GMM 的自由度为 $(K-1)+Kd+K\,d(d+1)/2$。实际比较可使用正则约束后的有限拟合量，需记录假设和度量。【刘远超·4 聚类问题 p47–50】【补充：多维参数计数】

谱聚类（spectral clustering）先建立对称非负相似度矩阵 $W$，令度 $D_{ii}=\sum_jW_{ij}$。未归一化拉普拉斯 $L=D-W$ 有
\[
\bm f^\top L\bm f
=\sum_iD_{ii}f_i^2-\sum_{i,j}W_{ij}f_if_j
=\frac12\sum_{i,j}W_{ij}(f_i-f_j)^2\ge0.
\]
第二个等号由 $W_{ij}=W_{ji}$ 得到，说明 $L$ 半正定且常向量具有零能量。若图由多个连通分量组成，分量内常数、分量间允许不同的向量均为零能量。
若无孤立点，也可计算对称归一化拉普拉斯 $L_{\mathrm{sym}}=I-D^{-1/2}WD^{-1/2}$。寻找最小若干特征值对应方向，将特征向量按列拼接，以每个样本对应的行向量作为谱特征。课内可举对嵌入行归一化后做 $K$ 均值的版本。谱嵌入使图上相连的点表示相近，最后的 $K$ 均值把连续向量变成离散簇号。【2021深度真题 第4题】【2022春真题 二(1)、二(3)】【补充：常见拉普拉斯性质】

\section{计算题型与解题模板}

\paragraph{题型一：由列联表做外部评价。}
先求行列和，再求四种配对计数，最后计算外部分类式指标与信息指标。原创例：行是类别、列是簇，列联表为
$\bigl(\begin{smallmatrix}5&1&0\\1&2&2\end{smallmatrix}\bigr)$。总数为 11，行和 $(6,5)$，列和 $(6,3,2)$。因此
\[
TP=12,\quad FP=19-12=7,\quad FN=25-12=13,\quad TN=55-12-7-13=23.
\]
纯度为 $9/11$，$\mathrm{RI}=35/55$，配对精确率为 $12/19$，召回率为 $12/25$，$F_1=24/44$。
取 $\log_2$，非零格依次给
\[
I=\tfrac5{11}\log_2\tfrac{55}{36}+\tfrac1{11}\log_2\tfrac{11}{18}
 +\tfrac1{11}\log_2\tfrac{11}{30}+\tfrac2{11}\log_2\tfrac{22}{15}
 +\tfrac2{11}\log_2\tfrac{22}{10}.
\]
代入边际熵 $H(Y)=H(6/11,5/11)$、$H(Z)=H(6/11,3/11,2/11)$，互信息约为 $0.3890$ bit。同质性 $h=I/H(Y)\approx0.3914$，完整性 $c=I/H(Z)\approx0.2710$，等权的 $V$ 约为 $0.3203$。【刘远超·4 聚类问题 p8–15】

\paragraph{题型二：中心、SSE、轮廓与两种初始化操作。}
原创例：$X=\{0,2,3,11,14\}$，初始中心为 0 和 11。逐点比较平方距离得两簇 $\{0,2,3\}$ 和 $\{11,14\}$，初始 SSE 为 $4+9+9=22$。更新中心为 $5/3,25/2$。保持当前分配，第一簇误差为 $14/3$，第二簇为 $9/2$，总和 $55/6$。若从 2 开始做 $K$-means++，到 2 的平方距离为 $4,0,1,81,144$，点 14 被抽中的概率为 $144/230=72/115$。这与按原样本均匀抽下一个中心不同。第一簇的点 2 的 $a=(2+1)/2=3/2$，去第二簇的 $b=(9+12)/2=21/2$，所以 $s=6/7$。
独立的 Mini-Batch 小问：旧中心为 3，累计计数 4，新分到两点 8、10，更新为 $(4\cdot3+8+10)/6=5$。写题时区分「本轮更新前按旧中心算 SSE」和「更新后仍用同一分配算 SSE」。【刘远超·4 聚类问题 p13、p18–22】

\paragraph{题型三：三种链接画树。}
原创例：一维点 $A=1,B=3,C=6.5,D=11.5$。先列出 $d_{AB}=2,d_{AC}=5.5,d_{AD}=10.5,d_{BC}=3.5,d_{BD}=8.5,d_{CD}=5$。第一轮三法都合并 $A,B$。
\begin{itemize}
\item 单链：$d(AB,C)=3.5<5=d(C,D)$，先并 $C$，再以高度 $5$ 并 $D$。
\item 全链：$d(AB,C)=5.5>5$，先并 $C,D$，再以高度 $10.5$ 合并两个双点簇。
\item 平均链：$d(AB,C)=(5.5+3.5)/2=4.5<5$，先并 $C$，末步距离 $(10.5+8.5+5)/3=8$。
\end{itemize}
按从低到高列出三棵带高度的树：
\[
\begin{array}{ll}
\text{单链}&(((A,B)_2,C)_{3.5},D)_5\\
\text{全链}&((A,B)_2,(C,D)_5)_{10.5}\\
\text{平均链}&(((A,B)_2,C)_{4.5},D)_8 .
\end{array}
\]
\begin{center}
\setlength{\unitlength}{1pt}
\begin{picture}(414,86)
\put(15,5){\makebox(0,0){$A$}}\put(39,5){\makebox(0,0){$B$}}
\put(63,5){\makebox(0,0){$C$}}\put(87,5){\makebox(0,0){$D$}}
\put(15,13){\line(0,1){10}}\put(39,13){\line(0,1){10}}
\put(15,23){\line(1,0){24}}\put(27,23){\line(0,1){7.5}}
\put(63,13){\line(0,1){17.5}}\put(27,30.5){\line(1,0){36}}
\put(45,30.5){\line(0,1){7.5}}\put(87,13){\line(0,1){25}}
\put(45,38){\line(1,0){42}}\put(51,76){\makebox(0,0){单链}}
\put(159,5){\makebox(0,0){$A$}}\put(183,5){\makebox(0,0){$B$}}
\put(207,5){\makebox(0,0){$C$}}\put(231,5){\makebox(0,0){$D$}}
\put(159,13){\line(0,1){10}}\put(183,13){\line(0,1){10}}
\put(159,23){\line(1,0){24}}\put(171,23){\line(0,1){42.5}}
\put(207,13){\line(0,1){25}}\put(231,13){\line(0,1){25}}
\put(207,38){\line(1,0){24}}\put(219,38){\line(0,1){27.5}}
\put(171,65.5){\line(1,0){48}}\put(195,76){\makebox(0,0){全链}}
\put(303,5){\makebox(0,0){$A$}}\put(327,5){\makebox(0,0){$B$}}
\put(351,5){\makebox(0,0){$C$}}\put(375,5){\makebox(0,0){$D$}}
\put(303,13){\line(0,1){10}}\put(327,13){\line(0,1){10}}
\put(303,23){\line(1,0){24}}\put(315,23){\line(0,1){12.5}}
\put(351,13){\line(0,1){22.5}}\put(315,35.5){\line(1,0){36}}
\put(333,35.5){\line(0,1){17.5}}\put(375,13){\line(0,1){40}}
\put(333,53){\line(1,0){42}}\put(339,76){\makebox(0,0){平均链}}
\end{picture}
\end{center}
画图时横向排列叶，竖线高度等于对应下标，达到该高度用横线连接。题目只要求得到两簇时，读第二个高度之后、最后一次合并之前的分支。【刘远超·4 聚类问题 p25】【刘扬·6 层次聚类 p21–28】

\paragraph{题型四：DBSCAN 逐点判型。}
原创例：$\{0,1,2,3.5,10\}$，$\varepsilon=1.5$，$\mathrm{MinPts}=3$，含自身及等于半径的点。逐点列出邻域为 $\{0,1\}$、$\{0,1,2\}$、$\{1,2,3.5\}$、$\{2,3.5\}$、$\{10\}$。1 和 2 各有 3 个邻居且互相可达，是同一个密度簇的核心。0 和 3.5 都在核心邻域内，是边界点。10 为噪声。先圈核心，之后才能区分非核心点中的边界点与噪声。【2025秋补缓考 简答2】【补充】

\paragraph{题型五：GMM 完成一轮 EM。}
原创例：一维数据 $(-2,0,2)$，初始均值 $-1,1$，等权，共同方差 $v=2/\ln 2$。正均值分量对负均值分量的后验赔率为
$\exp(((x+1)^2-(x-1)^2)/(2v))=2^x$，
故正分量责任度为 $(1/5,1/2,4/5)$，负分量为 $(4/5,1/2,1/5)$。
两者的 $N_k$ 都为 $3/2$，新权重仍各为 $1/2$。正分量新均值是 $[-2(1/5)+2(4/5)]/(3/2)=4/5$，负分量为 $-4/5$。各自加权二阶矩为 $4/(3/2)=8/3$，新方差为 $8/3-(4/5)^2=152/75$。这道题把旧权重与方差用于后验，M 步的中心化用的是新均值。若参数数为 5 与 8、极大对数似然为 $-40,-36$、样本数 40，则两个 AIC 为 $90,88$，两个 BIC 为 $80+5\ln40,72+8\ln40$，两种准则的选择不同。【刘远超·4 聚类问题 p38–45、p47–49】

\paragraph{题型六：写一步均值偏移与 SOM。}
原创例：在 0、2 两个一维样本上，从 0 开始做高斯加权均值偏移。带宽 $h^2=2/\ln3$，两个未归一化权重是 $1,1/3$，所以新点是 $(0+2/3)/(1+1/3)=1/2$。SOM 的另一小问中，输入为 4，两个网格邻位原型是 2 和 9，BMU 为 2，$\eta=1/2$，BMU 的邻域权重为 1，相邻原型权重为 $1/4$。更新后分别为 $2+(4-2)/2=3$，$9+(4-9)/8=67/8$。两题都检查起点、权重和向量差的方向。 【刘远超·4 聚类问题 p28、p33】【补充：高斯权重】

\paragraph{题型七：谱聚类和伯努利混合的极简检查。}
原创例：四个顶点只有权重分别为 2 与 3 的两条互不相交无向边 $A$--$B$、$C$--$D$。对称归一化后两边的非零邻接权重都为 1，$L_{\mathrm{sym}}$ 为两个 $\bigl(\begin{smallmatrix}1&-1\\-1&1\end{smallmatrix}\bigr)$ 块。最小两个特征方向可取两个连通分量的归一化指示向量，行嵌入分别为 $(1/\sqrt2,0)$ 和 $(0,1/\sqrt2)$，$K=2$ 的 $K$ 均值将两条边分开。
另例：某伯努利成分对三个样本 $(1,0),(1,1),(0,0)$ 的责任度依次为 $1/4,1/2,1/8$，$N_k=7/8$。该成分两坐标的新参数为 $(3/4,1/2)/(7/8)=(6/7,4/7)$。两维均使用相同的成分有效样本数。【2021深度真题 第4题】【刘扬·8GMM-EM p41–47】【补充】

\section{易混辨析}

\begin{tabularx}{\linewidth}{@{}p{.22\linewidth}XX@{}}
\toprule
比较 & 第一项 & 第二项\\
\midrule
$K$ 均值／GMM & 硬分配，均值作中心，优化平方误差 & 后验软分配，估均值、协方差与权重，优化似然\\
单链／全链 & 最近点对，容易链式连通 & 最远点对，倾向紧凑，可能受离群点影响\\
$K$ 均值／DBSCAN & 给定 $K$，一般每点都分到一个中心 & 给定邻域参数，可标噪声，适合连通的非凸簇\\
同质性／完整性 & 固定预测簇，问真实类别混不混 & 固定真实类别，问预测簇散不散\\
纯度／轮廓 & 用参考标签列联表 & 用样本距离，可无标签\\
E 步／M 步 & 旧参数求隐变量的后验 & 固定旧后验权重更新模型参数\\
\bottomrule
\end{tabularx}

\section{常见错误与陷阱}
\begin{itemize}
\item 距离比较和代价有时不能共享同一种量。$K$ 均值用平方欧氏代价，层次或轮廓题用题设的距离，$K$-means++ 用到最近已选中心的\textbf{平方距离}。
\item 聚类编号可以置换。直接与类别整数逐点比较前要先对齐。纯度、兰德、互信息则可以直接对列联表计算。训练时塞进提供的真实标签会改变任务。
\item 生成数据时把任意矩阵 $M$ 改成 $MM^\top$ 只能保证半正定。满秩或加 $\epsilon I$ 才能保证正定。GMM 里样本集中在低维子空间、有效簇样本过少或某分量方差塌缩，都可能造成奇异或病态。可加对角正则、设下界、重启低权重分量、用对角／共享协方差。计算密度时优先使用 Cholesky 分解和 log-sum-exp，避免显式求逆和概率下溢。【2019秋真题 第7题】【2021秋真题 第6题】【实验三·kmeans.py】
\item 「似然不减」是精确 E 步与不降低 $Q$ 的 M 步的结论。直接在 M 步后加正则属于修改目标或更新策略，不能不加条件沿用原无正则目标的单调性。分量参数为 0、协方差不正定时，还要检查有效性。
\item GMM 的一维 $\sigma$ 是标准差，方差是 $\sigma^2$。均值与协方差的分母为 $N_k$。DBSCAN 的半径等号、是否含自身、簇切树的平局都按题设。
\end{itemize}

\section{往年真题与参考答案}
\begin{description}
\item[2019 秋第 7 题] 写观测似然 $\ell=\sum_i\ln\sum_k\pi_k\mathcal N(x_i)$，区分 E 步责任度、M 步三类参数。协方差奇异答「加权离差秩不足或成分坍缩」并给可操作办法如 $\epsilon I$ 或重初始化。原稿无逐点分值。
\item[2020 春第 2.1 题] 地形新颖性任务先判是否偏离已知分布，再聚类新颖候选。可用已知类拟合密度并以验证集定阈值。对拒识样本用可解释特征聚类，按人工核查、簇稳定性和危险地形漏检率检验。属于方案题，评分关键是把检出、分组、验证三个环节说清。
\item[2021 秋第 4、6 题] 比较 $K$ 均值硬分配与 GMM 软分配，写 $\gamma_{ik}$、$Q=\sum\gamma(\ln\pi+\ln\mathcal N)$ 与更新量。指出 $K$ 均值以平方误差定义簇，GMM 能表达协方差方向。奇异题补充有效计数及修正。
\item[2021 深度第 4 题、2022 春二(3)] 谱聚类答「样本建图、计算图拉普拉斯、取特征向量嵌入、在嵌入行上做 $K$ 均值」。谱是矩阵的特征值和对应特征向量。最后一步把变换后的点转成清晰的簇号。直接对原空间用 $K$ 均值可能错分曲形簇。
\item[2022 春二(1)] 回忆题写「拉普拉斯性质 1」，没有抄出原性质。若题设为未归一化二次型，按「写 $W,D,L$—展开两项—对称合并—解释半正定和零能量」证明，推导见本章。原题具体指定的性质仍缺失。
\item[2022 秋第 3 题] $K$ 均值先初始化，交替最近中心分配和均值更新，直到目标或分配稳定。分配逐个样本取最小，更新是各簇二次误差最小化。答出平方误差、两步的理由和局部解。
\item[2024 秋大题 9] 参考回忆 DOCX 对应题号为 9，整理 PDF 主观题第 5 项有同题。得分点是隐变量与完全数据、用旧参数算责任度、固定责任度做极大化、权重均值协方差各是什么。不需要假设缺失的具体数表。
\item[2025 秋补缓考单选 2 与简答 2] 单选留存的「不同初始化结果相同」为错。其他选项未完整保留。简答部分先写流程与缺点，再针对客户数据说明条件选择。四组已由两轴业务定义且近似紧凑、强制人人归类时，设 $K=4$ 的 $K$ 均值便于解释。如果题目重点是识别噪声和密度簇，DBSCAN 可留下噪声。说明先标准化可比特征，按邻域密度判核心、边界、噪声。作选择时写清依据，不能仅凭「客户」断言数据形状。
\end{description}

\section{本章自测}
\begin{enumerate}
\item 三类标签的样本各自都被完全拆成单点簇。纯度、同质性如何？完整性一定如何？\\
\textbf{答} 纯度与同质性均为 1。若某个类别含多个样本，给定这一类仍有簇不确定性，完整性小于 1。所有类别都只有一个样本时则为 1。
\item 旧簇 $A$ 有 2 个点，$B$ 有 5 个点，到同一外簇的平均链距离为 4 和 11。合并后是多少？\\
\textbf{答} $(2\cdot4+5\cdot11)/7=9$。
\item 七个样本的类—簇列联表为 $\bigl(\begin{smallmatrix}2&2\\1&2\end{smallmatrix}\bigr)$。求 $TN$ 与 $F_1$。\\
\textbf{答} 同类和同簇各有 9 对，同类同簇有 3 对。$FP=FN=6$，$TN=21-3-6-6=6$，$F_1=6/18=1/3$。
\item 两个分量对某样本给出相等密度，权重分别为 $1/4,3/4$。两责任度各多少？\\
\textbf{答} 贝叶斯公式中相等的密度可约去，依次为 $1/4,3/4$。
\item 固定数据与半径，把 $\mathrm{MinPts}$ 从 4 增到 6，原噪声点会变成核心吗？\\
\textbf{答} 不会。任一点的邻居计数不变，成为核心的条件更难。原噪声点本已不满足计数 4。
\item 在 3 维数据上比较 2 成分、各自全协方差 GMM。计算其自由参数数。\\
\textbf{答} 1 个自由权重、6 个均值坐标、12 个独立协方差元素，共 19 个。
\item 某轮 E 步的 $q$ 等于旧后验。M 步只提高 $Q$，没有找到最大值。观察似然为何仍可不减？\\
\textbf{答} 固定 $q$ 时熵不变，增大 $Q$ 就增大下界。旧参数处下界等于观察似然，新参数处观察似然至少为下界。此为广义 EM 的基本情形。
\end{enumerate}

\chapter{PCA 与 Fisher 判别降维}
\section{本章考点清单}
\begin{itemize}[leftmargin=2em]
\item \(\bigstar \bigstar \bigstar\) 主成分分析的中心化、协方差、特征分解、得分与重构。最大投影方差和最小重构误差的等价推导。【刘扬·9PCA p5--12】【刘远超·4 聚类问题 p19】【2019秋真题 第6题】【模拟卷一 填空5】【模拟卷二 单选7】【模拟卷三 大题2(1)--(3)(7)】【模拟卷四 判断13、大题4(3)】
\item \(\bigstar \bigstar \bigstar\) PCA 方向和解释方差，符号、重根与丢弃特征值。【刘扬·9PCA p6--10】【模拟卷一 填空5、判断14】【模拟卷二 单选7】【模拟卷四 判断13】主子空间唯一性的明确条件【补充】。
\item \(\bigstar \bigstar \bigstar\) Fisher 判别分析（FDA/LDA）的类内与类间散布、二类方向、PCA 与 LDA 对照。【刘扬·9PCA p13】【2022春真题 二(2)】【2023真题 第6题】【模拟卷三 判断12】散布矩阵与广义 Rayleigh 商的完整推导【补充】。
\item \(\bigstar \bigstar\) Gram 矩阵的非零谱与高维低样本量计算。【刘扬·9PCA p14--15】【2021深度真题 第3题】任意矩阵奇异值角度与零谱计数【补充】。
\item \(\bigstar\) 概率 PCA、因子分析、独立成分分析及线性自编码器的联系。【刘扬·9PCA p16--45】【模拟卷四 大题4(1)--(3)】本章只用线性情形建立降维对应，生成式推断另见自编码器章。
\end{itemize}

\section{核心概念与定义}
\begin{description}[leftmargin=2em,style=nextline]
\item[降维与中心化] 降维（dimensionality reduction）使用较短坐标表示数据，以便压缩、可视化、去冗余或下游学习。中心化矩阵 \(X\in\mathbb R^{n\times d}\) 的第 \(i\) 行是 \(\bm z_i^\top=(\bm x_i-\bar{\bm x})^\top\)，\(\bar{\bm x}=n^{-1}\sum_i\bm x_i\)，本章协方差 \(S=X^\top X/n\)。若题设用 \(1/(n-1)\)，其特征值按比例改，方向及解释方差比例不改，平均重构损失仍按题设分母算。
\item[主成分] 主成分分析（principal component analysis, PCA）选正交方向，使保留坐标的总方差最大，等价于使平均平方正交重构误差最小。\(U=[\bm u_1,\ldots,\bm u_m]\) 的列为高特征值对应的单位特征向量。\(\bm t_i=U^\top(\bm x_i-\bar{\bm x})\) 是得分，\(\hat{\bm x}_i=\bar{\bm x}+U\bm t_i\) 是重构。刘扬课件的最大方差记号 \(u_1^\top S u_1\) 对应此处单列 \(U\)，最小误差用被舍弃方向的残差展开。
\item[Fisher 判别] Fisher 判别分析（Fisher discriminant analysis, FDA，常称线性判别分析 linear discriminant analysis, LDA）利用标签，让投影的类均值相对分开，同时压低每类在该方向的散布。求方向后仍需另定分类阈值，可由题设中点、类内投影模型或贝叶斯决策决定。这里的 Fisher 方向问题与生成模型中“共享协方差的高斯判别”有关，但论证和假设不能互换。
\item[随机潜变量] 概率 PCA（probabilistic PCA, PPCA）设隐变量服从标准正态，观测为线性映射加各向同性噪声。因子分析（factor analysis, FA）可给各观测维不同的噪声方差。独立成分分析（independent component analysis, ICA）强调恢复相互独立、通常非高斯的成分。PCA 的去相关只约束二阶矩，不一般保证独立。
\end{description}

\section{公式与完整推导}
\subsection*{最大投影方差}
对单位向量 \(\bm u\)，一维得分 \(t_i=\bm u^\top\bm z_i\) 的均值为零，方差
\[
 \frac1n\sum_i t_i^2=\bm u^\top
 \left(\frac1n\sum_i\bm z_i\bm z_i^\top\right)\bm u
 =\bm u^\top S\bm u .
\]
构造 \(\mathcal L(\bm u,\lambda)=\bm u^\top S\bm u-\lambda(\bm u^\top\bm u-1)\)。
对 \(\bm u\) 求导得到 \(2S\bm u-2\lambda\bm u=0\)，故 \(S\bm u=\lambda\bm u\)。再证明最大性：取正交单位特征基 \(\{\bm e_j\}\)，对应 \(\lambda_1\ge\cdots\ge\lambda_d\)，写 \(\bm u=\sum_j a_j\bm e_j\)，则 \(\sum_j a_j^2=1\) 且 \(\bm u^\top S\bm u=\sum_j\lambda_j a_j^2\le\lambda_1\)。取 \(\bm e_1\) 达上界。

对 \(m\) 个单位正交方向，得分协方差为 \(U^\top S U\)，总体方差为它的迹。把单位谱基代入，令 \(\omega_j=\|U^\top\bm e_j\|^2\in[0,1]\)，\(\sum_j\omega_j=m\)（因为等于投影矩阵的迹）。于是
\[
 \operatorname{tr}(U^\top S U)=\sum_j\lambda_j\omega_j .
\]
最大方差与上式的差
\[
 \sum_{j\le m}\lambda_j(1-\omega_j)-\sum_{j>m}\lambda_j\omega_j
 \ \ge\ \lambda_m\left(\sum_{j\le m}(1-\omega_j)-\sum_{j>m}\omega_j\right)=0 .
\]
取前 \(m\) 个特征向量可取等号。若 \(\lambda_m>\lambda_{m+1}\)，最优主子空间唯一。特征向量符号本身任意，重根的特征子空间内可换正交基。

\subsection*{最小平方重构误差}
对给定 \(U^\top U=I\)，若用坐标 \(\bm t\) 近似中心化样本 \(\bm z\)，展开
\[
 \|\bm z-U\bm t\|^2=\|\bm z\|^2-2\bm t^\top U^\top\bm z+\bm t^\top\bm t .
\]
求导得 \(2(\bm t-U^\top\bm z)\)，所以最优 \(\bm t=U^\top\bm z\)。令 \(P=UU^\top\)，\(P=P^\top=P^2\)，故
\[
 \frac1n\sum_i\|(I-P)\bm z_i\|^2
 =\operatorname{tr}\big((I-P)S\big)
 =\operatorname{tr}(S)-\operatorname{tr}(U^\top S U).
\]
固定第一项，误差最小相当于第二项最大，所以最小误差为 \(\sum_{j>m}\lambda_j\)。总方差非零时，保留比例为
\[
\frac{\sum_{j\le m}\lambda_j}{\sum_j\lambda_j}.
\]

中心化之所以重要，是因为若平移基准由 \(\bar{\bm x}\) 改为 \(\bm b\)，额外的误差为
\[
\|(I-P)(\bar{\bm x}-\bm b)\|^2 .
\]
交叉项因中心化残差之和为零而消失，因此基准点可取样本均值。

\subsection*{高维低样本量时用 Gram 矩阵}
令 \(S=X^\top X/n,\ G=XX^\top/n\)。若 \(S\bm u=\lambda\bm u,\lambda>0\)，则 \(X\bm u\ne0\)，并且 \(G(X\bm u)=X(S\bm u)=\lambda X\bm u\)。反之若 \(G\bm v=\lambda\bm v,\lambda>0\)，则 \(X^\top\bm v\ne0\)，\(S(X^\top\bm v)=X^\top(G\bm v)=\lambda X^\top\bm v\)。因此正特征值及其重数相同，零特征值数量随矩阵尺寸改变。若 \(\|\bm v\|=1\)，则 \(\|X^\top\bm v\|^2=\bm v^\top XX^\top\bm v=n\lambda\)，故原空间单位方向为 \(X^\top\bm v/\sqrt{n\lambda}\)。中心化使 \(X^\top\bm 1=0\)，非零特征值个数不超过 \(\min(d,n-1)\)。若 Gram 定义为不除 \(n\) 的 \(XX^\top\)，其非零谱是 \(S\) 的 \(n\) 倍。

\subsection*{Fisher 方向与相似图的答题说明}
类别 \(k\) 的样本均值 \(\bm\mu_k\)、数量 \(n_k\)，总体均值 \(\bm\mu\)。定义
\[
 \begin{aligned}
 S_W&=\sum_k\sum_{i:y_i=k}(\bm x_i-\bm\mu_k)(\bm x_i-\bm\mu_k)^\top,\\
 S_B&=\sum_k n_k(\bm\mu_k-\bm\mu)(\bm\mu_k-\bm\mu)^\top.
 \end{aligned}
\]
在每类内，\(\bm x_i-\bm\mu=(\bm x_i-\bm\mu_k)+(\bm\mu_k-\bm\mu)\)。各类内第一项之和为零，故总体散布 \(S_T=S_W+S_B\)。投影到 \(\bm u\) 后，类内、类间散布分别是 \(\bm u^\top S_W\bm u,\ \bm u^\top S_B\bm u\)。
 
二类时沿投影的均值差为 \(\bm u^\top\bm d\)，\(\bm d=\bm\mu_2-\bm\mu_1\)，而 \(S_B\) 和 \(\bm d\bm d^\top\) 只差正比例系数。最大化
\[
 \frac{(\bm u^\top\bm d)^2}{\bm u^\top S_W\bm u}.
\]
若 \(S_W\) 正定，令 \(\bm v=S_W^{1/2}\bm u\)、\(\bm g=S_W^{-1/2}\bm d\)，目标变成 \((\bm v^\top\bm g)^2/\|\bm v\|^2\)。由柯西--施瓦茨，上限 \(\|\bm g\|^2\) 在 \(\bm v\parallel\bm g\) 取得，故 \(\bm u\propto S_W^{-1}\bm d\)。多类时在 \(\bm u^\top S_W\bm u=1\) 约束写拉格朗日式，求导得 \(S_B\bm u=\lambda S_W\bm u\)。\(S_B\) 由 \(K\) 个加权和为零的中心差构成，秩至多 \(K-1\)，因此具有非零 Fisher 判别值的可用维数至多 \(K-1\)。\(S_W\) 奇异时，由题设决定加小的正则或在可逆子空间求解。

从相似图看，PCA 的总体散布等价于全体点对的差：
\[
 \sum_i(\bm x_i-\bm\mu)(\bm x_i-\bm\mu)^\top
 =\frac1{2n}\sum_{i,j}(\bm x_i-\bm x_j)(\bm x_i-\bm x_j)^\top .
\]
右侧展开，有两个自乘和各给 \(n\sum_i\bm x_i\bm x_i^\top\)，两个交叉和各给 \((\sum_i\bm x_i)(\sum_j\bm x_j)^\top\)。相除恰得左式。将同一恒等式逐类别应用得到 \(S_W\)，每类点对权重为 \(1/(2n_k)\)。PCA 对全体点对不看类别，最大化其投影散布。FDA 用标签构造同类点对，分母衡量同类散布，分子用各类中心差。这样回答不依赖记忆某张图的特定符号。【2022春真题 二(2)】【补充】

\subsection*{概率 PCA 与线性自编码器}
设 \(\bm x=\bm\mu+W\bm z+\bm\epsilon\)，\(\bm z\sim\mathcal N(0,I)\)、\(\bm\epsilon\sim\mathcal N(0,\sigma^2I)\) 独立。线性组合仍为高斯，取均值和协方差得 \(p(\bm x)=\mathcal N(\bm\mu,WW^\top+\sigma^2I)\)。将 \(\|(\bm x-\bm\mu)-W\bm z\|^2/\sigma^2+\|\bm z\|^2\) 按 \(\bm z\) 展开，二次项系数为 \(I+W^\top W/\sigma^2\)，一次项是 \(-2\bm z^\top W^\top(\bm x-\bm\mu)/\sigma^2\)。配方后后验均值为 \((W^\top W+\sigma^2I)^{-1}W^\top(\bm x-\bm\mu)\)，协方差为 \((I+W^\top W/\sigma^2)^{-1}\)。主方向已固定时，可直接对方向上的观测方差求似然。线性自编码器在平方重构、\(m\) 维瓶颈及训练达到全局最优的条件下可学习 PCA 的主子空间，隐坐标可作可逆基变换，未必逐个对齐主轴。【刘扬·9PCA p16--30、p44--45】

\subsection*{概率 PCA 的后验、极大似然和 EM}
令 $z\sim\mathcal N(0,I_m)$，$x=Wz+\mu+\epsilon$，独立噪声 $\epsilon\sim\mathcal N(0,\sigma^2I_d)$。条件分布为 $x\mid z\sim\mathcal N(Wz+\mu,\sigma^2I_d)$。由线性高斯封闭性与 $\operatorname{Cov}(Wz)=WW^\top$，
$p(x)=\mathcal N(\mu,C)$，$C=WW^\top+\sigma^2I$。
令 $r=x-\mu$，将 $\ln p(z\mid x)=\ln p(x\mid z)+\ln p(z)+C_0$ 展开。含 $z$ 的项是
\[
-\frac{1}{2\sigma^2}
\bigl[z^\top(W^\top W+\sigma^2I)z-2z^\top W^\top r\bigr].
\]
记 $M=W^\top W+\sigma^2 I$。配方得到均值 $M^{-1}W^\top r$、协方差 $\sigma^2M^{-1}$。协方差与观测值无关。若 $W$ 的列空间固定且满秩，$\sigma^2$ 趋零时 $W\mathbb E[z\mid x]$ 趋向原空间的正交投影，潜坐标仍受列尺度与旋转影响。【刘扬·9PCA p16–23、p29】

对 $n$ 个样本最大化边际似然，均值最优为 $\bar x$，因为固定 $C$ 时导数为 $\sum_i C^{-1}(x_i-\mu)$。令 $S=\sum_i r_ir_i^\top/n$。忽略常数，将负对数似然除以 $n/2$ 得 $\log|C|+\operatorname{tr}(C^{-1}S)$。记样本特征值从大到小为 $\lambda_j$。对固定的模型特征值，迹项最小可令协方差特征向量对齐，并把最大的模型协方差特征值配给最大样本特征值：两组谱基之间旋转矩阵元素的平方形成双随机权重，迹为 $\lambda_j/c_i$ 的加权和，最小值可在某个配对排列达到。交换反序两对会降低或保持迹，因为 $(\lambda_i-\lambda_j)(1/c_i-1/c_j)\le0$ 当两组谱同向排序。对前 $m$ 个模型谱值使用 $\partial(\ln c+\lambda/c)/\partial c=(c-\lambda)/c^2$，最优为 $c_j=\lambda_j$。其余谱值相同，令
$\partial[(d-m)\ln\sigma^2+\sum_{j>m}\lambda_j/\sigma^2]/\partial\sigma^2=0$，得到
\[
\widehat{\sigma}^2=\frac{\sum_{j>m}\lambda_j}{d-m},\qquad
\widehat W=U_m(\Lambda_m-\widehat{\sigma}^2I_m)^{1/2}R,
\]
其中 $m<d$，$R$ 为任意正交矩阵。特征值等于噪声方差时可没有对应活跃维度。若残余方差为零，应按零噪声极限理解。这说明主子空间可由 PPCA 恢复，$W$ 本身还有旋转不确定性。【刘扬·9PCA p24–29】

用 EM 参数估计时，E 步用旧 $W,\sigma^2$ 计算每个样本的 $m_i=\mathbb E[z_i]$ 和 $T_i=\mathbb E[z_i z_i^\top]=\sigma^2 M^{-1}+m_im_i^\top$。重构平方误差的期望为
$r_i^\top r_i-2r_i^\top Wm_i+\operatorname{tr}(W^\top WT_i)$。
对 $W$ 求导并令零，$W'(\sum_i T_i)=\sum_i r_im_i^\top$。固定新 $W'$，对方差求导得到 $\sigma^{2\prime}$ 为该期望平方误差总和除以 $nd$。贝叶斯 PCA 可给 $W$ 的各列不同的先验精度，精度较大的列在估计中被更强地收缩。因子分析把独立噪声扩展到对角 $\Psi$，边际协方差 $WW^\top+\Psi$。PCA 会偏向原空间的大方差方向，FA 可以把某个独立的大噪声方差分到 $\Psi$。两者若给 $W$ 乘潜空间正交矩阵，$WW^\top$ 不变。混合因子分析再引入离散成分标签，在每个簇内学低维表示。ICA 则以非高斯独立成分的分布假设解除只靠二阶统计无法识别的部分。【刘扬·9PCA p28–43】

\section{计算题型与解题模板}

\paragraph{补充例题：PPCA 投影与后验均值。}
取 $\mu=(0,0)^\top$、$W=(2,0)^\top$、$\sigma^2=1$。则边际协方差 $C=\operatorname{diag}(5,1)$。对样本 $(3,1)^\top$，$M=2^2+1=5$，$W^\top(x-\mu)=6$，所以 $z\mid x\sim\mathcal N(6/5,1/5)$。代回得后验重构均值 $(12/5,0)^\top$。几何上直接投到第一主轴为 $(3,0)^\top$，有限噪声时存在收缩。
\paragraph{例 1 二维 PCA 全流程。}
\textbf{步骤} 求均值与中心化矩阵，写协方差并求降序特征对，规定方向符号，按内积投影，最后加均值重构。给定 \((5,-4),(-1,2),(3,0),(1,-2)\)，选第一主轴第一分量为正。
\textbf{解} 均值为 \((2,-1)\)，中心化点是 \(\pm(3,-3)\) 和 \(\pm(1,1)\)。协方差
\[
 S=\tfrac12\begin{pmatrix}9&-9\\-9&9\end{pmatrix}
   +\tfrac12\begin{pmatrix}1&1\\1&1\end{pmatrix}
  =\begin{pmatrix}5&-4\\-4&5\end{pmatrix}.
\]
特征方程为 \(\det(S-\lambda I)=(5-\lambda)^2-16=0\)，根为 \(9\) 和 \(1\)。主轴分别为 \((1,-1)/\sqrt2\) 和 \((1,1)/\sqrt2\)。四个一维得分依次为 \(3\sqrt2\)、\(-3\sqrt2\)、\(0\)、\(0\)。保留方差 \(9/10\)，训练平均重构误差为 \(1\)。新点 \((5,0)\) 中心化为 \((3,1)\)，得分 \(\sqrt2\)。在均值 \((2,-1)\) 上加回 \((1,-1)\)，得到重构 \((3,-2)\)。

\paragraph{例 2 Gram 快速求原空间方向。}
\textbf{步骤} 写低维 Gram 谱，对正特征值的单位向量乘 \(X^\top\) 并用 \(\sqrt{n\lambda}\) 归一。两个三维中心化点为 \((2,-1,2)\)、\((-2,1,-2)\)。
\textbf{解} \(G=\frac12\begin{pmatrix}9&-9\\-9&9\end{pmatrix}\)，最大特征值 \(9\)，取 \(\bm v=(1,-1)^\top/\sqrt2\)。映射并归一得 \(\bm u=(2,-1,2)^\top/3\)。原空间 \(S\) 还另有两个零根。训练点均在这一轴线上，保留一个主成分可零误差重构。

\paragraph{例 3 Fisher 用标签求方向。}
\textbf{步骤} 逐类求均值与散布，合并 \(S_W\)，解 \(S_W\bm u=\bm\mu_2-\bm\mu_1\)，按题设确定阈值。两类的中心分别 \((-2,0),(2,1)\)，每类恰有相对自身中心的四个偏移 \(\pm(1,1),\pm(0,2)\)。规定用两均值投影的中点为阈值。
\textbf{解} 各类中心即给定值，合计
\[
 S_W=4\begin{pmatrix}1&1\\1&5\end{pmatrix},\qquad
 \bm\mu_2-\bm\mu_1=(4,1)^\top.
\]
解得 \(\bm u=(19/16,-3/16)^\top\)，等比例方向为 \((19,-3)\)。以这个方向投影时，两类中心得分 \(-38,35\)，阈值为 \(-3/2\)。新点 \((1,2)\) 得分 \(13\)，按给定阈值判类 2。得分方向乘正数时阈值也须同倍缩放。

\paragraph{例 4 固定主轴下的概率 PCA 方差。}
\textbf{步骤} 在样本协方差特征基里写生成协方差，分离两个标量似然并检查 \(a\ge b>0\)。取例 1 的 \(S\)，\(W=\rho\bm u_1\)，固定 \(\bar{\bm x}\) 为均值。
\textbf{解} 对应生成协方差特征值 \(a=\rho^2+\sigma^2,b=\sigma^2\)。高斯负对数似然去掉常数及正因子后最小化 \(\ln a+9/a+\ln b+1/b\)。前两项合并后的导数 \((a-9)/a^2\)，后两项的导数 \((b-1)/b^2\)，分别在 \(a=9,b=1\) 由负转正。两处全局最小同时满足协方差约束，故 \(\widehat{\rho}^{\,2}=8,\widehat{\sigma}^{\,2}=1\)。

\section{易混辨析}
\noindent\begin{tabularx}{\linewidth}{@{}p{.21\linewidth}XX@{}}
\toprule
概念 & PCA & LDA／FDA\\
\midrule
学习信号 & 不用类标，描述总体方差 & 用类标，对比类内与类间散布\\
典型方程 & \(S\bm u=\lambda\bm u\) & \(S_B\bm u=\lambda S_W\bm u\)\\
维数上限 & 按协方差非零秩，选保留维数 & 非零判别方向至多 \(K-1\)\\
输出目的 & 压缩、重构、可视化和前处理 & 促进分类。阈值需另由题设或概率模型给定\\
共同点 & 可把原向量线性投影到低维 & 投影方向是否正交以及中心化方式由各自约束决定\\
\bottomrule
\end{tabularx}

\section{常见错误与陷阱}
\begin{itemize}[leftmargin=2em]
\item 求协方差前忘了减均值，或用全数据均值帮测试点投影。新样本应使用训练阶段的均值与方向。
\item 求了特征值忘了归一化特征向量，或者把方差最大误写成选最小特征值。重构误差用舍弃特征值，解释方差用保留特征值。
\item 主轴换符号必须同时让投影得分换符号，\(UU^\top\) 和原尺度重构不变。特征值发生重合时应区分基不唯一和子空间是否唯一。
\item Gram 与协方差的尺度需统一，只保证其非零谱按正确比例对应。白化（whitening）用正特征值开方缩放坐标使协方差为单位阵，一般不由此获得独立性。
\item PCA 不看类标，高方差方向可能恰好是类内变化最大的方向。用 Fisher 商时分子是类间、分母是类内，不能直接最大化类间散布而忽略类内抵消。
\end{itemize}

\section{往年真题与参考答案}
\begin{enumerate}[leftmargin=2em]
\item \textbf{2019 秋第 6 题，推导 PCA 及应用。} 先中心化，再把单位方向的投影方差写为 \(\bm u^\top S\bm u\)。约束求极值或按谱分解得到最大特征方向，更高维按前 \(m\) 个方向组合。若从误差推导，先对给定平面求正交投影，再写等于总方差减保留方差。应用可各写一例，如图像低秩压缩和高维数据可视化。
\item \textbf{2021 深度第 3 题，内积矩阵与协方差矩阵。} 先解释矩阵尺寸，一个比较样本对，一个描述特征之间的中心化协变。中心化并统一 \(1/n\) 时，左乘 \(X\) 或 \(X^\top\) 在正特征空间相互映射，非零特征值及重数相同。原题问“特征值是否相同”，零根数量和缩放应分别说明。
\item \textbf{2022 春二(2)，相似图角度的 PCA/FDA。} 写全体点对的散布恒等式说明 PCA 没有引入标签，FDA 中同类点形成组，使用组内分母和类均值间分子，解释两者为什么会选不同低维方向。可配一组沿横向幅度大、沿竖直分两类的散点口头示意。
\item \textbf{2023 第 6 题。} 可辨认的是 PCA 与 FDA 对比，回忆者未确认相似图是否为硬性要求。按“是否有标签、目标函数、求解特征方程、输出维数”展开。若要求图角度，加上一问的同类边与全体点对分析。
\end{enumerate}

\section{本章自测}
\begin{enumerate}[leftmargin=2em]
\item 中心化样本协方差特征值 \(7,2,0\)，保留 1 维，求比例与平均重构误差。\textbf{答} \(7/9,2\)。
\item 若 \(S=\operatorname{diag}(2,6)\)，取第一主轴的第二坐标为正，新中心化点为 \((3,-2)\)。\textbf{答} 主轴 \((0,1)\)，得分 \(-2\)，中心化重构 \((0,-2)\)。
\item 一组中心化观测有 \(n=3,d=12\)，协方差最多有多少正特征值？\textbf{答} 至多 \(2\)，样本行相加为零。
\item \(S_W=\operatorname{diag}(2,8)\)，两类均值差 \((2,4)\)，求 Fisher 方向。\textbf{答} \(S_W^{-1}\bm d=(1,1/2)\)，任意非零同向倍数均可。
\item \(S=\operatorname{diag}(4,4,1)\)，只保留 1 维，最优方向唯一吗？\textbf{答} 两个最大特征值相同，前两坐标形成的平面内任一单位向量都最优。
\item 能否仅靠 PCA 的解释方差比例判断分类一定变好？\textbf{答} 不能。保留的高方差未必是类间差异，需按下游分类指标验证。
\item 若把中心化矩阵 \(X\) 的 Gram 写成 \(XX^\top\)，却把协方差写成 \(X^\top X/5\)，两者正特征值有什么关系？\textbf{答} Gram 的正谱为协方差正谱的 5 倍。
\end{enumerate}

\chapter{隐马尔可夫模型}

\section{本章考点清单}
\begin{itemize}
\item ★★★ 模型 $\lambda=(A,B,\pi)$、齐次一阶状态链和观测条件独立。从箭头写联合概率。区分概率计算、学习、解码。 【刘远超·5 隐马尔科夫模型 p4–12】【刘扬·12HMM p3–16】【2024秋真题 单选4、判断2】【模拟卷二 选择9、大题8(1–3)】
\item ★★★ 前向的定义、初始化、递推、终止和计算量。刘远超课件有盒球手算。刘扬强调状态表示与前向、后向的配合。 【刘远超·5 隐马尔科夫模型 p15–20】【刘扬·12HMM p18–26】【2022春真题 一(2)】【2023真题 第2题】【模拟卷一 大题28】【模拟卷二 大题8(3–8)】【模拟卷三 大题8(1–6)】【模拟卷四 大题8(1–6)】
\item ★★ 后向、滤波与平滑、$\gamma_t$ 与 $\xi_t$，单步预测观测。 【刘扬·12HMM p18–27、p35】【2024秋真题 判断2】【模拟卷三 大题8(7–10)】【模拟卷四 判断16、大题8(7–8)】
\item ★★ Viterbi 的取最大与回溯、Baum–Welch 的三类重估、正则化和数值缩放。 【刘扬·12HMM p28–34、p40–41】【模拟卷二 判断18、大题8(9)】【模拟卷三 大题8(11)】【模拟卷四 大题8(6–9)】【补充：缩放的统一约定】
\item ★ HMM 与 RNN、带输入或更高阶的状态模型、线性动态系统与粒子滤波概念。 【刘扬·12HMM p36–53】【刘远超·8 建模序列数据的神经网络】【2024秋真题 单选4】【模拟卷四 大题8(9)】
\end{itemize}

\section{核心概念与定义}
隐马尔可夫模型（Hidden Markov Model, HMM）描述一条看不见的随机状态链，以及从各时刻状态发出的可见观测。例：设备真实工况不可直接看到，传感读数可看到。刘远超记状态集 $Q=\{q_1,\ldots,q_N\}$、观测集 $V=\{v_1,\ldots,v_M\}$、状态序列 $I=(i_1,\ldots,i_T)$、观测序列 $O=(o_1,\ldots,o_T)$。本章沿用其 $\lambda=(A,B,\pi)$，写 $S_t=i$ 简记 $i_t=q_i$。刘扬以 $\bm z_t$ 的 one-hot 表示同一隐状态，$\bm x_t$ 是观测，$A_{ij}$ 仍为从 $i$ 到 $j$ 的条件概率，发射可写 $p(\bm x_t\mid \bm z_t)$。在连续观测时发射是密度。两套符号中的 $\alpha(\bm z_t)$、$\beta(\bm z_t)$ 对应逐状态的 $\alpha_t(i)$、$\beta_t(i)$。刘扬以 $\xi$ 描述相邻帧时，可能用落在后一个时刻的下标，本章用转移起点 $t$。【刘远超·5 隐马尔科夫模型 p5–8、p17】【刘扬·12HMM p8–12、p20–27】

$\pi_i=P(S_1=i)$，$a_{ij}=P(S_{t+1}=j\mid S_t=i)$，$b_j(o)=P(o_t=o\mid S_t=j)$。离散模型中 $A\in\mathbb R^{N\times N}$、$B\in\mathbb R^{N\times M}$，每行和为 1，$\pi$ 各项和为 1。齐次一阶马尔可夫假设是下一个状态在给定当前状态后，与更早历史条件独立，且 $A$ 在每步相同。观测独立性是假设给定全状态链，观测在时间上条件独立，每个观测只用自己的状态生成。边际上跨时刻观测仍可能相关。若状态本身可直接观察，就是普通的马尔可夫链（Markov chain），可以用相邻状态计数估参数。二阶链需给定最近两个状态。【刘远超·5 隐马尔科夫模型 p4–11】【刘扬·12HMM p3–12】

三个经典问题是观测似然 $P(O\mid\lambda)$、给定观测学习参数、给定观测找整体最佳状态路径。对应工具依次为前向或前后向、Baum–Welch（HMM 的 EM）、Viterbi。刘远超课件把第三类称为「预测问题」，并明说是解码（decoding），即推断已经观察到的序列背后的最可能状态链。未来观测预测则求 $p(o_{t+1}\mid o_{1:t})$，可以单独用滤波与一次转移发射得到。 【刘远超·5 隐马尔科夫模型 p12、p17】【刘扬·12HMM p26–35】【2024秋真题 判断2】

\section{公式与完整推导}

\subsection{联合概率与前向递推}
对任意状态路径 $s_{1:T}$，先从 $\pi$ 抽第一步状态，再在每一步从该状态发射观测，并用 $A$ 转移至下一步。链式法则结合两个条件独立假设给出
\[
p(O,S_{1:T}=s_{1:T}\mid\lambda)
=\pi_{s_1} b_{s_1}(o_1)\prod_{t=1}^{T-1} a_{s_t s_{t+1}} b_{s_{t+1}}(o_{t+1}).
\]
直接求 $P(O\mid\lambda)$ 要对 $N^T$ 条路径加总。把具有相同前缀终点的路径合在一起，定义
$\alpha_t(i)=P(o_{1:t},S_t=i\mid\lambda)$。初值 $\alpha_1(i)=\pi_i b_i(o_1)$。
从前一步全概率展开，有
\[
\begin{aligned}
\alpha_{t+1}(j)
&=\sum_i p(o_{1:t},S_t=i,S_{t+1}=j,o_{t+1}\mid\lambda)\\
&=\sum_i p(o_{1:t},S_t=i\mid\lambda)
   p(S_{t+1}=j\mid S_t=i)\,p(o_{t+1}\mid S_{t+1}=j)\\
&=b_j(o_{t+1})\sum_i \alpha_t(i)a_{ij}.
\end{aligned}
\]
终点状态互斥且穷尽，$P(O\mid\lambda)=\sum_i\alpha_T(i)$。每步用 $N^2$ 个转移项，时间 $O(TN^2)$。只求似然时只要上一步向量，额外内存可为 $O(N)$。行向量实现是 $\bm\alpha_{t+1}=(\bm\alpha_t A)\odot \bm b(o_{t+1})$。列向量实现用 $A^\top\bm\alpha_t$。发射因子属于到达的状态。 【刘远超·5 隐马尔科夫模型 p15–20】【刘扬·12HMM p20–22】

\subsection{后向、滤波、平滑与相邻转移}
令 $\beta_t(i)=P(o_{t+1:T}\mid S_t=i,\lambda)$，空后缀的概率为 1，因此 $\beta_T(i)=1$。从后一时刻展开
\[
\begin{aligned}
\beta_t(i)
&=\sum_j P(S_{t+1}=j,o_{t+1},o_{t+2:T}\mid S_t=i,\lambda)\\
&=\sum_j a_{ij}b_j(o_{t+1})\beta_{t+1}(j).
\end{aligned}
\]
给定 $S_t$，前缀与后缀条件独立，故在任意 $t$，
$P(O,S_t=i)=\alpha_t(i)\beta_t(i)$，
$P(O)=\sum_i\alpha_t(i)\beta_t(i)$。由此得到平滑概率（smoothing）
$\gamma_t(i)=\alpha_t(i)\beta_t(i)/P(O)$。
滤波（filtering）只拥有前缀，$P(S_t=i\mid o_{1:t})=\alpha_t(i)/\sum_j\alpha_t(j)$。
把前缀、一次转移和一次发射、后缀拼接得
\[
\xi_t(i,j)=P(S_t=i,S_{t+1}=j\mid O)=
\frac{\alpha_t(i)a_{ij} b_j(o_{t+1})\beta_{t+1}(j)}
{\sum_{r,s}\alpha_t(r)a_{rs}b_s(o_{t+1})\beta_{t+1}(s)} .
\]
分母等于整条观测似然。$\sum_j\xi_t(i,j)=\gamma_t(i)$，$\sum_i\xi_t(i,j)=\gamma_{t+1}(j)$。若已知前缀，下一观测 $v$ 的预测概率是 $\sum_{i,j}P(S_t=i\mid o_{1:t})a_{ij}b_j(v)$。过滤、平滑和未来预测用了不同的信息集合，答题时写清条件。【刘扬·12HMM p18–27、p35】

\subsection{维特比算法与回溯}
维特比算法（Viterbi algorithm）把「前缀中所有路径的总概率」换成「到达当前状态的最大一条路径的联合概率」。定义
$\delta_t(j)=\max_{s_{1:t-1}}p(o_{1:t},s_{1:t-1},S_t=j\mid\lambda)$。
初始化 $\delta_1(j)=\pi_jb_j(o_1)$。以当前步终点为 $j$，末前一状态若为 $i$，后续转移发射乘子与更早选择无关。先在所有落在 $i$ 的前缀中取最大，再在所有 $i$ 中取最大，得到
\[
\begin{aligned}
\delta_{t+1}(j)&= b_j(o_{t+1})\max_i\{\delta_t(i)a_{ij}\},\\
\psi_{t+1}(j)&=\arg\max_i\{\delta_t(i)a_{ij}\}.
\end{aligned}
\]
从 $s_T^*=\arg\max_j\delta_T(j)$ 开始回溯 $s_t^*=\psi_{t+1}(s_{t+1}^*)$。全观测固定，最大化联合概率与最大化除以同一个 $P(O)$ 的路径后验等价。若题目问路径后验值，则再除以 $P(O)$。逐时刻选 $\gamma_t(i)$ 最大的状态优化的是逐点正确数的期望，可能拼出不高概率甚至不合法的路径。【刘扬·12HMM p33–34】【模拟卷二 大题8(9)】

\subsection{Baum–Welch 的三项重估}
完全数据对数似然由初始、转移和发射的对数相加。令 $\mathbf1_{\{S_t=i\}}$ 是指示函数，$\mathbb E[\mathbf1_{\{S_t=i\}}\mid O]=\gamma_t(i)$，相邻状态指示乘积的期望是 $\xi_t(i,j)$。将指示量换成这些后验期望，得
\[
\begin{aligned}
Q={}&\sum_i\gamma_1(i)\ln\pi_i
 +\sum_{i,j}\left(\sum_{t=1}^{T-1}\xi_t(i,j)\right)\ln a_{ij}\\
&+\sum_{i,k}\left(\sum_{t=1}^{T}\gamma_t(i)[o_t=v_k]\right)\ln b_i(v_k).
\end{aligned}
\]
对于一行候选概率 $p_j$，子目标 $\sum_j c_j\ln p_j$。施加 $\sum_jp_j=1$，导数条件为 $c_j/p_j=\lambda$。求和得 $\lambda=\sum_jc_j$，故 $p_j=c_j/\sum_jc_j$。三处都用期望计数归一化：
\[
\begin{aligned}
\pi_i'&=\gamma_1(i),\\
a_{ij}'&=\frac{\sum_{t=1}^{T-1}\xi_t(i,j)}{\sum_{t=1}^{T-1}\gamma_t(i)},\\
b_i'(v_k)&=\frac{\sum_{t=1}^{T}\gamma_t(i)[o_t=v_k]}{\sum_{t=1}^{T}\gamma_t(i)}.
\end{aligned}
\]
转移只到 $T-1$，因为最后时刻没有出边。发射包括所有 $T$ 步。多序列时每条先做各自前后向，汇总计数，序列之间不制造虚假转移。连续高斯发射与 GMM 使用同款加权均值和加权协方差，权重改为 $\gamma_t(i)$。零概率路径、零参数和不可达状态需检查是否是结构限制，必要时加先验伪计数。【刘扬·12HMM p28–32、p40–41】

\subsection{逐步缩放与可选扩展}
定义 $c_1=\sum_i\pi_i b_i(o_1)$，$\hat\alpha_1(i)=\pi_i b_i(o_1)/c_1$。以归一的前一步产生
\[
\tilde\alpha_t(j)=b_j(o_t)\sum_i\hat\alpha_{t-1}(i)a_{ij},\qquad
c_t=\sum_j\tilde\alpha_t(j),\qquad
\hat\alpha_t(j)=\tilde\alpha_t(j)/c_t .
\]
归纳假设 $\hat\alpha_{t-1}=\alpha_{t-1}/P(o_{1:t-1})$。代入递推得 $\tilde\alpha_t=\alpha_t/P(o_{1:t-1})$，求和得 $c_t=P(o_t\mid o_{1:t-1})$。再除以 $c_t$ 即 $\hat\alpha_t=\alpha_t/P(o_{1:t})$，因此 $\ln P(O)=\sum_t\ln c_t$。

同一约定下令 $\hat\beta_t(i)=\beta_t(i)/\prod_{k=t+1}^T c_k$。从未缩放递推代入该定义，分母剩下 $c_{t+1}$，得 $\hat\beta_t(i)=\sum_j a_{ij}b_j(o_{t+1})\hat\beta_{t+1}(j)/c_{t+1}$，$\hat\beta_T=1$。乘积 $\hat\alpha_t(i)\hat\beta_t(i)=\gamma_t(i)$。若题目把缩放系数定义为这里的倒数，似然改为负的对数和。若发现 $c_t=0$，当前模型给前缀零似然，需处理不可能观测或初始模型。对数域的前向用 log-sum-exp，维特比用 $\max$ 加对数。【模拟卷四 大题8(1–7、9)】【补充】

若 $A$ 各行相同且初始分布也取该行，隐状态跨时刻独立，模型退化为独立混合模型。自回归 HMM（autoregressive HMM）允许发射在当前状态外依赖过去观测。输入输出 HMM（input-output HMM）让外部输入调节转移或发射。因子 HMM（factorial HMM）用多条潜在链和共同发射，联合状态数可能按链的数目乘积增长。标准 HMM 的隐藏状态离散，常以一阶链记忆。普通 RNN 的隐藏状态是可学习的连续递推向量，可编码历史信息。【刘扬·12HMM p10、p36–41】【2024秋真题 单选4】

刘扬课件还比较线性动态系统（Linear Dynamical System, LDS）与粒子滤波（particle filter）。线性高斯状态 $\bm z_t=F\bm z_{t-1}+\bm w_t$ 发射 $\bm x_t=H\bm z_t+\bm v_t$。若前一步后验为 $\mathcal N(m,P)$ 且状态噪声协方差为 $Q$，线性变换得预测均值 $m^-=Fm$、协方差 $P^-=FPF^\top+Q$。再对联合高斯 $(z_t,x_t)$ 做条件化，交叉协方差 $P^-H^\top$，观测协方差 $HP^-H^\top+R$，得到 Kalman 增益 $K=P^-H^\top(HP^-H^\top+R)^{-1}$ 及更新 $m=m^-+K(x-Hm^-)$、$P=P^--KHP^-$。卡尔曼滤波（Kalman filtering）是这种结构下的精确过滤。非线性或非高斯时，可用粒子推进、乘发射权重、归一化与重采样近似分布。权重集中过高时有效粒子数会下降。【刘扬·12HMM p42–53】

\section{计算题型与解题模板}

\paragraph{共用原创数据。}
记录装置的隐状态 1、2，观测字母 $U,V$，行表示当前状态，列表示下一状态或观测。令
\[
\pi=(2/3,1/3),\quad A=\begin{pmatrix}7/10&3/10\\1/5&4/5\end{pmatrix},\quad
B=\begin{pmatrix}3/4&1/4\\1/3&2/3\end{pmatrix},\quad O=(U,U,V).
\]
每条路径均可枚举验证。以下各表是精确分数，可留分母。

\paragraph{题型一：用前向求似然、归一化与预测。}
先初始乘发射，再做「上一行乘 $A$、到达点乘发射」，最后求和。
第一步 $\alpha_1=(1/2,1/9)$。
第二步的第一分量是 $[(1/2)(7/10)+(1/9)(1/5)](3/4)=67/240$，另一分量为 $43/540$。
第三步的第一分量是 $[(67/240)(7/10)+(43/540)(1/5)](1/4)=913/17280$，第二分量为 $637/6480$。终点和 $P(O)=1567/10368$。
缩放题可从每个前缀和的比值验算，$c=(11/18,155/264,1567/3720)$，三项乘积仍为 $1567/10368$，对数似然是三个自然对数之和。最后滤波向量 $(2739,5096)/7835$。再乘 $A$ 与 $U$ 的发射列，下一读数为 $U$ 的概率为 $18409/37608$。写出起始向量、每步向量、最后一行的和，能分别展示完整步骤。【刘远超·5 隐马尔科夫模型 p17–20】【模拟卷四 大题8(1–8)】

\paragraph{题型二：后向与平滑。}
把 $\beta_3$ 写为 $(1,1)$。$\beta_2(1)=(7/10)(1/4)+(3/10)(2/3)=3/8$，$\beta_2(2)=7/12$。再用第二个观测 $U$ 得 $\beta_1=(49/192,61/288)$。例如时刻 2 的状态 1 平滑概率为 $(67/240)(3/8)/(1567/10368)=5427/7835$，状态 2 为 $2408/7835$。从时刻 2 的状态 1 到时刻 3 的状态 2，分子乘前向、转移、发射、后向：
$(67/240)(3/10)(2/3)(1)$，除以整序列似然后为 $14472/39175$。回算 $\sum_j\xi_2(1,j)=5427/7835$。如果直接把两个边缘后验相乘，会漏掉真实的相邻时间依赖。【刘扬·12HMM p23–27】

\paragraph{题型三：最大路径。}
用同一数据，$\delta_1=(1/2,1/9)$。第二步各状态的最优父节点均为 1，乘发射后 $\delta_2=(21/80,1/20)$。第三步两个父节点仍都是 1，$\delta_3=(147/3200,21/400)$。从较大的第二分量回溯 $2\gets1\gets1$，路径为 $(1,1,2)$，该路径和观测的联合概率为 $21/400$。例如末步到 2 的两项分别为 $(21/80)(3/10)$、$(1/20)(4/5)$，先选前者再乘 $2/3$。误把这一答案写成观测的总概率会丢掉其他路径的权重。【刘扬·12HMM p33–34】

\paragraph{题型四：用前后向统计量做一次重估。}
同一数据下两组相邻后验矩阵与三个平滑向量为
\[
\begin{gathered}
\xi_1=\frac1{7835}\begin{pmatrix}5103&1512\\324&896\end{pmatrix},\quad
\xi_2=\frac1{39175}\begin{pmatrix}12663&14472\\1032&11008\end{pmatrix},\\
\gamma_1=(1323,244)/1567,\quad
\gamma_2=(5427,2408)/7835,\quad
\gamma_3=(2739,5096)/7835 .
\end{gathered}
\]
初值以新后验替代，$\pi'=\gamma_1$。例如 $a_{12}'$ 的分子是两个矩阵的右上项相加，分母是 $\gamma_1(1)+\gamma_2(1)$，结果为 $408/1115$。另一行按同法。发射 $U$ 分子加前两时刻，$V$ 分子只取第 3 时刻。各行分母均为该状态在三步中的后验占比之和。完整结果为
\[
A'=\begin{pmatrix}707/1115&408/1115\\663/4535&3872/4535\end{pmatrix},\qquad
B'=\begin{pmatrix}4014/4927&913/4927\\907/2181&1274/2181\end{pmatrix}.
\]
最后逐行检查概率和为 1。数值可通过每个时间点的 $\alpha_t^\top\beta_t$ 一致性检查。【刘扬·12HMM p28–32】

\paragraph{题型五：连续状态过滤的简单数值。}
原创一维 LDS：上步后验均值 1、方差 2，$F=H=1$，过程及测量噪声方差分别为 1 与 3。新观测为 3。预测均值仍为 1，预测方差 $2+1=3$，增益 $3/(3+3)=1/2$。更新均值 $1+(1/2)(3-1)=2$，方差 $3-(1/2)3=3/2$。只要求 HMM 手算时，应按离散前后向处理。【刘扬·12HMM p44–47】

\section{易混辨析}
\begin{tabularx}{\linewidth}{@{}p{.25\linewidth}XX@{}}
\toprule
比较 & 第一项 & 第二项\\
\midrule
前向／维特比 & 同终点路径概率求和，由末向量求似然 & 同终点保留一条最大路径，存父节点，末状态回溯\\
前向／后向 & 当前状态与含当前观测的前缀联合量 & 给定当前状态，计算当前之后的观测概率\\
滤波／平滑 & 截至当前的观测 & 全序列观测，结合前后两侧证据\\
学习／解码 & 给样本更新 $\pi,A,B$ & 给定参数找状态路径\\
HMM／RNN & 离散隐状态和概率转移，发射有明确定义 & 连续隐状态由可训练函数递推，任务可用分类或生成头\\
\bottomrule
\end{tabularx}

\section{常见错误与陷阱}
\begin{itemize}
\item 转移矩阵的行是出发、列是到达。第一步用 $\pi$，没有从不存在的第 0 步转移。递推的发射应与到达状态和新观测对齐。
\item 后向用 $o_{t+1}$，不是 $o_t$。$\beta_T=1$。归一化前向向量可做滤波，不能拿来替代仍需后缀的平滑。
\item 算维特比时记录每个状态的父节点，再自最后一时刻回溯。直觉上逐步选最大、或写每行最大的 $\delta$，都可能与最终回溯路径不一致。
\item 先固定缩放系数约定，再累计对数。用表的归一化行求和永远是 1，无法直接当 $P(O)$。记住转移 M 步到 $T-1$、发射到 $T$。
\item 确定性预测不能代替概率预测。HMM 既能给过去状态后验，也能给未来观测分布。这两种操作分别用全序列平滑和一次转移发射。
\end{itemize}

\section{往年真题与参考答案}
\begin{description}
\item[2022 春一(2)] 回忆只说课件中「朋友在外地」例题变式，已知转移矩阵第二行改为 $(0.5,0.5)$。其余参数及求解目标未存。评分关键词：先用该新行替换原行，写出 $A$ 的行列定义，再按题目所求分别做前向或 Viterbi。如果求似然，最后对状态求和。缺失的原卷结果不能从修改的一行唯一计算。
\item[2022 春一(3)] 以递归贝叶斯解释 Kalman 滤波。先由旧后验和转移对前一状态积分得到状态预测。新观测到来后乘观测似然再归一化。线性高斯的预测和条件化都给高斯，故递推只需均值和协方差。
\item[2023 第 2 题] 材料只记天气与行动例题。答题点：天气作状态，行动作观测，物理含义分开。用 $\pi,A,B$ 建模，若问序列概率按前向三步。若问最佳天气路径用 Viterbi 回溯。原参数和具体小问缺失。
\item[2023 第 3 题] Kalman 参数意义。$F$ 映射上一状态到下一状态，$H$ 把状态映射到观测空间，$Q$、$R$ 是过程与观测噪声协方差，$P$ 是状态估计误差协方差，$K$ 将创新 $x-Hm^-$ 变成状态修正量。原稿未列符号，按所给课件或试卷符号替换。
\item[2024 秋单选 4] 可辨认内容是 HMM 与 RNN 的区别。记住 HMM 是隐变量概率序列模型，标准 RNN 是用递推神经网络建模序列，二者都能表达序列依赖，未保留的选项不猜字母。
\item[2024 秋判断 2] 「HMM 不能预测序列」判错。用已有观测过滤当前状态，再通过 $A$ 得下一状态分布，乘 $B$ 得下一观测概率。若问过去路径，用条件观测的 Viterbi，而非未来预测公式。
\end{description}

\Needspace{18\baselineskip}
\section{本章自测}
\begin{enumerate}
\item 一个三状态四观测的离散 HMM，独立的同质参数自由度是多少？\\
\textbf{答} $A$ 三行各 2，$B$ 三行各 3，$\pi$ 有 2，自由度 $6+9+2=17$。
\item 当 $T=1$，$\pi=(1/5,4/5)$，发射当前观测概率为 $(1/2,1/4)$。求前向量、似然和第 2 状态平滑概率。\\
\textbf{答} $\alpha_1=(1/10,1/5)$，似然 $3/10$，末步 $\beta_1=(1,1)$，故 $\gamma_1(2)=2/3$。
\item 输入 $c=(1/2,1/4,1/5)$，这里的 $c_t$ 是分量和。求 $P(O)$ 与 $\ln P(O)$。\\
\textbf{答} $P(O)=1/40$，$\ln P(O)=-\ln40$。
\item 对某观测，两步的路径后验为 $AA=.36,AB=.34,BA=0,BB=.30$。逐步边缘最大和整体路径最大各是什么？\\
\textbf{答} 第一步 $A$ 的边缘概率 .70，第二步 $B$ 的边缘概率 .64，逐步给 $AB$。整体最佳为 $AA$。
\item 在 E 步后，期望从状态 1 转到状态 1、2 的次数为 1.2、0.8。更新 $A$ 第 1 行。\\
\textbf{答} 归一后 $(1.2,0.8)/2=(0.6,0.4)$。
\item 若所有状态对当前观测都给相同且非零的发射概率，这一步如何影响过滤更新？\\
\textbf{答} 它给所有状态乘同一常数，归一化后抵消。当步后验保持转移后的预测状态分布。
\item 已知 $\xi_t(2,1)=0.1$、$\xi_t(2,2)=0.3$，两状态模型下求 $\gamma_t(2)$。\\
\textbf{答} 对下一状态求和为 $0.4$。求 $\gamma_{t+1}(2)$ 则要对当前状态求和。
\end{enumerate}

\chapter{前馈神经网络与反向传播}

\section{本章考点清单}

下文的页码是 PDF 阅读器显示的物理页码。模拟卷按卷面原题号定位，三、四卷的“大题 1”与第一卷的大题 21 均指网络手算题。

\begin{itemize}
\item \textbf{\(\bigstar\bigstar\bigstar\) 前馈与参数量。} 从仿射变换写到激活，能确定任意层矩阵形状、权重数和偏置数。知道纯线性层堆叠仍是仿射变换。【刘远超·6 感知机模型 p22--23】【刘扬·10前馈神经网与反向传递算法 p3--15】【模拟卷一 大题21(1)】【模拟卷二 大题1(1--3)】
\item \textbf{\(\bigstar\bigstar\bigstar\) 激活函数。} 写出 Sigmoid、tanh、ReLU 及导数，比较饱和、零中心、负半轴。能写常见 ReLU 变体。【刘远超·7 网络优化 p46--55】【刘扬·10前馈神经网与反向传递算法 p4--7、p35】【2021深度真题 第5题】【2023真题 第4题】【模拟卷二 单选5】【本次范本 单选2】
\item \textbf{\(\bigstar\bigstar\bigstar\) 完整反向传播。} 掌握输出误差、隐藏误差、权重外积、偏置梯度四式。能按同一组旧参数求梯度再更新。【刘远超·6 感知机模型 p24--28】【刘扬·10前馈神经网与反向传递算法 p22--34】【模拟卷一 大题21(2--4)】【模拟卷二 大题1(4--8)】【模拟卷三 大题1(3--9)】
\item \textbf{\(\bigstar\bigstar\bigstar\) 输出层与损失配对。} 线性输出加平方误差，Sigmoid 加二元交叉熵，Softmax 加类别交叉熵。分别推导评分变量的梯度。【刘扬·10前馈神经网与反向传递算法 p10、p16--26】【刘远超·7 网络优化 p23--26】【2024秋真题 大题3】【2025秋补／缓考 简答4(2)】【模拟卷四 大题1(2--6)、大题7(6--8)】
\item \textbf{\(\bigstar\bigstar\) 初始化和计算图。} Xavier（Glorot）/He 初始化的适用条件、全部神经元同初始化的对称性问题、计算图的分支梯度累加。【刘扬·10前馈神经网与反向传递算法 p29--35】【刘扬·11卷积神经网络 p34（训练配置）】【模拟卷二 大题1(8)】【模拟卷三 单选8】【补充：方差传播推导与两种初始化的标准形式】
\end{itemize}

\section{核心概念与定义}

前馈神经网络（feedforward neural network）的计算图从输入到输出依次展开。多层感知机（multilayer perceptron, MLP）是由仿射层和非线性激活层组合的前馈网络。对一条样本采用列向量，$a^{(0)}=x$，
\[
 z^{(l)}=W^{(l)}a^{(l-1)}+b^{(l)},\qquad
 a^{(l)}=\phi_l(z^{(l)}).
\]
第 $l$ 层有 $n_l$ 个单元，故 $W^{(l)}\in\mathbb R^{n_l\times n_{l-1}}$，$b^{(l)},z^{(l)},a^{(l)}\in\mathbb R^{n_l}$。仿射层有 $n_l(n_{l-1}+1)$ 个可学习标量。激活（activation）通常逐分量作用，Softmax 则让各分量在同一分母中耦合。

刘远超课件把第 $l$ 层加权输入、激活、误差写作 $z^l,a^l,\delta^l$，权重 $w^l_{jk}$ 表示从第 $l-1$ 层第 $k$ 个节点到第 $l$ 层第 $j$ 个节点。刘扬课件以 $t$ 编层，用净输入 $z^{(t)}$ 和输出 $h^{(t)}$，分别对应本章的 $z^{(l)}$、$a^{(l)}$。课件里 $a(\cdot)$ 作为激活函数时对应本文 $\phi_l(\cdot)$，上下文不要把函数 $a(\cdot)$ 误读为激活向量。两者的误差都指对该层净输入求导。【刘远超·6 感知机模型 p22--26】【刘扬·10前馈神经网与反向传递算法 p22--28】

\begin{center}
\begin{tabularx}{\linewidth}{@{}lXX@{}}
\toprule
函数 & 定义与导数 & 位置和提醒\\
\midrule
Sigmoid & $\sigma(z)=(1+e^{-z})^{-1}$，$\sigma'(z)=\sigma(z)(1-\sigma(z))$ & 二分类的输出与门控。正负两端饱和，导数至多 $1/4$。\\
tanh & $\tanh z=(e^z-e^{-z})/(e^z+e^{-z})$，导数为 $1-\tanh^2 z$ & 值域 $(-1,1)$，零中心。两端仍饱和。\\
ReLU & $\max(0,z)$。$z>0$ 导数为 $1$，$z<0$ 为 $0$ & 常作隐藏激活。$z=0$ 按约定取次梯度，常取 $0$。\\
Leaky ReLU & $z\ge0$ 时为 $z$，$z<0$ 时为 $\alpha z$，常用固定 $0<\alpha<1$ & 负区间留有梯度。若 $\alpha$ 可训练，属于 PReLU【补充】。\\
ELU & $z\ge 0$ 时为 $z$，否则为 $\alpha(e^z-1)$ & 负区导数 $\alpha e^z$，正区导数 $1$。负侧远端仍饱和。\\
Softmax & $p_i=e^{z_i}/\sum_j e^{z_j}$ & 单标签多分类输出。需全向量归一化。\\
\bottomrule
\end{tabularx}
\end{center}
【刘远超·7 网络优化 p46--52】【刘扬·10前馈神经网与反向传递算法 p7、p10】

回归可用恒等输出和平方误差，二类分类常用一个 Sigmoid。互斥多类用 Softmax。多标签分类则为每个类别单独用 Sigmoid。两个仿射层相接时，$W_2(W_1x+b_1)+b_2=(W_2W_1)x+(W_2b_1+b_2)$，非线性激活才会改变这类复合映射的表达能力。【刘扬·10前馈神经网与反向传递算法 p3、p8--15】

\section{公式与完整推导}

\subsection*{反向传播的四个关系}

反向传播（backpropagation, BP）对已确定的计算图应用链式法则。定义 $\delta^{(l)}=\partial L/\partial z^{(l)}$。若输出激活逐元素且损失写成 $L(a^{(L)})$，则
\[
 \delta^{(L)}_i
 =\sum_j\frac{\partial L}{\partial a^{(L)}_j}
       \frac{\partial a^{(L)}_j}{\partial z^{(L)}_i}
 =\frac{\partial L}{\partial a^{(L)}_i}\phi_L'(z^{(L)}_i).
\]
隐藏层误差也由相邻层连接给出。由 $z_j^{(l+1)}=\sum_i W^{(l+1)}_{ji}\phi_l(z^{(l)}_i)+b_j^{(l+1)}$，
\begin{align*}
 \delta^{(l)}_i
 &=\sum_j \delta^{(l+1)}_j\,
    \frac{\partial z_j^{(l+1)}}{\partial z_i^{(l)}}\\
 &=\sum_j \delta^{(l+1)}_j W_{ji}^{(l+1)}\phi_l'(z_i^{(l)}),\\
 \delta^{(l)}
 &=\bigl((W^{(l+1)})^{\mathsf T}\delta^{(l+1)}\bigr)
       \odot\phi_l'(z^{(l)}).
\end{align*}
$\odot$ 表示逐元素乘积。再用 $z^{(l)}_j=\sum_k W^{(l)}_{jk}a_k^{(l-1)}+b_j^{(l)}$，
\[
 \frac{\partial L}{\partial W^{(l)}_{jk}}
 =\delta^{(l)}_j a_k^{(l-1)},\qquad
 \nabla_{W^{(l)}}L=\delta^{(l)}(a^{(l-1)})^{\mathsf T},\qquad
 \nabla_{b^{(l)}}L=\delta^{(l)}.
\]
输出误差、隐藏误差、权重、偏置分别对应刘远超课件的 BP1--BP4。若有跳连等多条后继路径，同一变量的梯度是各路径贡献之和。若小批量的样本按列堆叠，$A^{(l-1)}\in\mathbb R^{n_{l-1}\times m}$、$D^{(l)}\in\mathbb R^{n_l\times m}$ 各列为该样本的激活或误差，损失取平均时
\[
 \nabla_{W^{(l)}}J=\frac1mD^{(l)}(A^{(l-1)})^{\mathsf T},
 \qquad
 \nabla_{b^{(l)}}J=\frac1mD^{(l)}\mathbf 1_m .
\]
如果损失定义为样本损失之和，去掉 $1/m$。若求 $\frac1m\sum_i(\hat y_i-y_i)^2$ 的梯度而非 $\frac1{2m}\sum_i(\hat y_i-y_i)^2$，还须保留平方求导的系数 $2$。【刘远超·6 感知机模型 p24--28】【刘扬·10前馈神经网与反向传递算法 p22--34】

\subsection*{两个常用输出误差为何是预测减标签}

Sigmoid 配二元交叉熵（binary cross-entropy, BCE）时，
\begin{align*}
 p&=\sigma(z),\quad L=-y\ln p-(1-y)\ln(1-p),\\
 \frac{\partial L}{\partial z}
 &=\left(-\frac{y}{p}+\frac{1-y}{1-p}\right)p(1-p)
 =-y(1-p)+(1-y)p=p-y.
\end{align*}
多类交叉熵的标签为概率向量 $y$，满足 $\sum_i y_i=1$，独热是特例。由 $S=\sum_k e^{z_k}$，
\begin{align*}
 \frac{\partial p_i}{\partial z_j}
 &=\frac{\mathbf 1_{i=j}e^{z_i}S-e^{z_i}e^{z_j}}{S^2}
 =p_i(\mathbf 1_{i=j}-p_j),\\
 \frac{\partial L}{\partial z_j}
 &=\sum_i(-y_i/p_i)p_i(\mathbf 1_{i=j}-p_j)
 =-y_j+p_j\sum_i y_i=p_j-y_j.
\end{align*}
也可把两个依赖得分的项分开求导：
\[
 L=\ln \sum_k e^{z_k}-\sum_i y_i z_i.
\]
实现时设 $m=\max_i z_i$，算 $p_i=e^{z_i-m}/\sum_k e^{z_k-m}$，并用 $\ln S=m+\ln\sum_k e^{z_k-m}$ 防止指数上溢。线性输出配 $\tfrac12\|\hat y-y\|^2$ 也得 $\hat y-y$。Sigmoid 配平方误差时则是 $(p-y)p(1-p)$。【刘远超·7 网络优化 p23--26】【刘扬·10前馈神经网与反向传递算法 p17--18、p23--26】【补充：对数和指数的数值稳定写法】

\subsection*{初始化与信号方差}

设一层有 $n_{\rm in}$ 个输入。若权重独立同分布、均值为零，并近似独立于输入，偏置为零，则
\[
 {\rm Var}(z_j)
 =\sum_{i=1}^{n_{\rm in}} {\rm Var}(W_{ji}a_i)
 \simeq n_{\rm in}{\rm Var}(W_{ji})\,\mathbb E[a_i^2].
\]
对近线性的零中心激活，前向希望系数 $n_{\rm in}{\rm Var}(W)$ 接近 $1$，同理反向希望 $n_{\rm out}{\rm Var}(W)$ 接近 $1$。Xavier 以 $2/(n_{\rm in}+n_{\rm out})$ 折中。若均匀分布范围是 $[-a,a]$，其方差是 $a^2/3$，故 $a=\sqrt{6/(n_{\rm in}+n_{\rm out})}$。对零均值、分布对称的 ReLU 输入，$\mathbb E[\max(0,z)^2]=\mathbb E[z^2]/2$，He 取方差 $2/n_{\rm in}$ 抵消这一因子。用 ReLU 归一化的是二阶矩。ReLU 激活通常均值不为零。相同隐藏单元若参数初值相同且收到相同更新，训练仍保持对称。随机权重用于打破这种对称，偏置仍可初始化为零。【模拟卷三 单选8】【刘扬·11卷积神经网络 p34】【补充：独立与对称假设下的近似推导】

\section{计算题型与解题模板}

\paragraph{型一：完整前向、BP 与同步更新。}
依次列尺寸，算各层 $z,a$ 和损失。从最后一层回推误差。权重用外积，偏置直接取误差。最后用同一组旧梯度更新全部参数。

\textbf{例题。} 三维输入 $x=(2,1,0)^{\mathsf T}$，隐藏层为双单元 tanh，记 $r=\tfrac12\ln3$。取
\[
 W=\begin{pmatrix}0&r&0\\-r&r&1\end{pmatrix},
 \quad b=0,\quad v=(2,1)^{\mathsf T},\quad c=\tfrac14 .
\]
输出 $\hat y=v^{\mathsf T}h+c$，标签 $y=1/4$，损失 $L=\tfrac12(\hat y-y)^2$。用学习率 $1/5$ 更新一次。

\textbf{解。} 先有 $z=(r,-r)^{\mathsf T}$、$h=(1/2,-1/2)^{\mathsf T}$，故 $\hat y=3/4$、$L=1/8$。输出误差为 $1/2$，因而 $\nabla_v L=(1/4,-1/4)^{\mathsf T}$、$\partial L/\partial c=1/2$。隐藏激活导数各为 $3/4$，得到 $\delta=(3/4,3/8)^{\mathsf T}$。于是
\[
 \nabla_W L
 =\begin{pmatrix}3/2&3/4&0\\3/4&3/8&0\end{pmatrix},
 \qquad \nabla_b L=\delta .
\]
一次同步更新后，
\begin{align*}
 W'&=\begin{pmatrix}-3/10&r-3/20&0\\-r-3/20&r-3/40&1\end{pmatrix},\\
 b'&=(-3/20,-3/40)^{\mathsf T},\quad
 v'=(39/20,21/20)^{\mathsf T},\quad c'=3/20.
\end{align*}
输入第三维为零，所以这条样本对隐藏层第三列权重没有梯度，偏置并不受这个零值屏蔽。

\paragraph{型二：Softmax 与分支梯度。}
先算归一化概率，独热标签在正确下标减 $1$。需要参数梯度时再外积。例题给定得分 $(\ln4,0,\ln3)$、标签为第 $3$ 类、前一层输出 $h=(2,-1)^{\mathsf T}$，求概率、损失与输出仿射层梯度。

\textbf{解。} 权重和为 $4+1+3=8$，$p=(1/2,1/8,3/8)^{\mathsf T}$，$L=\ln(8/3)\approx0.9808$，$\delta=(1/2,1/8,-5/8)^{\mathsf T}$。偏置梯度为 $\delta$，权重梯度为
\[
 \delta h^{\mathsf T}
 =\begin{pmatrix}1&-1/2\\1/4&-1/8\\-5/4&5/8\end{pmatrix}.
\]
三类梯度之和为零，可用作自检。

\paragraph{型三：参数计数与初始化。}
逐层统计每个输出有多少连接和偏置。再按激活和该层扇入、扇出选初始化。例题为 $6$ 输入、$5$ 个 tanh 隐单元、$3$ 类输出。参数数为 $5(6+1)+3(5+1)=53$。隐藏层的 Xavier 均匀边界为 $\pm\sqrt{6/11}$。改成 ReLU 隐层时，He 正态的标准差为 $\sqrt{2/6}$。初始化公式求得的是分布参数，不能把每个权重都填成同一标准差。

\paragraph{实验对照。}
实验四把 MLP（多层感知机）和线性分类器放在同一任务上对比，程序分别实现 \texttt{forward}、\texttt{backward} 和参数更新。先在同一批数据上保存前向中间量，再核查两层矩阵形状、每列或每行的样本轴、损失归约和偏置。若代码采用“样本为行”，其权重保存布局可能与讲义的列向量布局互为转置，确认矩阵乘法后再比较公式。【实验4·lab4.pdf p4--9】

\section{易混辨析}

\begin{center}
\begin{tabularx}{\linewidth}{@{}lXX@{}}
\toprule
对象 & 公式或作用 & 常见对照\\
\midrule
前向传播 & 根据输入和参数计算净输入、激活与损失 & 反向传播计算导数，优化器再用导数更新。\\
权重梯度 & $\delta a^{\mathsf T}$，保留输入和输出两个轴 & 偏置梯度为 $\delta$。批量时沿样本轴求和或取平均。\\
Sigmoid 输出 & 单个值表示正类或各个独立标签的概率 & Softmax 的互斥类概率之和为 $1$。\\
死 ReLU & 某单元长期落入负区间，激活导数恒为零 & 个别样本激活为零，不足以判定该单元对所有样本都死亡。\\
Xavier 与 He & 前者近似折中两侧方差，后者补 ReLU 的二阶矩折损 & 都不是“随机数越大越好”，也不保证深网稳定。\\
\bottomrule
\end{tabularx}
\end{center}

\section{常见错误与陷阱}

\begin{itemize}
\item 输出误差是对净输入求导。Sigmoid 配 BCE 已经化成 $p-y$，再乘一次 $\sigma'$ 会重复计入局部导数。
\item 权重外积的行属于本层，列属于前层。把次序倒过来可能在方阵中悄悄通过尺寸检查，故还须核查 $W_{jk}$ 的连接含义。
\item 所有层的梯度用同一组参数计算。先改输出层再传回隐藏层，不是本次梯度下降。
\item Softmax 的雅可比包含非对角的 $-p_i p_j$。单独对各分量套用 Sigmoid 公式，漏掉了分母的耦合依赖。
\item 全零权重可使同层多个单元无法分工。输入到输出极简模型的零初始化可有意义，不能把对称性结论推广到所有模型。
\end{itemize}

\section{往年真题与参考答案}

\textbf{2021 深度学习与模式识别第 5 题（回忆）}问为何 ReLU 常替代 Sigmoid。得分点是：（1）写出两种激活与导数。（2）说明 Sigmoid 两端饱和且梯度上界 $1/4$，深层连乘可能较快衰减。（3）ReLU 正区间导数为 $1$、计算简单。（4）补上负区间梯度为零的代价。回忆未给逐项分值。

\textbf{2022 春四（1）}问从浅 MLP 扩到深网的困难与三种改进。得分点是：（1）从连乘说明梯度消失或爆炸，并补充深模型的优化或泛化问题。（2）分别给出合适激活与初始化、归一化、残差捷径等方法。（3）把每种方法和它作用的困难对应起来。题面没有唯一三种方法。

\textbf{2023 第 4 题}要求写 ReLU 及三种变体。可写普通 ReLU、Leaky ReLU、PReLU、ELU 的分段式。Leaky ReLU 和 PReLU 的负区都形如 $\alpha z$，须写明 $\alpha$ 前者固定、后者学习。另可选 ReLU6【补充】，其定义为 $\min\{\max(0,z),6\}$。

\textbf{2024 秋大题 3、2025 秋补／缓考简答 4（2）}的共用得分点是：（1）写 $\sigma(s)=(1+e^{-s})^{-1}$。（2）给出二元交叉熵及正类概率。（3）区分逻辑回归输出概率与线性回归输出数值。单层线性得分配 Sigmoid 的特例与本章输出层相同。2024 秋 PDF 的 Softmax 提示没有编号，不能据此补写选择题第 3 题。2025 秋简答 1 与本章相邻，感知机和逻辑回归的最大间隔辨析见感知机章。

\Needspace{18\baselineskip}
\section{本章自测}

\begin{enumerate}
\item 三层宽度依次为 $4,2,2$，每层仿射都有偏置，求参数总数。
\item $a=(1,-2)$、某层误差 $\delta=(3,1)$，用列向量约定求该层权重和偏置梯度。
\item Sigmoid 的输出为 $1/4$、标签为零。写 BCE 对评分的导数并求损失。
\item Softmax 输出 $(1/6,1/3,1/2)$，真类为第一类，求得分梯度并验证其分量和。
\item 一条梯度沿两个分支回到同一变量，局部贡献分别为 $5$、$-3$。若这个变量前一层的局部导数是 $4$，再往前一步的梯度是多少？
\item 说明为何“隐藏层用了 ReLU，所以梯度永远不会消失”不成立，给出一种能缓解死单元的替代激活。
\end{enumerate}
\textbf{答案。} （1）$2(4+1)+2(2+1)=16$。（2）$W$ 梯度为 $\begin{pmatrix}3&-6\\1&-2\end{pmatrix}$，偏置梯度为 $(3,1)^{\mathsf T}$。（3）$p-y=1/4$，$L=\ln(4/3)$。（4）$(-5/6,1/3,1/2)$，和为零。（5）先相加，再乘局部导数，$4(5-3)=8$。（6）负区导数为零，且梯度还受权重雅可比连乘影响。可选负区有斜率的 Leaky ReLU。

\chapter{网络优化与正则化}

\section{本章考点清单}

\begin{itemize}
\item \textbf{\(\bigstar\bigstar\bigstar\) 批量、随机、小批量梯度和学习率。} 区分整个训练集、单样本、少量样本，能算一次或两次更新，说明步长过大或过小的表现。【刘远超·7 网络优化 p29--36】【刘扬·10前馈神经网与反向传递算法 p19--21】【2024秋真题 PDF p1 无编号判断示例】【模拟卷一 判断13】【模拟卷三 判断14】【模拟卷四 大题1(7、10)】
\item \textbf{\(\bigstar\bigstar\bigstar\) 动量、自适应学习率与 Adam。} 会写 Momentum、NAG、AdaGrad、RMSProp、Adam 及偏差修正公式，能对比各自累计对象。【刘远超·7 网络优化 p37--43】【刘扬·11卷积神经网络 p34】【2021深度真题 第6题】【2022春真题 四(2)】【2023真题 第5题（主题）】【模拟卷三 大题1(7--9)】【模拟卷四 单选6】【补充：完整 Adam 与 AdamW 更新】
\item \textbf{\(\bigstar\bigstar\bigstar\) 梯度消失、爆炸和裁剪。} 从雅可比连乘说明现象，能算全局范数裁剪，区分计算稳定和泛化策略。【刘远超·7 网络优化 p58--59】【刘扬·10前馈神经网与反向传递算法 p35】【2022春真题 四(1)】【模拟卷三 判断15、大题1(10)】
\item \textbf{\(\bigstar\bigstar\bigstar\) BatchNorm 与 LayerNorm。} 能按轴手算均值、方差和标准化，区分训练、推理统计量来源。【刘远超·7 网络优化 p59】【刘远超·8 建模序列数据的神经网络 p63--65】【刘扬·10前馈神经网与反向传递算法 p36】【模拟卷一 填空7】【模拟卷四 大题1(8--9)】
\item \textbf{\(\bigstar\bigstar\bigstar\) 正则化与泛化。} 掌握 L1、L2、早停、数据增强、Dropout，能说明权重衰减和 L2 在自适应优化下的差别。【刘远超·7 网络优化 p20--22、p62--65】【刘扬·10前馈神经网与反向传递算法 p37--38】【2021秋真题 第1题】【2022秋真题 第4题】【本次范本 单选1】【模拟卷一 判断12】【模拟卷二 判断11、14】【模拟卷三 判断13、大题7(10)】
\end{itemize}

\section{核心概念与定义}

经验风险（empirical risk）常取 $J_{\rm data}(\theta)=N^{-1}\sum_{i=1}^N L_i(\theta)$。结构风险最小化（structural risk minimization）的常见训练目标加上复杂度惩罚 $J(\theta)=J_{\rm data}(\theta)+\lambda R(\theta)$。损失是单个预测的代价。模型选择看验证指标，最终泛化报告使用独立测试集。调学习率改变优化路径，调正则化主要改变训练目标或模型受到的约束。【刘远超·7 网络优化 p20--26、p62--65】

批量梯度下降（batch gradient descent, BGD）用所有样本计算精确经验梯度。随机梯度下降（stochastic gradient descent, SGD）的狭义形式每步用一条样本。小批量梯度下降（mini-batch gradient descent）每步平均一组样本的梯度。编程接口常把小批量更新也叫 SGD。训练要说明抽样方式和损失的归约方式。对均匀无偏抽样，批次平均梯度的期望是训练集梯度。批次越小通常噪声越强，实际方差也受样本相关性影响。【刘远超·7 网络优化 p29--33】【刘扬·10前馈神经网与反向传递算法 p20--21】

刘远超课件用学习率 $\eta$ 或 $\alpha$、当前梯度 $g_t$ 和累积量 $v_t$。刘扬课件用参数向量 $\theta$ 与目标的梯度，并有“小批量随机梯度”之说。本文将梯度均值写作 $g_t$，动量记 $u_t$，平方梯度统计记 $s_t$，避免与别章的输出参数 $v$、LayerNorm 方差混淆。刘远超称“批量标准化”，刘扬称“批量归一化”，均指 Batch Normalization（BN）。层归一化是 Layer Normalization（LN）。【刘远超·7 网络优化 p32--43、p59】【刘远超·8 建模序列数据的神经网络 p64】【刘扬·10前馈神经网与反向传递算法 p21、p36】

\section{公式与完整推导}

\subsection*{更新式、学习率和动量}

记 $g_t=|B_t|^{-1}\sum_{i\in B_t}\nabla L_i(\theta_t)$。纯梯度下降为 $\theta_{t+1}=\theta_t-\eta_t g_t$。在一维二次目标 $J(w)=\frac12 a(w-w^*)^2$、$a>0$ 上，
\[
 w_{t+1}-w^*
 =w_t-w^*-\eta a(w_t-w^*)=(1-\eta a)(w_t-w^*).
\]
因此稳定收敛需要 $0<\eta a<2$。$0<\eta a<1$ 时从同侧接近，$1<\eta a<2$ 时两侧震荡但幅度衰减，$\eta a=2$ 时非零起点来回跳，$\eta a>2$ 时发散。多维二次函数的每个特征方向按本式变化，最大曲率限制稳定步长。模型非凸时，这个推导说明局部近似，不能据此宣称一条全局收敛保证。【刘远超·7 网络优化 p29--37】【补充：二次局部模型的稳定条件】

动量（Momentum）的一种约定为
\[
 u_t=\beta u_{t-1}+g_t,\qquad
 \theta_{t+1}=\theta_t-\eta_t u_t,\quad u_0=0 .
\]
展开得 $u_t=g_t+\beta g_{t-1}+\beta^2g_{t-2}+\cdots$，方向一致的分量得到累积，交替变号的分量可部分抵消。另有写法用 $(1-\beta)g_t$，或把 $\eta_t$ 放进速度定义，手算以题设为准。Nesterov 加速梯度（Nesterov accelerated gradient, NAG）在临时位置 $\theta_t-\eta_t\beta u_{t-1}$ 取梯度，再合并旧动量与该梯度。【刘远超·7 网络优化 p37--41】

学习率调度（learning-rate schedule）可按步数分段、指数衰减或余弦退火，也可在验证表现停滞时减小学习率。若题目给余弦周期 $T$，常见一个周期内写
$\eta_t=\eta_{\min}+\tfrac12(\eta_{\max}-\eta_{\min})(1+\cos(\pi t/T))$。
Warmup（预热）在起始阶段逐步增大学习率，通常与后续衰减组合。【刘远超·7 网络优化 p36】【补充：调度公式及预热】

\subsection*{平方梯度与偏差修正}

在各坐标上分别计算平方、平方根和除法，$g_t^2$ 表示逐元素平方：
\begin{align*}
 \text{AdaGrad:}\quad& s_t=s_{t-1}+g_t^2,\qquad
 \theta_{t+1}=\theta_t-\eta\,g_t/(\sqrt{s_t}+\epsilon),\\
 \text{RMSProp:}\quad& s_t=\rho s_{t-1}+(1-\rho)g_t^2,\qquad
 \theta_{t+1}=\theta_t-\eta\,g_t/(\sqrt{s_t}+\epsilon),\\
 \text{Adam:}\quad& m_t=\beta_1 m_{t-1}+(1-\beta_1)g_t,\qquad
 s_t=\beta_2 s_{t-1}+(1-\beta_2)g_t^2,\\
 &\widehat m_t=m_t/(1-\beta_1^t),\qquad
 \widehat s_t=s_t/(1-\beta_2^t),\\
 &\theta_{t+1}=\theta_t-\eta\,\widehat m_t/(\sqrt{\widehat s_t}+\epsilon).
\end{align*}
AdaGrad（adaptive gradient algorithm）的累计量单调不减，某坐标早期平方梯度很大时，后面步长也可能偏小。RMSProp（root mean square propagation）给久远梯度较小权重，可随梯度尺度改变步长。Adam（adaptive moment estimation）还平滑带符号的一阶梯度。零初始化时若 $\mathbb E[g_j]=\mu$ 在这些步近似恒定，
\[
 \mathbb E[m_t]=(1-\beta_1)\sum_{j=1}^t\beta_1^{t-j}\mu
               =(1-\beta_1^t)\mu.
\]
同理，若 $\mathbb E[g_j^2]=\nu$，则 $\mathbb E[s_t]=(1-\beta_2^t)\nu$，于是两项各除对应因子。二阶矩是 $\mathbb E[g^2]$，方差才是 $\mathbb E[g^2]-\mathbb E[g]^2$。时间步从 $1$ 开始。$\epsilon$ 的放置按题面约定。【刘远超·7 网络优化 p42--43】【2021深度真题 第6题】【模拟卷三 大题1(7--9)】【补充：Adam 偏差修正推导】

\subsection*{梯度稳定与归一化的导数}

设 $A_j$ 为第 $j$ 步或第 $j$ 层的局部雅可比。梯度反传包含 $g_l=A_{l+1}^{\mathsf T}\cdots A_L^{\mathsf T}g_L$，因此
$\|g_l\|\le\|g_L\|\prod_{j=l+1}^L\|A_j\|$。若各范数有严格小于 $1$ 的共同上界，梯度衰减。连续放大某些方向时可爆炸，范数上界大于 $1$ 本身不足以证明某条梯度一定会爆炸。全局范数裁剪（global gradient norm clipping）把所有梯度拼成 $g$，置
\[
 g'=\min\{1,C/\|g\|_2\}\,g\quad(g\ne0),\qquad g'=0\quad(g=0).
\]
它保留拼接向量的方向。逐元素裁剪分别截断各坐标，通常会改变方向。【刘远超·7 网络优化 p58--59】【刘扬·10前馈神经网与反向传递算法 p35】

对于同一统计组的 $m$ 个数值，设
\[
 \mu=m^{-1}\sum_j x_j,\quad u_i=x_i-\mu,\quad
 v=m^{-1}\sum_j u_j^2,\quad s=\sqrt{v+\epsilon},
 \qquad y_i=\gamma_i u_i/s+\beta_i.
\]
BN 在特征矩阵（行是样本）中按列跨样本统计。常用卷积 BN 对同一通道跨批次和空间位置统计。训练更新运行均值、方差。推理通常使用积累的统计量。LN 在每个样本内沿特征统计，Transformer 中通常对每个词元的隐藏维统计。训练和推理都从当前样本计算。BN 的 $\gamma,\beta$ 通常按特征共享到各样本，LN 的 $\gamma_i,\beta_i$ 常按特征设置。由 $\sum_i u_i=0$，
\[
 \frac{\partial u_i}{\partial x_j}=\mathbf1_{i=j}-1/m,
 \quad
 \frac{\partial v}{\partial x_j}
 =\frac2m\sum_i u_i(\mathbf1_{i=j}-1/m)=\frac{2u_j}{m}.
\]
从而 $\partial s/\partial x_j=u_j/(ms)$，进而
\[
 \frac{\partial (u_i/s)}{\partial x_j}
 =\frac{\mathbf1_{i=j}-1/m}{s}-\frac{u_i u_j}{m s^3}.
\]
反传到某个 $x_j$ 时对组内全部 $y_i$ 的贡献求和。这里的 $1/m$ 是训练标准化的通常定义，题设若另定样本方差分母则按题设。【刘远超·7 网络优化 p59】【刘远超·8 建模序列数据的神经网络 p63--64】【刘扬·10前馈神经网与反向传递算法 p36】【补充：标准化雅可比及推理统计量说明】

\subsection*{训练阶段的随机正则化与权重衰减}

倒置 Dropout（inverted dropout）中，保留率 $q>0$，$M_i\sim{\rm Bernoulli}(q)$，训练激活 $\widetilde a_i=M_i a_i/q$。故 $\mathbb E[\widetilde a_i\mid a_i]=a_i$，推理直接用 $a_i$。仿射层输出的期望也相同，非线性层之后一般不能把期望移到激活内部。早停（early stopping）用验证集保存最佳检查点，数据增强（data augmentation）以任务允许的变换增加训练样本变化。增强须保持或同步调整标签，且训练验证集划分后再做增强，避免同源样本泄露。【刘远超·7 网络优化 p64--65】【刘扬·10前馈神经网与反向传递算法 p37--38】【模拟卷一 判断12】【模拟卷二 判断11】

若 $R(w)=\tfrac12\|w\|^2_2$，$\nabla R=w$，纯 SGD 对 $J_{\rm data}+\lambda R$ 的一步为
\[
 w' = w-\eta(g+\lambda w)=(1-\eta\lambda)w-\eta g.
\]
所以在此约定的 SGD 中，L2 正则项和乘性权重衰减（weight decay）等价。Adam 若把 $\lambda w$ 放到梯度中，动量和平方梯度统计一起改变。解耦权重衰减（decoupled weight decay, AdamW）用数据梯度更新矩，只对权重另外乘 $1-\eta\lambda$。L1 惩罚 $\sum_i|w_i|$ 在零点的次梯度区间为 $[-1,1]$，可产生精确零。高斯和拉普拉斯先验对应的正则化系数还取决于数据损失采用求和还是平均。【刘远超·7 网络优化 p22、p59、p64】【2021秋真题 第1题】【模拟卷二 判断14】【补充：解耦形式】

\section{计算题型与解题模板}

\paragraph{型一：SGD 与小批量更新的分歧。}
同一点求各梯度再取均值是一次小批量。逐条 SGD 的第二次梯度要代新参数。例题取 $L_i(w)=\frac12(w-a_i)^2$，$a_1=1,a_2=7$，$w_0=0,\eta=1/2$。先求两条初始梯度 $-1,-7$。小批量平均为 $-4$，一步到 $2$。逐条更新先到 $1/2$，再用 $1/2-7=-13/2$ 得到 $15/4$。比较时除了梯度评估点不同，还要说明两次单样本步骤的总步幅不必等于一次平均步幅。

\paragraph{型二：逐步更新各优化器的状态。}
先列梯度流和初始状态，再按优化器递推，Adam 最后修正。例题的两个梯度按给定流 $g_1=4,g_2=-1$，$w_0=1,\eta=0.1$，各状态初始为 $0$，$\epsilon=0$。Momentum 取 $\beta=1/2$，RMSProp 取 $\rho=3/4$，Adam 取 $\beta_1=1/2,\beta_2=3/4$。

\begin{center}
\begin{tabularx}{\linewidth}{@{}lXX@{}}
\toprule
方法 & 第一步 & 第二步\\
\midrule
Momentum & $u_1=4,\ w_1=0.6$ & $u_2=1,\ w_2=0.5$\\
AdaGrad & $s_1=16,\ w_1=0.9$ & $s_2=17,\ w_2=0.9+0.1/\sqrt{17}\approx0.9243$\\
RMSProp & $s_1=4,\ w_1=0.8$ & $s_2=13/4,\ w_2=0.8+0.1/\sqrt{13/4}\approx0.8555$\\
Adam & $m_1=2,s_1=4$。修正为 $4,16$。$w_1=0.9$ & $m_2=1/2,s_2=13/4$。修正为 $2/3,52/7$。$w_2\approx0.8755$\\
\bottomrule
\end{tabularx}
\end{center}
Adam 第二步的一阶矩仍为正，虽 $g_2=-1$，更新依然减小 $w$。这比只比较公式更能说明一阶动量与平方累计的区别。

\paragraph{型三：BN 与 LN 的轴。}
先标行列意义，再按指定轴减均值、除标准差、缩放和平移。例题为两样本三特征 $A=\begin{pmatrix}1&2&3\\5&6&9\end{pmatrix}$，分母取组大小，$\gamma=1,\beta=0,\epsilon=0$。BN 的列均值 $(3,4,6)$、列方差 $(4,4,9)$，输出两行为 $(-1,-1,-1)$、$(1,1,1)$。LN 的行均值为 $2,20/3$，行方差为 $2/3,26/9$，输出两行为
\[
 \frac{(-1,0,1)}{\sqrt{2/3}},\qquad
 \frac{(-5,-2,7)}{\sqrt{26}} .
\]
核查每个归一化组的均值为零、均方为一即可定位轴选错的问题。

\paragraph{型四：裁剪、权重衰减和 Dropout。}
例题给 $g=(6,-8)^{\mathsf T}$、裁剪阈值 $5$，参数 $w=(2,1)^{\mathsf T}$。先求 $\|g\|_2=10$，统一乘 $1/2$ 得 $(3,-4)^{\mathsf T}$。再按 $\eta=0.1,\lambda=0.2$ 作一次带解耦衰减的 SGD 更新，$(1-\eta\lambda)w-\eta g'=(1.66,1.38)^{\mathsf T}$。另对激活 $(2,-1,4)$ 用保留率 $1/2$、掩码 $(1,0,1)$，训练样本的激活变成 $(4,0,8)$，推理使用原向量。不能把“某一次训练掩码结果不同”当成期望不相等的证据。

\section{易混辨析}

\begin{center}
\begin{tabularx}{\linewidth}{@{}lXX@{}}
\toprule
易混项 & 判别依据 & 可写在答卷上的一句话\\
\midrule
Momentum 与 RMSProp & 前者看带符号梯度，后者看梯度平方 & 动量改变方向的时间平滑，RMSProp 调节各坐标的尺度。\\
平方梯度均值与方差 & 二阶原点矩与中心二阶矩 & Adam 的分母由二阶原点矩估计构成。\\
BN 与 LN & 前者常跨样本，后者在样本内 & LN 不依赖同批的其他样本，BN 的训练输出通常依赖。\\
学习率衰减与权重衰减 & 分别作用于步长和参数 & 降低学习率减少每步移动，权重衰减直接缩小参数。\\
Dropout 与残差 & 前者随机置零并按概率缩放，后者提供恒等路径 & Dropout 主要用于正则，残差主要改善优化。\\
\bottomrule
\end{tabularx}
\end{center}

\section{常见错误与陷阱}

\begin{itemize}
\item 对 Adam 第一步按 $t=0$ 修正将出现除零。从第一次用梯度更新矩开始计为 $t=1$。
\item 把 $(1-\beta)$ 的写法与题面没有这项的动量约定混用，可能使整条迭代轨迹错误。
\item 写 BatchNorm 时混用训练批统计与推理运行统计。标准化要求的方差分母是统计组大小，须确认不是无偏估计的 $m-1$。
\item 全局范数裁剪只约束裁剪后传入优化器的梯度。随后的动量和参数衰减仍会影响实际参数改变量。
\item 早停以独立验证数据为依据，不能用测试集反复调停止轮次。倒置 Dropout 不在测试时再乘一次保留概率。
\item “学习率过小使每步更新过大”方向反了。过小主要导致步子小、收敛慢。大学习率未必每次都发散，稳定性取决于目标局部曲率。
\end{itemize}

\section{往年真题与参考答案}

\textbf{2021 秋第 1 题}问正则化与先验。得分点是：（1）由 MAP 给出 $-\ln p(D\mid w)-\ln p(w)$。（2）零均值独立高斯先验对应平方惩罚，拉普拉斯先验对应绝对值。（3）说明控制容量、防止过拟合的作用，不能用先验惩罚替代训练数据似然。

\textbf{2021 深度第 6 题}回忆为 Adam 的两个矩。答题先给一阶 $m_t$、二阶 $s_t$ 的更新及统计量含义。再说明零初值偏差与修正。最后写参数更新并标 $\epsilon$。材料未记全具体问法，应先确保矩的定义正确。

\textbf{2022 春四（1）、四（2）}分别问深网训练的困难和使用动量的方法。四（1）得分点见前馈章，重点可写归一化、初始化和残差三种手段。四（2）可举 Momentum、NAG、Adam、Nadam【补充】。分别写累计历史梯度、在前瞻位置求新梯度、在带符号矩之外使用平方矩、将 Nesterov 式前瞻思想结合 Adam 矩估计。另可举带动量扩展的 RMSProp，需说明普通 RMSProp 的平方统计本身不是一阶动量。2023 第 5 题仅留下“优化算法”主题，不能补造题设。

\textbf{2022 秋第 4 题}问定义、实验现象与两种措施。得分点是：（1）训练误差低但独立数据误差较高，模型吸收了噪声或训练样本偶然结构。（2）结合实际实验，比如高阶多项式的训练与验证误差差距、边界波动。（3）给两项对症改进并解释，如降阶或加强正则、增加有效数据。2024 秋无编号判断示例按“过低学习率导致过大步幅”为错说明。

\textbf{2020 春 A 卷第 3.2、3.3 题}要求为图像去噪选择损失和训练优化。可按数据配对、像素重构损失、稳定性和实验验证组织。例如 MSE、合适的归一化和裁剪、比较验证损失与感知效果。网络部分另见 CNN 章。

\section{本章自测}

\begin{enumerate}
\item $J(w)=2(w-1)^2$，常数学习率使非最优点收敛到 $1$ 的范围是什么？
\item 全局梯度为 $(-5,12)$，裁剪阈值 $6.5$，求裁剪向量。
\item Adam 取初值矩为零，第一步某坐标的梯度为 $-3$，写出两项修正矩。
\item 当 $A$ 的行是样本时，改变同一批次的其他行，当前样本的 BN 和 LN 结果何者可能改变？
\item 倒置 Dropout 的保留率为 $0.8$，训练观察到某激活被保留，若原值为 $-2$，求本次训练值与推理值。
\item 某模型训练损失继续下降，而验证损失连续上升。说明应依据哪个检查点交付，并给出一种泛化改进。
\end{enumerate}
\textbf{答案。} （1）曲率为 $4$，故 $0<\eta<1/2$。（2）范数 $13$，结果 $(-2.5,6)$。（3）$\widehat m_1=-3,\widehat s_1=9$。（4）训练时 BN 可能变，LN 只用当前样本组内数值。（5）$-2/0.8=-2.5$，推理为 $-2$。（6）保留验证最佳检查点，并如增加正则、做标签一致的增强或补充数据，最终测试只用于报告。

\chapter{卷积神经网络}

\section{本章考点清单}

\begin{itemize}
\item \textbf{\(\bigstar\bigstar\bigstar\) 卷积尺寸和参数量。} 按互相关手算，区分核通道数、输出通道数、偏置、步长和填充。【刘远超·9 卷积神经网络 p6--16、p23、p30】【刘扬·11卷积神经网络 p4--22】【模拟卷二 单选4、大题5(1)】【模拟卷三 大题5(1--3)】
\item \textbf{\(\bigstar\bigstar\bigstar\) 池化和梯度。} 能算最大／平均池化、输出尺寸、重叠窗口梯度以及全局平均池化。【刘远超·9 卷积神经网络 p17--18、p30--33】【刘扬·11卷积神经网络 p24--26】【本次范本 判断4】【模拟卷一 大题25(1)】【模拟卷二 判断19、大题5(2)】【模拟卷四 大题5(2)】【补充：全局平均池化详式】
\item \textbf{\(\bigstar\bigstar\bigstar\) 感受野、有效核与共享梯度。} 逐层更新感受野和采样间距，能比较堆叠小核、空洞卷积和扩大步长。按所有用到同一核权重的位置累加导数。【刘远超·9 卷积神经网络 p14、p52--55】【刘扬·11卷积神经网络 p15--22】【2020春 A卷 第3.1题】【2022春真题 五(1)】【模拟卷一 判断17、大题25(1)】【模拟卷二 大题5(3)】【模拟卷四 大题5(1)】【补充：一般感受野递推与空洞卷积】
\item \textbf{\(\bigstar\bigstar\) 经典网络与捷径。} 知道 LeNet、AlexNet、VGG、GoogLeNet/Inception、ResNet 的设计要点，能写残差直接梯度路径。【刘远超·9 卷积神经网络 p19--29、p35--61】【刘扬·11卷积神经网络 p27--34】【模拟卷三 单选9、大题5(4)】【模拟卷四 单选7】
\item \textbf{\(\bigstar\bigstar\) 任务设计与实验。} 根据分类、去噪、文本分类或目标定位选择输出和损失，说明输入表示与 CNN 特征提取方式。【刘远超·9 卷积神经网络 p63--64】【刘扬·11卷积神经网络 p23、p26--29】【2020春 A卷 第3.2--3.4、5题】【2021深度真题 第8题（残片）】【2022春真题 五(2)、八(2)】【2023真题 第9题（旧范围）】
\end{itemize}

\section{核心概念与定义}

卷积神经网络（convolutional neural network, CNN）把局部模式检测应用到许多空间位置。局部连接（local connectivity）缩小每个单元所读入的区域，参数共享（weight sharing）让同一个核在所有位置使用，不同输出通道可以学习不同的核。若没有边界影响且步长为 $1$，平移输入会相应平移输出，这是平移等变（translation equivariance）。池化和分类汇聚可提高某些小平移下的稳健性，有限的填充、步长和裁剪会改变实际性质。

互相关（cross-correlation）直接把核与当前输入块对齐乘加，数学卷积（convolution）则有核的翻转。在 CNN 算题中，“卷积层”通常依题面按互相关计算。刘扬课件 p4--6 明确区分卷积和相关，再用 CNN 名称称呼该计算。刘远超用“滑动窗口、卷积核、特征图”描述同一操作。下文把高宽通道写成 $H,W,C$。刘扬图中的 width、height、depth 对应 $W,H,C$，其中 depth 通常是通道数，并非网络层数。刘远超 p15--16 的“一组卷积核”指覆盖输入各通道、求和成一张输出特征图的权重组。【刘远超·9 卷积神经网络 p6--16】【刘扬·11卷积神经网络 p4--17】

感受野（receptive field）指某个输出可能受影响的原输入空间范围。理论感受野来自网络结构，实际梯度或影响分布可以集中在其中部分区域。全局平均池化（global average pooling, GAP）分别对每一通道的全图位置取平均，得到长度为通道数的向量。空洞卷积（dilated convolution）用间隔大于 $1$ 的采样位置扩大跨度，已存在的采样权重数不变。【刘远超·9 卷积神经网络 p17--18、p51--55】【补充：GAP、空洞和感受野术语】

\section{公式与完整推导}

\subsection*{互相关、输出尺寸与参数数}

令输出第 $q$ 个通道使用核 $K^{(q)}$，空间索引从 $0$ 开始，填充后的输入超出原范围按零处理。取对称填充 $p_h,p_w$，则
\[
 Y_{ijq}=b_q+
 \sum_{c=0}^{C_{\rm in}-1}
 \sum_{r=0}^{k_h-1}\sum_{s=0}^{k_w-1}
 K^{(q)}_{rsc}\,X_{\,iS_h+rD_h-p_h,\ jS_w+sD_w-p_w,\ c}.
\]
$S$ 是步长（stride），$D$ 是空洞率（dilation），$D=1$ 时相邻核元素连续。高方向最远的两个核元素相差 $D_h(k_h-1)$ 格，跨度为 $K_h^{\rm eff}=D_h(k_h-1)+1$。输入填充后的高为 $H+p_{\rm top}+p_{\rm bottom}$，末核窗口的起点 $iS_h$ 不超过填充后高减跨度，故
\[
 H_{\rm out}=\left\lfloor
 \frac{H+p_{\rm top}+p_{\rm bottom}-K_h^{\rm eff}}{S_h}
 \right\rfloor+1,
\]
宽方向同理。先确认填充后空间能够放置一个有效核。零填充（padding）是边界规则。若框架使用 SAME，步长大于 $1$ 时通常令输出为 $\lceil H/S_h\rceil$，还需按规则计算可能不对称的两侧填充。标准全通道卷积的参数数为 $C_{\rm out}(k_hk_w C_{\rm in}+1)$（每个输出通道有一偏置）。没有偏置就去掉 $1$。每个核被重复使用，输出空间尺寸不乘入参数量。分组卷积【补充】若分组数 $G$ 整除输入与输出通道，每个输出核读取 $C_{\rm in}/G$ 个输入通道，公式相应改为 $C_{\rm out}(k_hk_w C_{\rm in}/G+1)$。【刘远超·9 卷积神经网络 p6--16、p23】【刘扬·11卷积神经网络 p4--8、p19--22】【补充：空洞和分组的一般化】

$1\times1$ 卷积在第 $(i,j)$ 个位置做 $y_{ij}=Wx_{ij}+b$，$W$ 形状为 $C_{\rm out}\times C_{\rm in}$。步长 $1$、不填充时保持高宽，混合通道，可降维或升维。它不会因核大小为 $1$ 而自动减少通道，通道数由卷积核个数决定。【刘扬·11卷积神经网络 p20】

\subsection*{共享参数与池化的反向传播}

设 $\delta_{ijq}=\partial L/\partial Y_{ijq}$。同一个核权重用于所有 $(i,j)$，所以
\begin{align*}
 \frac{\partial L}{\partial K^{(q)}_{rsc}}
 &=\sum_{i,j}\delta_{ijq}
   \frac{\partial Y_{ijq}}{\partial K^{(q)}_{rsc}}
 =\sum_{i,j}\delta_{ijq} X_{iS_h+rD_h-p_h,\ jS_w+sD_w-p_w,\ c},\\
 \frac{\partial L}{\partial b_q}&=\sum_{i,j}\delta_{ijq}.
\end{align*}
若另有多个样本，也按损失归约规则跨样本合并。输入梯度的规则是把每个上游 $\delta_{ijq}$ 乘窗口内对应核系数，累加散回被读取的位置。池化输出尺寸仍可用窗口核跨度计算。最大池化（max pooling）窗口内最大值唯一时，仅 argmax 位置接收该窗口的梯度。若有并列须按题设选择某个位置或分配次梯度。平均池化（average pooling）把窗口梯度均分给参与求平均的元素，同一输入被多个窗口覆盖时累加。带填充的平均池化是否把补零计入分母依给定实现。GAP 对整个空间均分。其输出 $y_c=(HW)^{-1}\sum_{ij}X_{ijc}$，导数为 $1/(HW)$。【刘远超·9 卷积神经网络 p14、p17--18、p30--33】【刘扬·11卷积神经网络 p24--26】【模拟卷一 判断17】

\subsection*{感受野与残差梯度}

记上一层相邻两个位置的输入步距为 $j_{l-1}$，一个位置的理论感受野长度为 $r_{l-1}$，初值 $j_0=r_0=1$。本层等效跨度为 $k_l^{\rm eff}$，故它最远两处读入位置在原输入相隔 $(k_l^{\rm eff}-1)j_{l-1}$，再加上一层单位置的宽度，有
\[
 r_l=r_{l-1}+(k_l^{\rm eff}-1)j_{l-1},\qquad
 j_l=j_{l-1}S_l .
\]
高宽分别递推。填充影响中心对齐和有效输入区域，理论长度公式中的跨度与步距计算不变。三个步长 $1$ 的 $3\times3$ 连续卷积有 $7\times7$ 感受野。若各层通道均为 $C$，忽略偏置，参数为 $27C^2$，单层 $7\times7$ 则为 $49C^2$。中间非线性使两种操作一般不等价。【刘远超·9 卷积神经网络 p51--55】【2020春 A卷 第3.1题】【补充：统一跳距记法】

残差连接（residual connection）令 $y=x+F(x)$，可令 $F\approx0$ 逼近恒等映射，使较深模型容易保留较浅网络的解。梯度为
\[
 \nabla_x L=(I+J_F(x))^{\mathsf T}\nabla_y L
 =\nabla_y L+J_F(x)^{\mathsf T}\nabla_y L.
\]
第一项提供不经残差分支内部若干层的路径。当尺度或通道不匹配时可用投影捷径 $y=P(x)+F(x)$，第一项相应是 $J_P^{\mathsf T}\nabla_yL$。恒等通路有助于传递梯度，不能推出梯度永远不衰减。两项可抵消。退化（degradation）是更深普通网训练误差反而更高的观察，不能只用测试误差变高就认定为退化。【刘远超·9 卷积神经网络 p58--61】【刘扬·11卷积神经网络 p30--34】

\subsection*{经典结构速比较}

\begin{center}
\begin{tabularx}{\linewidth}{@{}lXX@{}}
\toprule
网络 & 关键设计 & 题目容易追问\\
\midrule
LeNet-5 & 卷积、降采样、再卷积、全连接，面向手写字符 & 核大小和逐层参数。历史 S 层有带参数的降采样，不能一律等同现代无参 MaxPool。\\
AlexNet & 更大数据与更深网络，引入 ReLU、重叠池化、Dropout 和数据增强，使用局部响应归一化 LRN & 典型结构首层 $11\times11$、步长 $4$。LRN 跨邻近通道利用同一空间位置的响应。\\
VGG & 堆叠小 $3\times3$ 核，阶段末池化 & 相同感受野下参数量与非线性的比较。\\
GoogLeNet & Inception 模块并行融合不同尺寸的局部处理，常用 $1\times1$ 先压通道 & 并行分支沿通道拼接，需保持空间尺寸相容。细节【补充】。\\
ResNet & 跳跃连接学习残差，必要时投影匹配形状 & 区分容易逼近恒等映射和梯度直达两点。\\
\bottomrule
\end{tabularx}
\end{center}
【刘远超·9 卷积神经网络 p19--29、p35--61】【刘扬·11卷积神经网络 p27--34】刘远超 p53 的 VGG 对比用了 $7\times7$，作为三层 $3\times3$ 的感受野例子成立。AlexNet 的第一层卷积核以其本章 p37 和刘扬 p29 的结构为准。通道是图像 RGB 或特征维，并非空间维。文本 CNN 可先把句子做成词向量矩阵，用不同词窗口提特征，再池化和分类。【刘远超·9 卷积神经网络 p63--64】

\subsection*{经典结构中的连接量、LRN 与文本卷积}
LeNet-5 的卷积首层 C1 和降采样 S2 保留 6 个通道，边长由 32 到 28 再到 14。C3 对 S2 的输入采用部分连接，不能把这一层套成每个核连满所有输入通道。其原始输出用径向基（Radial Basis Function, RBF）原型距离 $\sum_j(x_j-w_{ij})^2$。特征图每个位置复用同一组权重，\textbf{连接数}可以比独立参数数大。刘远超用 $6\cdot(5^2+1)=156$ 的 C1 参数与 $28^2\cdot156=122304$ 的连接次数举例，包含偏置连接。【刘远超·9 卷积神经网络 p22–29】

AlexNet 的局部响应归一化（Local Response Normalization, LRN）在同一空间位置对邻近通道的响应使用
$b^i_{x,y}=a^i_{x,y}/(k+\alpha\sum_{j\in\mathcal N(i)}(a^j_{x,y})^2)^\beta$。
它用的是通道响应平方和，和按批统计均值方差的 BN 不同。VGG 的全连接权重可重排为空间卷积：假设首次全连接接收 $h\times w\times c$，其每一输出行可还原成同尺寸卷积核，在原尺寸输入时等价，在更大特征图上滑动生成一张得分图。后续对单个位置的通道混合可重排为 $1\times1$ 卷积。处理变长文本时，先把词嵌入排成「位置 $\times$ 向量维数」，常用横跨全部嵌入维数、沿位置滑动的核识别局部词组。对各核的时间响应取全局最大，再连接分类器。【刘远超·9 卷积神经网络 p47–48、p54–55、p63–64】

\section{计算题型与解题模板}

\paragraph{补充例题：局部归一化与文本卷积。}
同一位置三个相邻通道的激活为 $(1,2,3)$。取 $k=1,\alpha=1/2,\beta=1$，中间通道的邻域含三者，归一化结果为 $2/(1+14/2)=1/4$。若另有 8 个词、嵌入维数 5，使用 3 个高度 2 的核和 2 个高度 4 的核，均跨全嵌入宽度，步长 1、无填充、带偏置，时间输出长度分别为 7 和 5，参数总数 $3(2\cdot5+1)+2(4\cdot5+1)=75$。各核全局最大池化后为 5 维分类向量。

\paragraph{型一：多通道卷积、空洞和逐点卷积。}
先写有效核高宽，再套尺寸公式，随后读通道计算参数。例题输入 $8\times10\times4$，空间核 $3\times2$，空洞率 $(2,1)$、步长 $(2,3)$、两侧对称填充 $(2,1)$，输出 $6$ 通道且各有偏置。跨度为 $5,2$，故高为 $\lfloor(8+4-5)/2\rfloor+1=4$，宽为 $\lfloor(10+2-2)/3\rfloor+1=4$，输出 $4\times4\times6$。参数为 $6(3\cdot2\cdot4+1)=150$。续接步长 $1$、不填充的两核 $1\times1$ 卷积，输出 $4\times4\times2$，参数为 $2(6+1)=14$。若某输出位置两个通道的局部两点输入分别为 $(2,1)$、$(0,3)$，核系数分别为 $(1,-2)$、$(2,1)$，偏置为 $1$，则先逐通道算 $0,3$ 再加一次偏置，输出 $4$。

\paragraph{型二：重叠池化前向和梯度。}
输入
$X=\begin{pmatrix}0&3&1\\2&5&8\\6&4&7\end{pmatrix}$，
窗口 $2\times2$、步长 $1$、无填充。最大池化为
$\begin{pmatrix}5&8\\6&8\end{pmatrix}$。
若上游梯度为 $\begin{pmatrix}1&2\\-1&3\end{pmatrix}$，左上、左下窗口的梯度分别给值 $5$、$6$ 的格子。右上和右下都给值 $8$ 的格子，累加 $2+3$。所以
\[
 \nabla_X L=\begin{pmatrix}0&0&0\\0&1&5\\-1&0&0\end{pmatrix}.
\]
同图首个平均池化值为 $(0+3+2+5)/4=2.5$，整图 GAP 值为 $36/9=4$。

\paragraph{型三：结构感受野。}
例题依次经过 $3$ 核步长 $2$ 卷积、$2$ 窗步长 $2$ 池化、空洞率 $2$ 的 $3$ 核步长 $1$ 卷积。从 $(r,j)=(1,1)$ 开始，第一层为 $(3,2)$，第二层为 $(5,4)$，第三层等效核为 $5$，得到 $(21,4)$。最终两个临近输出对应原图采样起点间距为 $4$，即使理论感受野有大量重叠，也不能把间距写成 $21$。

\paragraph{型四：共享核梯度与参数更新。}
一维互相关 $x=(2,0,1,-1)$，核 $(1,2)$，步长 $1$、不填充、偏置为零，输出 $(2,2,-1)$。给定上游梯度 $(1,-2,3)$，三个窗口分别是 $(2,0)$、$(0,1)$、$(1,-1)$。各乘上对应上游得 $(2,0),(0,-2),(3,-3)$。逐列加，核梯度为 $(5,-5)$，偏置梯度 $1-2+3=2$。如学习率为 $0.1$，更新后核 $(0.5,2.5)$，偏置 $-0.2$。

\paragraph{实验对照。}
实验五的 AlexNet 改作 MNIST 灰度十分类，报告采用 $1\times224\times224$ 的单通道输入，首层为 $64$ 个 $11\times11$ 核、步长 $4$、填充 $2$，分类头输出 $10$ 类。对这套输入作题时首层高宽为 $\lfloor(224+4-11)/4\rfloor+1=55$，输出 $55\times55\times64$。实验代码另用自适应平均池化对齐分类头的空间尺寸。古典课件中的三通道、首层 $96$ 个核与千类输出都属于另一配置。【实验5·实验五报告 p2--4】

\section{易混辨析}

\begin{center}
\begin{tabularx}{\linewidth}{@{}lXX@{}}
\toprule
项目 & 特点 & 容易混淆的量\\
\midrule
卷积核 vs 输出特征图 & 一个全通道核生成一张输出图 & 核深度由输入决定，输出深度由核的个数决定。\\
步长 vs 空洞 & 步长挪窗口，空洞拉大核内采样间隔 & 前者改变输出步距，后者先改变等效核大小。\\
卷积 vs 池化 & 卷积有学习参数，常见最大/平均池化按固定规则聚合 & 无参数的池化仍传梯度。\\
局部连接 vs 权重共享 & 分别限制连接范围、复用同一检测器 & 权重共享会把多位置梯度累加到同一个参数。\\
残差 vs 拼接 & 同形残差按元素相加，Inception 拼接通道 & 相加不扩通道。拼接通常增加通道数。\\
\bottomrule
\end{tabularx}
\end{center}

\section{常见错误与陷阱}

\begin{itemize}
\item 有 $C_{\rm in}$ 个输入通道时，每个核大小为 $k_hk_wC_{\rm in}$。一个输出位置通常只加一次该输出通道偏置。
\item 手算先用题面确定是相关还是数学卷积。通常不翻转。若做 SAME 且核为偶数，切勿自行假定两边同宽。
\item 对空洞卷积，参数按采样点数 $k_hk_w$，空间输出按跨度 $K_h^{\rm eff},K_w^{\rm eff}$。二者不是同一个量。
\item 经典网络结构表中的历史参数与现代简化实现可不同，例如 LeNet 降采样与 AlexNet 输入尺寸。以题给层表为准。
\item 为扩大感受野而过早下采样会失去定位细节。去噪或检测题应具体说如何保留局部信息。
\end{itemize}

\section{往年真题与参考答案}

\textbf{2020 春 A 卷第 3.1 题（必答 15 分）}问步长 $1$ 的 $3\times3$ 串联达到感受野 $27$。得分点：（1）每层扩大 $2$，$1+2L=27$，所以 $L=13$。（2）至少给三种减少计算或扩大覆盖的设计，如空洞卷积、低分辨率分支、$1\times1$ 通道瓶颈后接空间卷积、深度可分离卷积【补充】。（3）说明去噪取舍，空洞采样较稀，下采样可能丢纹理，瓶颈过窄会限制表达。若画多尺度分支，要写清如何恢复原尺寸。

\textbf{2020 春 A 卷第 3.2 题（必答 15 分）}给配对噪声图像与干净图像。得分点：（1）输出和目标像素对齐，使用预测与干净目标的 MSE 或 MAE。（2）MSE惩罚大偏差，过度平滑时可比较边缘损失。（3）提出训练稳定方案和独立验证。第 3.3 题（选答 10 分）从网络、学习和集成角度改进，可写残差块、合适学习率及增强、独立训练模型的平均。第 3.4 题（选答 10 分）为视频，时间模块见序列章。第 5 题（选答 20 分）可按暗光、景深或运动模糊等增强需求分类，再用代价加权交叉熵处理错分代价不等，报告分场景的准确率及对后续增强的影响。2021 深度第 8 题只记与 2020 卷积题类似，沿本题复习即可。

\textbf{2022 春五（1）、五（2）}分别问四种扩大感受野的方法和 CNN 借鉴 Transformer。五（1）可给堆叠小核、空洞、步幅／池化后补高分辨率跳连及低分辨率大范围分支。应说明降算力方法与扩大覆盖方法的先后关系。五（2）可写用全局注意力联系远处位置，保留局部卷积，采用窗口或降采样限制注意力成本。【补充：工程设计举例】

\textbf{2022 春八（2）、2023 第 9 题}属于目标检测的旧问法。核心答题点是将单个整图类别扩展成多个类别、位置框及其匹配目标，使分类分支与框回归共同学习。Faster R-CNN 的损失可拆成候选区域网络和第二阶段各自的分类与框回归。2023 原文把另一算法写作 DERT。若按 DETR 理解，可答集合预测用匹配分配训练目标，损失含类别（包括无对象）及边框回归。该解读不当作原文已确认。【补充：检测比较】

\section{本章自测}

\begin{enumerate}
\item 输入长 $13$，核长 $4$，两端各补 $1$，步长 $3$，空洞率 $1$。求输出长度。
\item $1\times1$ 卷积把十个输入通道变成七个输出通道，每个输出有偏置，求参数量。
\item 某一维卷积核长 $5$、空洞率 $3$，求空间跨度与一条核在一通道上的可学习权重数。
\item 一个最大池化的两个窗口都选中同一个输入元，上游梯度为 $-2$、$5$。求该输入从它们收到的梯度。
\item 两层一维卷积的核长均为 $3$，第一层步长 $2$、第二层步长 $1$，求第二层感受野与跳距。
\item 写出同形残差层梯度的两个来源。若残差分支的雅可比为 $-I$，这点的梯度如何？
\end{enumerate}
\textbf{答案。} （1）$\lfloor(13+2-4)/3\rfloor+1=4$。（2）$7(10+1)=77$。（3）跨度 $13$，权重 $5$。（4）$3$。（5）第一层 $(3,2)$，第二层 $(7,2)$。（6）恒等与残差分支。此点二者抵消，局部输入梯度为零。

\chapter{序列模型与 Transformer}

\section{本章考点清单}

\begin{itemize}
\item \textbf{\(\bigstar\bigstar\bigstar\) RNN 的时序、不同输入输出模式、正反方向。} 会写共享状态递推、末状态分类和矩阵形状，还要防止双向结构在实时任务中读到未来。【刘远超·8 建模序列数据的神经网络 p7--15、p31】【刘扬·13RNN-10.25 p3--6】【2024秋真题 单选4、大题8】【2025秋补／缓考 判断4、简答6(1)】【本次范本 大题2】【模拟卷三 判断16、大题5(5)】【模拟卷四 大题5(3)】
\item \textbf{\(\bigstar\bigstar\bigstar\) 时间反向传播。} 写出局部梯度、来自未来的状态梯度和共享参数的跨时求和，能在双向图或逆向图上判断反传顺序。【刘扬·13RNN-10.25 p7--9】【刘远超·8 建模序列数据的神经网络 p11--15（前向图）】【模拟卷一 大题25(2)】【模拟卷三 大题5(6)】
\item \textbf{\(\bigstar\bigstar\bigstar\) LSTM 与 GRU。} 记忆状态、三门、逐元素乘积、长依赖，能算单步状态。理解耦合门与猫眼连接。【刘远超·8 建模序列数据的神经网络 p18--31、p37--43】【刘扬·13RNN-10.25 p9--16】【2022春真题 六】【2025秋补／缓考 判断2、7、简答6(2)】【本次范本 判断3】【模拟卷二 大题5(7--8)】【模拟卷四 大题5(4--5)】
\item \textbf{\(\bigstar\bigstar\) 注意力和 Seq2Seq。} 会区分固定上下文和随目标位置更新的上下文，能解释目标序列错位输入及训练推理差别。【刘远超·8 建模序列数据的神经网络 p34--35、p47--50、p66--69】【刘扬·13RNN-10.25 p17--20】【模拟卷四 大题5(6--8)】【补充：teacher forcing 的统一表述】
\item \textbf{\(\bigstar\bigstar\bigstar\) Transformer。} 考 $Q/K/V$、缩放、掩码、位置编码、多头、交叉注意力、残差、LayerNorm 和 FFN。能算形状、权重以及时间复杂度。【刘远超·8 建模序列数据的神经网络 p47--71】【刘扬·13RNN-10.25 p17--24】【2024秋真题 判断3（残片）】【2025秋补／缓考 判断5、简答7】【模拟卷一 判断18、大题25(4)】【模拟卷二 大题5(9--10)】【模拟卷三 大题5(8--10)】【模拟卷四 大题5(8--10)】
\item \textbf{\(\bigstar\) 递归网络及扩展。} 知道树形递归网络与时序循环网络的连接结构差别、双向与自回归的用途，分清哪些属于早期卷的应用题。【刘远超·8 建模序列数据的神经网络 p74--83】【2020春 A卷 第3.4题】【2021深度真题 第9题】【2022春真题 五(2)】【2023真题 第7题】
\end{itemize}

\section{核心概念与定义}

循环神经网络（recurrent neural network, RNN）通过重复使用同一个状态更新函数处理序列。参数在时间上共享，隐状态携带已读入的序列信息。本文用
\[
 a_t=W_xx_t+W_hh_{t-1}+b,\qquad h_t=f(a_t),\qquad
 s_t=W_yh_t+c,\qquad o_t=g(s_t).
\]
若 $x_t,h_t,o_t$ 的维数分别是 $d_x,d_h,d_y$，则 $W_x$、$W_h$、$W_y$ 依次为 $d_h\times d_x,d_h\times d_h,d_y\times d_h$。刘远超课件写 $U x_t+W h_{t-1}$、$o_t=\operatorname{softmax}(Vh_t)$，对应本章 $(W_x,W_h,W_y)=(U,W,V)$。刘扬课件将输入权重写 $W$、循环权重写 $U$，公式 $r_t=U h_{t-1}+W x_t+b$、$p_t=\operatorname{softmax}(Vh_t+c)$，所以本章 $(W_x,W_h,W_y)=(W,U,V)$。同一字母在不同讲义可能指不同连接，按乘到哪个变量辨认。【刘远超·8 建模序列数据的神经网络 p11--13】【刘扬·13RNN-10.25 p4--6】

多对一（many-to-one）以一串输入预测整句标签，一个常见做法是取沿箭头读完整序列的最后隐状态。一对多（one-to-many）如给一幅图后依次生成描述。多对多（many-to-many）可以按位输出标签，也可以由编码器和解码器进行不同长度的映射。反向 RNN 读序列从右向左，$h_t=f(W_xx_t+W_hh_{t+1}+b)$，边界改为 $h_{T+1}$。双向 RNN（bidirectional RNN）分别做正向和反向递推，通常拼接同一位置的两个状态。实时预测不能直接读取全序列的反向状态。多层 RNN（stacked RNN）让上层在同一时刻接收下层状态，各层在各自时间轴上共享参数。刘远超 p74--83 的递归神经网络（recursive neural network）是树上递归合成子节点表示，用于情感树等，树的父子接边不等同于时间相邻。【刘远超·8 建模序列数据的神经网络 p8--10、p31、p74--83】【刘扬·13RNN-10.25 p3--6】【补充：多层记法】

\section{公式与完整推导}

\subsection*{展开的反向传播与共享梯度}

时间反向传播（backpropagation through time, BPTT）是对时间展开图应用 BP。设 $L=\sum_{t=1}^T\ell_t(s_t)$，$\delta_t=\partial L/\partial a_t$，$\varepsilon_t=\partial\ell_t/\partial s_t$，在 Softmax 配交叉熵时 $\varepsilon_t=o_t-y_t$。从 $h_t$ 出发，有到当前输出 $s_t$ 和到下一步输入 $a_{t+1}$ 的两条路径：
\begin{align*}
 \frac{\partial L}{\partial h_t}
 &=W_y^{\mathsf T}\varepsilon_t+W_h^{\mathsf T}\delta_{t+1},\\
 \delta_t&=\left(W_y^{\mathsf T}\varepsilon_t+
                 W_h^{\mathsf T}\delta_{t+1}\right)\odot f'(a_t),
                 \qquad\delta_{T+1}=0.
\end{align*}
开始于最后时刻，倒序计算。局部地，
$\partial a_{t,i}/\partial(W_h)_{ij}=h_{t-1,j}$，
$\partial a_{t,i}/\partial(W_x)_{ij}=x_{t,j}$，所以
\begin{align*}
 \nabla_{W_h}L&=\sum_t\delta_t h_{t-1}^{\mathsf T},&
 \nabla_{W_x}L&=\sum_t\delta_t x_t^{\mathsf T},&
 \nabla_b L&=\sum_t\delta_t,\\
 \nabla_{W_y}L&=\sum_t\varepsilon_t h_t^{\mathsf T},&
 \nabla_c L&=\sum_t\varepsilon_t .
\end{align*}
若只在末时刻有监督，则其他时刻的 $\varepsilon_t$ 置零，状态间递推仍在。若前向从右向左，来自 $h_t$ 的后一计算节点是 $h_{t-1}$，递推中的 $\delta_{t+1}$ 换成 $\delta_{t-1}$，共享循环权重梯度用 $h_{t+1}$ 作输入。截断 BPTT【补充】在分段边界切掉梯度链，前向仍可把数值状态传到后段，节省内存但短化梯度可追溯的距离。【刘扬·13RNN-10.25 p7--8】【模拟卷三 大题5(5--6)】

对早期状态的影响要乘多个 $\partial h_t/\partial h_{t-1}=\operatorname{diag}(f'(a_t))W_h$。乘积可能指数衰减或增大，故“状态在图上连接”并不保证普通 RNN 学到长期依赖。【刘远超·8 建模序列数据的神经网络 p14--15】【刘扬·13RNN-10.25 p9】

\subsection*{LSTM、GRU 和门的约定}

长短期记忆网络（long short-term memory, LSTM）另维护细胞状态（cell state）$c_t$。对 $r_t=[h_{t-1};x_t]$，
\begin{align*}
 f_t&=\sigma(W_f r_t+b_f),&
 i_t&=\sigma(W_i r_t+b_i),\\
 o_t&=\sigma(W_o r_t+b_o),&
 \widetilde c_t&=\tanh(W_c r_t+b_c),\\
 c_t&=f_t\odot c_{t-1}+i_t\odot \widetilde c_t,&
 h_t&=o_t\odot\tanh(c_t).
\end{align*}
遗忘门（forget gate）控制旧状态保留，输入门（input gate）控制新内容的写入，输出门（output gate）控制细胞状态对隐藏输出的可见程度。沿直接记忆路径且固定各门和候选值，$\partial c_{t,j}/\partial c_{t-1,j}=f_{t,j}$。连续 $k$ 步是 $\prod_{\tau=t-k+1}^t f_{\tau,j}$。遗忘门接近 $1$ 时，这条路径衰减慢。实际总导数还包含经 $h$ 和门输入的间接路径。刘扬称“长短时记忆网络”“记忆元”，刘远超称“长短期记忆网络”“长期状态”，同指上述结构。猫眼连接（peephole connection）让某些门额外读取细胞状态。耦合遗忘、输入门的变体常用 $i_t=1-f_t$。【刘远超·8 建模序列数据的神经网络 p18--27】【刘扬·13RNN-10.25 p9--16】

门控循环单元（gated recurrent unit, GRU）通常没有单列的 $c_t$ 和 LSTM 式输出门。按刘远超 p28 的约定，重置门 $r_t$ 和更新门 $z_t$ 为
\begin{align*}
 r_t&=\sigma(W_r[h_{t-1};x_t]+b_r),&
 z_t&=\sigma(W_z[h_{t-1};x_t]+b_z),\\
 \widetilde h_t&=\tanh(W[r_t\odot h_{t-1};x_t]+b),&
 h_t&=z_t\odot\widetilde h_t+(1-z_t)\odot h_{t-1}.
\end{align*}
重置门调节旧状态对候选值的影响，更新门权衡旧状态和新候选。有的常见实现恰把 $z_t$ 定义成保留旧状态的权重，写成 $h_t=z_t\odot h_{t-1}+(1-z_t)\odot\widetilde h_t$。【补充：互补门命名】本文手算一律按题给式，通用结论与命名无关。

\subsection*{Seq2Seq 注意力、缩放点积与多头形状}

序列到序列（sequence-to-sequence, Seq2Seq）编码器产生源状态 $h_1,\dots,h_{n_s}$。解码器可用当前状态 $d_t$ 对每个输入状态打分 $e_{ti}=\mathrm{score}(d_t,h_i)$，然后
\[
 \alpha_{ti}=\frac{e^{e_{ti}}}{\sum_j e^{e_{tj}}},\qquad
 c_t=\sum_i\alpha_{ti}h_i.
\]
刘远超课件 p35 用 $v_i=\tanh(W_wh_i+b_w)$、$e_i=v_i^{\mathsf T}u_w$ 得逐词权重，然后对 $h_i$ 加权，体现相同的打分、归一化、加权流程。在 Seq2Seq 解码时，打分常随目标时刻改变。教师强制（teacher forcing）在训练时用真实前一个词元作当前解码输入，推理要用已预测词元。对目标 $(y_1,y_2,\mathrm{EOS})$，训练输入为 $(\mathrm{BOS},y_1,y_2)$，输出仍是目标序列。训练看了整个移位矩阵仍需因果掩码，防止某位置读取后面目标。预测时前错可影响后续，这是训练和推理条件不同带来的问题。【刘远超·8 建模序列数据的神经网络 p34--35、p47--50、p66--69】【刘扬·13RNN-10.25 p17--19】【补充：教师强制与暴露偏差】

Transformer 的缩放点积注意力（scaled dot-product attention）用行表示位置：
\[
 S=QK^{\mathsf T}/\sqrt{d_k}+M,\qquad
 P=\operatorname{softmax}_{\rm row}(S),\qquad O=PV .
\]
若各查询和键分量近似独立、零均值、单位方差，则点积 $\sum_{j=1}^{d_k}q_j k_j$ 的方差约为 $d_k$，除以 $\sqrt{d_k}$ 使波动尺度约为 $1$，有助避免 Softmax 过早饱和。因果掩码（causal mask）在 $j>i$ 的未来位置加 $-\infty$，指数后为零。可用键原值的位置加零。对填充标记也可以设掩码。若整行全不可见，普通 Softmax 无归一化分母，应按实现约定处理，训练题通常保证至少一个可用键。【刘远超·8 建模序列数据的神经网络 p55--59、p66--69】【刘扬·13RNN-10.25 p19--20】

自注意力（self-attention）让三者来自同一序列。交叉注意力（cross-attention）在编码器--解码器模型中，$Q$ 从目标侧得到，$K,V$ 从编码器的源侧得到。$n_t$ 个查询、$n_s$ 个键，$Q,K,V$ 的尺寸分别为 $n_t\times d_k,\ n_s\times d_k,\ n_s\times d_v$。因此分数是 $n_t\times n_s$，输出是 $n_t\times d_v$。多头注意力（multi-head attention）用多组不同线性投影得到各头的查询、键、值。每头的输出按通道拼接，再乘 $W^O$。在无偏置、同一个模型维度 $d$、$h$ 个头且 $d_k=d_v=d/h$ 的常见设定，$Q,K,V$ 共有 $3h\cdot d(d/h)=3d^2$ 个参数，输出投影另有 $d^2$，合计 $4d^2$。若四个组合投影全带偏置，再加 $4d$。每头对所有 $n^2$ 对位置做点积和加权求和，合计时间 $O(n^2d)$。显式保存所有注意力分数需 $O(hn^2)$ 空间，线性投影自身另有 $O(nd^2)$ 时间。交叉注意力将 $n^2$ 换成 $n_t n_s$。【刘远超·8 建模序列数据的神经网络 p55--62、p66--67】【刘扬·13RNN-10.25 p19--21】【补充：参数和复杂度分项】

\subsection*{位置编码与子层顺序}

正弦位置编码（sinusoidal positional encoding）常写作
\[
 {\rm PE}(p,2j)=\sin(p/10000^{2j/d}),\qquad
 {\rm PE}(p,2j+1)=\cos(p/10000^{2j/d}).
\]
与同维词嵌入相加可提供位置。普通全局自注意力若没有位置编码、掩码或其他位置结构，则对序列置换等变。设置换矩阵为 $\Pi$，$Q',K',V'=\Pi Q,\Pi K,\Pi V$，则 $S'=\Pi S\Pi^{\mathsf T}$。逐行 Softmax 与行列重排交换，结果为 $\Pi P\Pi^{\mathsf T}$，乘 $V'=\Pi V$ 得 $\Pi PV$。因此输出会随输入重排。不能称整个输出对置换不变。RNN 的递推顺序本身提供次序，通常不需要同样的绝对位置编码。【刘远超·8 建模序列数据的神经网络 p51--53】【刘扬·13RNN-10.25 p21】【模拟卷一 判断18】

课件采用后归一化（post-norm）结构，每个子层先算 $y=\operatorname{LN}(x+\operatorname{SubLayer}(x))$。编码层依次自注意力和逐位置前馈网络（position-wise feed-forward network, FFN），解码层依次掩码自注意力、源--目标交叉注意力、FFN。FFN 常为
$F(x)=W_2\operatorname{ReLU}(W_1x+b_1)+b_2$，同参数独立地用于每个位置。跨位置的信息混合主要发生在注意力，FFN 负责每一位置自己的非线性变换。先归一化（pre-norm）【补充】写 $x+\operatorname{SubLayer}(\operatorname{LN}(x))$，考试按题面结构交代。【刘远超·8 建模序列数据的神经网络 p63--69】【刘扬·13RNN-10.25 p19】

\subsection*{树形递归、矩阵向量递归与张量递归}
对二叉树的两个子节点词向量 $\bm b,\bm c\in\mathbb R^d$，标准递归用 $\bm p=\tanh(W[\bm b;\bm c]+\bm a)$，$W\in\mathbb R^{d\times2d}$。可在每个有标注节点用 $W_s\bm p$ 的 Softmax 预测情感。矩阵向量递归（Matrix-Vector Recursive Neural Network, MV-RNN）给子节点同时保留矩阵 $B,C$ 和向量 $\bm b,\bm c$，父向量由 $W[C\bm b;B\bm c]$ 得到，父矩阵用 $d\times2d$ 的组合矩阵乘纵向拼接的 $[B;C]$。递归神经张量网络（Recursive Neural Tensor Network, RNTN）令 $q=[\bm b;\bm c]$，为第 $j$ 个输出额外加二次项 $q^\top V_jq$，每片 $V_j$ 为 $2d\times2d$，要 $d$ 片才能给 $d$ 个输出。这样的交互允许一个子节点的意义依赖另一个子节点。节点标签损失累加，训练时沿树结构链式反传。【刘远超·8 建模序列数据的神经网络 p74–83】

\section{计算题型与解题模板}

\paragraph{补充例题：识别树结点的二次交互。}
设 $d=1$，左右子向量是 $1,-1$，线性合成的权重都为 1，偏置为 0。普通合成给 $\tanh(1-1)=0$。若张量唯一一片为 $\bigl(\begin{smallmatrix}0&-1/2\\-1/2&0\end{smallmatrix}\bigr)$，$q=(1,-1)^\top$ 给 $q^\top V_1q=1$，父状态变成 $\tanh 1$。用这一步区分线性拼接得分和双线性交互。

\paragraph{型一：RNN 方向、分类与共享权重的 BPTT。}
画出时间箭头，先走前向，再逆箭头传梯度，最后对各位置的局部权重梯度求和。例题用标量线性状态，正向 $h_t=\frac14h_{t-1}+2x_t$、$h_0=1$，输入 $(2,-1,3)$。得
\[
 h_1=17/4,\qquad h_2=-15/16,\qquad h_3=369/64.
\]
若分类有三类，输出得分为 $Vh_3+b$，此处 $V$ 为 $3\times1$。改从右向左读且 $h_4=1$，依次得 $h_3=25/4,h_2=-7/16,h_1=249/64$，整句表示应取现在的 $h_1$。回到正向，若 $\ell_t$ 对当步状态的局部导数是 $(2,0,-1)$，则汇总导数依次逆序为 $-1,-1/4,31/16$。循环系数的梯度为
\[
 (31/16)\cdot1+(-1/4)\cdot(17/4)+(-1)\cdot(-15/16)=29/16.
\]
输入系数的梯度为 $(31/16)\cdot2+(-1/4)\cdot(-1)+(-1)\cdot3=9/8$。

\paragraph{型二：门值既定的 LSTM 和 GRU。}
LSTM 例题固定 $c_{t-1}=-2,f_t=1/4,i_t=3/4,\widetilde c_t=0.4,o_t=1/2$。先算 $c_t=-0.5+0.3=-0.2$，再算 $h_t=\tfrac12\tanh(-0.2)\approx-0.09869$。固定各门和候选时，直接导数为 $1/4$。另设 GRU 的 $h_{t-1}=0.4,r_t=1/4,z_t=3/4$，候选采用 $\tanh(0.5+r_th_{t-1})$。先算候选 $\tanh(0.6)\approx0.5370$，按本章约定得 $h_t=0.25\cdot0.4+0.75\cdot0.5370\approx0.5028$。

\paragraph{型三：上下文、掩码与形状。}
源值依次为 $(1,1)$、$(2,-1)$ 和 $(0,3)$。解码位置的指数化分数为 $(1,4,2)$，归一化后权重为 $(1,4,2)/7$，上下文为 $(9/7,3/7)$。

另取因果自注意力的三步指数化分数矩阵为
\[
\begin{pmatrix}2&1&3\\1&4&2\\3&2&4\end{pmatrix}.
\]
值向量逐位置为 $(3,0,-2)^{\mathsf T}$。先把未来列的指数权重置零，再逐行归一：
\[
 P=\begin{pmatrix}1&0&0\\1/5&4/5&0\\1/3&2/9&4/9\end{pmatrix},
 \qquad O=(3,3/5,1/9)^{\mathsf T}.
\]
若模型维数为 $12$、三头、每头维数 $4$，则无偏置的注意力投影参数为 $4\cdot12^2=576$。源序列长 $4$，目标序列长 $3$ 时，每头交叉注意力的分数为 $3\times4$，输出为 $3\times4$（本例恰有 $d_v=n_s=4$，二者角色不同），拼接后是 $3\times12$。

\paragraph{型四：位置编码及维度。}
令 $d=4$、位置 $p=1$。两对坐标分别使用分母 $10000^0=1$、$10000^{1/2}=100$，故向量 $(\sin 1,\cos 1,\sin 0.01,\cos 0.01)$。先求 $2j/d$ 再取幂，正弦、余弦使用弧度。若词嵌入为 $(1,-2,0,1)$，按坐标相加，不做拼接。无位置编码但有因果掩码时，掩码已经引入可见范围随位置变化的结构。“输入重排输出等变”的证明需要无此位置依赖结构。

\section{易混辨析}

\begin{center}
\begin{tabularx}{\linewidth}{@{}lXX@{}}
\toprule
对象 & 机制或轴 & 典型易混项\\
\midrule
RNN 正向与逆向 & 分别读 $t-1$ 和 $t+1$ 的状态 & 反传方向与实际前向图相反，与时刻编号正负无必然关系。\\
双向与多层 & 双向沿两个时间方向。多层沿网络深度堆叠 & 可以同时存在。双向通常使用了后文。\\
LSTM 与 GRU & LSTM 有三门、独立 $c_t$。GRU 常用重置与更新两门 & GRU 更新门的两种约定互补，按题给式。\\
自注意力与交叉注意力 & 自注意力同源，交叉注意力的查询来自目标侧 & 源长度与目标长度可以不相同。\\
位置编码与因果掩码 & 前者补充顺序，后者禁止读到未来 & 两者可同时使用，作用不同。\\
注意力与 FFN & 注意力跨位置混合值向量。FFN逐位置变换通道 & FFN 参数共享并不意味着它能跨时刻传输入。\\
\bottomrule
\end{tabularx}
\end{center}

\section{常见错误与陷阱}

\begin{itemize}
\item 把“最后一个位置”机械地写作图最右侧。逆向遍历的最后计算节点可能在最左侧。
\item 将所有时刻的隐状态梯度直接相加当成循环权重梯度。每项还要与该时刻的输入隐状态做外积。
\item 写 LSTM 时把哈达玛逐元素乘法换成点积，或把 $c_t$ 等同于 $h_t$。输出门屏蔽隐输出时，细胞状态仍可传递。
\item 缩放使用每头键维数 $d_k$，不是序列长或所有头的总宽度。掩码加在分数上。将未来分数写为零仍会给非零指数权重。
\item Softmax 对每个查询的所有可见键按行归一，不能把三个头或所有查询一起归一。
\item 交叉注意力默认用目标 $Q$ 和源 $K,V$。Teacher forcing 在训练中要右移目标词元，不该让第 $t$ 步直接读入它要预测的同一个标签。
\end{itemize}

\section{往年真题与参考答案}

\textbf{2020 春 A 卷第 3.4 题（视频去噪，选答 10 分）。} 建议得分点：（1）每帧用共享 CNN 提取图像特征。（2）按时间用 RNN、卷积循环模块（ConvLSTM【补充】）或时序融合传播状态。（3）解码回各帧，写清配对的重构损失，必要时再考虑时序一致性。（4）实时只看当前和过去，离线可加入未来帧。

\textbf{2021 深度第 9 题（回忆）。} 比较 BERT 与 PixelRNN／PixelCNN。可按两点给分：BERT 在预训练的掩码词预测中用两侧文本上下文。PixelRNN 或 PixelCNN 对图像按像素顺序做自回归因子分解，预测依赖已生成的像素，循环或掩码卷积处理可见上下文。【补充】回忆未保留第三个比较维度。2022 春五（2）可从全局建模为 CNN 补足长距离关系，同时说明高分辨率注意力成本。2023 第 7 题问如何改进 Transformer（ViT），可举局部窗口或分层降采样节省成本、融合卷积分支提取局部纹理，并给出对照验证。【补充：开放题思路】

\textbf{2022 春第六题。} 比较 LSTM 与 GRU 门控。得分点：（1）LSTM 有遗忘、输入、输出门和独立细胞状态。（2）GRU 以重置和更新门在一个状态内合并旧值和候选。（3）分别解释何时保存旧信息、写新信息及对外输出的控制，通常 GRU 的门控参数更少，但具体性能需验证。

\textbf{2024 秋单选第 4 题和大题第 8 题。} 客观题只保留 HMM、RNN 辨析，HMM 用概率状态转移和发射建模，RNN 用隐藏向量的可学习递推与反向传播。大题得分点：（1）图箭头从右向左，所以依赖 $h_{t+1}$。（2）写 $h_t=f(Ux_t+Wh_{t+1}+b)$。（3）写 $o_t=g(Vh_t+c)$ 并定义参数，若要求分类概率可选 Softmax。判断第 3 题只留了自注意力实现残片，不能补其原命题。

\textbf{2025 秋补／缓考判断第 2、4、5、7 题。} 常规 LSTM 门用 Sigmoid 产生 $0$ 到 $1$ 的门值，乘法按元素。典型 Transformer 用注意力和前馈子层而非靠递归或卷积完成主干。RNN 的共享参数可复用在不同长度序列，不能由此推出对任意长依赖均能成功拟合。LSTM 用门控记忆缓解长期信息丢失。

\textbf{2025 秋补／缓考简答 6（1--2）。} 保留未记下数值的 $h_0$。第一步 $h_1=\tanh(1.6+0.1h_0)$，第二步 $h_2=\tanh(-0.8+0.1h_1)$。题解中的 $h_0=0$ 是额外示例，不能当原题条件。第二问按长期依赖困难、细胞状态直接路径、三个门各自功能组织。

\textbf{2025 秋补／缓考简答 7（1--3）。} 第一小问可只写文字过程：投影出查询、键、值，对每个查询与键打分并缩放，Softmax 得权重，再加权求值，多个头可并行。第二问定位编码给注意力序列顺序。第三问写 RNN 的状态由前一时刻和当前输入递推，顺序已由运算结构提供。课件范本大题 2 还要求用前向末状态做十分类，若 $h$ 长度为 $200$，$V$ 应为 $10\times200$。其作用是把隐藏向量映射成十个分数。

\section{本章自测}

\begin{enumerate}
\item 从右向左的四步 RNN 作多对一分类，若所有输入都应参与预测，取 $h_1$ 或 $h_4$ 的哪个？
\item 一个共享标量权重 $u$ 在三步中所乘状态为 $2$、$1$、$-1$。各净输入的汇总梯度分别为 $1/2$、$2$、$3$。求 $\partial L/\partial u$。
\item 对 LSTM 直接记忆路径，三个遗忘门值依次为 $1,1/2,1/2$，求三步的直接梯度因子。
\item 一个解码位置能看见三个键，缩放点积为 $0,\ln 3,0$，值为 $(-1,2,4)$。写注意力输出。
\item 输入模型维数为 $24$、四个头、各头宽度均为 $6$，无偏置。求四组投影的总参数数。
\item 目标生成序列为三个词元 $j,k,\mathrm{EOS}$。Teacher forcing 下第 2 步读什么？因果掩码为什么仍有必要？
\end{enumerate}
\textbf{答案。} （1）$h_1$。（2）$(1/2)\cdot2+2\cdot1+3\cdot(-1)=0$。（3）$1/4$。（4）权重 $(1,3,1)/5$，输出 $9/5$。（5）$4\cdot24^2=2304$。（6）读真实前词元 $j$。并行训练时后续列也在矩阵中，需要阻止第 2 步访问晚于自身的位置。

\chapter{自编码器与变分自编码器}

\section{本章考点清单}

\begin{itemize}
\item \textbf{\(\bigstar\bigstar\bigstar\) AE 的训练目标、重构维度与用途。} 分清输入、潜变量与输出。会写编码解码、重构误差，并用图像解释去噪或特征提取。【刘远超·10 生成式神经网络 p7--13、p21--22】【2024秋真题 大题7】【模拟卷一 大题24(1)】【模拟卷二 大题4(1--4)】【模拟卷四 大题4(1--5)】
\item \textbf{\(\bigstar\bigstar\) 自编码器变体及 PCA。} 堆栈、去噪、稀疏与受限线性 AE 的最优子空间，会分辨潜空间插值和输出插值。【刘远超·10 生成式神经网络 p8、p11--13】【刘扬·9PCA p1--23（PCA）】【刘扬·14VAE p16--18（隐空间图）】【模拟卷四 大题4(1--6)】【补充：稀疏约束、线性 AE 与 PCA 的证明】
\item \textbf{\(\bigstar\bigstar\bigstar\) VAE 的分布与生成流程。} 写先验、似然、真实后验、近似后验，知道训练重构用近似后验、生成可从先验抽样。【刘远超·10 生成式神经网络 p15--20、p23--24】【刘扬·14VAE p2--5、p10--15】【本次范本 单选4】【模拟卷一 大题24(1)】【模拟卷二 大题4(9)】
\item \textbf{\(\bigstar\bigstar\bigstar\) ELBO 的两种推导、取等和 Gaussian KL。} 能从 Jensen 不等式推下界，从贝叶斯公式推后验差距，手算一维或对角高斯 KL。【刘扬·14VAE p4--9】【模拟卷一 大题24(2--3)】【模拟卷二 判断20、大题4(5--7)】【模拟卷三 大题4(2--9)】
\item \textbf{\(\bigstar\bigstar\bigstar\) 重参数化与梯度。} 给 $\mu,\sigma$ 或 log-variance，求标准差、潜变量、重构项和固定噪声时的偏导。写随机梯度变分贝叶斯估计的角色。【刘远超·10 生成式神经网络 p15--17】【刘扬·14VAE p6--11】【模拟卷一 大题24(3--4)】【模拟卷二 大题4(8)】【模拟卷四 大题4(9)】
\item \textbf{\(\bigstar\bigstar\) 条件生成。\(\bigstar\) 权衡与后验坍塌。} 区分 AE、VAE、CVAE。能按 $\beta$ 比较重构和 KL，说明零 KL 不能独立推出重构好坏。【刘远超·10 生成式神经网络 p18--20、p23--24】【模拟卷二 大题4(9)】【模拟卷一 判断19】【模拟卷四 大题4(7--8)】【补充：$\beta$-VAE 和后验坍塌的诊断】
\end{itemize}

\section{核心概念与定义}

自编码器（autoencoder, AE）把样本 $x$ 经编码器 $f_\phi$ 得到表示 $z$，经解码器 $g_\theta$ 得到重构 $\widehat x$。用 $L_{\rm rec}(x,\widehat x)$ 衡量重构，连续数据常用平方误差。若输出表示一组二值变量的概率，可用逐坐标二元交叉熵。典型重构必须与目标同形，潜变量可因瓶颈较小，也可在稀疏等约束下高维。编码器的可用表示来源于瓶颈或其他约束，否则可能只学会复制输入。去噪自编码器（denoising autoencoder, DAE）收受损输入 $\widetilde x$，目标保持为干净样本 $x$。稀疏自编码器（sparse autoencoder）限制隐藏激活，常在激活上加 L1 或惩罚平均激活偏离稀疏目标【补充】。堆栈自编码器（stacked autoencoder）让多个编码块串联，解码端把高层表示映回输入空间，某些训练方式逐层预训练再微调。【刘远超·10 生成式神经网络 p7--13】

变分自编码器（variational autoencoder, VAE）用概率模型表示数据。定义先验 $p(z)$、解码似然 $p_\theta(x\mid z)$ 和近似后验 $q_\phi(z\mid x)$：
\[
 p_\theta(x,z)=p_\theta(x\mid z)p(z),\quad
 p_\theta(x)=\int p_\theta(x,z)\,dz,\quad
 p_\theta(z\mid x)=p_\theta(x,z)/p_\theta(x).
\]
编码器常输出对角高斯的均值 $\mu_\phi(x)$ 及对数方差 $\ell_\phi(x)$，而非一个固定的编码。训练时从 $q_\phi(z\mid x)$ 抽样以估计重构项。生成时可先抽 $z\sim p(z)$，再抽 $x\sim p_\theta(x\mid z)$，展示时也常显示解码均值。普通 AE 的重构训练不指定统一采样分布。若另建潜空间密度则仍可构造某种生成流程。刘远超 AE 图中 $\phi$、$\psi$ 表编码和解码映射，VAE 图以 $z_{\rm mean}$、$z_{\rm log\_var}$ 表均值与对数方差。刘扬的 Kingma--Welling 英文课件写 $q_\phi(z|x)$、$p_\theta(x|z)$，近似后验记号 $q_\phi$ 与均值／方差网络是整体与参数的关系。下文用 $\phi$ 表编码参数、$\theta$ 表解码参数，$\sigma$ 始终指标准差。【刘远超·10 生成式神经网络 p9--10、p15--24】【刘扬·14VAE p2--11】

\section{公式与完整推导}

\subsection*{最小重构误差与线性 AE 的主子空间}

对中心化数据 $x\in\mathbb R^d$、协方差 $\Sigma=\mathbb E[xx^{\mathsf T}]$，取线性 tied 权重（tied weights），令 $U\in\mathbb R^{d\times k}$、$U^{\mathsf T}U=I$，$z=U^{\mathsf T}x$、$\widehat x=UU^{\mathsf T}x$。令 $P=UU^{\mathsf T}$，则 $P^2=P=P^{\mathsf T}$，从而
\begin{align*}
 \|x-\widehat x\|^2
 &=x^{\mathsf T}(I-P)^{\mathsf T}(I-P)x
 =x^{\mathsf T}(I-P)x,\\
 \mathbb E\|x-\widehat x\|^2
 &=\operatorname{tr}(\Sigma)-\operatorname{tr}(U^{\mathsf T}\Sigma U).
\end{align*}
设 $\Sigma e_i=\lambda_i e_i$，$\lambda_1\ge\cdots\ge\lambda_d$。令 $s_i=\|U^{\mathsf T}e_i\|^2$，则 $0\le s_i\le1$、$\sum_i s_i=k$，而 $\operatorname{tr}(U^{\mathsf T}\Sigma U)=\sum_i\lambda_i s_i$。把权重从较小特征值挪到尚未取满的较大特征值不会降低目标，因此最大值 $\sum_{i\le k}\lambda_i$ 可由前 $k$ 个特征向量取得，最小重构误差是 $\sum_{i>k}\lambda_i$。这就是 PCA 的主子空间。正交基可以翻转或旋转而不改变其投影。非线性、自编码损失非平方、带稀疏约束时，没有这种一般等价保证。【刘扬·9PCA p1--23（主方差背景）】【模拟卷四 大题4(1--3)】【补充：线性 AE 推导】

\subsection*{证据下界的两种推导}

证据下界（evidence lower bound, ELBO）用 $\mathcal L(\theta,\phi;x)$ 记。第一种由 Jensen 不等式得到。假设 $q_\phi(z\mid x)$ 在 $p_\theta(x,z)$ 有质量的区域取正值，
\begin{align*}
 p_\theta(x)
  &=\int q_\phi(z\mid x)\frac{p_\theta(x,z)}{q_\phi(z\mid x)}\,dz
     =\mathbb E_q[p_\theta(x,z)/q_\phi(z\mid x)],\\
 \ln p_\theta(x)
  &\ge \mathbb E_q[\ln p_\theta(x,z)-\ln q_\phi(z\mid x)]\\
  &=\mathbb E_q[\ln p_\theta(x\mid z)]
     -D_{\mathrm{KL}}\!\left(q_\phi(z\mid x)\,\|\,p(z)\right)
     \equiv\mathcal L(\theta,\phi;x).
\end{align*}
由于 $\ln$ 严格凹，有效支撑上 $p_\theta(x,z)/q_\phi(z\mid x)$ 恒定时取等。若 $q$ 缺失证据的部分支撑，第一行的期望只算 $q$ 支撑内的质量，不能直接写成全部证据。在这类情况下，适当情况下仍可先用 $\ln p_\theta(x)\ge\ln\int_{q>0}p_\theta(x,z)\,dz$，再对该部分质量用 Jensen。

第二种由真实后验与近似后验之间的 KL 得：
\begin{align*}
 D_{\mathrm{KL}}(q_\phi\|p_\theta(z\mid x))
 &=\mathbb E_q[\ln q_\phi(z\mid x)-\ln p_\theta(z\mid x)]\\
 &=\mathbb E_q[\ln q_\phi(z\mid x)-\ln p_\theta(x,z)]
   +\ln p_\theta(x)\\
 &=\ln p_\theta(x)-\mathcal L(\theta,\phi;x).
\end{align*}
这说明差距是对真实后验的 KL。它与 ELBO 内部的先验 KL 不同。$p_\theta(x)>0$ 且各项定义良好时，差距为零当且仅当 $q_\phi(z\mid x)=p_\theta(z\mid x)$ 几乎处处成立。固定 $\theta$ 时，对 $\phi$ 最大化 ELBO 等价于最小化该后验 KL。若 $\theta$ 也改变，$\ln p_\theta(x)$ 会改变，不能把联合训练只看成最小化后验 KL。实现常最小化 $-\mathcal L=-\mathbb E_q\ln p_\theta(x\mid z)+D_{\rm KL}(q\|p)$。【刘扬·14VAE p4--6】【模拟卷一 大题24(2)】【模拟卷二 大题4(5--6)、判断20】【模拟卷三 大题4(2--6、9)】

\subsection*{高斯 KL 与重参数化}

设 $q=\mathcal N(\mu,\sigma^2)$、$p=\mathcal N(0,1)$，$\sigma>0$。代入密度，
\begin{align*}
 \ln q(z)-\ln p(z)
 &=-\tfrac12\ln\sigma^2-\frac{(z-\mu)^2}{2\sigma^2}+\tfrac12 z^2,\\
 D_{\mathrm{KL}}(q\|p)
 &=\mathbb E_q[\ln q(z)-\ln p(z)]\\
 &=\tfrac12\left(\mu^2+\sigma^2-1-\ln\sigma^2\right).
\end{align*}
最后一步使用 $\mathbb E_q[(z-\mu)^2]=\sigma^2$、$\mathbb E_q[z^2]=\sigma^2+\mu^2$。多个对角高斯坐标的对数密度比可逐维相加，得到
\[
 D_{\mathrm{KL}}=\tfrac12\sum_j(\mu_j^2+\sigma_j^2-1-\ln\sigma_j^2).
\]
若换成先验 $\mathcal N(m_0,s_0^2)$，标准化 $z$ 到 $(z-m_0)/s_0$，结果为
$\tfrac12[\ln(s_0^2/\sigma^2)+(\sigma^2+(\mu-m_0)^2)/s_0^2-1]$。【刘扬·14VAE p5--11】【模拟卷三 大题4(7--8)】【补充：一般一维和对角高斯推导】

重参数化（reparameterization）将随机性和可导参数分开：
\[
 \epsilon\sim\mathcal N(0,I),\qquad
 \sigma=\exp(\ell/2),\qquad z=\mu+\sigma\odot\epsilon .
\]
因为各坐标条件均值为 $\mu_j$、方差为 $\sigma_j^2$，得到所需的近似后验。重构期望由
\[
 \mathbb E_{q_\phi}[\ln p_\theta(x\mid z)]
 =\mathbb E_{\epsilon}[\ln p_\theta(x\mid\mu_\phi(x)+\sigma_\phi(x)\odot\epsilon)]
\]
改写。在满足可以交换求导与积分的条件下，固定噪声的函数路径可用链式法则求梯度，再对噪声取期望。固定一个噪声样本时，$\partial z_j/\partial\mu_j=1$、$\partial z_j/\partial\sigma_j=\epsilon_j$、$\partial z_j/\partial\ell_j=\frac12\sigma_j\epsilon_j$。用少量噪声均值得无偏的重构项估计，方差随模型和抽样数而异。已知高斯 KL 时常直接计算，不必给它再抽样。刘扬课件称该类估计为随机梯度变分贝叶斯（stochastic gradient variational Bayes, SGVB）。原始随机采样节点不直接给出这条路径导数。另有分数函数估计等办法【补充】，所以“有随机性就无法优化编码器”并不成立。【刘远超·10 生成式神经网络 p15--17】【刘扬·14VAE p6--11】

\subsection*{条件生成、插值和坍塌}

条件变分自编码器（conditional variational autoencoder, CVAE）把类别或属性 $c$ 作为条件。典型训练近似后验为 $q_\phi(z\mid x,c)$，解码似然为 $p_\theta(x\mid z,c)$，较一般的先验为 $p(z\mid c)$。刘远超的生成示例是在指定标签后仍从标准正态取 $z$，即 $p(z\mid c)=p(z)$ 的情况。对固定 $c$ 应用同样的 Jensen 推导：
\[
 \ln p_\theta(x\mid c)\ge
 \mathbb E_{q_\phi(z\mid x,c)}[\ln p_\theta(x\mid z,c)]
 -D_{\rm KL}(q_\phi(z\mid x,c)\|p(z\mid c)).
\]
训练用 $(x,c)$，生成先指定 $c$，抽取相应 $z$ 后以 $(z,c)$ 解码。训练时编码器看真实 $x$，条件生成时无需先有该 $x$。【刘远超·10 生成式神经网络 p18--20】【模拟卷二 大题4(9)】【补充：条件先验的一般形式】

潜空间插值（latent interpolation）可令 $z_\tau=(1-\tau)z_a+\tau z_b$，$0\le\tau\le1$，再解码各点。非线性解码器的中间输出一般不等于输出插值，且直线经过潜空间低概率区域时不保证样本自然。$\beta$-VAE【补充】最大化 $R-\beta K$，其中 $R$ 是对数重构似然期望，$K$ 是先验 KL。$\beta=1$ 是普通 ELBO。$\beta>1$ 加强对近似后验的先验约束，结果仍不高于原 ELBO。$0<\beta<1$ 的结果不一定仍是证据下界。后验坍塌（posterior collapse）表现之一是编码器给大多数样本的后验接近先验，区别不同 $x$ 的信息很弱。强解码器可能在少用潜变量时仍重构较好。可结合每维 KL、重构、改变 $z$ 后的输出敏感性共同判断。零 KL 仅说明后验与先验匹配，不能单独推出重构误差或解码行为。【刘扬·14VAE p15--19（生成与潜空间图）】【模拟卷一 判断19】【模拟卷四 大题4(6--8)】

\section{计算题型与解题模板}

\paragraph{型一：离散证据、Jensen 差距和最优后验。}
例题用 $z\in\{a,b\}$，先验 $(2/3,1/3)$，观测 $x$ 的似然在两状态依次为 $3/4,1/4$，且 $q=(1/2,1/2)$。先算联合概率 $(1/2,1/12)$ 和证据 $7/12$。比值 $p(x,z)/q(z\mid x)$ 为 $1,1/6$，所以 $\mathcal L=-\tfrac12\ln6$，差距为 $\ln(7\sqrt6/12)>0$。把两个联合量除以证据，得真实后验 $(6/7,1/7)$，若取此 $q$，两个比值均是 $7/12$，达到等号。出题若指定候选 $q$ 的分量为零，先检查能否写完整证据的期望。

\paragraph{型二：对角高斯 KL、单样本重构与路径梯度。}
例题设 $\mu=(-1/2,1)$、$\sigma=(1/2,\sqrt2)$，先验 $\mathcal N(0,I)$，样本噪声 $\epsilon=(2,0)$。解码分布 $x\mid z\sim\mathcal N(z_1+2z_2,1)$，观测 $x=2$。步骤如下。
\begin{enumerate}
\item 由 $\sigma^2=(1/4,2)$，逐维 KL 相加得 $\tfrac12[(1/4+1/4-1-\ln(1/4))+(1+2-1-\ln2)]=\tfrac34+\tfrac12\ln2$。
\item 用固定噪声得到 $z=(1/2,1)$，均值 $5/2$。重构对数似然为 $-\tfrac12\ln(2\pi)-\tfrac18$。单样本 ELBO 为 $-\tfrac12\ln(2\pi)-\tfrac18-\tfrac34-\tfrac12\ln2\approx-2.141$。
\item 对 $z$ 的重构梯度是 $(2-5/2)(1,2)=(-1/2,-1)$。对 $\mu$ 的重构梯度相同。对 $\sigma$ 为该向量逐坐标乘噪声，得 $(-1,0)$。
\item KL 对 $\mu$ 的梯度为 $\mu$，对 $\sigma$ 为 $\sigma-1/\sigma=(-3/2,1/\sqrt2)$。由于求的是最大化目标，路径梯度分别减去 KL 导数，得到 $\nabla_\mu\widehat{\mathcal L}=(0,-2)$、$\nabla_\sigma\widehat{\mathcal L}=(1/2,-1/\sqrt2)$。若实现参数是 $\ell$，还要逐坐标乘 $\sigma/2$。
\end{enumerate}

\paragraph{型三：最优线性瓶颈、单样本重构及插值。}
例题取 $\Sigma=\operatorname{diag}(9,4,2)$，中心化三维数据投影到二维瓶颈。选前两个坐标方向，最优总体误差是剩余特征值 $2$。具体点 $x=(-1,2,4)^{\mathsf T}$ 变成编码 $(-1,2)^{\mathsf T}$，重构 $(-1,2,0)^{\mathsf T}$，该点的平方误差为 $16$。总体平均误差与单个样本误差不可混用。另给一维潜变量端点 $-1,3$，取 $\tau=1/2$，$g(z)=z^2$。先插值再解码是 $g(1)=1$，先解码再插值是 $(1+9)/2=5$。

\paragraph{型四：给定类别生成及 $\beta$ 权衡。}
题设条件 $c=1$ 下，先验 $\mathcal N(1,1/16)$，解码器 $p(x=1\mid z,c)=\sigma(z+c)$。如抽到标准正态噪声 $\epsilon=-2$，先算 $z=1+\frac14(-2)=1/2$，再算概率 $\sigma(3/2)\approx0.8176$。只需概率时保留该值，若实际生成一个二值样本还需按该概率抽样。另比较 $(R_A,K_A)=(-2.4,0.9)$、$(R_B,K_B)=(-3.6,0.2)$。两者目标差 $1.2-0.7\beta$，$\beta<12/7$ 时 $A$ 更高，超过此阈值 $B$ 更高，相等时得分相同。

\paragraph{实验对照。}
实验六 MNIST 代码分别输出 \texttt{mu}、\texttt{logvar}，用 \texttt{exp(0.5*logvar)} 取标准差，再累加重构 BCE 和解析 KL。该实现两项都在同一批次求和。若为了记录而按样本平均，两个损失项都用同一分母。一个求和、另一个按像素平均会改变二者的权衡。【实验6·VAE.py 第120--134行】

\section{易混辨析}

\begin{center}
\begin{tabularx}{\linewidth}{@{}lXX@{}}
\toprule
对象 & 输入及训练输出 & 生成或重构时的区别\\
\midrule
普通 AE & 编码为确定的 $z$，从 $z$ 重构 $x$ & 无默认先验，也能在另建潜变量分布后扩展生成流程。\\
去噪 AE & 输入是 $\widetilde x$，目标是干净的 $x$ & 目标形状与干净输入一致，训练拟合清除扰动。\\
VAE & 编码为分布参数，近似后验抽样估计 $R$，另加先验 KL & 无条件生成使用 $p(z)$。\\
CVAE & 编码和解码均可读取 $c$ & 先给 $c$，再抽 $z$，解码依赖二者。\\
ELBO 的先验 KL & $D_{\rm KL}(q_\phi(z\mid x)\|p(z))$ & 惩罚编码分布偏离先验。\\
证据与 ELBO 的差距 & $D_{\rm KL}(q_\phi(z\mid x)\|p_\theta(z\mid x))$ & 约束真实后验的近似质量。\\
\bottomrule
\end{tabularx}
\end{center}

\section{常见错误与陷阱}

\begin{itemize}
\item 误把 $\widehat x$ 与潜变量 $z$ 的维数互换。重构目标通常和输入同形，而瓶颈维数依模型设定，不必总小于输入。
\item 用 $\ln\mathbb E[R]$ 替换 $\mathbb E[\ln R]$，或将 Jensen 方向写反。应明确 $\ln$ 是凹函数。
\item 对高斯 log-variance $\ell$ 写成标准差 $e^\ell$。正确为方差 $e^\ell$、标准差 $e^{\ell/2}$。
\item 最大化 ELBO 时先验 KL 前是负号，最小化负 ELBO 时先验 KL 前是正号。高斯 KL 带 $\ln\sigma^2$，不能任意改对数底数。
\item 梯度估计应说明固定当次噪声。对标准差、对数方差、方差求导是不同结果。
\item 从近似后验采样是条件于已给的 $x$ 的重构路线。新样本生成阶段若无 $x$，应从模型先验采样。
\item 单次采样像某一类别，或平均 KL 较低，都不足以评价生成覆盖与重构质量。应分别检查生成和重构。
\end{itemize}

\section{往年真题与参考答案}

\textbf{2024 秋大题第 7 题（5 分）。} 题目问自编码器的输出长度和图像重构。参考得分点：（1）输入 $x$ 和重构 $x'$ 同长度，目标是逐元素比较。（2）写 $z=f_\phi(x)$、$x'=g_\theta(z)$。（3）用图像举例说明低维编码保留主要特征、解码回到图像空间，以重构误差训练。（4）列出去噪、压缩、特征提取等用途。PDF 回忆漏了这道题，2024 秋回忆 DOCX 和配套参考答案均保留了题干。

\textbf{邻近题型与范本。} 2020 春 A 卷第 3.2 题是配对图像去噪。在类比去噪 AE 时，可把噪声图作输入、干净图作重构目标。原题指定的是卷积网络，并未要求它一定具有编码器瓶颈。本次范本单选第 4 题考 VAE 的用途，得分点是“有先验、可抽样生成新数据”，重构旧图不是生成流程的唯一目标。2025 秋补／缓考的可辨认简答未收录独立的 VAE 题，第 5 道简答缺失，不能推断其主题。

\section{本章自测}

\begin{enumerate}
\item 输入 $32\times32$ 的单通道图像，瓶颈向量长为 $8$，问重构的图像尺寸以及训练去噪时的数据角色。
\item $q=\mathcal N(0,4)$、$p=\mathcal N(0,1)$，求 $D_{\rm KL}(q\|p)$。
\item 编码器给出 $\mu=2$、$\ell=\ln9$，重参数噪声为 $-1/3$。求 $\sigma,z$ 和 $\partial z/\partial\ell$。
\item 在固定 $\theta$ 下，若证据对数为 $-2.8$、ELBO 为 $-3.5$，计算后验差距。它是否给出 ELBO 中的先验 KL？
\item 某个无偏置的线性三维 AE，协方差特征值为 $7,3,1/2$，若保留一维最优瓶颈，求期望平方重构误差。
\item CVAE 若训练使用 $q(z\mid x,c)$ 且先验独立于 $c$，指定标签生成时该从哪个分布抽取潜变量，标签给到哪里？
\end{enumerate}
\textbf{答案。} （1）$32\times32$ 的单通道输出。输入噪声图、目标干净图。（2）$(3-\ln4)/2$。（3）$\sigma=3,z=1,\partial z/\partial\ell=-1/2$。（4）差距为 $0.7$，对应真实后验 KL，不能据此确定先验 KL。（5）$3+1/2=3.5$。（6）从与类别无关的 $p(z)$ 取样，标签与 $z$ 一起输入解码器。

\chapter{课程实验与结合实验作答}

\section{本章考点清单}
\begin{itemize}
\item ★★★ 实验一正弦拟合：构造幂特征，写出带或不带 $L_2$ 的损失和梯度，用验证集比较阶数与惩罚。 【实验一·sin\_loss1.py】【实验一·lab1.pdf p2–6】【2022秋真题 第4题】【2024秋真题 大题2】【2025秋补缓考 简答4】【模拟卷三 大题7(1–7)】【模拟卷四 大题7(1–5)】
\item ★★★ 实验二逻辑回归：预处理、后验与阈值、对数损失、正则、验证与测评。区分特征相关和朴素贝叶斯的条件独立。 【实验二·logit.py】【实验二·lab2.pdf p2–9】【2024秋真题 大题3】【2025秋补缓考 简答1、4】【模拟卷二 大题7(6–7)】
\item ★★★ 实验三聚类：距离、硬分配、后验责任度、加权协方差、奇异性和轮次记录，训练结束后再和类别对照。 【实验三·kmeans.py】【实验三·lab3.pdf p2–14】【2019秋真题 第7题】【2021秋真题 第4、6题】【2024秋真题 大题9】【模拟卷四 大题3(3–9)】
\item ★★★ 实验四 MLP：两种模型的可表示边界，矩阵 BP，改样本量、宽度、学习率对性能的影响。 【实验四·MLP.py】【实验四·实验四报告.pdf p2–9】【2022春真题 四(1–2)】【模拟卷三、四 大题1】
\item ★★ 实验五 CNN/AlexNet：卷积输出尺寸、张量布局、训练与推理、比较激活与 Dropout。 【实验五·AlexNet.py、卷积.py】【实验五·实验五报告.pdf p2–7】【2022春真题 五(1)】【模拟卷三 大题5(1–4)】
\item ★★ 实验六 VAE：编码、重参数化、解码、重构损失与 KL、随机生成和潜空间编辑。 【实验六·VAE.py】【实验六·实验六报告.pdf p3–8】【刘远超·模拟试题范本 p1 选择4】【模拟卷三、四 大题4】
\end{itemize}

\section{核心概念与定义}
\subsection{实验一：用高阶多项式拟合含噪正弦}
目标是把标量 $x$ 映射成 $(1,x,\ldots,x^d)$，用线性参数近似非线性真函数。源码 \texttt{sin\_loss1.py} 用均匀采样生成输入，将输入缩到 $[-1,1]$，目标是 $\sin(2\pi x)$ 加标准差 0.1 的高斯噪声。比较阶数、样本量和 $L_2$ 惩罚。代码含梯度下降、方程求解和共轭梯度（Conjugate Gradient, CG）。同目录的报告另有分区拟合和 RMSProp 讨论，答题时说明自己选用的实现版本。建议真实比较先固定分组，再从训练组估缩放参数，避免整组样本极值参与预处理。【实验一·sin\_loss1.py 的 data、loss1、loss1\_solve、CG】【实验一·lab1.pdf p2–7】

\subsection{实验二：逻辑回归分类}
\texttt{logit.py} 先人工构造特征相关的二维样本，另用威斯康星乳腺癌（Wisconsin Diagnostic Breast Cancer, WDBC）数据作实际分类。训练统计量标准化测试特征。二元逻辑回归把线性得分变成后验概率，用交叉熵训练，阈值决定预测。源码此版本比较 $L_2$，偏置单独更新。另一份 \texttt{lab2.py} 与报告版本讨论 $L_1$、皮肤分类数据和手工构造样本。两份文件的正则类型和数据集不可混说。特征相关性是对朴素贝叶斯条件独立的检验背景，逻辑回归可以直接学习联合特征的条件后验。准确率配合同类召回、精确率和混淆矩阵阅读。【实验二·logit.py 的 train\_logistic\_regression、standardize、入口】【实验二·lab2.py】【实验二·lab2.pdf p2–9】

\subsection{实验三：合成簇与鸢尾花聚类}
\texttt{kmeans.py} 的流程包括高斯簇数据、平均链层次初始化、$K$ 均值、用聚类结果初始化 GMM 和 EM。中心以簇内均值更新，GMM 按责任度更新，协方差初始化与 M 步加小的对角量。数据有真实标签时，训练过程仍不需要用到，最后才把聚类编号与真实类别对照。报告另比较随机或 $K$ 均值初始化，观察样本椭圆、似然和鸢尾花数据上的 ARI（Adjusted Rand Index，调整兰德指数）。原始类别编号与算法簇号可以重排。【实验三·kmeans.py 的 hierarchical\_init\_centroids、gmm\_fit】【实验三·lab3.pdf p2–14】

\subsection{实验四：二维非线性四分类}
\texttt{MLP.py} 生成二维螺旋四分类样本，实现线性 Softmax 分类器与一隐层 MLP。线性头只能沿仿射分界，MLP 的 ReLU 隐层可先变换表示，用多个激活区域形成分段边界。输出为 4 个类别得分，Softmax 与交叉熵相配。代码用行样本矩阵，求平均梯度，扫描样本量、隐藏宽度、学习率。另一版 \texttt{lab4.py} 用 \texttt{make\_blobs} 与 Sigmoid。其评分分析应和各自数据与损失配置绑定。若某版程序的损失与输出梯度不一致，考试写所声明损失的数学导数。【实验四·MLP.py 的 generate\_spiral、MLP、train\_model】【实验四·lab4.py】【实验四·实验四报告.pdf p2–9】

\subsection{实验五：手写数字卷积分类}
\texttt{AlexNet.py} 将 MNIST 单通道图放到 $224\times224$，使用 $1\to64\to192\to384\to256\to256$ 的卷积通道和池化。展平后接全连接，输出 10 个 logits，用交叉熵和 Adam。对比激活函数、Dropout、样本量及学习率。一个更浅的 \texttt{卷积.py} 采用 Keras，报告相应给简单模型。考卷问参数和尺寸时，用题设核大小、步长、填充与是否含偏置，避免把原始 AlexNet 的 RGB 输入直接套到 MNIST 单通道版本。【实验五·AlexNet.py 的 AlexNet、train\_model】【实验五·卷积.py】【实验五·实验五报告.pdf p2–7】

\subsection{实验六：手写数字卷积 VAE}
编码器从 $1\times28\times28$ 到三层降采样特征，展平后两个头输出 $\mu,\log v$。采样 $\epsilon\sim\mathcal N(0,I)$，用 $z=\mu+\exp(\log v/2)\odot\epsilon$ 送入解码器。三层转置卷积从 $4\times4$ 解码至 $32\times32$，再裁剪成 $28\times28$，最后用二元交叉熵的像素和加先验 KL。重建使用给定图片的后验采样。随机生成从先验抽 $z$。潜空间编辑沿某方向调整 $z$，记录视觉变化。实验比较数据量、维数、Dropout 与激活。评价宜分开呈现损失、重建和随机生成。源码是连续灰度目标上的 BCE 工程约定，严格 Bernoulli 概率模型的原始观测为 0/1。【实验六·VAE.py 的 Encoder、Decoder、reparameterize、loss\_function】【实验六·实验六报告.pdf p3–8】

\section{公式与完整推导}

\paragraph{回归的和式、均值与正则。}
设 $X$ 已含第一列 1。源码的一个版本用 $J=\frac12\|Xw-y\|^2+\frac\lambda2\|w\|^2$。对 $w$ 作微扰 $dw$，残差 $e=Xw-y$ 的微扰为 $X\,dw$。因此
$dJ=e^\top X\,dw+\lambda w^\top dw$，故梯度 $X^\top(Xw-y)+\lambda w$。
令它为 0 得 $(X^\top X+\lambda I)w=X^\top y$。若按训练样本数取平均，只有平方误差项被除以 $N$ 时，等效惩罚相对变强。若偏置不罚，将 $I$ 的首个对角元换成 0。写 MSE 时使用 $\|e\|^2/N$，除以 $2N$ 是另定义的损失。【实验一·sin\_loss1.py 的 loss1、loss1\_gradient、loss1\_solve】

\paragraph{二分类与矩阵 BP。}
单样本二元交叉熵 $L=-y\ln p-(1-y)\ln(1-p)$。代入 $p=\sigma(a)$，
\[
\frac{\partial L}{\partial a}
=\Bigl(-\frac{y}{p}+\frac{1-y}{1-p}\Bigr)p(1-p)=p-y .
\]
从 $a=x^\top w+b$ 得 $w$ 梯度 $(p-y)x$、$b$ 梯度 $p-y$。代码用批平均加 $\lambda w$。对于多类行样本，$Z_1=XW_1+\mathbf1b_1$，$H=\phi(Z_1)$，$Z_2=HW_2+\mathbf1 b_2$。行 Softmax 的交叉熵可化为每行的 $\ln\sum_k e^{z_k}-\sum_ky_kz_k$，所以 $\Delta_2=P-Y$。根据 $dZ_2=dH\,W_2+H\,dW_2+\mathbf1db_2$，
\[
\nabla_{W_2}J=H^\top\Delta_2/N,\qquad
\Delta_1=(\Delta_2W_2^\top)\odot\phi'(Z_1),\qquad
\nabla_{W_1}J=X^\top\Delta_1/N .
\]
偏置梯度沿样本行求平均。权重矩阵梯度的形状应与权重本身相同。把列向量课件公式直接搬到行样本实现前要转置相应项。【实验二·logit.py】【实验四·MLP.py 的 forward、backward】

\paragraph{配对模型的 GMM 与卷积维度。}
GMM 的 E 步把「混合权重乘高斯密度」逐行归一化。M 步对中心的加权平方损失求导，得到加权均值。加权离差外积的平均得到协方差。对任意单位方向 $u$，$u^\top S u$ 是离差投影的加权二阶矩，故非负。加 $\epsilon I$ 后对任何非零 $u$ 增加 $\epsilon\|u\|^2$，因此协方差变成正定。卷积的有效核宽为 $k_{\rm eff}=d(k-1)+1$，第一个落点到最后一个满足 $ms+k_{\rm eff}\le H+2p$，故输出高度 $H'=\lfloor(H+2p-k_{\rm eff})/s\rfloor+1$。宽度同算。普通卷积每个核跨输入通道，参数总数是 $C_{\rm out}(k_hk_wC_{\rm in}+1)$，其中 $1$ 是每输出通道的偏置。【实验三·kmeans.py】【刘远超·9 卷积神经网络 p23、p30】【实验五·AlexNet.py】

\paragraph{VAE 的两个输出头与 KL。}
每维对角后验记 $q=\mathcal N(\mu,v)$。写两种一维高斯的对数比，得
$\ln(q/p)=-\frac12\ln v-\frac{(z-\mu)^2}{2v}+\frac12 z^2$。
取 $q$ 下期望，$\mathbb E_q(z-\mu)^2=v$、$\mathbb E_q z^2=v+\mu^2$，所以
$D_{\mathrm{KL}}=\frac12(\mu^2+v-1-\ln v)$。对角维相加。程序输出 $\ell=\ln v$，故标准差为 $e^{\ell/2}$。重建损失对所有像素求和，再加这个 KL。记录每样本损失时可在全批或总集上除以样本数，报告必须说明分母。要最小化的是负 ELBO，KL 前是加号。【实验六·VAE.py 的 loss\_function】【刘扬·14VAE】

\section{计算题型与解题模板}

\paragraph{实验一对应题：从测试表分开算 MSE 与训练目标。}
原创模型 $1+2x-x^2+x^3$，三个目标点 $(-1,-2),(0,0),(2,7)$。逐点代入得预测 $(-3,1,9)$，残差 $(-1,1,2)$，MSE 为 $(1+1+4)/3=2$。若同一批用于定义带正则损失且四个参数都罚，$\lambda=0.4$，损失是 $\frac12\cdot6+\frac{0.4}{2}(1+4+1+1)=4.4$。考试先圈「测试 MSE」还是「训练目标」。

\paragraph{实验二对应题：标准化后读概率并算一次梯度。}
训练统计量 $m=(3,10)$，$s=(2,5)$，输入 $(7,5)$。先得标准化特征 $(2,-1)$。若 $w=(\ln2,0)$、$b=-\ln2$，得 $a=\ln2$、$p=2/3$。真标签为 0 时，未加惩罚的 $w$ 梯度是 $(2/3)(2,-1)=(4/3,-2/3)$，$b$ 梯度为 $2/3$。加 $L_2$ 时给 $w$ 梯度附加 $\lambda w$。

\paragraph{实验三对应题：定位奇异协方差。}
三个等权向量 $(0,1),(2,1),(5,1)$。均值为 $(7/3,1)$，第一方向方差 $29/3-(7/3)^2=38/9$，第二方向方差及交叉项为 0。矩阵 $\operatorname{diag}(38/9,0)$ 不可逆。加 $0.01I$ 得 $\operatorname{diag}(38/9+0.01,0.01)$，两特征值均正。

\paragraph{实验四对应题：ReLU 门控与输出层的误差。}
一个样本的隐层净输入为 $(-1,2)$，ReLU 后 $h=(0,2)$。取列向量约定，二类输出权重 $W_o=\bigl(\begin{smallmatrix}0&0\\0&\ln2/2\end{smallmatrix}\bigr)$，无偏置。得到 logits $(0,\ln2)$、概率 $(1/3,2/3)$。真类为第 1 类，$\delta_o=(-2/3,2/3)$，
$\nabla W_o=\delta_o h^\top=\bigl(\begin{smallmatrix}0&-4/3\\0&4/3\end{smallmatrix}\bigr)$。再把 $W_o^\top\delta_o$ 按激活斜率 $(0,1)$ 逐项相乘，得到隐层净输入梯度 $(0,\ln2/3)$。每一步都写清向量方向。

\paragraph{实验五对应题：从尺寸表推展平特征。}
依据 \texttt{AlexNet.py}：输入边长 224，首核 11、步长 4、填充 2，边长 $1+\lfloor(224+4-11)/4\rfloor=55$。首池化核 3、步长 2，$55\to27$。第二卷积核 5、步长 1、填充 2，仍为 27。第二次池化得 13。之后三层 $3\times3$ 卷积保持 13，第三次池化得 6。末层通道 256，展平长度 $256\cdot6^2=9216$。首卷积有 $64(11^2+1)=7808$ 参数。

\paragraph{实验六对应题：一次重参数与损失核对。}
单维 $\mu=-1/2,v=4$，抽到 $\epsilon=3/4$，则 $z=-1/2+2(3/4)=1$。KL 为 $\frac12(1/4+4-1-\ln4)=13/8-\ln2$。
若示例的两个二值像素为 $(1,0)$，输出概率为 $(0.6,0.2)$，重建和式 $-\ln(0.6)-\ln(0.8)=-\ln0.48$。总损失为 $13/8-\ln2-\ln0.48$。若题设取像素平均，需要明确同时采用的 KL 相对权重。

\section{易混辨析}
\begin{tabularx}{\linewidth}{@{}p{.24\linewidth}XX@{}}
\toprule
容易并写的量 & 第一项 & 第二项\\
\midrule
训练曲线／验证曲线 & 更新使用的样本上的误差 & 独立数据上的误差，引导选阶数和早停\\
逻辑回归得分／概率 & $x^\top w+b$，可为任意实数 & $\sigma$ 后在 $[0,1]$，可再做阈值\\
聚类标签／监督标签 & 簇编号可重排 & 监督训练时类别含义固定\\
真实试验参数／教材原架构 & MNIST 改成单通道、$224^2$ & 原架构的通道和输入尺寸遵循其自己的约定\\
VAE 重建／随机生成 & 从给定图片近似后验取潜变量 & 从训练时约定的先验取潜变量\\
\bottomrule
\end{tabularx}

\section{常见错误与陷阱}
\begin{itemize}
\item 实验报告和同目录代码有多个版本时先说明版本。不要把 \texttt{logit.py} 的 WDBC 与 \texttt{lab2.py} 的 $L_1$／Skin 场景写成同一次固定实验。
\item Lab1 的 \texttt{data} 先从全部生成输入取最大最小再切分，是实际代码顺序。严格复现实验框架时可预先声明理想抽样区间。严格评估应从训练组估变换，才方便解释独立验证。
\item \texttt{logit.py} 的一段正则系数比较直接用测试准确率选优，\texttt{AlexNet.py} 保存各轮最高测试准确率。将其当作探索指标或另留测试集。作答「如何设计公平验证」时用验证集选择、测试集只做最终评估。
\item Lab3 中概率很小导致行和接近零时，仅给分母加极小量可能使责任度行和偏离 1。用对数和规范化更清楚。空簇回退、协方差正定检查、比较多起点也都值得写。
\item CNN 必须区分 logits 与 Softmax 概率。交叉熵实现如果已含 log-softmax，输入直接用 logits。评估时关闭训练态 Dropout，必要时固定种子并报告重复运行的均值。
\item VAE 图像好看和单个训练损失低不能简单等同。按相同归约、相同数据划分看 KL 与重建，也应查随机先验样本是否合理，评估隐变量是否被利用。
\end{itemize}

\section{往年真题与参考答案}
\begin{description}
\item[2019 秋第 7 题] 实验中协方差奇异。答出 GMM 目标、责任度、加权均值与协方差。解释秩不足与分量塌缩。给加 $\epsilon I$、重置小簇、换对角协方差等操作，并说明用似然和重复试验检查。
\item[2019 秋第 8、9 题] 特征映射、内积的例子。结合实验一幂基对输入建非线性模型但对参数线性，或实验四隐层把原二维螺旋映射为更适合线性打分的表示。内积可举线性得分、卷积窗口、核方法中映射后内积。题目若问是否在自做实验中实际使用，区分自己实现过的变换与理论例子。
\item[2020 春第 1.1、2.1、3 题] 工程题分别涉及小样本地形识别、新颖性检测、拍照去噪。思路按数据与任务、模型与损失、独立验证三个层次写。地形遮挡需要单独验证危险误分。新颖性先识别再分组。去噪用配对干净目标与对应像素损失。
\item[2021 秋第 4、6 题] 对比实验中 $K$ 均值硬分配和 GMM 责任度。后者可改变协方差，出现协方差奇异时解释低秩并写修正。比较同一数据、同 $K$ 的多次运行和外部指标。
\item[2022 秋第 4 题] 结合实验解释过拟合，给两种办法。模板：写实验一的正弦训练集及幂特征。说高阶在有限含噪样本上追随噪声，训练残差小但验证误差大。加 $L_2$ 收缩大系数、验证选低阶、增加覆盖充分的样本三选二，说出机制与验证曲线怎么变化。
\item[2024 秋大题 2、3、9] 第 2 题区分拟合与测试误差，已知预测 $12.8,20$、真实值 $14.8,22$，MSE 为 $4$。第 3 题说明逻辑回归后验和线性判别，阈值可按损失改变。第 9 题用实验三代码顺序解释 E 步逐样本、M 步逐成分。
\item[2025 秋补缓考简答 4、8] 第 4(1) 问原计算数据缺失，应写残差平方和除以样本数。第 4(2) 问答逻辑回归用二元交叉熵 $-y\ln p-(1-y)\ln(1-p)$。第 8 题写 MSE 及梯度 $2X^\top(Xw-y)/n$，然后用实验一解释高阶模型把噪声拟合进去，写验证曲线与增加样本、降阶或加惩罚的解决机制。
\end{description}

\Needspace{18\baselineskip}
\section{本章自测}
\begin{enumerate}
\item 自己报告「增加训练样本使泛化提高」。要固定哪些比较条件？\\
\textbf{答} 固定测试划分、模型容量、主要预处理、训练策略与评价指标。训练样本子集可用相同抽样规则，重复试验并给误差波动。
\item 实验一训练曲线一直下降，独立验证曲线先下后上。怎么选轮次和最终评估？\\
\textbf{答} 用验证最低附近的轮次与参数，固定选择后只在测试集评估一次。
\item 实验二权重无正则时，$x=(3,-2)$、$p=0.25$、$y=1$，偏置梯度是多少？\\
\textbf{答} $p-y=-0.75$。
\item 实验三某 GMM 的责任度矩阵是 $200\times3$。第 2 列和为 80，新权重是多少？\\
\textbf{答} $80/200=0.4$。先检查各样本责任度行和约为 1。
\item 实验五步长 2 的池化前一层边长为 19，核 3，无填充。输出边长与参数量？\\
\textbf{答} 边长 $\lfloor(19-3)/2\rfloor+1=9$，普通最大池化没有可训练参数。
\item 实验六编码器输出 $\mu=0,\log v=0$。单维 KL 和标准差是多少？\\
\textbf{答} $v=1,\sigma=1$，后验等于标准正态先验，KL 为 0。
\end{enumerate}

\chapter{考前速查与答题模板}

\section{本章考点清单}
\begin{itemize}
\item ★★★ 在相同记号下默写贝叶斯后验和条件风险、树分裂、回归梯度、支持向量条件、EM 三类更新、HMM 前后向、链式梯度和 ELBO。 【2024秋真题 大题2–9】【2025秋补缓考 简答1–4、6–8】【模拟卷一至四 大题】
\item ★★★ 看到数据就建立中间表：真实标签与预测格、SVM 约束、核窗口位置、时间状态和路径、深层张量维度。 【刘远超·模拟试题范本 p2】【模拟卷一 大题22–25、28】【模拟卷二 大题1–3、5、8】【模拟卷三 大题1–3、5、8】【模拟卷四 大题1–3、5、8】
\item ★★ 区分名称相似的方法并写适用条件。典型是硬软间隔、L1/L2、滤波和解码、AE 和 VAE、归一化轴。 【2024秋真题 大题4–5、8】【2025秋补缓考 简答1、7–8】【模拟卷四 大题1(8–9)、大题4】
\item ★★ 按「定义—机制—结果—局限」写简答，结合实验时增加可对照的证据。 【2019秋真题 第7–9题】【2022秋真题 第4题】【2024秋真题 大题9】
\end{itemize}

\section{核心概念与定义}
这章用每种方法的\textbf{输入、目标、决定式和检查量}串起全课程。预测问题最终要交付数值或标签。估计问题要给出训练时求到的参数。解释题要给出机制成立的条件。例如「最小风险」的输入除了概率还要有损失矩阵，「KNN 密度」需要搜索窗的体积，「VAE 生成」需要模型先验和解码器。把补充条件写在公式前面，避免把同一符号下不同任务的对象接起来。出处随每组公式列在表中，推导和两版记号可回查前章。

\section{公式与完整推导}
\subsection{全课程公式表}
取每表内的符号为该行定义，跨行不强制同名。默认 $\ln$，如题给底 2 或 10，熵和 KL 须一起换底。$n$ 通常指样本数，$d$ 是特征维数，$K$ 是类或簇数。

\begin{longtable}{@{}p{.19\linewidth}p{.76\linewidth}@{}}
\toprule
知识块 & 核心公式与成立条件\\
\midrule
\endfirsthead
\toprule 知识块 & 核心公式与成立条件\\ \midrule
\endhead
\bottomrule
\endfoot
分组与指标 & $z=(x-\mu_{\rm tr})/\sigma_{\rm tr}$，min-max 为 $(x-a_{\rm tr})/(b_{\rm tr}-a_{\rm tr})$。归一化的拟合只见训练段。$\mathrm{Acc}=(TP+TN)/n$，$P=TP/(TP+FP)$，$R=TP/(TP+FN)$，$F_1=2PR/(P+R)$。ROC 横坐标 $FP/(FP+TN)$，纵坐标 $TPR=R$。有并列时，AUC 正负配对的同分计半。MSE 是平方和除以 $n$，MAE 是绝对值和除以 $n$。平方误差期望分解为偏差平方、方差和不可约噪声。 【第1章】【刘远超·2 机器学习概述 p28–34】【刘远超·3 分类及回归问题 p9–18、p53–58】\\[.4em]
概率与信息 & $P(y\mid x)\propto p(x\mid y)P(y)$。MLE 最大化 $\sum_i\ln p(x_i\mid\theta)$。MAP 加 $\ln p(\theta)$。$H(P)=-\sum p_i\log p_i$，$H(P,Q)=-\sum p_i\log q_i$，$D_{\rm KL}(P\Vert Q)=H(P,Q)-H(P)$，$I(X;Y)=H(X)-H(X\mid Y)$。 【第2章】【刘扬·2 概率论速通】【2024秋真题 大题6】【刘远超·模拟试题范本 p2 第1题】\\[.4em]
决策树 & $H(D)=-\sum_kp_k\log p_k$。给分支 $D_j$，增益 $H(D)-\sum_j|D_j|H(D_j)/|D|$。增益比再除以分支比例的熵。基尼 $1-\sum_kp_k^2$，分裂基尼是各子节点加权和。CART 代价复杂度剪枝常写 $R(T)+\alpha|T_{\rm leaves}|$。 【第3章】【刘扬·1 决策树】\\[.4em]
贝叶斯决策 & $R(a\mid x)=\sum_jL(a,y_j)P(y_j\mid x)$，决策取 $\arg\min_a$。高斯对数判别 $g_k=-\tfrac12(x-\mu_k)^\top\Sigma_k^{-1}(x-\mu_k)-\tfrac12\ln|\Sigma_k|+\ln\pi_k$（略共同常数）。协方差相同消去共同二次项。 【第4章】【刘扬·3 贝叶斯判别】【2022秋真题 第5题】\\[.4em]
非参数 & 核密度 $\hat p(x)=\frac1{nh^d}\sum_iK((x-x_i)/h)$。矩形核数点后除以固定体积，KNN 密度 $k/(nV_k)$ 由近邻决定体积。KNN 分类按投票或加权投票，$k$ 与距离度量属于超参数。 【第5章】【刘扬·3 贝叶斯判别与KNN】\\[.4em]
朴素贝叶斯 & $\hat y=\arg\max_y(\ln P(y)+\sum_i\ln p(x_i\mid y))$。多项式词概率 $(N_{yi}+\alpha)/(N_y+\alpha n)$ 并按词频乘对数。伯努利每位为 $q_i^{x_i}(1-q_i)^{1-x_i}$，$q_i=(M_{yi}+1)/(M_y+2)$。高斯 $\hat\mu=\sum x/n_y$，$\hat v=\sum(x-\hat\mu)^2/n_y$。 【第6章】【刘远超·3 分类及回归问题】【刘扬·4 朴素贝叶斯 p25–28】\\[.4em]
回归与正则 & $f=Xw$，半平方和梯度 $X^\top(Xw-y)$。岭回归 $(X^\top X+\lambda I)w=X^\top y$（是否罚偏置依设定）。逻辑回归 $\sigma(a)=1/(1+e^{-a})$，BCE 对 $a$ 的梯度 $p-y$。多类中类别权重按行写时 $a=Wx+b$，$p_k=e^{a_k}/\sum e^{a_j}$，交叉熵对得分的梯度 $\bm p-\bm y$。Lasso 加 $\lambda\|w\|_1$。 【第7章】【刘扬·5 逻辑回归与SVM】【2024秋真题 大题2–3】\\[.4em]
感知机、SVM & 感知机分错时 $w\gets w+\eta yx,\ b\gets b+\eta y$。以更新次数展开成对偶和。硬间隔 $y_i(w^\top x_i+b)\ge1$，总宽度 $2/\|w\|$。软间隔最小化 $\|w\|^2/2+C\sum_i[1-y_if_i]_+$。核对偶 $\sum_i\alpha_i-\tfrac12\sum_{ij}\alpha_i\alpha_jy_iy_jK_{ij}$，$\sum_i\alpha_iy_i=0$、盒上界 $C$，预测分数 $\sum_i\alpha_i y_iK(x_i,x)+b$。 【第8、9章】【刘远超·6 感知机】【刘扬·5 逻辑回归与SVM】【2024秋真题 大题4–5】\\[.4em]
聚类与验证 & $K$ 均值最小化 $\sum_i\|x_i-\mu_{z_i}\|^2$。中心取簇均值。$K$-means++ 按 $D_i^2$ 归一化。纯度 $\sum_k\max_c n_{ck}/n$，$\mathrm{RI}=(TP+TN)/\binom n2$（按配对），轮廓 $s=(b-a)/\max(a,b)$。链接取跨簇点对的 min/max/mean。DBSCAN 以 $\varepsilon$ 和 MinPts 定义核心、边界、噪声。 【第10章】【刘远超·4 聚类问题 p8–25】【2025秋补缓考 简答2】\\[.4em]
GMM 与 EM & $\gamma_{ik}=\pi_k\mathcal N_k(x_i)/\sum_j\pi_j\mathcal N_j(x_i)$，$N_k=\sum_i\gamma_{ik}$，$\pi_k'=N_k/n$，$\mu_k'=\sum_i\gamma_{ik}x_i/N_k$，$\Sigma_k'=\sum_i\gamma_{ik}(x_i-\mu_k')(x_i-\mu_k')^\top/N_k$。$Q=\mathbb E_{\rm old}\ln p(X,Z\mid \theta)$。AIC 为 $2d_f-2\hat\ell$，BIC 为 $d_f\ln n-2\hat\ell$。 【第10章】【刘远超·4 聚类问题 p36–49】【2024秋真题 大题9】\\[.4em]
降维 & $S=X_c^\top X_c/n$（题定分母优先），最大化 $u^\top Su$ 且 $\|u\|=1$ 得 $Su=\lambda u$。一维投影 $u^\top x_c$，重构 $\bar x+uu^\top x_c$，最小误差为舍弃特征值之和。LDA 两类方向 $\propto S_W^{-1}(\mu_1-\mu_2)$，区别在是否使用类标签和类内散布。 【第11章】【刘扬·9PCA】【模拟卷三 大题2(1–3)】\\[.4em]
HMM & $\alpha_1(i)=\pi_i b_i(o_1)$。$\alpha_{t+1}(j)=b_j(o_{t+1})\sum_i\alpha_t(i)a_{ij}$。似然 $\sum_i\alpha_T(i)$。$\beta_T=1$。$\beta_t(i)=\sum_j a_{ij}b_j(o_{t+1})\beta_{t+1}(j)$。$\gamma_t(i)\propto\alpha_t(i)\beta_t(i)$，$\xi_t(i,j)\propto\alpha_t(i)a_{ij}b_j(o_{t+1})\beta_{t+1}(j)$。Viterbi 把内层求和换成最大并存回溯指针。 【第12章】【刘远超·5 隐马尔科夫模型 p17–20】【刘扬·12HMM p20–34】\\[.4em]
HMM 数值与学习 & 令 $c_t$ 为前向逐步缩放前向量之和，则 $\ln P(O)=\sum_t\ln c_t$。Baum–Welch：$\pi_i'=\gamma_1(i)$，$a_{ij}'=\sum_{t<T}\xi_t(i,j)/\sum_{t<T}\gamma_t(i)$，$b_i'(v)=\sum_t\gamma_t(i)[o_t=v]/\sum_t\gamma_t(i)$。 【第12章】【刘扬·12HMM p28–32】【模拟卷四 大题8】\\[.4em]
前馈与 BP & $z^{(l)}=W^{(l)}a^{(l-1)}+b^{(l)}$，$\delta^{(l)}=(W^{(l+1)\top}\delta^{(l+1)})\odot\phi'(z^{(l)})$。梯度 $\delta^{(l)}a^{(l-1)\top}$ 与 $\delta^{(l)}$。批次取平均时都要除以批大小。ReLU 导数正区为 1，负区为 0。Xavier 方差常取 $2/(n_{\rm in}+n_{\rm out})$，He 常取 $2/n_{\rm in}$。 【第13章】【刘扬·10 前馈神经网与反向传递算法】【模拟卷三、四 大题1】\\[.4em]
网络优化 & 动量可用 $u_t=\rho u_{t-1}+g_t$，$\theta\gets\theta-\eta u_t$。AdaGrad 加各步平方。RMSProp 用指数平方均值。Adam 有 $\hat m=m/(1-\beta_1^t)$、$\hat v=v/(1-\beta_2^t)$，更新 $\theta\gets\theta-\eta\hat m/(\sqrt{\hat v}+\epsilon)$。梯度裁剪乘 $\min(1,c/\|g\|)$。归一化 $\gamma(x-\mu)/\sqrt{\sigma^2+\epsilon}+\beta$，BN 按特征统计批，LN 按单样本内特征。inverted Dropout 训练乘 $m/(1-p)$，推理直接用原激活。 【第14章】【刘远超·7 网络优化】【模拟卷三、四 大题1】\\[.4em]
CNN & $k_{\rm eff}=d(k-1)+1$，$H'=\lfloor(H+2p-k_{\rm eff})/s\rfloor+1$。全通道 $k_h\times k_w$ 卷积参数 $C_{\rm out}(k_hk_wC_{\rm in}+1)$。全局平均池化在空间上取均值，最大池化梯度回获胜位置。串联感受野 $j_l=j_{l-1}s_l,\ r_l=r_{l-1}+(k_{\rm eff}-1)j_{l-1}$。 【第15章】【刘扬·11 卷积神经网络 p18–22】【刘远超·9 卷积神经网络 p23、p30】\\[.4em]
序列与注意力 & $h_t=f(Ux_t+Wh_{t-1}+b)$。逆向需用 $h_{t+1}$，双向并行拼接。LSTM $c_t=f_t\odot c_{t-1}+i_t\odot\tilde c_t$, $h_t=o_t\odot\tanh c_t$。GRU 示例约定 $h_t=(1-z_t)h_{t-1}+z_t\tilde h_t$。注意力 $\operatorname{softmax}(QK^\top/\sqrt{d_k}+M)V$，标准无偏置多头投影总参数 $4d^2$（总头维均为 $d$），稠密注意力主项 $O(n^2d)$。 【第16章】【刘远超·8 建模序列数据的神经网络】【2025秋补缓考 简答6–7】\\[.4em]
AE 与 VAE & AE 选 $\hat x=d(e(x))$ 降重构误差。VAE 下界 $\mathbb E_q[\ln p_\theta(x\mid z)]-D_{\rm KL}(q\Vert p(z))$。高斯对标准正态 KL 为 $\frac12\sum_j(\mu_j^2+\exp\ell_j-1-\ell_j)$，$z=\mu+e^{\ell/2}\odot\epsilon$。先验采样生成、近似后验采样重建。条件 VAE 将条件纳入编码、解码及约定的先验。 【第17章】【刘扬·14VAE】【刘远超·10 生成式神经网络】【2024秋真题 大题7】\\
\end{longtable}

\subsection{考场推导的最小闭环}
有三个式子经常一问分多步。对 Softmax 先把损失写成 $\ln\sum e^{a_j}-\sum y_ja_j$，再求 $\partial/\partial a_k$。对 EM 先写含指示的完全数据对数似然，再用条件期望替换指示，最后固定责任度极大化。对 PCA 先中心化并把方差写成 $u^\top S u$，再引入单位约束推出特征方程。每条链都写出中间等式，最后才给出结论。【刘扬·8GMM-EM p19–38】【刘扬·9PCA】【刘扬·10 前馈神经网与反向传递算法】【见第10、11、13章】

\section{计算题型与解题模板}
\paragraph{选择公式前先问方向。}
原创例：正样本得分为 $.83,.42$，负样本为 $.72,.41,.11$，高分判正。扫阈值时顺序为正、负、正、负、负。AUC 可按六个正负对算，$.83$ 赢三次，$.42$ 赢两次，所以 $\mathrm{AUC}=5/6$。阈值 $.5$ 时 $TP=1,FP=1,FN=1,TN=2$，因此 $P=1/2,R=1/2,F_1=1/2$。同一组分数可求整条 ROC 或指定阈值下的分类质量，先判断所求。

\paragraph{验算形状时从输出反推。}
原创例：输入 $9\times10\times4$，用 7 个 $3\times2$ 核，双向步长 2，填充高 1、宽 0，无空洞。高为 $\lfloor(9+2-3)/2\rfloor+1=5$，宽 $\lfloor(10-2)/2\rfloor+1=5$，输出 $5\times5\times7$。若有偏置，参数 $7(3\cdot2\cdot4+1)=175$。偏置是每输出通道一项，不随滑动位置增加。

\paragraph{简答题模板。}
\begin{description}
\item[解释某模型] 「对象是什么」给一句定义。「怎么用」写输入到输出的对应式。「为什么」指出目标或假设。「适用」写一种优势和一种失败条件。例如软间隔 SVM 解读违反约束的点及 $C$ 的权衡。
\item[比较两种方法] 按监督信号、假设、训练目标、输出、计算或适用条件各写一项。比如 LSTM 与 GRU 先写状态和门，之后写更新式，最后才谈参数、依赖长度和实验选择。
\item[结合实验] 拿第 18 章的「数据和划分—可重现操作—观察现象—推理依据—对照验证」。若无原始图上的实测数，描述趋势和判断依据。
\item[为何有下界或单调] 先写两个量的精确差值，指出非负项。再写等号成立条件和一次更新后的不等式链。避免凭经验声称一定收敛到全局。
\end{description}

\section{易混辨析}
\begin{tabularx}{\linewidth}{@{}p{.23\linewidth}XX@{}}
\toprule
配对概念 & 第一项 & 第二项\\
\midrule
交叉熵／KL & 用 $P$ 加权 $-\log Q$ & 从交叉熵扣 $H(P)$，有方向\\
最小错误／最小风险 & 误分代价一样，取最大后验 & 给每个误分不同代价，比较期望损失\\
多项式／伯努利 NB & 词频重复增加证据 & 只记录出现与否，缺词也贡献\\
L1／L2 & 可产生精确零权重，拉普拉斯先验 & 平滑收缩，高斯先验\\
感知机／逻辑回归／SVM & 感知机看错分，逻辑回归拟合条件概率 & SVM 优化几何间隔或合页与范数\\
$K$ 均值／GMM & 最近中心的硬分配 & 责任度的软分配和协方差\\
PCA／LDA & 无监督最大方差 & 有监督类间与类内散布的比值\\
前向／Viterbi & 对路径求和得到似然 & 最大联合概率加回溯\\
BN／LN & 同一通道跨样本（CNN 常含空间） & 每个样本在特征维度内统计\\
LSTM／GRU & 有独立细胞状态和输出门 & 将记忆并入隐状态，更新门与重置门\\
AE／VAE & 确定性编码、重构训练 & 概率后验、先验约束，可从先验生成\\
\bottomrule
\end{tabularx}

\section{常见错误与陷阱}
\begin{itemize}
\item 分类概率先看条件，$\Pr(y\mid x)$ 和 $p(x\mid y)$ 方向不同。KL 的两个分布顺序同样不可交换。
\item 核宽、通道、有效感受野与偏置分别计数。Sigmoid 与交叉熵组合时可化简 $p-y$。若模型换成 Sigmoid 与 MSE，要保留额外导数。
\item 维度、平均、缩放、归一化对每个数据集和实现都可能不同。训练和推理的 BatchNorm、Dropout 分开写。
\item 聚类外部评价对样本对或列联表，PCA 重构先中心化后加回均值，EM 坐标更新先旧后验再新参数。
\end{itemize}

\clearpage
\thispagestyle{fancy}
\noindent{\LARGE\bfseries 考前一页纸}\par
\medskip
\noindent 考前最后一遍，遮住右栏，按左栏闭卷复述并写公式。

\bigskip
\begin{tabularx}{\linewidth}{@{}p{.22\linewidth}X@{}}
\toprule
看到题干 & 先写的式子或条件\\
\midrule
条件概率、代价 & 先验乘类条件再归一。风险等于后验加权的损失。贝叶斯对数差检查协方差。\\[.6em]
熵、树划分 & 底数、方向和子节点占比。$D_{\mathrm{KL}}=H(P,Q)-H(P)$。增益减条件熵，增益比再除分支熵。\\[.6em]
近邻、NB & 明确距离和窗口的体积。词频反复计，伯努利缺席乘补概率，高斯方差按指定估计。\\[.6em]
回归、感知机 & 二次损失的分母。先写 $\hat y,y,e$，再梯度。二元概率 sigmoid，误差 $p-y$，感知机用 $y$ 更新错分点。\\[.6em]
SVM & $y_if_i\ge1-\xi_i$，$\xi\ge0$。几何距离除以 $\|w\|$。核进入内积，支持向量检查 $\alpha$ 和 KKT。\\[.6em]
聚类、PCA & 分配后更新均值。GMM 先责任度后权重、均值、协方差。PCA 去均值，投影，重构加回均值。\\[.6em]
HMM、序列 & $\alpha$ 对路径求和，Viterbi 取最大并回溯。后向用下一观测。RNN 按实际箭头写前一计算得到的状态。\\[.6em]
BP、正则 & 先前向缓存 $z$ 和 $a$。输出残差、反传、逐项激活导数、外积和偏置，批次梯度记分母。\\[.6em]
卷积、注意力 & 明确 $H,W,C$ 与有效核。一个输出通道一个偏置。$QK^\top/\sqrt{d_k}$，掩码在 softmax 前。\\[.6em]
AE、VAE、实验 & 对照重建对象、先验和后验。负 ELBO 为重构损失加 KL。实验答数据、模型、现象、解释、验证。\\
\bottomrule
\end{tabularx}

\bigskip
\noindent\textbf{交卷前三项：} 数值的方向与分母，矩阵的行列与偏置，缺失条件与答题假设。
\clearpage

\section{往年真题与参考答案}
\begin{description}
\item[2024 秋大题 2、6、9] MSE 要把两行样本分别算完再平均。KL 先写 $P/Q$ 顺序和底数。EM 先定义隐变量与完全数据，再写 E 步责任度和 M 步的三组参数。对应第 2、7、10 章。
\item[2024 秋大题 4–5、7–8] 第 4–5 题按分离线、支持向量、几何距离、软间隔约束、核依次写。第 7 题说清 AE 的输出长度等于重构输入的特征数。第 8 题是 RNN，确认题面的箭头从右向左再写前向。对应第 9、16、17 章。
\item[2025 秋补缓考简答 2、6–7] 客户聚类说明噪声、簇形与各自优缺点。序列题区分词、状态、得分。LSTM 给状态和门。Transformer 给注意力而非仅罗列子模块。对应第 10、16 章。
\item[2026 试题范本第 2 页] 第 1 题 KL 按底 10 保留两位。第 2 题整句取最后读到输入时的 $s_3$，再计算 softmax，列向量条件下 $V$ 是 $10\times200$。第 3 题 Iris 的先验每类 $1/3$，从第 2 类第 3 列样本用 MLE 估均值和方差，再代入 $2.5$ 的高斯密度。 【刘远超·模拟试题范本 p2】
\end{description}

\section{本章自测}
\begin{enumerate}
\item 把两个分布一起用对数底 $b>1$，从自然对数换底时 KL 如何变化？\\
\textbf{答} $D_b=D_e/\ln b$。相同次序的大小比较保持。
\item CNN 的两个输出通道各有偏置，输出空间变为原来一半，偏置参数数目是多少？\\
\textbf{答} 仍为 2，每输出通道共享一个。
\item 若 $z_t$ 与 $z_{t+1}$ 是 HMM 两时刻隐状态，$\xi_t$ 的行和对应哪个后验？\\
\textbf{答} $\sum_j\xi_t(i,j)=\gamma_t(i)$，即较早时刻的后验。
\item 一维线性 AE 以中心化输入训练并使用平方误差，最优编码方向与第一主轴相反可以吗？\\
\textbf{答} 可以，方向取反只改变投影坐标符号，投影矩阵和重构不变。
\item 给 Softmax 三类的所有 logit 同加 3，概率怎样变？\\
\textbf{答} 分子和分母都乘 $e^3$，概率不变。
\item 标签先验 $\pi_1=0.6,\pi_2=0.4$，类条件密度相等。最小错误率选哪类？若唯一关注第 2 类的漏判，而漏判代价比误报大四倍，最小风险选哪类？\\
\textbf{答} 后验还等于先验，最小错误率选 1。设置误报代价为 1，判 1 的风险 $4\cdot0.4=1.6$，判 2 的风险 $0.6$，故选 2。
\end{enumerate}

\end{document}
